% Options for packages loaded elsewhere
\PassOptionsToPackage{unicode}{hyperref}
\PassOptionsToPackage{hyphens}{url}
\PassOptionsToPackage{dvipsnames,svgnames,x11names}{xcolor}
\documentclass[
  english,
  11pt,
  a4paper,
]{article}
\usepackage{xcolor}
\usepackage[margin=20mm]{geometry}
\usepackage{amsmath,amssymb}
\setcounter{secnumdepth}{-\maxdimen} % remove section numbering
\usepackage{iftex}
\ifPDFTeX
  \usepackage[T1]{fontenc}
  \usepackage[utf8]{inputenc}
  \usepackage{textcomp} % provide euro and other symbols
\else % if luatex or xetex
  \usepackage{unicode-math} % this also loads fontspec
  \defaultfontfeatures{Scale=MatchLowercase}
  \defaultfontfeatures[\rmfamily]{Ligatures=TeX,Scale=1}
\fi
\usepackage{lmodern}
\ifPDFTeX\else
  % xetex/luatex font selection
\fi
% Use upquote if available, for straight quotes in verbatim environments
\IfFileExists{upquote.sty}{\usepackage{upquote}}{}
\IfFileExists{microtype.sty}{% use microtype if available
  \usepackage[]{microtype}
  \UseMicrotypeSet[protrusion]{basicmath} % disable protrusion for tt fonts
}{}
\makeatletter
\@ifundefined{KOMAClassName}{% if non-KOMA class
  \IfFileExists{parskip.sty}{%
    \usepackage{parskip}
  }{% else
    \setlength{\parindent}{0pt}
    \setlength{\parskip}{6pt plus 2pt minus 1pt}}
}{% if KOMA class
  \KOMAoptions{parskip=half}}
\makeatother
\usepackage{longtable,booktabs,array}
\newcounter{none} % for unnumbered tables
\usepackage{calc} % for calculating minipage widths
% Correct order of tables after \paragraph or \subparagraph
\usepackage{etoolbox}
\makeatletter
\patchcmd\longtable{\par}{\if@noskipsec\mbox{}\fi\par}{}{}
\makeatother
% Allow footnotes in longtable head/foot
\IfFileExists{footnotehyper.sty}{\usepackage{footnotehyper}}{\usepackage{footnote}}
\makesavenoteenv{longtable}
\ifLuaTeX
\usepackage[bidi=basic,shorthands=off]{babel}
\else
\usepackage[bidi=default,shorthands=off]{babel}
\fi
\ifLuaTeX
  \usepackage{selnolig} % disable illegal ligatures
\fi
\setlength{\emergencystretch}{3em} % prevent overfull lines
\providecommand{\tightlist}{%
  \setlength{\itemsep}{0pt}\setlength{\parskip}{0pt}}
\usepackage{amsmath,amssymb}
\usepackage{mathrsfs}
\usepackage{newunicodechar}
\newunicodechar{₂}{\textsubscript{2}}
\newunicodechar{∎}{\ensuremath{\blacksquare}}
\newunicodechar{□}{\ensuremath{\square}}
\allowdisplaybreaks
\setlength{\emergencystretch}{3em}
\sloppy
\relpenalty=100
\binoppenalty=100
\usepackage{fvextra}
\usepackage{needspace}
\setmainfont{LibertinusSerif-Regular.otf}[BoldFont=LibertinusSerif-Bold.otf,ItalicFont=LibertinusSerif-Italic.otf,BoldItalicFont=LibertinusSerif-BoldItalic.otf]
\setmonofont{DejaVuSansMono.ttf}
\makeatletter\renewcommand{\@pnumwidth}{3em}\renewcommand{\@tocrmarg}{4em}\makeatother
\RecustomVerbatimEnvironment{verbatim}{Verbatim}{breaklines=true,breakanywhere=true,fontsize=\small}
\AtBeginDocument{\pdfstringdefDisableCommands{\def\varepsilon{epsilon}}}
\usepackage{bookmark}
\IfFileExists{xurl.sty}{\usepackage{xurl}}{} % add URL line breaks if available
\urlstyle{same}
\hypersetup{
  pdftitle={Haar measure and quotient integration},
  pdfauthor={Original authors and AI contributors credited within each chapter},
  pdflang={en},
  colorlinks=true,
  linkcolor={Maroon},
  filecolor={Maroon},
  citecolor={Blue},
  urlcolor={Blue},
  pdfcreator={LaTeX via pandoc}}

\title{Haar measure and quotient integration}
\usepackage{etoolbox}
\makeatletter
\providecommand{\subtitle}[1]{% add subtitle to \maketitle
  \apptocmd{\@title}{\par {\large #1 \par}}{}{}
}
\makeatother
\subtitle{4 course lessons, 7 prerequisite chapters and 1 supplementary
proof · Published four-lesson companion; source authors self-checked; no
independent AI or human review claimed}
\author{Original authors and AI contributors credited within each
chapter}
\date{3 October 2026}

\begin{document}
\maketitle

{
\setcounter{tocdepth}{2}
\tableofcontents
}
\section*{About this edition}\label{about-edition}
\addcontentsline{toc}{section}{About this edition}

This downloadable edition contains the 4 course lessons, 7 prerequisite
chapters and 1 supplementary proof in the selected reader edition of
\emph{Haar measure and quotient integration}. It preserves the included
lesson text, mathematics, references, exercises and supplied solutions.

Source status: Published four-lesson companion; source authors
self-checked; no independent AI or human review claimed. Includes all
four lessons, all seven bundled prerequisite chapters and the
Fremlin-derived supplementary proof. It is not the complete
harmonic-analysis course, and further reading listed by the source
remains outside this companion.

Original writing provenance: Claude Opus 5.5 (Anthropic; original Haar
lesson, effort not recorded); OpenAI Codex --- GPT-6 Astra, Ultra
(quotient lesson); GPT-6.1 Sol, Ultra (measure tools, arithmetic
integration, corrections and self-check). Each bundled chapter retains
its own author statement.. Source-preserving format conversion and
packaging: OpenAI Codex --- GPT-6 Astra, Ultra effort (source-preserving
format conversion). The source-specific author, contributor and review
statements remain authoritative. Format validation adds no claim of
human editing, independent mathematical review or course completion.
This is a source-preserving edition, not an everyday-English rewrite.

\href{https://kokunoyumeto.github.io/open-math-courses/courses/harmonic-analysis-on-locally-compact-groups/reader/}{Read
the course online} ·
\href{https://kokunoyumeto.github.io/program-matematika-indonesia/en/}{Return
to the mathematics programme}.

The four course lessons and seven bundled programme chapters are CC0
1.0, with all original human source credits and AI author statements
retained. The supplementary proof Haar uniqueness through shrinking
neighborhoods adapts D. H. Fremlin and remains under the Design Science
License. Its source and change notices accompany the proof, and the
complete licence is included. This mixed collection does not relicense
Fremlin's work as CC0.

The PDF and EPUB share the same ordered content selection. Links within
the included course selection work offline; external scholarly and
prerequisite links require access to their destinations. The source
manifest records the exact lesson and source-edition identities.

\section{Measure and Hilbert space tools for Haar
integration}\label{measure-and-hilbert-space-tools}

\emph{Written by GPT-6.1 Sol (OpenAI), Ultra reasoning effort, October
2026. Self-checked by the writing AI, GPT-6.1 Sol. Independent AI review
remains separate. Original exposition is public domain (CC0).}

Haar integration needs convergence theorems even when the group has
uncountably many open components. It also needs a precise meaning for a
tensor product of Hilbert spaces. This lesson proves the measure results
used in the next lesson and explains which parts require countability.
The finite-measure arguments are deliberately separated from the
arguments valid on every measure space.

The earlier programme lesson
\hyperref[hilbert-spaces-and-compact-operators]{Hilbert spaces and
compact operators}, Theorems 2.1--2.3, 4.1 and 8.1, proves orthogonal
projection, representation of a bounded Hilbert-space functional,
orthonormal bases and the Hilbert tensor product. Those complete proofs
are the Hilbert-space prerequisites below. Elementary set theory,
including the axiom of choice, is assumed. No topology is needed until
the following Haar lesson. The core course
\href{https://kokunoyumeto.github.io/program-matematika-indonesia/en/\#course-D10}{Measure
and Integration} uses Fremlin's human text. The core
\href{https://kokunoyumeto.github.io/program-matematika-indonesia/en/\#course-C10}{Real
Analysis I} and
\href{https://kokunoyumeto.github.io/program-matematika-indonesia/en/\#course-C20}{Real
Analysis II} use J. Lebl's open \emph{Basic Analysis}. Our Lemma 6.1
supplies the one-variable integral identities used here. Basic
references for measure theory are {[}Fremlin{]} and {[}Lebesgue{]}.

\subsection{1. From an outer measure to a
measure}\label{measure-and-hilbert-space-tools--1-from-an-outer-measure-to-a-measure}

A sigma-algebra on a set \(X\) is a family containing \(X\) and closed
under complements and countable unions. A measure on it is a function
into \([0,\infty]\) taking the empty set to zero and countably additive
on disjoint sets. An outer measure \(m^*\) is defined on every subset of
\(X\), is zero at the empty set, is monotone, and satisfies
\[m^*\left(\bigcup_n A_n\right)\leq\sum_n m^*(A_n).\]

\textbf{Theorem 1.1 (Carathéodory).} The sets \(E\subseteq X\)
satisfying \[m^*(A)=m^*(A\cap E)+m^*(A\setminus E)\qquad(A\subseteq X)
\tag{1.1}\] form a sigma-algebra. The restriction of \(m^*\) to it is a
complete measure. To verify (1.1) it suffices to prove its
greater-than-or-equal inequality for \(A\) with finite outer measure.

\emph{Proof.} The other inequality is subadditivity; when
\(m^*(A)=\infty\), the greater-than-or-equal inequality is automatic.
Complements preserve the condition. If \(E,F\) satisfy it, split \(A\)
first by \(E\), and then split \(A\setminus E\) by \(F\). The sum of the
first two resulting terms is at least \(m^*(A\cap(E\cup F))\), by
subadditivity. This proves the condition for \(E\cup F\); differences
and finite intersections follow.

Let \(E_n\) be disjoint sets satisfying the condition, and put
\(E=\bigcup_nE_n\). Repeated finite splitting gives
\[m^*(A)=\sum_{n=1}^N m^*(A\cap E_n)+m^*\left(A\setminus\bigcup_{n=1}^N E_n\right)
\geq\sum_{n=1}^N m^*(A\cap E_n)+m^*(A\setminus E).\] Let \(N\) increase.
Subadditivity gives \(\sum_n m^*(A\cap E_n)\geq m^*(A\cap E)\), so \(E\)
satisfies the condition. An arbitrary countable union is reduced to this
case by replacing \(E_n\) with \(E_n\setminus\bigcup_{j<n}E_j\). Thus
the family is a sigma-algebra. Taking \(A=E\) in the finite-splitting
inequality proves countable additivity on the disjoint \(E_n\), together
with subadditivity. If \(N\) has outer measure zero, every subset of
\(N\) satisfies (1.1): monotonicity makes the first term zero, while
monotonicity and subadditivity force the second term to equal
\(m^*(A)\). This also proves completeness. \(\square\)

\textbf{Example.} Assign to a subset of the real line the infimum of the
sums of lengths of countable open interval covers. This is an outer
measure: combine covers with errors \(\varepsilon2^{-n}\) to prove
subadditivity. Splitting every covering interval at a fixed point splits
its total length between the two half-lines; enlarge the split intervals
by errors whose sum is arbitrarily small. This gives the reverse
Carathéodory inequality for each half-line. Half-lines generate the
Borel sigma-algebra, so Theorem 1.1 makes every Borel set measurable.

The measure of \([a,b]\) is \(b-a\). An enlarged interval gives the
upper bound. For any countable open cover of \([a,b]\), compactness
selects a finite subcover. The sum of its interval lengths is at least
\(b-a\): arrange its endpoints in order, and each successive subinterval
of \([a,b]\) is covered by at least one member. Summing lengths proves
the lower bound. Singletons have measure zero, so the same length
formula holds for open and half-open bounded intervals. Translating or
dilating interval covers shows directly that \(m^*(E+t)=m^*(E)\) and
\(m^*(aE)=|a|m^*(E)\) for \(a\ne0\), using the inverse operation for the
reverse inequalities. The resulting Borel Lebesgue measure, or its
completion, is therefore translation invariant. The next lesson instead
constructs an outer measure using continuous functions on an arbitrary
LCH space.

\subsection{2. Integration and convergence without countability
assumptions}\label{measure-and-hilbert-space-tools--2-integration-and-convergence-without-countability-assumptions}

Fix a measure space \((X,\Sigma,\mu)\). A measurable real function is
one for which \(\{f>t\}\in\Sigma\) for every real \(t\). Countable
suprema and infima are measurable: use
\(\{\sup_n f_n>t\}=\bigcup_n\{f_n>t\}\), and negate for infima. Limits
are measurable because \(\liminf f_n=\sup_N\inf_{n\geq N}f_n\). Complex
functions are measurable when their real and imaginary parts are.

For a nonnegative simple function \(s=\sum_{j=1}^r a_j1_{E_j}\), with
disjoint measurable \(E_j\) and \(a_j\geq0\), define
\(\int s=\sum_j a_j\mu(E_j)\), with \(0\cdot\infty=0\). Refining two
partitions proves independence of the presentation, monotonicity, and
additivity for simple functions. For measurable \(f\geq0\), define
\[\int f=\sup\{\int s:0\leq s\leq f,\ s\text{ simple}\}.\] There are
increasing simple \(s_n\to f\): truncate \(f\) at \(n\) and round down
to multiples of \(2^{-n}\). The truncation bounds and the refining
dyadic grids ensure monotonicity.

Countable additivity implies continuity from below: if
\(E_n\uparrow E\), write \(E\) as the disjoint union of
\(E_1,E_2\setminus E_1,\ldots\). Then \(\mu(E_n)\uparrow\mu(E)\). If
\(\mu(E_1)<\infty\) and \(E_n\downarrow E\), apply this to
\(E_1\setminus E_n\) to obtain continuity from above.

\textbf{Theorem 2.1 (Monotone convergence).} If \(0\leq f_n\uparrow f\)
almost everywhere, then \(\int f_n\uparrow\int f\).

\emph{Proof.} Discard one measurable null set to obtain pointwise
inequalities; changing values there changes none of the integrals.
Monotonicity gives \(\lim\int f_n\leq\int f\). For simple \(s\leq f\)
and \(0<t<1\), the sets \(A_n=\{f_n\geq ts\}\) increase and cover
\(\{s>0\}\). Consequently continuity from below on the finitely many
level sets of \(s\) gives \(\int s1_{A_n}\uparrow\int s\). Since
\(f_n\geq ts1_{A_n}\), the limit of their integrals is at least
\(t\int s\). Let \(t\uparrow1\) and then take the supremum over \(s\).
The same argument works when \(\int s=\infty\). \(\square\)

Applying this to increasing simple approximations of \(f,g\geq0\) proves
\(\int(f+g)=\int f+\int g\). Applying it to partial sums proves
\[\int\sum_n f_n=\sum_n\int f_n\qquad(f_n\geq0).
\tag{2.1}\] These identities justify defining the integral of an
integrable real function as \(\int f^+-\int f^-\), and then treating
complex functions by real and imaginary parts. Here integrable means
\(\int|f|<\infty\). Linearity follows from nonnegative additivity by
moving all negative parts to the opposite side of the desired equality.
Choosing a phase so that \(e^{i\theta}\int f\) is nonnegative real gives
\(|\int f|\leq\int|f|\).

\textbf{Theorem 2.2 (Fatou and dominated convergence).} For measurable
\(f_n\geq0\), \[\int\liminf_n f_n\leq\liminf_n\int f_n.\] If complex
measurable \(u_n\to u\) almost everywhere and \(|u_n|\leq g\) almost
everywhere for one integrable \(g\geq0\), then \(u\) is integrable and
\(\int|u_n-u|\to0\). In particular \(\int u_n\to\int u\).

\emph{Proof.} The functions \(v_N=\inf_{n\geq N}f_n\) increase to the
liminf, and \(\int v_N\leq\inf_{n\geq N}\int f_n\). Monotone convergence
proves Fatou. For the second assertion, \(|u|\leq g\) almost everywhere.
Fatou applied to the nonnegative functions \(2g-|u_n-u|\), whose limit
is \(2g\), gives \[2\int g\leq2\int g-\limsup_n\int|u_n-u|.\] All
integrals here are finite, so the limsup is zero. The integral
inequality proved above gives the last assertion. \(\square\)

\textbf{Example.} On an uncountable set with counting measure, an
integrable function has countable support: for each positive integer
\(n\), only finitely many points can have \(|f|>1/n\). Nevertheless the
whole space is not sigma-finite. Theorems 2.1 and 2.2 remain valid. A
later interchange of integrals will need a separate theorem; convergence
alone does not supply it.

\subsection{3. The complete spaces of integrable
functions}\label{measure-and-hilbert-space-tools--3-the-complete-spaces-of-integrable-functions}

For \(1\leq p<\infty\), let \(L^p(\mu)\) be measurable functions modulo
almost-everywhere equality with norm \(\|f\|_p=(\int|f|^p)^{1/p}\). Let
\(L^\infty(\mu)\) consist of essentially bounded functions, with
essential supremum norm. These definitions use the given sigma-algebra,
with no assumption that the measure is complete.

\textbf{Theorem 3.1 (Hölder and Minkowski).} If \(1/p+1/q=1\), including
the endpoints, then \(\int|fg|\leq\|f\|_p\|g\|_q\). For
\(1\leq p\leq\infty\), \(\|f+g\|_p\leq\|f\|_p+\|g\|_p\).

\emph{Proof.} For \(1<p<\infty\), minimization of \(a^p/p-ab+b^q/q\) in
\(a\geq0\) gives Young\textquotesingle s inequality
\(ab\leq a^p/p+b^q/q\). Apply it to \(a=|f|/\|f\|_p\), \(b=|g|/\|g\|_q\)
and integrate. A zero norm means the corresponding function vanishes
almost everywhere. The endpoints follow by bounding the essentially
bounded factor outside its exceptional null set.

For Minkowski with \(1<p<\infty\), the elementary convexity inequality
\((a+b)^p\leq2^{p-1}(a^p+b^p)\) first proves \(f+g\in L^p\). Put
\(h=|f+g|\). Hölder gives \[\int h^p\leq\int(|f|+|g|)h^{p-1}
\leq(\|f\|_p+\|g\|_p)\left(\int h^p\right)^{1/q}.\] Divide when
\(\|h\|_p>0\). The cases \(p=1,\infty\) are the pointwise triangle
inequality. \(\square\)

\textbf{Theorem 3.2 (Completeness and simple approximation).} Every
\(L^p(\mu)\), \(1\leq p\leq\infty\), is complete. For finite \(p\),
simple functions with support of finite measure are dense, and every
\(f\in L^p\) vanishes outside a sigma-finite measurable set.

\emph{Proof.} Let \(f_n\) be Cauchy in \(L^p\), \(p<\infty\). Choose a
subsequence \(f_{n_j}\) with \(\|f_{n_{j+1}}-f_{n_j}\|_p\leq2^{-j}\),
and measurable representatives. The partial sums of
\[v=\sum_{j\geq1}|f_{n_{j+1}}-f_{n_j}|\] have \(L^p\) norm at most one
by Minkowski. Monotone convergence applied to their \(p\)-th powers
gives \(\int v^p\leq1\). Thus \(v<\infty\) outside a measurable null
set, and the subsequence converges there to a measurable \(f\), defined
as zero on that null set. Its tail difference satisfies
\[|f-f_{n_j}|\leq\sum_{k\geq j}|f_{n_{k+1}}-f_{n_k}|,
\qquad \|f-f_{n_j}\|_p\leq\sum_{k\geq j}2^{-k}.\] The latter follows by
Minkowski for finite sums and monotone convergence. Hence \(f\in L^p\)
and the subsequence converges in norm; the original Cauchy sequence does
too.

For \(p=\infty\), choose a subsequence with the same summable bounds.
Remove a measurable null set on which its first term fails its essential
bound, as well as the countably many null sets on which successive
differences fail their bounds. On the complement the first term is
bounded and the series of differences converges uniformly. Define the
limit as zero on the removed measurable set. It is bounded and
measurable, and the same tail estimate proves norm convergence.

For finite \(p\), the sets \(E_n=\{1/n<|f|\leq n\}\) increase to
\(\{f\ne0\}\), and \(\mu(E_n)\leq n^p\|f\|_p^p<\infty\). Dominated
convergence shows \(f1_{E_n}\to f\) in \(L^p\). Approximate the bounded
complex function \(f1_{E_n}\) uniformly by simple functions supported on
\(E_n\), with error at most \(\varepsilon/(1+\mu(E_n))^{1/p}\). This
proves density and the support assertion. \(\square\)

In particular \(L^2\) is a Hilbert space for
\(\langle f,g\rangle=\int f\overline g\). Hölder proves that the inner
product exists, and the theorem proves completeness. This is the point
at which the earlier Hilbert-space results may be applied to measure
spaces.

\subsection{4. Densities and bounded
functionals}\label{measure-and-hilbert-space-tools--4-densities-and-bounded-functionals}

Absolute continuity \(\nu\ll\mu\) means that every \(\mu\)-null
measurable set is \(\nu\)-null. We use the classical Hilbert-space
method of J. von Neumann, also presented in S. Axler's
\href{https://measure.axler.net/MIRA.pdf}{\emph{Measure, Integration \&
Real Analysis}}. The following proof first produces a bounded density
relative to \(\mu+\nu\); division then gives the density relative to
\(\mu\). It avoids assuming a Radon--Nikodym theorem inside a
representation theorem.

\textbf{Theorem 4.1 (Radon--Nikodym).} If \(\mu,\nu\) are sigma-finite
measures on the same sigma-algebra and \(\nu\ll\mu\), there is a
measurable \(h\geq0\), unique \(\mu\)-almost everywhere, such that
\(\nu(E)=\int_Eh\,d\mu\) for every measurable \(E\).

\emph{Proof.} First assume both measures finite and put
\(\lambda=\mu+\nu\). The functional \(u\mapsto\int u\,d\nu\) on
\(L^2(\lambda)\) is bounded: Cauchy--Schwarz gives a bound
\(\nu(X)^{1/2}\|u\|_{L^2(\lambda)}\). The Riesz--Fréchet theorem in the
earlier Hilbert lesson, Theorem 2.3, supplies \(g\in L^2(\lambda)\) with
\(\int u\,d\nu=\int ug\,d\lambda\); conjugate the representing vector
because our inner product is linear in its first variable.

Testing measurable indicators shows that \(g\) is real with
\(0\leq g\leq1\) almost everywhere. Indeed \(\int_Eg\) is real and
between zero and \(\lambda(E)\) for every \(E\); the level sets where
the imaginary part has one sign, or where \(g<0\) or \(g>1\), would
violate these inequalities. Redefine \(g\) on the exceptional measurable
null set. Then \(\nu=g\lambda\) and \(\mu=(1-g)\lambda\) on all sets. On
\(A=\{g=1\}\), \(\mu(A)=0\), so \(\nu(A)=0\), hence \(\lambda(A)=0\).
Set \(h=g/(1-g)\) outside \(A\), and \(h=0\) on \(A\). The indicator
identities extend to nonnegative functions by simple approximation and
monotone convergence. Consequently
\[\int_Eh\,d\mu=\int_Eh(1-g)\,d\lambda=\int_Eg\,d\lambda=\nu(E).\]

For sigma-finite measures choose a disjoint countable measurable
partition into pieces finite for both measures: intersect their two
finite-measure covers and disjointify. Apply the finite argument on each
piece and join the densities. Countable additivity and monotone
convergence give the formula on \(X\). For uniqueness, a candidate
density is finite almost everywhere on each piece where the target
measure is finite: its infinite level set must be null for the reference
measure. On a finite-measure piece where both candidate densities are
bounded, integrate over \(\{h\geq k+1/n\}\). Equality of the integrals
forces its measure to vanish. Countably many such pieces cover the set
where \(h>k\), up to null sets; interchange them for the reverse
inequality. \(\square\)

\textbf{Theorem 4.2 (The dual of \(L^1\)).} On a sigma-finite measure
space, every bounded complex linear functional \(\Phi\) on \(L^1(\mu)\)
has the form \(\Phi(u)=\int bu\,d\mu\) for a unique
\(b\in L^\infty(\mu)\), with \(\|\Phi\|=\|b\|_\infty\).

\emph{Proof.} If \(\mu(X)<\infty\), restrict \(\Phi\) to \(L^2\), using
\(\|u\|_1\leq\mu(X)^{1/2}\|u\|_2\). Hilbert representation gives a
measurable \(b\in L^2\) with \(\Phi(u)=\int bu\) there. For
\(c>\|\Phi\|\), take \(E=\{|b|>c\}\) and \(u=1_E\overline b/|b|\), with
zero value when \(b=0\). This bounded function belongs to \(L^2\), and
\[c\mu(E)\leq\int_E|b|=|\Phi(u)|\leq\|\Phi\|\mu(E).\] Thus \(\mu(E)=0\),
and \(\|b\|_\infty\leq\|\Phi\|\). Simple functions are in \(L^2\) and
dense in \(L^1\), by Theorem 3.2, so the identity extends to \(L^1\).

In the sigma-finite case partition \(X\) into finite-measure pieces
\(E_n\), apply the argument to \(u\mapsto\Phi(u1_{E_n})\), and join the
bounded densities. They have the uniform bound \(\|\Phi\|\). The series
\(u=\sum_nu1_{E_n}\) converges in \(L^1\), so the joined density
represents \(\Phi\). Its uniqueness follows by testing the phase of the
difference on finite-measure pieces. Hölder gives
\(\|\Phi\|\leq\|b\|_\infty\), proving norm equality. \(\square\)

The sigma-finite hypothesis in this last theorem is explicit. The Haar
lesson later decomposes a general group into open sigma-compact cosets
and applies the theorem on each coset; it will not assume the hypothesis
globally.

\subsubsection{Why sigma-finiteness matters for a
density}\label{measure-and-hilbert-space-tools--radon-nikodym-countability}

\textbf{Example 4.3.} Let \(X=[0,1]\) with its Borel sigma-algebra, let
\(\mu\) be counting measure, and let \(\nu\) be Lebesgue measure. Then
\(\nu\ll\mu\), but no nonnegative measurable function \(h\) satisfies
\(\nu(E)=\int_Eh\,d\mu\) for every Borel \(E\).

\emph{Proof.} Counting measure has no nonempty null set, so absolute
continuity is automatic. If a density existed, then for every
\(x\in[0,1]\), \[0=\nu(\{x\})=\int_{\{x\}}h\,d\mu=h(x).\] Thus \(h\)
would vanish everywhere, contradicting \(\nu([0,1])=1\). A
finite-counting-measure set is finite, and a countable union of finite
sets is countable, whereas \([0,1]\) is uncountable. Consequently
\(\mu\) is not sigma-finite. This example pinpoints the failure of the
reference-measure hypothesis in Theorem 4.1, even though the target
measure is finite. \(\square\)

\subsection{5. Tensor products as Hilbert--Schmidt
operators}\label{measure-and-hilbert-space-tools--5-tensor-products-as-hilbert-schmidt-operators}

For Hilbert spaces \(H,K\), a tensor is built in
\hyperref[hilbert-spaces-and-compact-operators]{Hilbert spaces and
compact operators}, Theorem 8.1: choose orthonormal bases indexed by
\(I,J\), take \(\ell^2(I\times J)\), and send \((\xi,\eta)\) to the
array \(\langle\xi,e_i\rangle\langle\eta,f_j\rangle\). The theorem
proves basis independence through its isometric universal property. No
separability is assumed.

There is a second useful interpretation. Write \(\overline H\) for the
conjugate Hilbert space: \(\overline\xi\) denotes a vector with scalar
rule \(a\overline\xi=\overline{\overline a\xi}\) and inner product
\(\langle\overline\xi,\overline\eta\rangle=\langle\eta,\xi\rangle\). A
bounded operator \(T:H\to K\) is Hilbert--Schmidt if
\[\|T\|_{\mathrm{HS}}^2=\sum_{i\in I}\|Te_i\|^2<\infty,\] where an
arbitrary nonnegative sum means the supremum of its finite subsums.

\textbf{Theorem 5.1.} This definition and norm do not depend on the
orthonormal basis. The Hilbert--Schmidt operators form a Hilbert space
canonically isometric to \(K\otimes\overline H\), through
\[k\otimes\overline h\longmapsto[\xi\mapsto\langle\xi,h\rangle k].
\tag{5.1}\] Every Hilbert--Schmidt operator is compact, and
\(\|T\|\leq\|T\|_{\mathrm{HS}}\).

\emph{Proof.} Fix bases \((e_i),(f_j)\). Parseval gives
\[\sum_i\|Te_i\|^2=\sum_{i,j}|\langle Te_i,f_j\rangle|^2
=\sum_j\|T^*f_j\|^2.
\tag{5.2}\] The interchange is valid for arbitrary sums of nonnegative
numbers: both iterated sums equal the supremum over finite rectangles,
and each finite subset of \(I\times J\) lies in a finite rectangle.
Fixing the \(f_j\) in (5.2) proves independence of \((e_i)\); the left
expression then proves independence of \((f_j)\).

A square-summable array \((a_{ji})\) determines an operator. On finite
linear combinations set \(Te_i=\sum_ja_{ji}f_j\). For finite-support
\((c_i)\), Cauchy--Schwarz gives \[\left\|\sum_i c_iTe_i\right\|^2
\leq\left(\sum_i|c_i|^2\right)\left(\sum_{i,j}|a_{ji}|^2\right).\] It
extends uniquely to a bounded operator on \(H\), with matrix
\((a_{ji})\) and Hilbert--Schmidt norm equal to its \(\ell^2\) norm.
Conversely every Hilbert--Schmidt operator has this matrix, so its space
is complete and has the inner product of these arrays.

For the rank-one operators in (5.1), summing their matrix inner products
gives \(\langle k,k'\rangle\langle h',h\rangle\). This is exactly the
elementary-tensor inner product on \(K\otimes\overline H\).
Finite-support matrices are finite sums of these rank-one operators and
are dense among the square-summable arrays. The universal property of
the tensor product therefore extends (5.1) to a unitary. Finally
finite-support matrices approximate every such array in Hilbert--Schmidt
norm, and hence in operator norm by the displayed estimate. They give
finite-rank operators. A norm limit of finite-rank operators is compact:
approximate its image of the unit ball by a finite net for the image
under one finite-rank approximant. \(\square\)

\textbf{Example.} A diagonal operator on \(\ell^2(I)\) is
Hilbert--Schmidt exactly when its diagonal \((a_i)\) is square summable.
For uncountable \(I\), at most countably many \(a_i\) can be nonzero.
The identity is Hilbert--Schmidt exactly when \(I\) is finite. This
explains why an uncountable basis causes no uncountable convergence
problem in (5.2).

\subsection{6. One-variable integral
identities}\label{measure-and-hilbert-space-tools--6-one-variable-integral-identities}

The explicit Haar and arithmetic examples use elementary calculus. Here
are the integral identities needed, with their proofs, so no general
multidimensional change-of-variables theorem is implicit.

\textbf{Lemma 6.1.} If \(h\) is continuous on a compact interval, its
Lebesgue integral equals its Riemann integral. For \(H\in C^1([a,b])\),
\(\int_a^b H'(t)\,dt=H(b)-H(a)\). If \(\phi\in C^1([a,b])\) and \(h\) is
continuous on an interval containing its image, then
\[\int_{\phi(a)}^{\phi(b)}h(u)\,du
=\int_a^b h(\phi(t))\phi'(t)\,dt.
\tag{6.1}\] In particular, the integral of a compactly supported
continuous derivative on the real line is zero. Differentiating a
parameter integral is valid when its difference quotients converge
pointwise and have a common integrable bound.

\emph{Proof.} Recall the elementary differential facts involved. A
continuous function on a compact interval attains its extrema. If its
endpoints agree, either it is constant or one extremum occurs in the
interior, where its derivative is zero: the left and right difference
quotients have opposite weak signs. This is Rolle\textquotesingle s
theorem. Subtract the line joining the endpoint values to obtain the
mean value theorem. The chain rule follows by substituting
\(\phi(t+h)-\phi(t)=\phi'(t)h+o(h)\) into the corresponding first-order
expansion of the outer function; differentiability makes the inner
increment \(O(h)\).

Uniform continuity makes the difference between the upper and lower step
functions on a sufficiently fine partition arbitrarily small in
integral; both bound \(h\). Their Lebesgue integrals are the same
interval-length sums that define the Riemann integral. For \(H\), the
mean value theorem on each partition interval expresses its increment as
\(H'(\xi)\) times the interval length. Summing telescopes to
\(H(b)-H(a)\); uniform continuity of \(H'\) makes the sums converge to
its integral.

Define \(A(v)=\int_{v_0}^v h(u)\,du\), with oriented integrals. The
difference quotient of \(A\) is an interval average of \(h\), tending to
\(h(v)\) by continuity. Thus \((A\circ\phi)'=(h\circ\phi)\phi'\). Apply
the preceding identity to \(A\circ\phi\) to obtain (6.1). The
zero-integral assertion follows by choosing endpoints outside the
support. Finally apply dominated convergence, Theorem 2.2, to the
difference quotients to justify the parameter derivative. ∎

\subsection{7. Exercises with complete
solutions}\label{measure-and-hilbert-space-tools--7-exercises-with-complete-solutions}

\textbf{1 (easy): finite-measure support.} Prove that \(f\in L^p\),
\(p<\infty\), has sigma-finite support without assuming \(X\)
sigma-finite. Explain why the same assertion fails for \(p=\infty\).

\emph{Solution.} The support in the measure-theoretic sense
\(\{f\ne0\}\) is \(\bigcup_n\{|f|>1/n\}\). On each set,
\(n^{-p}\mu(\{|f|>1/n\})\leq\int|f|^p\), so its measure is finite. On an
uncountable counting space, \(1\in L^\infty\) has all of \(X\) as
support, which cannot be covered by countably many finite sets.

\textbf{2 (medium): a summable series in \(L^p\).} If
\(\sum_n\|u_n\|_p<\infty\), \(p<\infty\), prove that \(\sum_nu_n\)
converges almost everywhere and in \(L^p\), and estimate every tail.

\emph{Solution.} Minkowski bounds the norm of \(\sum_{n=1}^N|u_n|\) by
\(\sum_n\|u_n\|_p\). Monotone convergence of its \(p\)-th power shows
that \(\sum_n|u_n|<\infty\) almost everywhere. Thus the series converges
there. Apply the same argument to a tail, obtaining
\(\|\sum_{n\geq N}u_n\|_p\leq\sum_{n\geq N}\|u_n\|_p\to0\). Define the
sum as zero on the measurable exceptional set.

\textbf{3 (medium): failure of domination.} On \((0,1)\) with Lebesgue
measure put \(u_n=n1_{(0,1/n)}\). Find its pointwise limit and its
integral. Prove that no integrable common dominator exists.

\emph{Solution.} Every \(x>0\) eventually lies outside \((0,1/n)\), so
\(u_n(x)\to0\). Yet \(\int u_n=1\). If \(|u_n|\leq g\) for an integrable
\(g\), dominated convergence would force these integrals to tend to
zero, a contradiction. The shrinking support does not replace
domination.

\textbf{4 (hard): weighted finite reduction.} For sigma-finite \(\mu\),
construct a strictly positive measurable \(w\) with
\(\int w\,d\mu<\infty\). Prove that \(w\mu\) has the same null sets, and
that multiplication by \(w^{-1/2}\) gives a unitary from \(L^2(\mu)\) to
\(L^2(w\mu)\).

\emph{Solution.} Partition \(X\) into finite-measure measurable \(E_n\)
and take \(w=\sum_n2^{-n}(1+\mu(E_n))^{-1}1_{E_n}\). Then \(w>0\) and
\(\int w\leq1\). For any measurable \(E\), \(\int_Ew=0\) implies
\(\mu(E\cap\{w\geq1/n\})=0\) for each \(n\), hence \(\mu(E)=0\). The
converse is immediate. The norm identity is
\(\int|w^{-1/2}f|^2w\,d\mu=\int|f|^2\,d\mu\), and the inverse is
multiplication by \(w^{1/2}\).

\textbf{5 (hard): which tensor represents an operator?} Explain why
\(K\otimes H\) is not the canonical tensor for complex-linear rank-one
maps \(H\to K\), and verify (5.1) on two rank-one maps.

\emph{Solution.} The map \(h\mapsto[\xi\mapsto\langle\xi,h\rangle k]\)
is conjugate-linear in \(h\), while elementary tensors are bilinear in
their factors. Passing to \(\overline H\) makes it linear. For
\(T\xi=\langle\xi,h\rangle k\) and \(S\xi=\langle\xi,h'\rangle k'\),
Parseval gives
\(\sum_i\langle Te_i,Se_i\rangle=\langle k,k'\rangle\sum_i\langle e_i,h\rangle\overline{\langle e_i,h'\rangle}=\langle k,k'\rangle\langle h',h\rangle\),
the required tensor inner product.

\subsection{What this
enables}\label{measure-and-hilbert-space-tools--what-this-enables}

The next lesson uses Theorem 1.1 to construct Radon measures, Theorems
2.1--3.2 to pass to limits and complete \(L^p\), and Theorems 4.1--4.2
on the sigma-finite pieces of a general group. Its Radon-product theorem
proves the additional hypotheses required for iterated integration; none
are implicit in this lesson. Hilbert tensors then identify \(L^2\) of
that product, and Theorem 5.1 interprets square-integrable kernels as
Hilbert--Schmidt operators.

\subsection{References}\label{measure-and-hilbert-space-tools--references}

\begin{itemize}
\item
  {[}Erdman{]} J. M. Erdman, \emph{Functional Analysis and Operator
  Algebras: An Introduction}, version 4 October 2015,
  \href{https://web.pdx.edu/~erdman/}{author\textquotesingle s text and
  editable sources}. The core Functional Analysis course uses this text.
  Its Hilbert--Schmidt exercises concern separable spaces; Theorem 5.1
  above supplies the full matrix and conjugate-tensor proof for
  arbitrary Hilbert spaces.
\item
  {[}Paschke{]} W. L. Paschke, \emph{Lecture Notes in Functional
  Analysis}, edition 0.9,
  \href{https://orthogonalpublishing.com/lnfa/html/hilbertschmidt.html}{author\textquotesingle s
  Hilbert--Schmidt chapter}. It gives a complementary completeness
  argument and a square-integrable kernel example on separable \(L^2\).
\item
  {[}Hilbert tools{]} \emph{Hilbert spaces and compact operators},
  programme lesson by Claude Opus 5.5 (Anthropic), October 2026, CC0,
  Theorems 2.1--2.3, 4.1 and 8.1. These are preceding internal
  full-proof providers.
\item
  {[}Fremlin{]} D. H. Fremlin, \emph{Measure Theory}, Volumes 1 and 2,
  Torres Fremlin, 2000--2001.
  \href{https://www1.essex.ac.uk/maths/people/fremlin/mt.htm}{Freely
  accessible text and supplied editable TeX}. Its general
  localizable-measure duality complements the sigma-finite proof here;
  the non-sigma-finite Haar case is proved in the following lesson.
\item
  {[}Lebesgue{]} H. Lebesgue, \emph{Leçons sur
  l\textquotesingle intégration et la recherche des fonctions
  primitives}, Gauthier-Villars, 1904,
  \href{https://fr.wikisource.org/wiki/Le\%C3\%A7ons_sur_l\%E2\%80\%99int\%C3\%A9gration_et_la_recherche_des_fonctions_primitives_(premi\%C3\%A8re_\%C3\%A9dition)}{freely
  accessible historical edition}. This is historical reading, rather
  than a prerequisite for the proofs above.
\item
  {[}Lebl{]} J. Lebl, \emph{Basic Analysis I} and \emph{Basic Analysis
  II}, \href{https://www.jirka.org/ra/}{author's open reader}. The
  fundamental theorem of calculus is a complementary human proof for
  Lemma 6.1.
\item
  {[}Axler{]} S. Axler,
  \href{https://measure.axler.net/MIRA.pdf}{\emph{Measure, Integration
  \& Real Analysis}}, author's freely accessible edition dated 12 June
  2026. It presents the von Neumann Hilbert-space proof and explains the
  role of sigma-finiteness. The proof and counterexample above are
  independently written.
\end{itemize}

\section{Haar measure on locally compact
groups}\label{haar-measure-on-locally-compact-groups}

\emph{Originally written by Claude Opus 5.5 (Anthropic), September 2026,
with a separate historical AI spot-check of that text. Repaired and
self-checked by GPT-6.1 Sol (OpenAI), Ultra reasoning effort, October
2026. The current edition has no separately recorded independent AI or
human review. Original programme content is dedicated under CC0.}

This lesson constructs Haar measure on an arbitrary locally compact
group and develops the measure theory that goes with it. Sections 2--6
work on locally compact spaces. They prove the Riesz representation
theorem for Radon measures, treat image measures and densities, build
the product of two Radon measures with its Tonelli and Fubini theorems,
and identify \(L^2\) of a product with a Hilbert tensor product.
Sections 7--11 turn to groups: topological groups, existence and
uniqueness of Haar measure, the modular function, and the inversion
formula. Sections 12--15 treat products of groups, groups that are not
\(\sigma\)-compact, convolution, and approximate identities. Section 16
has exercises with solutions.

No countability assumption is made anywhere. The group need not be
\(\sigma\)-compact, second countable or unimodular, and its Haar measure
need not be \(\sigma\)-finite. This generality needs care: a Borel set
of infinite measure can have only null compact subsets, and
Fubini\textquotesingle s theorem can fail. Section 13 shows what goes
wrong, and the rest of the lesson is arranged so that it does no harm.

Haar measure makes a locally compact group \(G\) into a measure space on
which \(G\) acts by translations. Then \(L^1(G)\) is a Banach
\(*\)-algebra under convolution, and \(L^2(G)\) carries the left and
right regular representations. These objects are the starting point of
abstract harmonic analysis, and of the group von Neumann algebras
studied in the course on modular theory and weights.

The lesson builds on four sources, cited where they are used:

\begin{itemize}
\tightlist
\item
  \hyperref[measure-and-hilbert-space-tools]{Measure and Hilbert space
  tools for Haar integration}, Theorems 1.1--4.2, for the full
  measure-theory proofs; the measure convention is compared with Fremlin
  below;
\item
  \hyperref[the-stone-weierstrass-theorem-for-functions-vanishing-at-infinity]{The
  Stone--Weierstrass theorem for functions vanishing at infinity}, for
  locally compact spaces and Urysohn\textquotesingle s lemma;
\item
  \hyperref[weak-topologies-tychonoff-banach-alaoglu-mazur-bipolars-krein-milman-and-eberlein-smulian]{Weak
  topologies: Tychonoff, Banach--Alaoglu, Mazur, bipolars, Krein--Milman
  and Eberlein--Šmulian}, for Tychonoff\textquotesingle s theorem;
\item
  \hyperref[hilbert-spaces-and-compact-operators]{Hilbert spaces and
  compact operators}, for Hilbert tensor products and multiplication
  operators.
\end{itemize}

The exact statements are collected at the end, under "Results used from
other lessons". Core preparation is
\href{https://kokunoyumeto.github.io/program-matematika-indonesia/en/\#course-C90}{Point-Set
Topology};
\href{https://kokunoyumeto.github.io/program-matematika-indonesia/en/\#course-D10}{Measure
and Integration} provides Fremlin's broader measure-theory setting. The
arbitrary-product and Hilbert-space proofs used here are the advanced
programme results named above.

Freely accessible complementary treatments are Fremlin's
\href{https://www1.essex.ac.uk/maths/people/fremlin/mtcont.htm}{Measure
Theory} and F. van Doorn's
\href{https://arxiv.org/abs/2102.07636v1}{Formalized Haar Measure}.
Fischer and Ruzhansky's freely accessible
\href{https://biblio.ugent.be/publication/8585474}{\emph{Quantization on
Nilpotent Lie Groups}} supplies further convolution and Heisenberg-group
examples. The historical development from Haar and von Neumann to Weil
and Cartan is discussed in F. van Doorn's paper; the original accessible
von Neumann and Cartan papers are linked below.

\subsection{1.
Conventions}\label{haar-measure-on-locally-compact-groups--1-conventions}

\emph{Spaces.} An \emph{LCH space} is a Hausdorff space that is locally
compact. Its compact subsets are closed. \(C_c(X)\) is the space of
continuous complex functions with compact support, \(C_c^+(X)\) is the
set of \(f\in C_c(X)\) with \(f\geq0\) and \(f\neq0\), and
\(\|f\|_{\sup}=\sup_x|f(x)|\). For an open set \(U\) we write
\(f\prec U\) when \(f\in C_c(X)\), \(0\leq f\leq1\) and
\(\operatorname{supp}f\subseteq U\). For a compact set \(K\) we write
\(K\prec f\) when \(f\in C_c(X)\), \(0\leq f\leq 1\) and \(f=1\) on
\(K\). A function \(h\colon X\to[0,\infty]\) is \emph{lower
semicontinuous} (lsc) if every set \(\{h>c\}\) is open; lsc functions
are Borel.

\emph{Four tools from topology.} The first two tools are proved in the
lesson on the Stone--Weierstrass theorem; we derive the other two from
them.

\begin{itemize}
\tightlist
\item
  \textbf{(T1)} (\emph{Shrinking}) If \(K\subseteq U\) with \(K\)
  compact and \(U\) open, there is an open \(V\) with compact closure
  such that \(K\subseteq V\subseteq\overline V\subseteq U\). This is
  \hyperlink{the-stone-weierstrass-theorem-for-functions-vanishing-at-infinity--oa-fnd-sw-15}{Proposition
  4.2(2) of the Stone--Weierstrass lesson}.
\item
  \textbf{(T2)} (\emph{Urysohn\textquotesingle s lemma}) If
  \(K\subseteq U\) as in (T1), there is \(f\) with \(K\prec f\prec U\).
  \emph{Proof.} Take \(V\) from (T1).
  \hyperlink{the-stone-weierstrass-theorem-for-functions-vanishing-at-infinity--oa-fnd-sw-16}{Corollary
  5.2 of the Stone--Weierstrass lesson}, applied to the compact set
  \(K\) and the closed set \(X\setminus V\), gives \(f\in C_c(X)\) with
  \(0\leq f\leq1\), \(f=1\) on \(K\) and \(f=0\) off \(V\). Then
  \(\operatorname{supp}f\subseteq\overline V\subseteq U\). \(\square\)
\item
  \textbf{(T3)} (\emph{Partitions of unity}) If \(K\) is compact and
  \(K\subseteq U_1\cup\dots\cup U_n\) with each \(U_i\) open, there are
  \(h_i\prec U_i\) with \(\sum_ih_i=1\) on \(K\) and \(\sum_ih_i\leq1\)
  everywhere. \emph{Proof.} For each \(x\in K\) pick an index \(i\) with
  \(x\in U_i\), and by (T1) an open \(V_x\ni x\) whose compact closure
  lies in \(U_i\). Finitely many \(V_x\) cover \(K\). Let \(F_i\) be the
  union of the closures of the chosen \(V_x\) that were assigned to
  \(i\). Then \(F_i\) is compact, \(F_i\subseteq U_i\), and
  \(K\subseteq\bigcup_iF_i\). Choose \(F_i\prec g_i\prec U_i\) by (T2)
  and set \(h_1=g_1\) and \(h_k=(1-g_1)\cdots(1-g_{k-1})g_k\). By
  induction, \(\sum_{k\leq m}h_k=1-\prod_{k\leq m}(1-g_k)\). This lies
  in \([0,1]\), and it equals \(1\) on each \(F_i\). \(\square\)
\item
  \textbf{(T4)} (\emph{Tube lemma}) If \(K\subseteq X\) and
  \(L\subseteq Y\) are compact and \(K\times L\subseteq W\) with \(W\)
  open in \(X\times Y\), there are open \(U\supseteq K\) and
  \(V\supseteq L\) with \(U\times V\subseteq W\). \emph{Proof.} Fix
  \(x\in K\). Each \((x,y)\), \(y\in L\), lies in an open box inside
  \(W\); finitely many boxes \(P_j\times Q_j\) of this kind have \(Q_j\)
  covering \(L\). Put \(P^x=\bigcap_jP_j\) and \(Q^x=\bigcup_jQ_j\);
  then \(P^x\times Q^x\subseteq W\). Cover \(K\) by finitely many
  \(P^{x_i}\), and let \(U=\bigcup_iP^{x_i}\) and
  \(V=\bigcap_iQ^{x_i}\). \(\square\)
\end{itemize}

\emph{Measures and functions.} A measure is positive and countably
additive. \(\mathcal B(X)\) denotes the Borel sets of \(X\), the
\(\sigma\)-algebra generated by the open sets; functions on \(X\) are
Borel unless said otherwise. A Borel set is \emph{\(\sigma\)-finite} for
a measure when countably many Borel sets of finite measure cover it. For
\(1\leq p<\infty\), \(L^p(\mu)\) consists of the Borel functions with
\(\int|f|^p\,d\mu<\infty\), modulo functions that vanish almost
everywhere. Hilbert spaces are complex, and the inner product
\(\langle\xi,\eta\rangle=\int\xi\bar\eta\,d\mu\) is linear in \(\xi\).
For a Borel \(f\colon X\to[0,\infty]\) we use the \emph{layer functions}
\[s_n=2^{-n}\sum_{i=1}^{n2^n}1_{\{f>i2^{-n}\}} .
\tag{1.1}\] They are simple Borel functions, and they vanish where \(f\)
does. They increase, \(s_n\leq s_{n+1}\), because the \(i\)-th term of
\(s_n\) is at most the sum of the terms of \(s_{n+1}\) with indices
\(2i-1\) and \(2i\). And \(s_n(x)\to f(x)\) at every point: if
\(f(x)<n\), then \(f(x)-2^{-n}\leq s_n(x)\leq f(x)\), and if
\(f(x)=\infty\), then \(s_n(x)=n\). The monotone and dominated
convergence theorems, Hölder\textquotesingle s inequality, and the
completeness of \(L^p\) are proved in
\hyperref[measure-and-hilbert-space-tools]{Measure and Hilbert space
tools for Haar integration}, Theorems 2.1, 2.2, 3.1 and 3.2. The exact
uses are listed under "Results used from other lessons".

\emph{Groups.} A group \(G\) is a \emph{locally compact group} when it
carries a locally compact Hausdorff topology for which
\((x,y)\mapsto xy\) and \(x\mapsto x^{-1}\) are continuous; \(e\) is its
identity. For a function \(f\) on \(G\) and \(y\in G\) put
\[L_yf(x)=f(y^{-1}x),\qquad R_yf(x)=f(xy),\qquad \check f(x)=f(x^{-1}).\]
Then \(L_{yz}=L_yL_z\) and \(R_{yz}=R_yR_z\). The \emph{left and right
regular representations} of \(G\) on \(L^2(G)\) are \(\lambda(g)=L_g\)
and \(\rho(g)=\Delta(g)^{1/2}R_g\), where \(\Delta\) is the modular
function of Section 10; Exercise 16.3 shows that they are unitary
representations.

\subsection[2. Radon measures and the Riesz representation
theorem]{\texorpdfstring{\protect\hypertarget{haar-measure-on-locally-compact-groups--oa-fnd-hm-01}{}{}2.
Radon measures and the Riesz representation
theorem}{2. Radon measures and the Riesz representation theorem}}\label{haar-measure-on-locally-compact-groups--OA-FND-HM-01}

\textbf{Definition 2.1} (Radon measure). A \emph{Radon measure} on an
LCH space \(X\) is a measure \(\mu\) on \(\mathcal B(X)\) with the
following three properties.

\begin{itemize}
\tightlist
\item
  \textbf{(R1)} \(\mu(K)<\infty\) for every compact \(K\).
\item
  \textbf{(R2)} (\emph{outer regularity})
  \(\mu(E)=\inf\{\mu(U):U\supseteq E\ \text{open}\}\) for every Borel
  set \(E\).
\item
  \textbf{(R3)} (\emph{inner regularity on open sets})
  \(\mu(U)=\sup\{\mu(K):K\subseteq U\ \text{compact}\}\) for every open
  \(U\).
\end{itemize}

This is the notion of regularity used throughout the lesson. Inner
regularity is not required on all Borel sets, because it can fail
(Example 13.4). Proposition 2.3 shows that it holds on every
\(\sigma\)-finite set.

\textbf{Theorem 2.2} (Riesz representation theorem). Take a positive
linear functional \(I\) on \(C_c(X)\). There is exactly one Radon
measure \(\mu\) on \(X\) with \(I(f)=\int f\,d\mu\) for every
\(f\in C_c(X)\). Moreover, every Radon measure \(\mu\) satisfies, with
\(I(f)=\int f\,d\mu\),
\[\mu(U)=\sup\{I(f):f\prec U\}\quad(U\ \text{open}),\qquad
\mu(K)=\inf\{I(f):f\in C_c(X),\ 1_K\leq f\}\quad(K\ \text{compact}).
\tag{2.1}\]

\emph{Regularity warning:} The stronger requirement that a measure be
inner and outer regular on every Borel set is inappropriate here. In
that form it fails for the Haar integral of
\(\mathbb R\times\mathbb R_d\) (Example 13.4).

\textbf{Proof.} \emph{Step 1: formula (2.1) and uniqueness.} Let \(\mu\)
be Radon. If \(f\prec U\), then \(\int f\,d\mu\leq\mu(U)\). If
\(K\subseteq U\) is compact, Urysohn\textquotesingle s lemma (T2) gives
\(K\prec f\prec U\), and then \(\mu(K)\leq\int f\,d\mu\). With (R3) this
proves the first formula. If \(1_K\leq f\), then
\(\mu(K)\leq\int f\,d\mu\); if \(U\supseteq K\) is open, (T2) gives
\(K\prec f\prec U\) with \(\int f\,d\mu\leq\mu(U)\). With (R2) this
proves the second formula. By the first formula, two Radon measures that
give every \(f\in C_c(X)\) the same integral agree on open sets, and
then on all Borel sets by (R2).

\emph{Step 2: an outer measure.} For open \(U\) define
\(\mu(U)=\sup\{I(f):f\prec U\}\), and for any \(A\subseteq X\) put
\(\mu^*(A)=\inf\{\mu(U):U\supseteq A\ \text{open}\}\). Then
\(\mu^*=\mu\) on open sets, and \(\mu^*\) is monotone. Let
\(U=\bigcup_jU_j\) with \(U_j\) open, and let \(f\prec U\). The compact
set \(\operatorname{supp}f\) lies in \(U_1\cup\dots\cup U_n\) for some
\(n\). The partition of unity (T3) gives \(h_j\prec U_j\) with
\(\sum_{j\leq n}h_j=1\) on \(\operatorname{supp}f\). So
\(f=\sum_{j\leq n}fh_j\) with \(fh_j\prec U_j\), and
\(I(f)\leq\sum_j\mu(U_j)\). Hence \(\mu(U)\leq\sum_j\mu(U_j)\). For
arbitrary sets \(A_j\), choose open \(U_j\supseteq A_j\) with
\(\mu(U_j)\leq\mu^*(A_j)+\varepsilon2^{-j}\); this gives
\(\mu^*(\bigcup_jA_j)\leq\sum_j\mu^*(A_j)+\varepsilon\). So \(\mu^*\) is
an outer measure.

\emph{Step 3: open sets are measurable.} Let \(U\) be open. By
Carathéodory\textquotesingle s theorem
\hyperref[measure-and-hilbert-space-tools--1-from-an-outer-measure-to-a-measure]{the
preceding tools lesson, Theorem 1.1} it is enough to show
\(\mu^*(A)\geq\mu^*(A\cap U)+\mu^*(A\setminus U)\) whenever
\(\mu^*(A)<\infty\). First let \(A=V\) be open. Given \(\varepsilon>0\),
choose \(f\prec V\cap U\) with \(I(f)>\mu(V\cap U)-\varepsilon\). The
set \(V\setminus\operatorname{supp}f\) is open; choose
\(g\prec V\setminus\operatorname{supp}f\) with
\(I(g)>\mu(V\setminus\operatorname{supp}f)-\varepsilon\). The supports
of \(f\) and \(g\) are disjoint, so \(f+g\prec V\) and
\[\mu(V)\geq I(f+g)>\mu(V\cap U)+\mu(V\setminus\operatorname{supp}f)-2\varepsilon\geq\mu^*(V\cap U)+\mu^*(V\setminus U)-2\varepsilon ,\]
since \(V\setminus U\subseteq V\setminus\operatorname{supp}f\). For
general \(A\), take an open \(V\supseteq A\) with
\(\mu(V)\leq\mu^*(A)+\varepsilon\), apply the open case, and use
monotonicity. By Carathéodory\textquotesingle s theorem the
\(\mu^*\)-measurable sets form a \(\sigma\)-algebra on which \(\mu^*\)
is countably additive. It contains the open sets, hence
\(\mathcal B(X)\). Let \(\mu\) be the restriction of \(\mu^*\) to
\(\mathcal B(X)\). It satisfies (R2) by construction.

\emph{Step 4: compact sets.} Let \(K\) be compact and
\(1_K\leq f\in C_c(X)\); in particular \(f\geq0\). For \(0<c<1\) the
open set \(U_c=\{f>c\}\) contains \(K\), and every \(g\prec U_c\)
satisfies \(g\leq f/c\). So \(\mu(K)\leq\mu(U_c)\leq I(f)/c\), and
letting \(c\to1\) gives \(\mu(K)\leq I(f)\). Some \(f\) with
\(K\prec f\) exists by (T2) with \(U=X\), so \(\mu(K)<\infty\), which is
(R1). Given \(\varepsilon>0\), choose an open \(U\supseteq K\) with
\(\mu(U)\leq\mu(K)+\varepsilon\) and then \(f\) with
\(K\prec f\prec U\): \(I(f)\leq\mu(U)\leq\mu(K)+\varepsilon\). This
proves the second formula of (2.1).

\emph{Step 5: (R3).} If \(f\prec U\), then \(f\prec V\) for every open
\(V\supseteq\operatorname{supp}f\), so \(I(f)\leq\mu(V)\), and
\(I(f)\leq\mu(\operatorname{supp}f)\) by (R2). Taking the supremum over
\(f\prec U\) gives
\(\mu(U)\leq\sup\{\mu(K):K\subseteq U\ \text{compact}\}\). The reverse
inequality is monotonicity.

\emph{Step 6: \(\mu\) represents \(I\).} By linearity it suffices to
treat real \(f\in C_c(X)\) with \(0\leq f\leq1\). Fix \(N\geq1\). Put
\(K_0=\operatorname{supp}f\), \(K_j=\{f\geq j/N\}\) for
\(1\leq j\leq N\) (compact sets), and
\(f_j=\min(\max(f-\frac{j-1}N,0),\frac1N)\). Then \(f=\sum_{j=1}^Nf_j\),
each \(f_j\) is in \(C_c(X)\) with
\(\operatorname{supp}f_j\subseteq K_{j-1}\), and
\(\frac1N1_{K_j}\leq f_j\leq\frac1N1_{K_{j-1}}\). Integrating,
\(\frac1N\mu(K_j)\leq\int f_j\,d\mu\leq\frac1N\mu(K_{j-1})\). On the
other side, \(\mu(K_j)\leq I(Nf_j)\) by Step 4. Also \(Nf_j\prec U\) for
every open \(U\supseteq K_{j-1}\), so \(I(Nf_j)\leq\mu(U)\), and
\(I(Nf_j)\leq\mu(K_{j-1})\) by (R2). Summing over \(j\), both \(I(f)\)
and \(\int f\,d\mu\) lie between \(\frac1N\sum_{j=1}^N\mu(K_j)\) and
\(\frac1N\sum_{j=0}^{N-1}\mu(K_j)\). These bounds differ by at most
\(\mu(\operatorname{supp}f)/N\). Let \(N\to\infty\). \(\square\)

\textbf{Proposition 2.3} (Regularity on \(\sigma\)-finite sets). Every
Radon measure \(\mu\) on \(X\) has the following properties.

\begin{enumerate}
\tightlist
\item
  If \(E\) is Borel and \(\mu(E)<\infty\), then
  \(\mu(E)=\sup\{\mu(K):K\subseteq E\ \text{compact}\}\).
\item
  The same holds for every Borel set that is \(\sigma\)-finite for
  \(\mu\).
\item
  Every \(\sigma\)-finite Borel set \(E\) is the union of a
  \(\sigma\)-compact set and a Borel null set.
\end{enumerate}

\textbf{Proof.} (1) Let \(\varepsilon>0\). By (R2) choose an open
\(U\supseteq E\) with \(\mu(U)<\mu(E)+\varepsilon\), so
\(\mu(U\setminus E)<\varepsilon\). Again by (R2) choose an open
\(W\supseteq U\setminus E\) with \(\mu(W)<\varepsilon\). By (R3) choose
a compact \(F\subseteq U\) with \(\mu(F)>\mu(U)-\varepsilon\). The
compact set \(K=F\setminus W\) lies in \(E\): a point of \(F\) outside
\(E\) is in \(U\setminus E\subseteq W\). And
\(\mu(K)\geq\mu(F)-\mu(W)>\mu(E)-2\varepsilon\). (2) Write \(E\) as an
increasing union of Borel sets \(E_n\) of finite measure; then
\(\mu(E)=\lim\mu(E_n)\), and (1) applies to each \(E_n\). (3) With
\(E_n\) as in (2), choose compact \(K_{n,m}\subseteq E_n\) with
\(\mu(E_n\setminus K_{n,m})<1/m\). The \(\sigma\)-compact set
\(S=\bigcup_{n,m}K_{n,m}\) lies in \(E\), and \(\mu(E_n\setminus S)=0\)
for every \(n\), so \(\mu(E\setminus S)=0\). \(\square\)

\emph{Where the hypotheses enter.} Local compactness and the Hausdorff
property are used only through (T1)--(T3): they supply enough functions
\(f\prec U\). No countability is used, and the measure \(\mu\) need not
be \(\sigma\)-finite. A measure that is inner and outer regular on every
Borel set would be more convenient, but Example 13.4 shows that on the
group \(\mathbb R\times\mathbb R_d\) no measure of that kind represents
Haar integration. So the lesson never uses inner regularity on arbitrary
Borel sets.

\subsubsection{\texorpdfstring{Bounded functionals on
\(C_0(X)\)}{Bounded functionals on C\_0(X)}}\label{haar-measure-on-locally-compact-groups--bounded-functionals-on-c_0-x}

The operator-algebra lessons also use the bounded, not necessarily
positive, linear functionals on \(C_0(X)\), the space of continuous
functions vanishing at infinity with the norm \(\|\cdot\|_{\sup}\). Call
a functional \(\varphi\) \emph{hermitian} if \(\varphi(f)\) is real for
every real \(f\). Two facts are used below.

\begin{itemize}
\tightlist
\item
  \emph{\(C_c(X)\) is dense in \(C_0(X)\).} For \(f\in C_0(X)\) and
  \(\varepsilon>0\), the set \(K=\{|f|\geq\varepsilon\}\) is compact.
  Choose \(K\prec u\) by (T2). Then \(uf\in C_c(X)\) and
  \(\|f-uf\|_{\sup}\leq\varepsilon\).
\item
  \emph{Norms of positive functionals.} A positive linear functional
  \(\psi\) on \(C_0(X)\) is hermitian, since every real \(f\) is
  \(f^+-f^-\). If \(\psi\) is bounded, then
  \(\|\psi\|=\sup\{\psi(f):0\leq f\leq1\}\), and by (2.1) this is
  \(\nu(X)\) for the Radon measure \(\nu\) that represents \(\psi\) on
  \(C_c(X)\). Indeed, for \(\|f\|_{\sup}\leq1\) choose \(\theta\) with
  \(e^{i\theta}\psi(f)\geq0\), and write \(e^{i\theta}f=u+iv\) with
  \(u,v\) real. Then \(|\psi(f)|=\psi(u)+i\psi(v)\) forces
  \(\psi(v)=0\), and \(|\psi(f)|=\psi(u)\leq\psi(u^+)\) with
  \(0\leq u^+\leq1\). The functions \(f\prec X\) suffice, by the density
  just proved.
\end{itemize}

\textbf{Theorem 2.4} (Bounded functionals on \(C_0(X)\)). Let \(X\) be
an LCH space and \(\varphi\) a bounded linear functional on \(C_0(X)\).

\begin{enumerate}
\tightlist
\item
  (\emph{Jordan decomposition.}) Let \(\varphi\) be hermitian. For
  \(f\geq0\) put
  \(\varphi_+(f)=\sup\{\varphi(g):g\in C_0(X),\ 0\leq g\leq f\}\). Then
  \(\varphi_+\) extends to a bounded positive linear functional,
  \(\varphi_-=\varphi_+-\varphi\) is positive, and
  \(\|\varphi\|=\|\varphi_+\|+\|\varphi_-\|\). Every bounded \(\varphi\)
  is \(\varphi_1+i\varphi_2\) with \(\varphi_1,\varphi_2\) bounded and
  hermitian.
\item
  (\emph{Representation.}) There are finite Radon measures
  \(\nu_0,\dots,\nu_3\) with
  \(\varphi(f)=\sum_{k=0}^3i^k\int f\,d\nu_k\) for every
  \(f\in C_0(X)\). The set function \(\mu=\sum_ki^k\nu_k\) on
  \(\mathcal B(X)\) does not depend on the choice of the \(\nu_k\). It
  is the \emph{complex Radon measure} of \(\varphi\), and we write
  \(\varphi(f)=\int f\,d\mu\). Conversely, every such \(\mu\) defines a
  bounded functional in this way.
\item
  (\emph{Norm.}) \(\|\varphi\|=|\mu|(X)\), where
  \(|\mu|(E)=\sup\sum_j|\mu(E_j)|\), the supremum over the finite
  partitions of \(E\) into Borel sets.
\item
  (\emph{Positivity.}) \(\varphi\) is positive if and only if
  \(\mu\geq0\); then \(\mu\) is the measure of Theorem 2.2. If
  \(\varphi\) is hermitian, then \(\mu=\nu_+-\nu_-\), where \(\nu_\pm\)
  are the Radon measures of \(\varphi_\pm\), and \(|\mu|=\nu_++\nu_-\).
  In particular \(|\mu|\) is a finite Radon measure.
\end{enumerate}

\textbf{Proof.} (1) \emph{Additivity.} For \(f_1,f_2\geq0\), clearly
\(\varphi_+(f_1+f_2)\geq\varphi_+(f_1)+\varphi_+(f_2)\). Conversely, let
\(0\leq g\leq f_1+f_2\). Put \(g_1=\min(g,f_1)\) and \(g_2=g-g_1\). Then
\(0\leq g_1\leq f_1\) and \(0\leq g_2\leq f_2\), and
\(\varphi(g)=\varphi(g_1)+\varphi(g_2)\leq\varphi_+(f_1)+\varphi_+(f_2)\).

\begin{itemize}
\tightlist
\item
  \(\varphi_+\) is positively homogeneous, and
  \(0\leq\varphi_+(f)\leq\|\varphi\|\,\|f\|_{\sup}\).
\item
  So \(\varphi_+(f^+)-\varphi_+(f^-)\) for real \(f\), and then real and
  imaginary parts, extend \(\varphi_+\) to a positive linear functional.
  The extension is well defined: if \(f=a-b\) with \(a,b\geq0\), then
  \(a+f^-=b+f^+\).
\item
  \(\varphi_+(f)\geq\varphi(f)\) for \(f\geq0\), so \(\varphi_-\) is
  positive.
\end{itemize}

\emph{The norm.} \(\|\varphi\|\leq\|\varphi_+\|+\|\varphi_-\|\).
Conversely, let \(\varepsilon>0\).

\begin{itemize}
\tightlist
\item
  Choose \(0\leq f,h\leq1\) with
  \(\varphi_+(f)>\|\varphi_+\|-\varepsilon\) and
  \(\varphi_-(h)>\|\varphi_-\|-\varepsilon\).
\item
  Choose \(0\leq g\leq f\) with \(\varphi(g)>\varphi_+(f)-\varepsilon\).
\item
  Since \(\varphi_-(h)=\sup\{-\varphi(g'):0\leq g'\leq h\}\) (substitute
  \(g'=h-g''\) in the definition), choose \(0\leq g'\leq h\) with
  \(-\varphi(g')>\varphi_-(h)-\varepsilon\).
\item
  Then \(g-g'\) takes values in \([-1,1]\), and
  \(\|\varphi\|\geq\varphi(g-g')>\|\varphi_+\|+\|\varphi_-\|-4\varepsilon\).
\end{itemize}

\emph{General \(\varphi\).} Put
\(\varphi^\dagger(f)=\overline{\varphi(\bar f)}\),
\(\varphi_1=\frac12(\varphi+\varphi^\dagger)\) and
\(\varphi_2=\frac1{2i}(\varphi-\varphi^\dagger)\). For real \(f\),
\(\varphi_1(f)=\operatorname{Re}\varphi(f)\) and
\(\varphi_2(f)=\operatorname{Im}\varphi(f)\).

(2) Apply (1) to \(\varphi_1\) and \(\varphi_2\), and Theorem 2.2 to the
four positive parts restricted to \(C_c(X)\). This gives Radon measures
\(\nu_0,\nu_2\) for \(\varphi_{1,+},\varphi_{1,-}\) and \(\nu_1,\nu_3\)
for \(\varphi_{2,+},\varphi_{2,-}\). They are finite, with \(\nu_k(X)\)
the norm of the corresponding positive part. The identity
\(\varphi(f)=\sum_ki^k\int f\,d\nu_k\) holds on \(C_c(X)\), and both
sides are continuous for \(\|\cdot\|_{\sup}\), so it holds on
\(C_0(X)\).

\emph{Independence.} Let
\(\sum_ki^k\int f\,d\nu_k=\sum_ki^k\int f\,d\nu_k'\) on \(C_0(X)\), with
finite Radon measures.

\begin{itemize}
\tightlist
\item
  Taking real parts at real \(f\), the finite measures \(\nu_0+\nu_2'\)
  and \(\nu_0'+\nu_2\) give the same integrals on \(C_c(X)\).
\item
  A sum of two finite Radon measures is Radon: (R1) is clear; for (R2)
  intersect the two open sets; for (R3) take the union of the two
  compact sets.
\item
  By the uniqueness in Theorem 2.2 the two sums are equal, so
  \(\nu_0-\nu_2=\nu_0'-\nu_2'\) as set functions. The imaginary parts
  are treated in the same way.
\end{itemize}

\emph{Integrals.} For a bounded Borel \(h\), put
\(\int h\,d\mu=\sum_ki^k\int h\,d\nu_k\). This depends only on \(\mu\):
it does for simple functions, and every bounded Borel function is a
uniform limit of simple ones. Also \(|\int h\,d\mu|\leq\int|h|\,d\nu\)
with \(\nu=\sum_k\nu_k\). The converse statement of (2) is clear, with
\(\|\varphi\|\leq\nu(X)\).

(3) \emph{\(\|\varphi\|\leq|\mu|(X)\).} For a Borel simple function
\(s=\sum_jc_j1_{E_j}\), with the \(E_j\) a partition of \(X\),
\(|\int s\,d\mu|\leq\max_j|c_j|\sum_j|\mu(E_j)|\leq\|s\|_{\sup}|\mu|(X)\).
Let \(f\in C_0(X)\). Cutting the range of \(f\) into small Borel pieces,
and taking on each piece a value of \(f\), gives simple \(s_n\to f\)
uniformly with \(\|s_n\|_{\sup}\leq\|f\|_{\sup}\). So
\(|\varphi(f)|\leq\|f\|_{\sup}|\mu|(X)\).

\emph{\(|\mu|(X)\leq\|\varphi\|\).} Let \(E_1,\dots,E_n\) be a Borel
partition of \(X\) and \(\varepsilon>0\). The measure
\(\nu=\sum_k\nu_k\) is a finite Radon measure.

\begin{itemize}
\tightlist
\item
  By Proposition 2.3(1), choose compact \(K_j\subseteq E_j\) with
  \(\nu(E_j\setminus K_j)<\varepsilon/n\).
\item
  Disjoint compact sets in a Hausdorff space have disjoint open
  neighbourhoods: separate each point of one from the other set by
  compactness of the other, then use compactness of the first. Doing
  this for each pair and intersecting, we get disjoint open
  \(W_j\supseteq K_j\). By (R2) choose open \(U_j\) with
  \(K_j\subseteq U_j\subseteq W_j\) and
  \(\nu(U_j\setminus K_j)<\varepsilon/n\).
\item
  By (T2) choose \(K_j\prec g_j\prec U_j\). Let
  \(c_j=\overline{\mu(E_j)}/|\mu(E_j)|\), or \(c_j=0\) if
  \(\mu(E_j)=0\), and \(f=\sum_jc_jg_j\). The supports are disjoint, so
  \(\|f\|_{\sup}\leq1\).
\item
  \(|g_j-1_{E_j}|\leq1_{U_j\setminus K_j}+1_{E_j\setminus K_j}\), so
  \(|\int g_j\,d\mu-\mu(E_j)|<2\varepsilon/n\).
\item
  Hence
  \(\sum_j|\mu(E_j)|=\sum_jc_j\mu(E_j)\leq|\varphi(f)|+2\varepsilon\leq\|\varphi\|+2\varepsilon\).
\end{itemize}

The variation is itself a finite Radon measure, also for complex
\(\mu\). Refining finite partitions proves finite additivity of
\(|\mu|\) on disjoint Borel sets, and \(|\mu|(E)\leq\nu(E)\) with the
finite Radon measure \(\nu=\sum_k\nu_k\). For a disjoint sequence, the
\(\nu\)-measure of the union of the tail tends to zero. The same
domination makes its variation tend to zero, turning finite additivity
into countable additivity. To check regularity, approximate a Borel
\(E\) from within by a compact \(K\) and from outside by an open \(U\)
in \(\nu\)-measure, using Proposition 2.3(1) and (R2). Then
\(|\mu|(E\setminus K)\leq\nu(E\setminus K)\) and
\(|\mu|(U\setminus E)\leq\nu(U\setminus E)\), proving inner and outer
regularity.

(4) If \(\varphi\) is positive, then \(\varphi_+=\varphi\),
\(\varphi_-=0\) and \(\varphi_2=0\), so \(\mu\) is the measure of
Theorem 2.2. If \(\mu\geq0\), then \(\int s\,d\mu\geq0\) for simple
\(s\geq0\), and every \(f\geq0\) in \(C_0(X)\) is a uniform limit of
such \(s\); so \(\varphi\) is positive.

Let \(\varphi\) be hermitian, so \(\varphi_2=0\) and
\(\mu=\nu_+-\nu_-\).

\begin{itemize}
\tightlist
\item
  \(|\mu|(E)\leq\nu_+(E)+\nu_-(E)\) for every Borel \(E\), since
  \(|\mu(E_j)|\leq\nu_+(E_j)+\nu_-(E_j)\).
\item
  \(|\mu|\) is additive on disjoint Borel sets: partitions of \(E\) and
  \(F\) combine to a partition of \(E\cup F\), and a partition of
  \(E\cup F\) can be refined by \(E\) and \(F\) without decreasing the
  sum.
\item
  By (3) and (1),
  \(|\mu|(X)=\|\varphi\|=\|\varphi_+\|+\|\varphi_-\|=\nu_+(X)+\nu_-(X)\).
\item
  So
  \(|\mu|(E)+|\mu|(X\setminus E)=(\nu_++\nu_-)(E)+(\nu_++\nu_-)(X\setminus E)\),
  with \(\leq\) in each term. Hence \(|\mu|(E)=(\nu_++\nu_-)(E)\).
  \(\square\)
\end{itemize}

For compact \(X\), \(C_0(X)=C(X)\), and the finite Radon measures are
regular: outer regular by (R2), and inner regular on all Borel sets by
Proposition 2.3(1). So Theorem 2.4 identifies the dual of \(C(X)\) with
the complex regular Borel measures on \(X\), normed by total variation.

\subsection[3. Image measures, densities, and approximation by
continuous
functions]{\texorpdfstring{\protect\hypertarget{haar-measure-on-locally-compact-groups--oa-fnd-hm-02}{}{}3.
Image measures, densities, and approximation by continuous
functions}{3. Image measures, densities, and approximation by continuous functions}}\label{haar-measure-on-locally-compact-groups--OA-FND-HM-02}

This section collects four operations on a Radon measure that the rest
of the lesson uses constantly: moving it by a homeomorphism, multiplying
it by a positive continuous density, integrating lower semicontinuous
functions, and approximating \(L^p\) functions by continuous ones.

\textbf{Proposition 3.1.} Fix a Radon measure \(\mu\) on an LCH space
\(X\).

\begin{enumerate}
\tightlist
\item
  (\emph{Homeomorphisms}) Let \(\theta\colon X\to X'\) be a
  homeomorphism onto an LCH space. Then
  \(\theta_*\mu(E)=\mu(\theta^{-1}(E))\) defines a Radon measure on
  \(X'\), and \(\int f\,d\theta_*\mu=\int f\circ\theta\,d\mu\) for every
  Borel \(f\geq0\) and every \(f\in L^1(\theta_*\mu)\).
\item
  (\emph{Positive continuous densities}) Let
  \(\varphi\colon X\to(0,\infty)\) be continuous. Then
  \(\varphi\mu(E)=\int_E\varphi\,d\mu\) defines a Radon measure;
  \(\int f\,d(\varphi\mu)=\int f\varphi\,d\mu\) for every Borel
  \(f\geq0\); and \(\varphi\mu\) has the same null sets and the same
  \(\sigma\)-finite sets as \(\mu\). Consequently, if a Radon measure
  \(\nu\) satisfies \(\int f\,d\nu=\int f\varphi\,d\mu\) for all
  \(f\in C_c(X)\), then \(\nu=\varphi\mu\).
\item
  (\emph{Lower semicontinuous functions}) For every lsc
  \(h\colon X\to[0,\infty]\),
  \[\int h\,d\mu=\sup\Big\{\int g\,d\mu:g\in C_c(X),\ 0\leq g\leq h\Big\}.
  \tag{3.1}\]
\item
  (\emph{Density}) Let \(1\leq p<\infty\). Every \(f\in L^p(\mu)\)
  vanishes outside a \(\sigma\)-finite Borel set, and hence, after a
  change on a null set, outside a \(\sigma\)-compact set. The space
  \(C_c(X)\) is dense in \(L^p(\mu)\).
\end{enumerate}

\textbf{Proof.} (1) Since \(\theta\) is a homeomorphism, \(\theta^{-1}\)
maps compact, open and Borel sets of \(X'\) onto sets of the same kind
in \(X\), and it preserves inclusions. So (R1)--(R3) transfer from
\(\mu\) to \(\theta_*\mu\). The integral formula holds for indicators by
definition, and by linearity for simple functions. The layer functions
(1.1) and monotone convergence extend it to Borel \(f\geq0\), and
splitting \(f\) into four nonnegative parts extends it to integrable
\(f\).

(2) The formula for integrals holds for indicators by definition, hence
for Borel \(f\geq0\) by the layer functions (1.1) and monotone
convergence; in particular \(\varphi\mu\) is countably additive, so it
is a measure. Since \(\varphi>0\), \(\int_E\varphi\,d\mu=0\) exactly
when \(\mu(E)=0\). The sets \(\{1/n\leq\varphi\leq n\}\) cover \(X\),
and on each of them the two measures are comparable, so the
\(\sigma\)-finite sets agree. It remains to check (R1)--(R3) for
\(\varphi\mu\). (R1) holds because \(\varphi\) is bounded on compact
sets. (R2): Let \(E\) be Borel with \(\varphi\mu(E)<\infty\) (otherwise
take \(U=X\)), and let \(\varepsilon>0\). For \(j\in\mathbb Z\), the
open sets \(V_j=\{2^{j-1}<\varphi<2^{j+1}\}\) cover \(X\). Put
\(E_j=E\cap V_j\). Then \(\mu(E_j)\leq2^{1-j}\varphi\mu(E_j)<\infty\),
so by (R2) for \(\mu\) there is an open \(U_j\) with
\(E_j\subseteq U_j\subseteq V_j\) and
\(\mu(U_j\setminus E_j)<\varepsilon2^{-|j|-j-3}\). Then
\(\varphi\mu(U_j\setminus E_j)\leq2^{j+1}\mu(U_j\setminus E_j)<\varepsilon2^{-|j|-2}\).
The set \(U=\bigcup_jU_j\) is open, contains \(E\), and
\(U\setminus E\subseteq\bigcup_j(U_j\setminus E_j)\). So
\(\varphi\mu(U\setminus E)<\varepsilon\sum_j2^{-|j|-2}<\varepsilon\).
(R3): Let \(U\) be open and \(c<\varphi\mu(U)\). If
\(\mu(U\cap\{\varphi>s\})=\infty\) for some \(s>0\), then (R3) for
\(\mu\) gives a compact \(K\subseteq U\cap\{\varphi>s\}\) with
\(\mu(K)>c/s\), and \(\varphi\mu(K)\geq s\mu(K)>c\). Otherwise all these
sets have finite \(\mu\)-measure. Put
\(D_j=U\cap\{2^j<\varphi\leq2^{j+1}\}\) for \(j\in\mathbb Z\). These
Borel sets are disjoint, cover \(U\), and have finite \(\mu\)-measure.
Choose a finite set \(J\subseteq\mathbb Z\) with
\(\sum_{j\in J}\varphi\mu(D_j)>c+\eta\) for some \(\eta>0\). By inner
regularity on sets of finite measure (Proposition 2.3(1)), choose
compact \(K_j\subseteq D_j\) with
\(2^{j+1}\mu(D_j\setminus K_j)<\eta/|J|\). The compact set
\(K=\bigcup_{j\in J}K_j\subseteq U\) satisfies
\(\varphi\mu(K)\geq\sum_{j\in J}\big(\varphi\mu(D_j)-2^{j+1}\mu(D_j\setminus K_j)\big)>c\).
The last sentence of (2) follows from the uniqueness in the Riesz
theorem 2.2.

(3) The inequality \(\geq\) is clear. For \(\leq\), let
\(c<\int h\,d\mu\). The layer functions (1.1) of \(h\) have the form
\(\sum_i(t_i-t_{i-1})1_{\{h>t_i\}}\) with \(t_i=i2^{-n}\), and their
integrals tend to \(\int h\,d\mu\) by monotone convergence. So there is
a partition \(0=t_0<t_1<\dots<t_k\) with
\(\sum_i(t_i-t_{i-1})\mu(\{h>t_i\})>c\). Each set \(\{h>t_i\}\) is open,
so (R3) gives compact \(K_i\subseteq\{h>t_i\}\) with
\(\sum_i(t_i-t_{i-1})\mu(K_i)>c\). Choose
\(K_i\prec\varphi_i\prec\{h>t_i\}\) by (T2) and put
\(g=\sum_i(t_i-t_{i-1})\varphi_i\in C_c(X)\). At a point \(x\),
\(\varphi_i(x)\neq0\) forces \(h(x)>t_i\), so \(g(x)\) is at most the
sum of \(t_i-t_{i-1}\) over the indices \(i\) with \(t_i<h(x)\), which
is less than \(h(x)\) (or is \(0\)). So \(0\leq g\leq h\) and
\(\int g\,d\mu\geq\sum_i(t_i-t_{i-1})\mu(K_i)>c\).

(4) The set \(\{f\neq0\}\) is the union of the sets \(\{|f|>1/n\}\), and
by Chebyshev\textquotesingle s inequality each has measure at most
\(n^p\|f\|_p^p\). So \(\{f\neq0\}\) is \(\sigma\)-finite, and
Proposition 2.3(3) gives the \(\sigma\)-compact set. Simple functions
are dense in \(L^p(\mu)\), and a simple function in \(L^p\) is a
combination of indicators of Borel sets of finite measure. So it
suffices to approximate \(1_E\) with \(\mu(E)<\infty\). Given
\(\varepsilon>0\), inner regularity (Proposition 2.3(1)) and outer
regularity (R2) give a compact \(K\) and an open \(U\) with
\(K\subseteq E\subseteq U\) and \(\mu(U\setminus K)<\varepsilon\). Take
\(K\prec f\prec U\). Then \(|f-1_E|\leq1_{U\setminus K}\), so
\(\|f-1_E\|_p\leq\varepsilon^{1/p}\). \(\square\)

Part (4) needs no \(\sigma\)-finiteness of \(\mu\): an \(L^p\) function
with \(p<\infty\) automatically lives on a \(\sigma\)-finite set, and
there inner regularity holds. For a Haar measure (Section 8) and
\(p=2\), it says that \(C_c(G)\) is dense in \(L^2(G)\). In (2) the
density must be strictly positive; Exercise 16.1 shows what goes wrong
with a density that vanishes on a closed set.

\subsection[4. Products of Radon
measures]{\texorpdfstring{\protect\hypertarget{haar-measure-on-locally-compact-groups--oa-fnd-hm-03}{}{}4.
Products of Radon
measures}{4. Products of Radon measures}}\label{haar-measure-on-locally-compact-groups--OA-FND-HM-03}

Let \(X,Y\) be LCH spaces with Radon measures \(\mu,\nu\). For
\(F\in C_c(X\times Y)\) write \(F_x=F(x,\cdot)\), and let \(K_F\) and
\(L_F\) be the projections of \(\operatorname{supp}F\) to \(X\) and
\(Y\); they are compact. The product of \(\mu\) and \(\nu\) is built
from iterated integrals of such \(F\), and the first step is to
approximate \(F\) by sums of products.

\textbf{Lemma 4.1} (Tensor approximation). Let \(F\in C_c(X\times Y)\).

\begin{enumerate}
\tightlist
\item
  The map \(x\mapsto F_x\) is continuous from \(X\) to \(C_c(Y)\) with
  the norm \(\|\cdot\|_{\sup}\), and \(F_x\) vanishes outside \(L_F\).
\item
  Let \(U\supseteq K_F\) be open with compact closure. For every
  \(\varepsilon>0\) there are \(\varphi_1,\dots,\varphi_n\prec U\) and
  \(g_1,\dots,g_n\in C_c(Y)\) vanishing outside \(L_F\) such that
  \(|F(x,y)-\sum_i\varphi_i(x)g_i(y)|\leq\varepsilon\) for all
  \((x,y)\).
\end{enumerate}

\textbf{Proof.} (1) Fix \(x_0\) and \(\varepsilon>0\). For each
\(y\in L_F\), continuity of \(F\) at \((x_0,y)\) gives open
\(P_y\ni x_0\) and \(Q_y\ni y\) with
\(|F(x,y')-F(x_0,y)|<\varepsilon/4\) on \(P_y\times Q_y\). Finitely many
\(Q_{y_k}\) cover \(L_F\); let \(P=\bigcap_kP_{y_k}\). For \(x\in P\)
and \(y'\in L_F\), pick \(k\) with \(y'\in Q_{y_k}\); comparing both
\(F(x,y')\) and \(F(x_0,y')\) with \(F(x_0,y_k)\) gives
\(|F(x,y')-F(x_0,y')|<\varepsilon/2\). Off \(L_F\) both values are
\(0\). (2) By (1), each \(x\in K_F\) has an open neighbourhood
\(P_x\subseteq U\) with \(\|F_{x'}-F_x\|_{\sup}<\varepsilon\) for
\(x'\in P_x\). Cover \(K_F\) by \(P_{x_1},\dots,P_{x_n}\), take a
partition of unity \(\varphi_i\prec P_{x_i}\) from (T3), and let
\(g_i=F_{x_i}\). Then
\[F(x,y)-\sum_i\varphi_i(x)F(x_i,y)=\Big(1-\sum_i\varphi_i(x)\Big)F(x,y)+\sum_i\varphi_i(x)\big(F(x,y)-F(x_i,y)\big).\]
The first term vanishes: \(\sum_i\varphi_i=1\) on \(K_F\), and \(F_x=0\)
off \(K_F\). In the second, \(\varphi_i(x)\neq0\) forces
\(x\in P_{x_i}\), so the term is at most
\(\varepsilon\sum_i\varphi_i(x)\leq\varepsilon\) in absolute value.
\(\square\)

\textbf{Proposition 4.2} (Iterated integrals). Let
\(F\in C_c(X\times Y)\).

\begin{enumerate}
\tightlist
\item
  The function \(x\mapsto\int F(x,y)\,d\nu(y)\) is in \(C_c(X)\) and
  vanishes off \(K_F\); similarly with the roles of the factors
  exchanged.
\item
  The two iterated integrals agree:
  \[\int\Big(\int F(x,y)\,d\nu(y)\Big)d\mu(x)=\int\Big(\int F(x,y)\,d\mu(x)\Big)d\nu(y).
  \tag{4.1}\]
\item
  \(F\) is measurable for the product \(\sigma\)-algebra
  \(\mathcal B(X)\otimes\mathcal B(Y)\).
\end{enumerate}

\textbf{Proof.} (1)
\(\big|\int F(x,y)\,d\nu(y)-\int F(x_0,y)\,d\nu(y)\big|\leq\|F_x-F_{x_0}\|_{\sup}\,\nu(L_F)\),
which tends to \(0\) as \(x\to x_0\) by Lemma 4.1(1). (2) Choose an open
\(U\supseteq K_F\) with compact closure by (T1). For a sum
\(F'=\sum_i\varphi_i\otimes g_i\) as in Lemma 4.1(2), both sides equal
\(\sum_i\int\varphi_i\,d\mu\int g_i\,d\nu\). All these functions vanish
outside the compact set \(\overline U\times L_F\), and
\(|F-F'|\leq\varepsilon\), so each side of (4.1) for \(F\) differs from
the same side for \(F'\) by at most
\(\varepsilon\,\mu(\overline U)\,\nu(L_F)\). Let \(\varepsilon\to0\).
(3) The approximants are
\(\mathcal B(X)\otimes\mathcal B(Y)\)-measurable, and \(F\) is their
pointwise limit as \(\varepsilon=1/m\to0\). \(\square\)

\textbf{Definition 4.3} (Radon product). The \emph{Radon product}
\(\mu\hat\times\nu\) is the Radon measure on \(X\times Y\) that the
Riesz theorem 2.2 gives for the positive linear functional
\(F\mapsto\int\big(\int F(x,y)\,d\nu(y)\big)d\mu(x)\) on
\(C_c(X\times Y)\). By (4.1) it also represents the other iterated
integral.

\textbf{Proposition 4.4} (Rectangles). For open \(U\subseteq X\),
\(V\subseteq Y\) and compact \(K\subseteq X\), \(L\subseteq Y\), with
the convention \(0\cdot\infty=0\),
\[(\mu\hat\times\nu)(U\times V)=\mu(U)\,\nu(V),\qquad(\mu\hat\times\nu)(K\times L)=\mu(K)\,\nu(L).
\tag{4.2}\]

\textbf{Proof.} If \(F\prec U\times V\), then
\(x\mapsto\int F(x,y)\,d\nu(y)\) vanishes off \(U\) and is at most
\(\nu(V)\), so \(\int F\,d(\mu\hat\times\nu)\leq\mu(U)\nu(V)\); when
\(\mu(U)=0\) the integral is \(0\). If \(f\prec U\) and \(g\prec V\),
then \(f\otimes g\prec U\times V\) and its integral is
\(\int f\,d\mu\int g\,d\nu\). Taking suprema and using (2.1) for
\(\mu\), \(\nu\) and \(\mu\hat\times\nu\) gives the first formula. For
the second, if \(F\in C_c(X\times Y)\) and \(1_{K\times L}\leq F\), then
\(\int F(x,y)\,d\nu(y)\geq\nu(L)1_K(x)\), so
\(\int F\,d(\mu\hat\times\nu)\geq\mu(K)\nu(L)\). If \(1_K\leq f\) and
\(1_L\leq g\), then \(1_{K\times L}\leq f\otimes g\), whose integral is
\(\int f\,d\mu\int g\,d\nu\). Taking infima and using (2.1) gives the
second formula. \(\square\)

\subsection[5. Tonelli\textquotesingle s and Fubini\textquotesingle s
theorems for Radon
products]{\texorpdfstring{\protect\hypertarget{haar-measure-on-locally-compact-groups--oa-fnd-hm-10}{}{}5.
Tonelli\textquotesingle s and Fubini\textquotesingle s theorems for
Radon
products}{5. Tonelli\textquotesingle s and Fubini\textquotesingle s theorems for Radon products}}\label{haar-measure-on-locally-compact-groups--OA-FND-HM-10}

Let \(X,Y\) be LCH spaces with Radon measures \(\mu,\nu\), and let
\(\pi=\mu\hat\times\nu\) be their Radon product (Definition 4.3). For
\(E\subseteq X\times Y\) put \(E_x=\{y:(x,y)\in E\}\) and
\(E^y=\{x:(x,y)\in E\}\); if \(E\) is Borel, so are all sections,
because \(y\mapsto(x,y)\) is continuous. A function on \(X\) is
\emph{\(\mu\)-a.e. Borel} if it agrees with a Borel function outside a
Borel \(\mu\)-null set; its integral is that of the Borel function.

\textbf{Theorem 5.1.} Everything below also holds with the roles of
\(X\) and \(Y\) exchanged, since \(\pi\) represents both iterated
integrals (4.1).

\begin{enumerate}
\tightlist
\item
  For every open \(W\subseteq X\times Y\), the function
  \(x\mapsto\nu(W_x)\) is lsc and \(\pi(W)=\int\nu(W_x)\,d\mu(x)\).
\item
  For every compact \(C\subseteq X\times Y\), the function
  \(x\mapsto\nu(C_x)\) is Borel and \(\pi(C)=\int\nu(C_x)\,d\mu(x)\).
\item
  For every Borel \(E\) with \(\pi(E)<\infty\), the function
  \(x\mapsto\nu(E_x)\) is \(\mu\)-a.e. Borel and
  \(\pi(E)=\int\nu(E_x)\,d\mu(x)\). If \(\pi(E)=0\), then \(\nu(E_x)=0\)
  for \(\mu\)-almost every \(x\).
\item
  (\emph{Tonelli}) Let \(F\colon X\times Y\to[0,\infty]\) be Borel and
  vanish outside a set that is \(\sigma\)-finite for \(\pi\). Then
  \(x\mapsto\int F(x,y)\,d\nu(y)\) is \(\mu\)-a.e. Borel,
  \(y\mapsto\int F(x,y)\,d\mu(x)\) is \(\nu\)-a.e. Borel, and
  \[\int F\,d\pi=\int\Big(\int F(x,y)\,d\nu(y)\Big)d\mu(x)=\int\Big(\int F(x,y)\,d\mu(x)\Big)d\nu(y).
  \tag{5.1}\]
\item
  (\emph{Fubini}) Let \(F\in L^1(\pi)\). Then \(F(x,\cdot)\in L^1(\nu)\)
  for \(\mu\)-almost every \(x\), the \(\mu\)-a.e. defined function
  \(x\mapsto\int F(x,y)\,d\nu(y)\) is \(\mu\)-integrable, and (5.1)
  holds.
\item
  (\emph{Rectangles}) Let \(A\subseteq X\) and \(B\subseteq Y\) be Borel
  and \(\sigma\)-finite for \(\mu\) and \(\nu\). Then \(A\times B\) is
  \(\sigma\)-finite for \(\pi\), and \(\pi(A\times B)=\mu(A)\nu(B)\)
  with \(0\cdot\infty=0\). In particular \(A\times B\) is \(\pi\)-null
  when \(A\) is \(\mu\)-null.
\item
  If \(X\) and \(Y\) are \(\sigma\)-compact,
  \(\pi(E)=\int\nu(E_x)\,d\mu(x)\) for every
  \(E\in\mathcal B(X)\otimes\mathcal B(Y)\), so \(\pi\) extends the
  usual product measure. If \(X\) and \(Y\) are second countable, then
  \(\mathcal B(X\times Y)=\mathcal B(X)\otimes\mathcal B(Y)\).
\end{enumerate}

\emph{Scope warning:} Functions vanishing outside a \(\sigma\)-compact
set meet the support hypothesis, but (4) and (5) need only a
\(\sigma\)-finite set. Omitting this hypothesis and requiring a product
measure that is inner and outer regular on every Borel set would be
incorrect; Example 13.5 shows that no such product measure exists on
\(\mathbb R_d\times\mathbb R\).

\textbf{Proof.} (1) \emph{Lower semicontinuity.} Let \(c<\nu(W_{x_0})\).
The section \(W_{x_0}\) is open, so (R3) gives a compact
\(L\subseteq W_{x_0}\) with \(\nu(L)>c\). By the tube lemma (T4) applied
to \(\{x_0\}\times L\subseteq W\), there is an open \(P\ni x_0\) with
\(P\times L\subseteq W\). For \(x\in P\), \(L\subseteq W_x\), so
\(\nu(W_x)>c\). \emph{The inequality \(\leq\).} If \(F\prec W\), then
\(\int F(x,y)\,d\nu(y)\leq\nu(W_x)\), and integrating gives
\(\int F\,d\pi\leq\int\nu(W_x)\,d\mu(x)\). Take the supremum over \(F\)
and use (2.1). \emph{The inequality \(\geq\).} Put \(h(x)=\nu(W_x)\) and
let \(c<\int h\,d\mu\). By (3.1) choose \(g\in C_c(X)\) with
\(0\leq g\leq h\) and \(\int g\,d\mu>c\). Let
\(K=\operatorname{supp}g\), and choose \(\eta>0\) with
\(\int g\,d\mu-\eta\,\mu(K)>c\). For each \(x\in K\), (R3) gives a
compact \(L_x\subseteq W_x\) with \(\nu(L_x)\geq g(x)-\eta/2\) (empty if
\(g(x)\leq\eta/2\)). By (T4) there are open \(P_x\ni x\) and
\(Q_x\supseteq L_x\) with \(P_x\times Q_x\subseteq W\); shrinking
\(P_x\), we may also assume \(g<g(x)+\eta/2\) on \(P_x\). Cover \(K\) by
\(P_{x_1},\dots,P_{x_n}\), take a partition of unity
\(\alpha_i\prec P_{x_i}\) with \(\sum_i\alpha_i=1\) on \(K\) (T3), and
\(L_{x_i}\prec\beta_i\prec Q_{x_i}\) (T2). The function
\(F=\sum_i\alpha_i\otimes\beta_i\) satisfies \(0\leq F\leq1\) and
\(\operatorname{supp}F\subseteq\bigcup_i P_{x_i}\times Q_{x_i}\subseteq W\),
so \(F\prec W\). For \(x\in K\),
\[\int F(x,y)\,d\nu(y)=\sum_i\alpha_i(x)\int\beta_i\,d\nu\geq\sum_i\alpha_i(x)\,\nu(L_{x_i})\geq g(x)-\eta ,\]
because \(\alpha_i(x)\neq0\) forces \(x\in P_{x_i}\), and then
\(\nu(L_{x_i})\geq g(x_i)-\eta/2>g(x)-\eta\). Hence
\(\pi(W)\geq\int F\,d\pi\geq\int_K(g-\eta)\,d\mu>c\). (2) Let \(K\) and
\(L\) be the projections of \(C\). By (T1) choose open \(U\supseteq K\)
and \(V\supseteq L\) with compact closures, and put \(W=U\times V\).
Then \(\nu(W_x)=1_U(x)\nu(V)<\infty\), \(\pi(W)=\mu(U)\nu(V)<\infty\) by
(4.2), and \(W\setminus C\) is open. So
\(\nu(C_x)=\nu(W_x)-\nu((W\setminus C)_x)\) is a difference of finite
lsc functions, hence Borel, and by (1),
\(\pi(C)=\pi(W)-\pi(W\setminus C)=\int\nu(C_x)\,d\mu(x)\). (3) By outer
regularity (R2) and inner regularity on sets of finite measure
(Proposition 2.3(1)), choose open \(W_n\supseteq E\) and compact
\(C_n\subseteq E\) with \(\pi(W_n)<\pi(E)+1/n\) and
\(\pi(C_n)>\pi(E)-1/n\). Replacing \(W_n\) by \(W_1\cap\dots\cap W_n\)
and \(C_n\) by \(C_1\cup\dots\cup C_n\), we may assume \(W_n\) decreases
and \(C_n\) increases. Let \(a(x)=\lim_n\nu((W_n)_x)\) and
\(b(x)=\lim_n\nu((C_n)_x)\); both are Borel, and
\(b(x)\leq\nu(E_x)\leq a(x)\). By (1),
\(\int\nu((W_1)_x)\,d\mu=\pi(W_1)<\infty\). Its infinite-value set is
therefore Borel and null; define all section bounds as zero on that
exceptional set before taking differences. Dominated convergence now
gives \(\int a\,d\mu=\lim\pi(W_n)=\pi(E)\). By (2) and monotone
convergence, \(\int b\,d\mu=\lim\pi(C_n)=\pi(E)\). The finite-a.e.
nonnegative difference \(a-b\) has integral zero, so \(a=b\) outside a
Borel \(\mu\)-null set, and there \(\nu(E_x)=b(x)\). If \(\pi(E)=0\),
then \(\int b\,d\mu=0\), so \(b=0\) almost everywhere. (4) By (3), (5.1)
holds for \(F=1_E\) when \(\pi(E)<\infty\). If \(E\) is
\(\sigma\)-finite, write it as an increasing union of Borel sets of
finite measure and use monotone convergence on all three terms;
countably many null sets have a null union. By linearity (5.1) holds for
nonnegative simple Borel functions vanishing outside a \(\sigma\)-finite
set. For general \(F\), the layer functions (1.1) of \(F\) are such
simple functions and increase to \(F\); apply monotone convergence
again. (5) Split \(F\) into the four nonnegative functions
\((\operatorname{Re}F)^\pm\) and \((\operatorname{Im}F)^\pm\), and apply
(4) to each. They vanish outside \(\{F\neq0\}\), which is
\(\sigma\)-finite by Proposition 3.1(4), and their iterated integrals
are finite, so their inner integrals are finite almost everywhere.
Subtract. (6) By Proposition 2.3(3), write \(A=S_A\cup N_A\) and
\(B=S_B\cup N_B\) with \(S_A,S_B\) \(\sigma\)-compact and \(N_A,N_B\)
null. For compact \(L\subseteq Y\) and \(\varepsilon>0\), choose an open
\(U\supseteq N_A\) with \(\mu(U)<\varepsilon\) and an open
\(V\supseteq L\) with compact closure; by (4.2),
\(\pi(N_A\times L)\leq\mu(U)\nu(V)\leq\varepsilon\,\nu(\overline V)\).
So \(N_A\times L\) is null, and so is \(N_A\times S_B\). Choosing also
an open \(V'\supseteq N_B\) with \(\nu(V')<\varepsilon\), we get
\(\pi(N_A\times N_B)\leq\mu(U)\nu(V')<\varepsilon^2\). In the same way
\(S_A\times N_B\) is null. So \(A\times B\) differs from the
\(\sigma\)-compact set \(S_A\times S_B\) by a null set, and it is
\(\sigma\)-finite by (R1). Now (4) with \(F=1_{A\times B}\) gives
\(\pi(A\times B)=\int1_A(x)\nu(B)\,d\mu(x)=\mu(A)\nu(B)\). (7) If \(X\)
and \(Y\) are \(\sigma\)-compact, so is \(X\times Y\), and every Borel
set is \(\sigma\)-finite by (R1); apply (4) to \(1_E\), noting
\(\mathcal B(X)\otimes\mathcal B(Y)\subseteq\mathcal B(X\times Y)\)
because the projections are continuous. For \(E\) in the product
\(\sigma\)-algebra, \(\int\nu(E_x)\,d\mu(x)\) is the usual product
measure of \(\sigma\)-finite factors. If \(X\) and \(Y\) have countable
bases, every open \(W\subseteq X\times Y\) is the union of the countably
many products of basic sets contained in \(W\), so
\(W\in\mathcal B(X)\otimes\mathcal B(Y)\). \(\square\)

\emph{Where the hypotheses enter, and why they are needed.} Parts
(1)--(3) use only the Radon properties. \(\sigma\)-finiteness enters in
(4)--(6), and it cannot be dropped. Example 13.5 gives a closed set in
\(\mathbb R_d\times\mathbb R\) whose two iterated integrals are \(0\)
and \(\infty\), and Example 13.6 gives a set \(\{e\}\times G\) with
\(\pi(\{e\}\times G)=\infty\) although \(\mu(\{e\})=0\). In particular,
(4) applies to every Borel \(F\geq0\) that vanishes outside a
\(\sigma\)-compact set, since \(\sigma\)-compact sets are
\(\sigma\)-finite by (R1).

\textbf{Example 5.2} (Borel sets of products). If \(X\) and \(Y\) are
second countable,
\(\mathcal B(X\times Y)=\mathcal B(X)\otimes\mathcal B(Y)\) (Theorem
5.1(7)). Without second countability this fails even for discrete
groups. Let \(G\) be a discrete group of cardinality greater than
\(2^{\aleph_0}\), for instance a free abelian group on such a set. The
diagonal \(\Delta_G\subseteq G\times G\) is open, but it is not in
\(\mathcal B(G)\otimes\mathcal B(G)\), which here is the
\(\sigma\)-algebra generated by all rectangles. Proof: a set in the
\(\sigma\)-algebra generated by a family \(\mathcal C\) lies in the
\(\sigma\)-algebra generated by some countable subfamily, because the
sets with this property form a \(\sigma\)-algebra containing
\(\mathcal C\). So \(\Delta_G\) lies in the \(\sigma\)-algebra generated
by countably many rectangles \(A_n\times B_n\). Put
\(\varphi(x)=(1_{A_n}(x))_n\) and \(\psi(y)=(1_{B_n}(y))_n\) in
\(\{0,1\}^{\mathbb N}\). The sets \((\varphi\times\psi)^{-1}(S)\),
\(S\subseteq\{0,1\}^{\mathbb N}\times\{0,1\}^{\mathbb N}\), form a
\(\sigma\)-algebra containing each \(A_n\times B_n\), so
\(\Delta_G=(\varphi\times\psi)^{-1}(S)\) for some \(S\). If
\(\varphi(x)=\varphi(x')\), then \((x,x')\in\Delta_G\) because
\((\varphi(x),\psi(x'))=(\varphi(x'),\psi(x'))\in S\); so \(x=x'\).
Hence \(\varphi\) is injective and \(|G|\leq2^{\aleph_0}\), a
contradiction. So the Radon product (here counting measure on
\(G\times G\)) lives on a \(\sigma\)-algebra strictly larger than the
product \(\sigma\)-algebra. This is why Theorem 5.1 and Theorem 6.2 are
proved for Borel functions on \(X\times Y\), not only for
product-measurable ones.

\subsection[6. Hilbert tensor products and L² of a
product]{\texorpdfstring{\protect\hypertarget{haar-measure-on-locally-compact-groups--oa-fnd-hm-11}{}{}6.
Hilbert tensor products and L² of a
product}{6. Hilbert tensor products and L² of a product}}\label{haar-measure-on-locally-compact-groups--OA-FND-HM-11}

\emph{Hilbert tensor products.} For Hilbert spaces \(H,K\),
\(H\otimes K\) is the Hilbert tensor product of
\hyperref[hilbert-spaces-and-compact-operators]{Hilbert spaces and
compact operators}, Theorem 8.1. It carries a bilinear map
\((\xi,\eta)\mapsto\xi\otimes\eta\) with
\(\langle\xi\otimes\eta,\xi'\otimes\eta'\rangle=\langle\xi,\xi'\rangle\langle\eta,\eta'\rangle\),
and the elementary tensors span a dense subspace.

\textbf{Lemma 6.1} (Isometric extension). Suppose
\(B\colon H\times K\to\mathcal K\) is bilinear, where \(\mathcal K\) is
a third Hilbert space, and
\(\langle B(\xi,\eta),B(\xi',\eta')\rangle=\langle\xi,\xi'\rangle\langle\eta,\eta'\rangle\).
There is a unique isometry \(U\colon H\otimes K\to\mathcal K\) with
\(U(\xi\otimes\eta)=B(\xi,\eta)\).

\textbf{Proof.} For finite sums,
\(\|\sum_iB(\xi_i,\eta_i)\|^2=\sum_{i,j}\langle\xi_i,\xi_j\rangle\langle\eta_i,\eta_j\rangle=\|\sum_i\xi_i\otimes\eta_i\|^2\).
So \(\sum_i\xi_i\otimes\eta_i\mapsto\sum_iB(\xi_i,\eta_i)\) is well
defined (a sum of tensors that is \(0\) goes to a vector of norm \(0\))
and isometric on the span of the elementary tensors. It extends by
continuity to the closure, which is \(H\otimes K\). Uniqueness holds
because the elementary tensors are total. \(\square\)

\textbf{Theorem 6.2} (\(L^2\) of the Radon product). Let \(\mu,\nu\) be
Radon measures on LCH spaces \(X,Y\), and \(\pi=\mu\hat\times\nu\). For
functions \(\xi\) on \(X\) and \(\eta\) on \(Y\) write
\((\xi\boxtimes\eta)(x,y)=\xi(x)\eta(y)\).

\begin{enumerate}
\tightlist
\item
  For \(\xi\in L^2(\mu)\) and \(\eta\in L^2(\nu)\),
  \(\xi\boxtimes\eta\in L^2(\pi)\); its class depends only on the
  classes of \(\xi\) and \(\eta\); and
  \(\langle\xi\boxtimes\eta,\xi'\boxtimes\eta'\rangle=\langle\xi,\xi'\rangle\langle\eta,\eta'\rangle\).
\item
  There is a unique unitary
  \(U\colon L^2(\mu)\otimes L^2(\nu)\to L^2(\pi)\) with
  \(U(\xi\otimes\eta)=\xi\boxtimes\eta\).
\item
  Every \(\zeta\in L^2(\pi)\) satisfies \(\zeta(x,\cdot)\in L^2(\nu)\)
  for \(\mu\)-almost every \(x\), and
  \(\|\zeta\|^2=\int\|\zeta(x,\cdot)\|_{L^2(\nu)}^2\,d\mu(x)\); likewise
  in the other variable.
\item
  For a bounded Borel function \(a\) on \(X\), the operator
  \(U^*M_{a\circ\mathrm{pr}_1}U\), where \(M\) denotes multiplication on
  \(L^2(\pi)\), is the unique bounded operator on
  \(L^2(\mu)\otimes L^2(\nu)\) that sends every \(\xi\otimes\eta\) to
  \((a\xi)\otimes\eta\). It is the operator \(m_a\otimes1\). The same
  holds for the second factor.
\end{enumerate}

\textbf{Proof.} (1) Take Borel representatives. The sets
\(S_\xi=\{\xi\neq0\}\) and \(S_\eta=\{\eta\neq0\}\) are
\(\sigma\)-finite by Proposition 3.1(4). The function
\(\xi\boxtimes\eta\) is Borel on \(X\times Y\) and vanishes outside
\(S_\xi\times S_\eta\), which is \(\sigma\)-finite by Theorem 5.1(6).
Tonelli\textquotesingle s theorem 5.1(4) gives
\(\int|\xi\boxtimes\eta|^2\,d\pi=\|\xi\|^2\|\eta\|^2\). If \(\xi=\xi'\)
outside a \(\mu\)-null Borel set \(N\), then
\(\xi\boxtimes\eta=\xi'\boxtimes\eta\) outside \(N\times S_\eta\), which
is \(\pi\)-null by Theorem 5.1(6); similarly in the second variable. The
function \((\xi\bar\xi')\boxtimes(\eta\bar\eta')\) is in \(L^1(\pi)\) by
the Cauchy--Schwarz inequality, and Fubini\textquotesingle s theorem
5.1(5) gives the inner product formula. (2) Lemma 6.1, with
\(B(\xi,\eta)=\xi\boxtimes\eta\), gives the isometry \(U\). Its range is
closed, because \(U\) is an isometry on a complete space, and it
contains \(\varphi\boxtimes g\) for \(\varphi\in C_c(X)\) and
\(g\in C_c(Y)\). These span a dense subspace of \(L^2(\pi)\). Indeed,
let \(\zeta\in L^2(\pi)\) and \(\varepsilon>0\). Since \(C_c\) is dense
in \(L^2\) (Proposition 3.1(4)), there is \(F\in C_c(X\times Y)\) with
\(\|\zeta-F\|_2<\varepsilon\). Choose an open \(O\supseteq K_F\) with
compact closure (T1). The tensor approximation (Lemma 4.1(2)), applied
with \(O\), gives \(F'=\sum_i\varphi_i\boxtimes g_i\) with
\(|F-F'|\leq\delta\), where \(F\) and \(F'\) vanish outside the compact
set \(\overline O\times L_F\). Then
\(\|F-F'\|_2\leq\delta\,\pi(\overline O\times L_F)^{1/2}=\delta\,(\mu(\overline O)\nu(L_F))^{1/2}\)
by (4.2), which is less than \(\varepsilon\) for small \(\delta\). So
\(U\) is onto. Uniqueness holds because elementary tensors are total.
(3) Apply Tonelli\textquotesingle s theorem 5.1(4) to \(|\zeta|^2\). It
vanishes outside the set \(\{\zeta\neq0\}\), which is \(\sigma\)-finite
by Proposition 3.1(4) applied to \(\pi\). (4)
\(M_{a\circ\mathrm{pr}_1}(\xi\boxtimes\eta)=(a\xi)\boxtimes\eta\), so
\(U^*M_{a\circ\mathrm{pr}_1}U(\xi\otimes\eta)=(a\xi)\otimes\eta\). Two
bounded operators that agree on a total set are equal. \(\square\)

In words:
\textbf{\(L^2(\mu)\otimes L^2(\nu)\cong L^2(X\times Y,\mu\hat\times\nu)\)
through \(\xi\otimes\eta\mapsto\xi\boxtimes\eta\), for arbitrary Radon
measures, with the Radon product and not the ordinary product measure.}
Every element of \(L^2(\pi)\) has \(\sigma\)-finite support, so the
pathological sets of Examples 13.5 and 13.6 never carry \(L^2\) mass.

\subsection[7. Topological
groups]{\texorpdfstring{\protect\hypertarget{haar-measure-on-locally-compact-groups--oa-fnd-hm-04}{}{}7.
Topological
groups}{7. Topological groups}}\label{haar-measure-on-locally-compact-groups--OA-FND-HM-04}

\textbf{Proposition 7.1.} Let \(G\) be a group, topologized so that
multiplication \(G\times G\to G\) and inversion are continuous. The
Hausdorff property is not assumed in this proposition.

\begin{enumerate}
\tightlist
\item
  Left and right translations and inversion are homeomorphisms of \(G\).
  If \(U\) is open and \(A\subseteq G\) is any set, then \(AU\) and
  \(UA\) are open.
\item
  Every neighbourhood \(U\) of \(e\) contains a neighbourhood \(V\) of
  \(e\) with \(V=V^{-1}\) and \(VV\subseteq U\).
\item
  The closure of a subgroup is a subgroup. An open subgroup is closed.
\item
  If \(A\) and \(B\) are compact, so is \(AB\).
\item
  Let \(H\) be a subgroup, \(G/H\) the set of left cosets with the
  quotient topology, and \(q\colon G\to G/H\) the quotient map. Then
  \(q\) is open. If \(H\) is closed, \(G/H\) is Hausdorff. If \(G\) is
  locally compact, every point of \(G/H\) has a compact neighbourhood.
  If \(H\) is normal, \(G/H\) is a group with continuous multiplication
  and inversion.
\item
  If \(\{e\}\) is closed, \(G\) is Hausdorff. In general
  \(N=\overline{\{e\}}\) is a normal subgroup as well as closed, and
  \(G/N\) is Hausdorff.
\end{enumerate}

\textbf{Proof.} (1) Translation by \(x\) has the continuous inverse
translation by \(x^{-1}\), and inversion is its own inverse. Also
\(AU=\bigcup_{a\in A}aU\), and each \(aU\) is open. (2) Continuity of
multiplication at \((e,e)\) gives neighbourhoods \(W_1,W_2\) of \(e\)
with \(W_1W_2\subseteq U\). Put \(W=W_1\cap W_2\) and
\(V=W\cap W^{-1}\). (3) The continuous map \((x,y)\mapsto xy^{-1}\)
sends \(H\times H\) into \(H\), hence sends
\(\overline H\times\overline H=\overline{H\times H}\) into
\(\overline H\). If \(H\) is open, its complement is a union of cosets
\(xH\), which are open by (1). (4) Multiplication is continuous and maps
the compact set \(A\times B\) onto \(AB\). (5) For open
\(V\subseteq G\), \(q^{-1}(q(V))=VH\) is open by (1), so \(q(V)\) is
open. Let \(H\) be closed and \(y\notin xH\). The coset \(xH\) is
closed, so some neighbourhood \(U\) of \(e\) has
\(Uy\cap xH=\varnothing\); by (2) take \(V\) symmetric with
\(VV\subseteq U\). If \(q(Vx)\) and \(q(Vy)\) met, we would have
\(v_1xh_1=v_2yh_2\) with \(v_i\in V\) and \(h_i\in H\), so
\(v_1^{-1}v_2\,y=xh_1h_2^{-1}\in xH\) with
\(v_1^{-1}v_2\in VV\subseteq U\), which is impossible. So \(q(Vx)\) and
\(q(Vy)\) are disjoint neighbourhoods of \(xH\) and \(yH\) (\(q\) is
open). If \(C\) is a compact neighbourhood of \(e\) in \(G\), then
\(q(xC)\) is compact and contains the open set
\(q(x\operatorname{int}C)\ni xH\). If \(H\) is normal and \(W\) is an
open neighbourhood of \(q(xy)\), continuity of multiplication in \(G\)
gives open \(V_1\ni x\) and \(V_2\ni y\) with
\(V_1V_2\subseteq q^{-1}(W)\), so \(q(V_1)q(V_2)\subseteq W\) with
\(q(V_i)\) open. Inversion is handled the same way. (6) If \(\{e\}\) is
closed, apply (5) with \(H=\{e\}\). In general \(N\) is a subgroup by
(3), and it lies in every closed subgroup. For \(z\in G\), \(zNz^{-1}\)
is a closed subgroup (conjugation is a homeomorphism and an
automorphism), so it contains \(N\). Applying this to \(z^{-1}\) gives
\(zNz^{-1}=N\). Now apply (5) with \(H=N\). \(\square\)

By (6), dividing by \(\overline{\{e\}}\) removes any failure of the
Hausdorff property, which is why we may assume it throughout.

From now on \(G\) is a locally compact group. Two more facts are needed:
\(G\) splits into \(\sigma\)-compact pieces, and compactly supported
continuous functions are uniformly continuous.

\textbf{Proposition 7.2} (An open \(\sigma\)-compact subgroup). Some
subgroup \(G_0\) of \(G\) is at once open, closed and
\(\sigma\)-compact. Its left cosets \(xG_0\) partition \(G\) into open,
closed, \(\sigma\)-compact pieces.

\textbf{Proof.} Let \(C\) be a compact neighbourhood of \(e\) and
\(U=C\cap C^{-1}\), a compact symmetric neighbourhood of \(e\). The sets
\(U^n=U\cdots U\) (\(n\) factors) are compact by Proposition 7.1(4), and
\(G_0=\bigcup_nU^n\) is the subgroup generated by \(U\). It contains the
neighbourhood \(xU\) of each of its points \(x\), since
\(xU\subseteq U^{n+1}\) when \(x\in U^n\). So it is open, hence closed
by Proposition 7.1(3), and it is \(\sigma\)-compact. If \(G\) is
connected, the open and closed subgroup \(G_0\) is all of \(G\), so
every connected locally compact group is \(\sigma\)-compact. \(\square\)

\textbf{Proposition 7.3} (Uniform continuity). For every
\(f\in C_c(G)\), \(\|R_yf-f\|_{\sup}\to0\) and \(\|L_yf-f\|_{\sup}\to0\)
as \(y\to e\).

\textbf{Proof.} Let \(K=\operatorname{supp}f\) and \(\varepsilon>0\).
For \(x\in K\) choose a neighbourhood \(U_x\) of \(e\) with
\(|f(xz)-f(x)|<\varepsilon/2\) for \(z\in U_x\), and a symmetric
neighbourhood \(V_x\) with \(V_xV_x\subseteq U_x\); note
\(V_x\subseteq U_x\). Finitely many sets \(x_jV_{x_j}\) cover \(K\); let
\(V=\bigcap_jV_{x_j}\), a symmetric neighbourhood of \(e\). Let
\(y\in V\). If \(x\in K\), write \(x=x_ju\) with \(u\in V_{x_j}\). Then
\(xy=x_j(uy)\) with \(uy\in V_{x_j}V_{x_j}\subseteq U_{x_j}\), so
\(f(xy)\) and \(f(x)\) are both within \(\varepsilon/2\) of \(f(x_j)\),
and \(|f(xy)-f(x)|<\varepsilon\). If \(xy\in K\), the same argument
applied to the point \(xy\) and to \(y^{-1}\in V\) gives
\(|f(x)-f(xy)|<\varepsilon\). Otherwise both values are \(0\). So
\(\|R_yf-f\|_{\sup}\leq\varepsilon\) for \(y\in V\). For left
translations, \(L_yf(x)=\check f(x^{-1}y)=R_y\check f(x^{-1})\), so
\(\|L_yf-f\|_{\sup}=\|R_y\check f-\check f\|_{\sup}\) with
\(\check f\in C_c(G)\). \(\square\)

Proposition 7.3 underlies the continuity of translation in \(L^p\)
(Theorem 14.2(6)), and with it the strong continuity of the regular
representations (Exercise 16.3). A bounded function \(f\) with
\(\|L_yf-f\|_{\sup}\to0\) as \(y\to e\) is called \emph{left uniformly
continuous}, and one with \(\|R_yf-f\|_{\sup}\to0\) \emph{right
uniformly continuous}. The two names are also used the other way round,
so this lesson always writes the condition out.

\subsection[8. Existence of Haar
measure]{\texorpdfstring{\protect\hypertarget{haar-measure-on-locally-compact-groups--oa-fnd-hm-05}{}{}8.
Existence of Haar
measure}{8. Existence of Haar measure}}\label{haar-measure-on-locally-compact-groups--OA-FND-HM-05}

\textbf{Definition 8.1.} Fix a locally compact group \(G\). We call a
Radon measure \(\mu\neq0\) on \(G\) a \emph{left} (\emph{right})
\emph{Haar measure} when \(\mu(xE)=\mu(E)\) (\(\mu(Ex)=\mu(E)\)) for
every Borel set \(E\) and every \(x\in G\). A \emph{left Haar integral}
is a left-invariant positive linear functional \(I\neq0\) on \(C_c(G)\),
that is, \(I(L_yf)=I(f)\) for all \(f\) and \(y\).

\textbf{Proposition 8.2} (Invariance through integrals). Fix a Radon
measure \(\mu\) on \(G\).

\begin{enumerate}
\tightlist
\item
  \(\mu\) is a left Haar measure exactly when
  \(\tilde\mu(E)=\mu(E^{-1})\) is a right Haar measure.
\item
  \(\mu(yE)=\mu(E)\) for all Borel \(E\) and all \(y\) exactly when
  \(\int L_yf\,d\mu=\int f\,d\mu\) for all \(f\in C_c(G)\) and all
  \(y\). In that case \(\int f(yx)\,d\mu(x)=\int f\,d\mu\) for all
  \(y\), all Borel \(f\geq0\), and all \(f\in L^1(\mu)\).
\end{enumerate}

So, through the Riesz theorem 2.2, left Haar measures and left Haar
integrals are the same objects.

\textbf{Proof.} (1) Inversion \(\iota\) is a homeomorphism, so
\(\tilde\mu=\iota_*\mu\) is Radon by Proposition 3.1(1). Also
\(\tilde\mu(Ex)=\mu(x^{-1}E^{-1})\), and this equals
\(\tilde\mu(E)=\mu(E^{-1})\) for all \(E\) and \(x\) exactly when
\(\mu\) is left invariant. (2) Let \(\ell_z(x)=zx\) and
\(\mu_y=(\ell_{y^{-1}})_*\mu\), so \(\mu_y(E)=\mu(yE)\). By Proposition
3.1(1), \(\mu_y\) is Radon and
\(\int f\,d\mu_y=\int f(y^{-1}x)\,d\mu(x)=\int L_yf\,d\mu\). If
\(\int L_yf\,d\mu=\int f\,d\mu\) on \(C_c(G)\), the Radon measures
\(\mu_y\) and \(\mu\) agree by the uniqueness in Theorem 2.2. The
converse and the last sentence follow from the same integral formula,
applied to \(y^{-1}\). \(\square\)

\textbf{Theorem 8.3} (Existence). Every locally compact group has a left
Haar measure.

\emph{Historical reading:} F. van Doorn's freely accessible paper
describes Haar's original construction and the later generalization by
Weil. Fremlin supplies a complementary construction on general LCH
groups.

The proof builds a left Haar integral as a limit of normalized "covering
numbers", which compare a function \(f\) with translates of a function
\(\varphi\) of small support.

\emph{Covering numbers.} For \(f,\varphi\in C_c^+(G)\), let
\((f:\varphi)\) be the infimum of \(\sum_jc_j\) over all finite families
\(c_1,\dots,c_n\geq0\) and \(x_1,\dots,x_n\in G\) with
\(f\leq\sum_jc_jL_{x_j}\varphi\). This number is finite. Indeed, the
open set \(W=\{\varphi>\|\varphi\|_{\sup}/2\}\) is nonempty; fix
\(w\in W\). The open sets \(xw^{-1}W\), \(x\in\operatorname{supp}f\),
cover \(\operatorname{supp}f\), so finitely many of them, \(x_jw^{-1}W\)
for \(j\leq n\), do. On \(x_jw^{-1}W\) we have
\(L_{x_jw^{-1}}\varphi>\|\varphi\|_{\sup}/2\). Hence
\(f\leq\sum_j\frac{2\|f\|_{\sup}}{\|\varphi\|_{\sup}}L_{x_jw^{-1}}\varphi\),
and \((f:\varphi)\leq2n\|f\|_{\sup}/\|\varphi\|_{\sup}\). For
\(f,f_1,f_2,\varphi,\psi\in C_c^+(G)\), \(y\in G\) and \(c>0\):
\[\begin{gathered}
(L_yf:\varphi)=(f:\varphi),\quad (f_1+f_2:\varphi)\leq(f_1:\varphi)+(f_2:\varphi),\quad (cf:\varphi)=c(f:\varphi),\\
f_1\leq f_2\Rightarrow(f_1:\varphi)\leq(f_2:\varphi),\quad (f:\varphi)\geq\frac{\|f\|_{\sup}}{\|\varphi\|_{\sup}},\quad (f:\psi)\leq(f:\varphi)(\varphi:\psi).
\end{gathered}
\tag{8.1}\] For the first rule, \(f\leq\sum_jc_jL_{x_j}\varphi\) holds
exactly when \(L_yf\leq\sum_jc_jL_{yx_j}\varphi\). The next three follow
by adding, scaling or comparing dominations. For the fifth, evaluate a
domination at a point \(t_0\) where \(f(t_0)=\|f\|_{\sup}\):
\(\|f\|_{\sup}\leq\sum_jc_j\varphi(x_j^{-1}t_0)\leq\|\varphi\|_{\sup}\sum_jc_j\).
For the last, if \(f\leq\sum_ic_iL_{x_i}\varphi\) and
\(\varphi\leq\sum_kd_kL_{z_k}\psi\), then
\(f\leq\sum_{i,k}c_id_kL_{x_iz_k}\psi\), and
\(\sum_{i,k}c_id_k=(\sum_ic_i)(\sum_kd_k)\).

Fix \(f_0\in C_c^+(G)\) and put
\(I_\varphi(f)=(f:\varphi)/(f_0:\varphi)\). The last rule of (8.1), used
twice, gives \[(f_0:f)^{-1}\leq I_\varphi(f)\leq(f:f_0).
\tag{8.2}\] Each \(I_\varphi\) is left invariant, subadditive,
positively homogeneous and monotone. It is additive only approximately:

\textbf{Lemma 8.4} (Almost additivity). Let \(f_1,f_2\in C_c^+(G)\) and
\(\varepsilon>0\). There is a neighbourhood \(V\) of \(e\) such that
\(I_\varphi(f_1)+I_\varphi(f_2)\leq I_\varphi(f_1+f_2)+\varepsilon\)
whenever \(\operatorname{supp}\varphi\subseteq V\).

\textbf{Proof.} By Urysohn\textquotesingle s lemma (T2) choose
\(g\in C_c^+(G)\) with \(g=1\) on \(\operatorname{supp}(f_1+f_2)\). Let
\(\delta>0\), to be fixed at the end, and put \(h=f_1+f_2+\delta g\).
Let \(h_i=f_i/h\) where \(h>0\), and \(h_i=0\) elsewhere (\(i=1,2\)).
Since \(h\geq\delta\) on \(\operatorname{supp}f_i\), \(h_i\) is
continuous on the open set \(\{h>0\}\) and vanishes on the open set
\(G\setminus\operatorname{supp}f_i\); these two open sets cover \(G\).
So \(h_i\in C_c(G)\), and \(h_1+h_2\leq1\). By uniform continuity
(Proposition 7.3) there is a neighbourhood \(V\) of \(e\) with
\(|h_i(xv)-h_i(x)|<\delta\) for all \(x\in G\), \(v\in V\) and
\(i=1,2\). Let \(\operatorname{supp}\varphi\subseteq V\) and
\(h\leq\sum_jc_jL_{x_j}\varphi\). If \(\varphi(x_j^{-1}x)\neq0\), then
\(x_j^{-1}x\in V\), so \(|h_i(x)-h_i(x_j)|<\delta\). Hence
\[f_i(x)=h(x)h_i(x)\leq\sum_jc_j\varphi(x_j^{-1}x)h_i(x)\leq\sum_jc_j\varphi(x_j^{-1}x)\big(h_i(x_j)+\delta\big).\]
So \((f_i:\varphi)\leq\sum_jc_j(h_i(x_j)+\delta)\), and
\((f_1:\varphi)+(f_2:\varphi)\leq(1+2\delta)\sum_jc_j\). Taking the
infimum over dominations of \(h\) and using (8.1),
\[(f_1:\varphi)+(f_2:\varphi)\leq(1+2\delta)(h:\varphi)\leq(1+2\delta)\big((f_1+f_2:\varphi)+\delta(g:\varphi)\big).\]
Divide by \((f_0:\varphi)\) and use (8.2):
\(I_\varphi(f_1)+I_\varphi(f_2)\leq I_\varphi(f_1+f_2)+2\delta(f_1+f_2:f_0)+\delta(1+2\delta)(g:f_0)\).
Choose \(\delta\) so that the last two terms add up to less than
\(\varepsilon\). \(\square\)

\textbf{Proof of Theorem 8.3.} Let
\(\Omega=\prod_{f\in C_c^+(G)}\big[(f_0:f)^{-1},(f:f_0)\big]\). It is
compact Hausdorff by Tychonoff\textquotesingle s theorem, and by (8.2)
each \(I_\varphi\) is a point of \(\Omega\). For a neighbourhood \(V\)
of \(e\), let \(\Phi_V\) be the closure in \(\Omega\) of
\(\{I_\varphi:\varphi\in C_c^+(G),\ \operatorname{supp}\varphi\subseteq V\}\).
It is nonempty by (T2), and
\(\Phi_{V_1}\cap\dots\cap\Phi_{V_n}\supseteq\Phi_{V_1\cap\dots\cap V_n}\neq\varnothing\).
By compactness some point \(I\) lies in every \(\Phi_V\). By the
definition of the product topology this means: for every neighbourhood
\(V\) of \(e\), all \(f_1,\dots,f_m\in C_c^+(G)\) and all
\(\varepsilon>0\), there is \(\varphi\) with
\(\operatorname{supp}\varphi\subseteq V\) and
\(|I(f_k)-I_\varphi(f_k)|<\varepsilon\) for \(k\leq m\).

Let \(f_1,f_2\in C_c^+(G)\) and \(\varepsilon>0\). Take \(V\) from Lemma
8.4 and such a \(\varphi\) for \(f_1\), \(f_2\) and \(f_1+f_2\). Then
\[I(f_1)+I(f_2)<I_\varphi(f_1)+I_\varphi(f_2)+2\varepsilon\leq I_\varphi(f_1+f_2)+3\varepsilon<I(f_1+f_2)+4\varepsilon,\]
and
\(I(f_1+f_2)<I_\varphi(f_1+f_2)+\varepsilon\leq I_\varphi(f_1)+I_\varphi(f_2)+\varepsilon<I(f_1)+I(f_2)+3\varepsilon\).
So \(I\) is additive on \(C_c^+(G)\). The same approximation, without
the lemma, gives \(I(L_yf)=I(f)\) and \(I(cf)=cI(f)\) for \(c>0\),
because each \(I_\varphi\) has these properties exactly. Also
\(I(f)\geq(f_0:f)^{-1}>0\) and \(I(f_0)=1\).

Put \(I(0)=0\). A real \(f\in C_c(G)\) is \(f^+-f^-\) with
\(f^\pm\in C_c^+(G)\cup\{0\}\); set \(I(f)=I(f^+)-I(f^-)\). If also
\(f=g-h\) with \(g,h\in C_c^+(G)\cup\{0\}\), then \(g+f^-=h+f^+\), and
additivity gives \(I(g)-I(h)=I(f^+)-I(f^-)\). So \(I\) is well defined
and linear on real functions; for complex \(f\) put
\(I(f)=I(\operatorname{Re}f)+iI(\operatorname{Im}f)\). Then \(I\) is a
left Haar integral with \(I(f_0)=1\). By the Riesz theorem 2.2 and
Proposition 8.2, its Radon measure is a left Haar measure. \(\square\)

\emph{Where the hypotheses enter.} Compact supports make every
\((f:\varphi)\) finite. Local compactness and the Hausdorff property
supply the small functions \(\varphi\), through
Urysohn\textquotesingle s lemma (T2), and they make the Riesz theorem
available. Continuity of the group operations gives the uniform
continuity used in Lemma 8.4. The axiom of choice enters through
Tychonoff\textquotesingle s theorem; Cartan\textquotesingle s proof,
mentioned in the introduction, avoids it, but it is not given here.
Right Haar measures exist too: by Proposition 8.2(1), \(\tilde\mu\) is a
right Haar measure.

\subsection[9. Positivity and uniqueness of Haar
measure]{\texorpdfstring{\protect\hypertarget{haar-measure-on-locally-compact-groups--oa-fnd-hm-06}{}{}9.
Positivity and uniqueness of Haar
measure}{9. Positivity and uniqueness of Haar measure}}\label{haar-measure-on-locally-compact-groups--OA-FND-HM-06}

\textbf{Proposition 9.1} (Positivity). A left Haar measure \(\mu\)
satisfies \(\mu(U)>0\) for every nonempty open set \(U\), and
\(\int f\,d\mu>0\) for every \(f\in C_c^+(G)\).

\textbf{Proof.} Suppose \(U\neq\varnothing\) is open and \(\mu(U)=0\).
Fix \(u\in U\). For compact \(K\), the null open sets \(xu^{-1}U\),
\(x\in K\), cover \(K\), and finitely many suffice; so \(\mu(K)=0\).
Then \(\mu(G)=0\) by (R3), which contradicts \(\mu\neq0\). For
\(f\in C_c^+(G)\), the open set \(\{f>\|f\|_{\sup}/2\}\) is nonempty, so
\(\int f\,d\mu\geq\frac12\|f\|_{\sup}\,\mu(\{f>\|f\|_{\sup}/2\})>0\).
\(\square\)

\textbf{Theorem 9.2} (Uniqueness). Any two left Haar measures
\(\mu,\nu\) on \(G\) are proportional: \(\mu=c\nu\) for some
\(c\in(0,\infty)\). Equivalently, every left Haar integral is a positive
multiple of \(f\mapsto\int f\,d\mu\).

\emph{Historical reading:} The accessible original paper {[}von Neumann
1936{]} treats uniqueness for second countable groups. Fremlin's freely
accessible text proves general LCH uniqueness; the separately licensed
supplement gives that argument in full.

\textbf{Proof.} Fix \(f,g\in C_c^+(G)\). It suffices to show
\(\int f\,d\mu/\int f\,d\nu=\int g\,d\mu/\int g\,d\nu\); the
denominators are positive by Proposition 9.1. Fix a compact symmetric
neighbourhood \(V_0\) of \(e\), and let
\(A=(\operatorname{supp}f)V_0\cup V_0(\operatorname{supp}f)\) and
\(B=(\operatorname{supp}g)V_0\cup V_0(\operatorname{supp}g)\), which are
compact by Proposition 7.1(4). For \(y\in V_0\), the function
\(x\mapsto f(xy)-f(yx)\) vanishes off \(A\): \(f(xy)\neq0\) forces
\(x\in(\operatorname{supp}f)y^{-1}\), and \(f(yx)\neq0\) forces
\(x\in y^{-1}\operatorname{supp}f\). Let \(\varepsilon>0\). Since
\(|f(xy)-f(yx)|\leq|R_yf(x)-f(x)|+|f(x)-L_{y^{-1}}f(x)|\), uniform
continuity (Proposition 7.3) gives a symmetric neighbourhood
\(V\subseteq V_0\) of \(e\) with \(|f(xy)-f(yx)|<\varepsilon\) and
\(|g(xy)-g(yx)|<\varepsilon\) for all \(x\in G\) and \(y\in V\). Choose
\(h_1\in C_c^+(G)\) with \(\operatorname{supp}h_1\subseteq V\) and put
\(h=h_1+\check h_1\). Then \(h\in C_c^+(G)\),
\(\operatorname{supp}h\subseteq V\), and \(h(x^{-1})=h(x)\).

Every integrand below is a continuous function with compact support on
\(G\times G\), so (4.1) allows the order of integration to be exchanged.
By left invariance of \(\mu\) in \(x\),
\[\int h\,d\nu\int f\,d\mu=\iint h(y)f(yx)\,d\mu(x)\,d\nu(y).\] By left
invariance of \(\mu\) in \(x\), then an exchange and the symmetry of
\(h\), then left invariance of \(\nu\) in \(y\), then another exchange,
\[\begin{gathered}
\int h\,d\mu\int f\,d\nu\\
=\iint h(y^{-1}x)f(y)\,d\mu(x)\,d\nu(y)\\
=\iint h(x^{-1}y)f(y)\,d\nu(y)\,d\mu(x)\\
=\iint h(y)f(xy)\,d\nu(y)\,d\mu(x)\\
=\iint h(y)f(xy)\,d\mu(x)\,d\nu(y).
\end{gathered}\]
Subtract. For \(y\in\operatorname{supp}h\subseteq V\) the inner
integrand \(h(y)(f(yx)-f(xy))\) is at most \(\varepsilon h(y)\) in
absolute value and vanishes for \(x\notin A\). Hence
\[\Big|\int h\,d\nu\int f\,d\mu-\int h\,d\mu\int f\,d\nu\Big|\leq\varepsilon\,\mu(A)\int h\,d\nu .\]
Dividing by \(\int h\,d\nu\int f\,d\nu\) gives
\(\big|\frac{\int f\,d\mu}{\int f\,d\nu}-\frac{\int h\,d\mu}{\int h\,d\nu}\big|\leq\frac{\varepsilon\mu(A)}{\int f\,d\nu}\).
The same holds for \(g\) with \(B\). So the two ratios for \(f\) and
\(g\) differ by at most
\(\varepsilon\big(\mu(A)/\int f\,d\nu+\mu(B)/\int g\,d\nu\big)\), where
\(A\) and \(B\) do not depend on \(\varepsilon\). Hence the ratio is a
constant \(c>0\) on \(C_c^+(G)\). By linearity
\(\int f\,d\mu=c\int f\,d\nu\) on \(C_c(G)\), and \(\mu=c\nu\) by the
uniqueness in Theorem 2.2. \(\square\)

Only compactly supported continuous integrands on \(G\times G\) occur,
so this proof needs nothing beyond (4.1); in particular it needs no
\(\sigma\)-finiteness. From now on, \(\mu\) denotes a fixed left Haar
measure on \(G\), and we write \(\int f(t)\,dt=\int f\,d\mu\).

Fremlin gives a second proof that compares masses of shrinking symmetric
neighborhoods. The \hyperref[fremlin-neighbourhood-uniqueness]{complete
neighborhood-ratio proof} uses the open-set assertion of Theorem 5.1 and
reaches the same conclusion for every LCH group.

\subsection[10. The modular
function]{\texorpdfstring{\protect\hypertarget{haar-measure-on-locally-compact-groups--oa-fnd-hm-07}{}{}10.
The modular
function}{10. The modular function}}\label{haar-measure-on-locally-compact-groups--OA-FND-HM-07}

A left Haar measure need not be right invariant. The modular function
measures how far it is from right invariance.

\textbf{Theorem 10.1.} Let \(\mu\) be a left Haar measure on \(G\).

\begin{enumerate}
\tightlist
\item
  For each \(x\in G\), \(E\mapsto\mu(Ex)\) is a left Haar measure. So
  there is a unique \(\Delta(x)\in(0,\infty)\) with
  \(\mu(Ex)=\Delta(x)\mu(E)\) for every Borel \(E\). Another left Haar
  measure in place of \(\mu\) gives the same \(\Delta(x)\).
\item
  \(\Delta\) is a continuous homomorphism from \(G\) to the
  multiplicative group \((0,\infty)\).
\item
  For every Borel \(f\geq0\), and for every \(f\in L^1(\mu)\),
  \[\int f(tx)\,dt=\Delta(x)^{-1}\int f(t)\,dt .
  \tag{10.1}\] In particular \(\|R_xf\|_p=\Delta(x)^{-1/p}\|f\|_p\) for
  \(f\in L^p(\mu)\), \(1\leq p<\infty\).
\item
  \(\Delta=1\) on every compact subgroup. Every compact, abelian or
  discrete group is \emph{unimodular} (\(\Delta\equiv1\)). \(G\) is
  unimodular exactly when \(\mu\) is right invariant as well. If
  \(G/[G,G]\) is compact, where \([G,G]\) is the closure of the subgroup
  generated by the commutators \(xyx^{-1}y^{-1}\), then \(G\) is
  unimodular.
\end{enumerate}

\textbf{Proof.} (1) Right translation by \(x\) is a homeomorphism, so
\(\mu_x(E)=\mu(Ex)\) is Radon by Proposition 3.1(1). It is nonzero and
left invariant, since \(y(Ex)=(yE)x\). By uniqueness (Theorem 9.2),
\(\mu_x=\Delta(x)\mu\) with \(\Delta(x)>0\). Replacing \(\mu\) by
\(c\mu\) does not change \(\Delta\). (2) Let \(E\) be a compact
neighbourhood of \(e\), so \(0<\mu(E)<\infty\) by Proposition 9.1 and
(R1). Then
\(\Delta(xy)\mu(E)=\mu(Exy)=\Delta(y)\mu(Ex)=\Delta(y)\Delta(x)\mu(E)\).
So \(\Delta\) is a homomorphism, and \(\Delta(x^{-1})=\Delta(x)^{-1}\).
Continuity is proved after (3). (3) For \(f=1_E\),
\(\int1_E(tx)\,dt=\mu(Ex^{-1})=\Delta(x)^{-1}\mu(E)\). Linearity, then
the layer functions (1.1) with monotone convergence, then splitting into
four nonnegative parts extend this to simple functions, to Borel
\(f\geq0\), and to \(L^1\). Apply it to \(|f|^p\) for the norm formula.
\emph{Continuity of \(\Delta\).} Fix \(f\in C_c^+(G)\); by (3),
\(\Delta(x)^{-1}=\int R_xf\,dt/\int f\,dt\), where the denominator is
positive by Proposition 9.1. Let \(W\) be a compact neighbourhood of
\(e\) and \(x\in x_0W\). Then \(R_xf\) vanishes off the compact set
\((\operatorname{supp}f)W^{-1}x_0^{-1}\), and
\(\|R_xf-R_{x_0}f\|_{\sup}=\|R_{x_0}(R_{x_0^{-1}x}f-f)\|_{\sup}=\|R_{x_0^{-1}x}f-f\|_{\sup}\),
which tends to \(0\) as \(x\to x_0\) by uniform continuity (Proposition
7.3). So \(\int R_xf\,dt\to\int R_{x_0}f\,dt\). (4) \(\Delta(K)\) is a
compact subgroup of \((0,\infty)\) for a compact subgroup \(K\). If it
contained some \(r\neq1\), it would contain all powers \(r^n\),
\(n\in\mathbb Z\), which form an unbounded set. So \(\Delta(K)=\{1\}\);
take \(K=G\) for compact \(G\). For abelian \(G\),
\(\mu(Ex)=\mu(xE)=\mu(E)\). On a discrete group, counting measure is
invariant on both sides, and it is Radon: compact sets are finite and
every set is open. The equivalence with right invariance is the
definition of \(\Delta\). Finally, \(\Delta\) takes values in an abelian
group, so it is \(1\) on every commutator, hence on the subgroup they
generate, and by continuity on its closure \(N=[G,G]\). This \(N\) is
closed, and it is normal because conjugates of commutators are
commutators and conjugation is a homeomorphism. So
\(\Delta=\bar\Delta\circ q\) for a homomorphism \(\bar\Delta\) on the
locally compact group \(G/N\) (Proposition 7.1(5)), and \(\bar\Delta\)
is continuous because \(q\) is open:
\(\bar\Delta^{-1}(W)=q(\Delta^{-1}(W))\). If \(G/N\) is compact,
\(\Delta(G)=\bar\Delta(G/N)\) is a compact subgroup of \((0,\infty)\),
hence \(\{1\}\). \(\square\)

In shorthand, (10.1) reads \(d\mu(tx)=\Delta(x)\,d\mu(t)\).

\subsection[11. Inversion and right Haar
measure]{\texorpdfstring{\protect\hypertarget{haar-measure-on-locally-compact-groups--oa-fnd-hm-08}{}{}11.
Inversion and right Haar
measure}{11. Inversion and right Haar measure}}\label{haar-measure-on-locally-compact-groups--OA-FND-HM-08}

\textbf{Theorem 11.1} (Inversion). For every Borel
\(f\colon G\to[0,\infty]\),
\[\int f(t^{-1})\,dt=\int f(t)\,\Delta(t)^{-1}\,dt ,
\tag{11.1}\] and the same holds for every Borel \(f\) with
\(\int|f|\Delta^{-1}\,d\mu<\infty\). Equivalently, the right Haar
measure \(\tilde\mu(E)=\mu(E^{-1})\) equals \(\Delta^{-1}\mu\), that is,
\(\tilde\mu(E)=\int_E\Delta(t)^{-1}\,dt\). Moreover
\(\int f(t^{-1})\Delta(t)^{-1}\,dt=\int f(t)\,dt\) for every Borel
\(f\geq0\).

\textbf{Proof.} For \(f\in C_c(G)\) put
\(f^\natural(t)=f(t^{-1})\Delta(t)^{-1}\), again in \(C_c(G)\), and
\(J(f)=\int f^\natural\,d\mu\). The functional \(J\) is linear, and
\(J(f)>0\) for \(f\in C_c^+(G)\) by Proposition 9.1. It is left
invariant: with \(k=f^\natural\) and \(y\in G\),
\[(L_yf)^\natural(t)=f\big((ty)^{-1}\big)\Delta(t)^{-1}=f\big((ty)^{-1}\big)\Delta(ty)^{-1}\Delta(y)=\Delta(y)\,k(ty),\]
so \(J(L_yf)=\Delta(y)\Delta(y)^{-1}\int k\,d\mu=J(f)\) by (10.1). So
\(J\) is a left Haar integral, and by Proposition 8.2 and uniqueness
(Theorem 9.2), \(J=c\int\cdot\,d\mu\) on \(C_c(G)\) for some \(c>0\).
Now
\((f^\natural)^\natural(t)=f^\natural(t^{-1})\Delta(t)^{-1}=f(t)\Delta(t)\Delta(t)^{-1}=f(t)\).
So for \(f\in C_c^+(G)\),
\(\int f\,d\mu=J(f^\natural)=c\int f^\natural\,d\mu=cJ(f)=c^2\int f\,d\mu\),
and \(c=1\). Apply \(J=\int\cdot\,d\mu\) to
\(g=f\Delta^{-1}\in C_c(G)\): since
\(g(t^{-1})\Delta(t)^{-1}=f(t^{-1})\), this gives (11.1) for
\(f\in C_c(G)\).

The measure \(\tilde\mu=\iota_*\mu\) is Radon (Proposition 3.1(1)) with
\(\int f\,d\tilde\mu=\int f(t^{-1})\,dt\). The measure
\(\Delta^{-1}\mu\) is Radon by Proposition 3.1(2), because
\(\Delta^{-1}\) is continuous and positive. By the previous paragraph
they integrate every \(f\in C_c(G)\) alike, so they are equal by the
uniqueness in Theorem 2.2. Integrating a Borel \(f\geq0\) against both
(Proposition 3.1(1) and (2)) gives (11.1). For \(f\) with
\(\int|f|\Delta^{-1}d\mu<\infty\), apply this to the four nonnegative
parts of \(f\). For the last formula apply (11.1) to the Borel function
\(t\mapsto f(t)\Delta(t)\), using \(\Delta(t^{-1})=\Delta(t)^{-1}\).
\(\square\)

Formula (11.1) holds for all nonnegative Borel functions, whatever their
support.

\Needspace{5\baselineskip}
\textbf{Corollary 11.2.}

\begin{enumerate}
\tightlist
\item
  \(\mu\) and \(\tilde\mu\) have the same null sets and the same
  \(\sigma\)-finite sets.
\item
  For \(1\leq p<\infty\), \((S_pf)(t)=\Delta(t)^{-1/p}f(t^{-1})\)
  defines a linear isometry of \(L^p(\mu)\) onto itself with
  \(S_p^2=1\). For \(p=2\),
  \((J\xi)(t)=\Delta(t)^{-1/2}\overline{\xi(t^{-1})}\) is a
  conjugate-linear isometric involution of \(L^2(G)\).
\item
  \(f\mapsto\check f\) and \(f\mapsto\Delta^{1/p}f\) are isometries of
  \(L^p(\mu)\) onto \(L^p(\tilde\mu)\).
\item
  If \(G\) is not unimodular, \(\Delta\) is unbounded above and below.
\item
  The measure \((1+\Delta)\mu\) is a Radon measure, and \(C_c(G)\) is
  dense in \(L^2((1+\Delta)\mu)\). Consequently \(C_c(G)\) is a core for
  multiplication by \(\Delta^{1/2}\) on its maximal domain
  \(\{\xi\in L^2(G):\Delta^{1/2}\xi\in L^2(G)\}\), whose graph norm is
  \(\big(\int|\xi|^2(1+\Delta)\,d\mu\big)^{1/2}\).
\end{enumerate}

\textbf{Proof.} (1) Since \(\tilde\mu=\Delta^{-1}\mu\) with a continuous
positive density, this is Proposition 3.1(2). (2) By the last formula of
the theorem applied to \(|f|^p\),
\(\int\Delta(t)^{-1}|f(t^{-1})|^p\,dt=\int|f|^p\,dt\). Also
\(S_p(S_pf)(t)=\Delta(t)^{-1/p}\Delta(t^{-1})^{-1/p}f(t)=f(t)\), so
\(S_p\) is onto. Complex conjugation commutes with \(S_2\) and is a
conjugate-linear isometry. (3)
\(\int|\check f|^p\,d\tilde\mu=\int|f(t)|^p\,dt\) by the definition of
\(\tilde\mu\), and \(\int\Delta|f|^p\,d\tilde\mu=\int|f|^p\,d\mu\) by
the theorem; the inverse maps have the same form. (4) \(\Delta(G)\) is a
subgroup of \((0,\infty)\) containing some \(r\neq1\), hence all
\(r^n\). (5) Apply Proposition 3.1(2) with the positive continuous
density \(1+\Delta\), and then the density of \(C_c\) in \(L^2\)
(Proposition 3.1(4)). The graph norm identity is immediate, and a
function \(\xi\) lies in the maximal domain exactly when it lies in
\(L^2((1+\Delta)\mu)\). \(\square\)

\textbf{Example 11.3} (Explicit Haar measures on \(\mathbb R^\times\)
and on the \(ax+b\) group). Integrals of continuous compactly supported
functions of one real variable below are Riemann integrals, and
substitutions use the change-of-variables rule for Riemann integrals
(see "Results used from other lessons"). (a) On the multiplicative group
\(\mathbb R^\times=\mathbb R\setminus\{0\}\), \(I(f)=\int f(x)\,dx/|x|\)
is a Haar integral: the substitution \(u=cx\) gives
\(\int f(cx)\,dx/|x|=\int f(u)\,du/|u|\) for \(c\neq0\). The group is
abelian, so \(\Delta\equiv1\). (b) Let
\(G=\{(a,b):a>0,\ b\in\mathbb R\}\) with \((a,b)(a',b')=(aa',ab'+b)\),
identity \((1,0)\), inverse \((a,b)^{-1}=(1/a,-b/a)\), and the topology
of the open half-plane. Put
\(I(f)=\int_0^\infty\big(\int_{\mathbb R}f(a,b)\,db\big)a^{-2}\,da\).
For \(x_0=(a_0,b_0)\), the substitutions \(u=a_0b+b_0\) (inner) and
\(v=a_0a\) (outer) give
\[\int_0^\infty\!\!\int_{\mathbb R}f(a_0a,a_0b+b_0)\,db\,\frac{da}{a^2}=\int_0^\infty\frac1{a_0}\int_{\mathbb R}f(a_0a,u)\,du\,\frac{da}{a^2}=\int_0^\infty\!\!\int_{\mathbb R}f(v,u)\,du\,\frac{dv}{v^2},\]
so \(I\) is a left Haar integral, \(d\mu=a^{-2}\,da\,db\). The
substitutions \(u=ab_0+b\) and \(v=aa_0\) give
\(\int f(tx_0)\,d\mu(t)=a_0\int f\,d\mu\), so by (10.1)
\[\Delta(a,b)=1/a .\] The inversion formula (11.1) can be checked
directly: the substitutions \(u=-b/a\) and \(v=1/a\) give
\(\int f(t^{-1})\,d\mu(t)=\int_0^\infty\int_{\mathbb R}f(v,u)\,du\,v^{-1}\,dv\),
which is \(\int f\Delta^{-1}\,d\mu\). So the right Haar measure is
\(a^{-1}\,da\,db\). The commutators
\((a,b)(a',b')(a,b)^{-1}(a',b')^{-1}\) have first coordinate \(1\), and
they fill the normal subgroup \(\{1\}\times\mathbb R\); the quotient
\(G/(\{1\}\times\mathbb R)\cong(0,\infty)\) is not compact, in line with
Theorem 10.1(4). This group is a good test case for the signs in (10.1)
and (11.1).

\subsection{12. Products of
groups}\label{haar-measure-on-locally-compact-groups--12-products-of-groups}

The Radon product of Haar measures is a Haar measure. For two copies of
\(G\) this gives the measure on \(G\times G\) behind convolution
(Section 14); for infinitely many compact groups it gives the Haar
measure of an infinite product.

\textbf{Proposition 12.1} (The group \(G\times G\)). Keep \(G\) and its
left Haar measure \(\mu\), and let \(\pi=\mu\hat\times\mu\) on
\(G\times G\).

\begin{enumerate}
\tightlist
\item
  On the locally compact group \(G\times G\), \(\pi\) is invariant under
  left translations, so it is a left Haar measure; its modular function
  is \((s,t)\mapsto\Delta(s)\Delta(t)\). In the same way, if each
  \(G_j\) (\(j=1,\dots,n\)) carries a left Haar measure \(\mu_j\), the
  iterated Radon product
  \((\cdots(\mu_1\hat\times\mu_2)\cdots)\hat\times\mu_n\) is one on
  \(G_1\times\dots\times G_n\), and it integrates functions in \(C_c\)
  by iterated integrals.
\item
  For every continuous \(a\colon G\to G\), the map
  \(\Phi_a(s,t)=(s,a(s)t)\) is a homeomorphism of \(G\times G\) that
  preserves \(\pi\). So \((W_a\zeta)(s,t)=\zeta(s,a(s)t)\) is a unitary
  on \(L^2(G\times G,\pi)\), with inverse
  \(\zeta\mapsto\zeta(s,a(s)^{-1}t)\).
\item
  Under the unitary \(U\) of Theorem 6.2, the operator
  \(1\otimes\lambda(g)\) is \(W_a\) for the constant \(a\equiv g^{-1}\):
  \(\zeta(s,t)\mapsto\zeta(s,g^{-1}t)\). For \(a(s)=s\), \(W_a\) is the
  operator \((W\zeta)(s,t)=\zeta(s,st)\), with inverse
  \(\zeta(s,s^{-1}t)\).
\end{enumerate}

\textbf{Proof.} (1) For \(F\in C_c(G\times G)\) and
\((g,h)\in G\times G\),
\(\int\big(\int F(gs,ht)\,dt\big)ds=\int\big(\int F(gs,t)\,dt\big)ds=\int\big(\int F(s,t)\,dt\big)ds\),
by left invariance applied first to the inner function
\(t\mapsto F(gs,t)\) and then to the outer function
\(s\mapsto\int F(s,t)\,dt\), which lies in \(C_c(G)\) by Proposition
4.2(1). So \(\pi\) is a left Haar measure by Proposition 8.2 (it is
nonzero by (4.2)). In the same way, (10.1) gives
\(\iint F(sg,th)\,dt\,ds=\Delta(g)^{-1}\Delta(h)^{-1}\iint F\,dt\,ds\),
which identifies the modular function by Theorem 10.1(3). The statement
for \(n\) factors follows by induction, using Proposition 4.2 for the
iterated integral. (2) \(\Phi_a\) is continuous, with continuous inverse
\((s,t)\mapsto(s,a(s)^{-1}t)\). For \(F\in C_c(G\times G)\), also
\(F\circ\Phi_a\in C_c(G\times G)\), and
\(\int F\circ\Phi_a\,d\pi=\int\big(\int F(s,a(s)t)\,dt\big)ds=\int\big(\int F(s,t)\,dt\big)ds\)
by left invariance in \(t\) for each fixed \(s\). So the Radon measures
\((\Phi_a)_*\pi\) (Proposition 3.1(1)) and \(\pi\) agree by the
uniqueness in Theorem 2.2. Hence \(\zeta\mapsto\zeta\circ\Phi_a\)
preserves null sets and the \(L^2(\pi)\) norm, and its inverse is
composition with \(\Phi_a^{-1}\). (3)
\(U(\xi\otimes L_g\eta)(s,t)=\xi(s)\eta(g^{-1}t)=U(\xi\otimes\eta)(s,g^{-1}t)\).
Both operators are bounded and agree on elementary tensors. \(\square\)

\textbf{Example 12.2} (Infinite products of compact groups, and
\((\mathbb Z_2)^{\mathbb N}\) versus \([0,1]\)). Let
\((G_\alpha)_{\alpha\in A}\) be compact groups with Haar measures
\(\mu_\alpha(G_\alpha)=1\) (compact groups are unimodular by Theorem
10.1(4)). The product \(G=\prod_\alpha G_\alpha\), with the product
topology and the group law taken coordinate by coordinate, is a compact
group by Tychonoff\textquotesingle s theorem. Write \(C_F(G)\) for the
continuous functions that depend on only finitely many coordinates,
\(f(x)=f_0(x_{\alpha_1},\dots,x_{\alpha_n})\) with \(f_0\) continuous.
It is a self-adjoint algebra, it contains the constant functions, and it
separates the points of \(G\): if \(x_\alpha\neq x'_\alpha\), (T2) on
\(G_\alpha\) gives a continuous function of the coordinate \(x_\alpha\)
that separates them. By the Stone--Weierstrass theorem, \(C_F(G)\) is
dense in \(C(G)\). Put
\(I(f)=\int f_0\,d(\mu_{\alpha_1}\hat\times\cdots\hat\times\mu_{\alpha_n})\).
The value is independent of how \(f\) is written: by (4.1) the order of
the factors is irrelevant, and a coordinate on which \(f_0\) does not
depend contributes a factor \(\mu_\beta(G_\beta)=1\). So \(I\) is linear
and positive on \(C_F(G)\), \(|I(f)|\leq\|f\|_{\sup}\), \(I(1)=1\), and
\(I\) is left invariant because each finite product measure is a left
Haar measure (Proposition 12.1(1)). It extends uniquely to a positive,
left-invariant, norm-one functional on \(C(G)\) (positivity passes to
the limit because \(\operatorname{Re}f_n+\|f-f_n\|_{\sup}\geq0\) when
\(f\geq0\)). Its Radon measure is the Haar measure of \(G\) with total
mass \(1\).

For \(G=(\mathbb Z_2)^{\mathbb N}\), each factor with mass \(\frac12\)
on each point, let \(\Phi(a)=\sum_ja_j2^{-j}\in[0,1]\). This map is
continuous and onto, and it is one-to-one except over the dyadic
rationals \(j2^{-k}\) in \((0,1)\), each of which has two preimages. It
carries Haar measure to Lebesgue measure \(\lambda_1\) on \([0,1]\):
\(\lambda_1(E)=\mu(\Phi^{-1}(E))\) for every Borel \(E\). Proof:
\(\Phi_*\mu\) is a finite Borel measure on \([0,1]\). It is outer
regular: by (R2) for \(\mu\) there is an open \(O\supseteq\Phi^{-1}(E)\)
with \(\mu(O\setminus\Phi^{-1}(E))<\varepsilon\); then
\(U=[0,1]\setminus\Phi(G\setminus O)\) is open (\(\Phi\) is a closed
map, since \(G\) is compact), contains \(E\), and satisfies
\(\Phi^{-1}(U)\subseteq O\), so \(\Phi_*\mu(U\setminus E)<\varepsilon\).
A finite outer regular measure on a compact space is also inner regular
on open sets (take complements), so \(\Phi_*\mu\) is Radon. For
continuous \(f\) on \([0,1]\), the cylinder sets fixing
\(a_1,\dots,a_k\) have measure \(2^{-k}\) and are mapped onto the
intervals \([j2^{-k},(j+1)2^{-k}]\), so \(\int f\circ\Phi\,d\mu\)
differs from the Riemann sum \(\sum_j2^{-k}f(j2^{-k})\) by at most the
oscillation of \(f\) on intervals of length \(2^{-k}\). Letting
\(k\to\infty\), \(\int f\,d\Phi_*\mu=\int_0^1f\,dx\). By the uniqueness
in Theorem 2.2, \(\Phi_*\mu=\lambda_1\).

\subsubsection{The Heisenberg group and its volume
scaling}\label{haar-measure-on-locally-compact-groups--heisenberg-haar-volume}

\textbf{Example 12.3.} Let \(N=\mathbb R^3\) with multiplication
\[(x,y,z)(x',y',z')=(x+x',y+y',z+z'+xy').
\tag{12.3}\] This is the group of matrices
\(\begin{pmatrix}1&x&z\\0&1&y\\0&0&1\end{pmatrix}\), with the ordinary
Euclidean topology. Its identity is \((0,0,0)\) and its inverse is
\[(x,y,z)^{-1}=(-x,-y,-z+xy).\] Matrix multiplication proves
associativity, or one can check that the third coordinate of either
association is \(z+z'+z''+xy'+xy''+x'y''\). The polynomial
multiplication and inverse are continuous, so this is an LCH group.

Lebesgue measure \(dx\,dy\,dz\) is both left and right Haar measure.
Indeed, left and right multiplication by \((a,b,c)\) are respectively
\[(x,y,z)\longmapsto(a+x,b+y,c+z+ay),\qquad
(x,y,z)\longmapsto(x+a,y+b,z+c+xb).\] For \(f\in C_c(\mathbb R^3)\),
integrate in \(z\) first. Each added third-coordinate term is a
translation depending only on \(x,y\), so it leaves that inner integral
unchanged. Translate \(x\) and \(y\) next. Fubini and the one-variable
substitution rule give invariance of the full integral in both cases.
The transported Radon measures consequently agree on all Borel sets by
Riesz uniqueness. Thus \(\Delta_N=1\). The group is nonabelian: with the
commutator convention \([u,v]=uvu^{-1}v^{-1}\),
\[[(a,0,0),(0,b,0)]=(0,0,ab).\]

For \(r>0\), set \(D_r(x,y,z)=(rx,ry,r^2z)\). Substitution in (12.3)
proves \(D_r(uv)=D_r(u)D_r(v)\), and \(D_{1/r}\) is its inverse. Three
one-variable substitutions give
\[\int_N f(D_ru)\,du=r^{-4}\int_Nf(u)\,du,
\qquad \mu(D_rE)=r^4\mu(E)
\tag{12.4}\] for every Borel \(E\); the second identity follows from the
first by equality of Radon measures. The powers \(r,r,r^2\) explain why
the \textbf{homogeneous dimension} is \(4\), even though the underlying
manifold has dimension \(3\).

Fischer and Ruzhansky use symmetric coordinates \(t=z-xy/2\). The
coordinate change is a shear preserving Lebesgue measure, by integration
in the third coordinate, and converts (12.3) into
\[(x,y,t)(x',y',t')=(x+x',y+y',t+t'+(xy'-yx')/2).\] These two
descriptions therefore have the same Haar normalization. Integer matrix
coordinates give the lattice computed in
\hyperref[finite-covolume-and-arithmetic-quotients--heisenberg-lattice-volume]{the
arithmetic lesson}; in the symmetric coordinates its third coordinate is
\(k-mn/2\), rather than always an integer. This distinction prevents a
spurious factor of two in a fundamental domain.

\subsection[13. Groups that are not
σ-compact]{\texorpdfstring{\protect\hypertarget{haar-measure-on-locally-compact-groups--oa-fnd-hm-09}{}{}\protect\hypertarget{haar-measure-on-locally-compact-groups--OA-FND-HM-14}{}{}\protect\hypertarget{haar-measure-on-locally-compact-groups--oa-fnd-hm-14}{}{}13.
Groups that are not
σ-compact}{13. Groups that are not σ-compact}}\label{haar-measure-on-locally-compact-groups--OA-FND-HM-09}

When \(G\) is not \(\sigma\)-compact, its Haar measure is not
\(\sigma\)-finite, and several familiar facts need care: inner
regularity, the duality between \(L^1\) and \(L^\infty\), and
Fubini\textquotesingle s theorem. This section describes the measure
through the cosets of an open \(\sigma\)-compact subgroup, sets up the
right notion of \(L^\infty\), and ends with the examples that show what
fails.

Fix an open, closed, \(\sigma\)-compact subgroup \(G_0\) (Proposition
7.2) and a set \(Y\subseteq G\) that contains exactly one point of each
left coset of \(G_0\). So the open \(\sigma\)-compact sets \(yG_0\),
\(y\in Y\), partition \(G\), and \(\mu\) restricted to any one of them
is \(\sigma\)-finite by (R1). The restriction of \(\mu\) to the Borel
subsets of \(G_0\) is a left Haar measure of the locally compact group
\(G_0\): the restriction of a Radon measure to an open set is Radon
there, it is nonzero by Proposition 9.1, and it is invariant under left
translation by elements of \(G_0\).

\textbf{Proposition 13.1} (Cosets). Let \(E\subseteq G\) be Borel.

\begin{enumerate}
\tightlist
\item
  If \(E\) lies in countably many cosets \(y_jG_0\), then
  \(\mu(E)=\sum_j\mu(E\cap y_jG_0)\).
\item
  If \(E\) meets uncountably many cosets, then \(\mu(E)=\infty\).
\item
  Hence a Borel set is \(\sigma\)-finite exactly when it lies in
  countably many cosets. In particular every Borel set of finite measure
  does. Moreover \(G\) is \(\sigma\)-compact \(\iff\) \(\mu\) is
  \(\sigma\)-finite \(\iff\) \(Y\) is countable.
\end{enumerate}

\textbf{Proof.} (1) is countable additivity. (2) By (R2) it suffices to
show \(\mu(U)=\infty\) for every open \(U\supseteq E\). Such \(U\) meets
uncountably many cosets, and each nonempty open set \(U\cap yG_0\) has
positive measure by Proposition 9.1. For some \(n\), uncountably many of
these measures exceed \(1/n\); adding \(k\) of them gives
\(\mu(U)\geq k/n\) for every \(k\). (3) A \(\sigma\)-finite set is
covered by countably many sets of finite measure, each lying in
countably many cosets by (2); conversely each coset is
\(\sigma\)-finite. If \(Y\) is countable, countably many
\(\sigma\)-compact cosets cover \(G\); if \(G\) is \(\sigma\)-compact,
\(\mu\) is \(\sigma\)-finite by (R1); if \(\mu\) is \(\sigma\)-finite,
\(Y\) is countable by the first claim. \(\square\)

\textbf{Definition 13.2.} A set \(E\subseteq G\) is \emph{locally Borel}
when its intersection with each Borel set of finite measure is Borel,
and \emph{locally null} when, in addition, all these intersections are
null. A property holds \emph{locally almost everywhere} when the points
where it fails form a locally null set. Call \(f\colon G\to\mathbb C\)
\emph{locally measurable} when the preimage of each Borel subset of
\(\mathbb C\) is locally Borel. \(L^\infty(G)\) consists of the locally
measurable \(f\) with \(|f|\leq c\) locally almost everywhere for some
constant \(c\); two such functions are identified when they agree
locally almost everywhere; and \(\|f\|_\infty\) is the infimum of those
constants \(c\).

Why the local notions? With the ordinary definitions, \(L^\infty\) need
not be the dual of \(L^1\) when the measure is not \(\sigma\)-finite.
For Haar measure on \(\mathbb R\times\mathbb R_d\), the indicator of the
locally null set \(Y\) of Example 13.4 is a nonzero element of the
ordinary \(L^\infty\), yet it integrates every \(L^1\) function to
\(0\). The local versions repair this.

\Needspace{5\baselineskip}
\textbf{Theorem 13.3.}

\begin{enumerate}
\tightlist
\item
  \(E\) is locally Borel exactly when every \(E\cap yG_0\) is Borel, and
  locally null exactly when every \(E\cap yG_0\) is null. \(f\) is
  locally measurable exactly when every restriction \(f|_{yG_0}\) is
  Borel. A Borel set is locally null exactly when all its compact
  subsets are null. A locally null Borel set is null exactly when it is
  \(\sigma\)-finite. Countable unions of locally null sets are locally
  null, and \(|f|\leq\|f\|_\infty\) locally almost everywhere.
\item
  For \(f\in L^\infty(G)\) and \(g\in L^p(G)\), \(1\leq p<\infty\), the
  product \(fg\) is Borel, its class depends only on the classes of
  \(f\) and \(g\), and \(\|fg\|_p\leq\|f\|_\infty\|g\|_p\).
\item
  (\emph{\(L^1{}^*=L^\infty\)}) For every bounded linear functional
  \(\Phi\) on \(L^1(G)\) there is a unique \(f\in L^\infty(G)\) with
  \(\Phi(g)=\int fg\,d\mu\) for all \(g\in L^1(G)\), and
  \(\|\Phi\|=\|f\|_\infty\).
\item
  The map \(f\mapsto m_f\), \(m_f\xi=f\xi\), is an isometric
  \(*\)-isomorphism of \(L^\infty(G)\) onto a von Neumann algebra on
  \(L^2(G)\) that is maximal abelian.
\item
  (\emph{Radon--Nikodym, restricted form}) Let \(\nu\) and \(\mu'\) be
  Radon measures on an LCH space \(X\), with \(\nu\) \(\sigma\)-finite
  and \(\nu(E)=0\) whenever \(\mu'(E)=0\). Then
  \(\nu(E)=\int_Eh\,d\mu'\) for all Borel \(E\), for some Borel
  \(h\geq0\).
\end{enumerate}

\textbf{Proof.} (1) Each coset \(yG_0\) is the union of countably many
compact sets \(yK_n\), each of finite measure, so
\(E\cap yG_0=\bigcup_nE\cap yK_n\) is Borel (null) if \(E\) is locally
Borel (locally null). Conversely, a Borel \(F\) of finite measure lies
in countably many cosets \(y_jG_0\), so
\(E\cap F=\bigcup_j(E\cap y_jG_0)\cap F\) is Borel (null). The statement
about functions follows. If every compact subset of a Borel set \(E\) is
null and \(\mu(F)<\infty\), then \(\mu(E\cap F)=0\) by inner regularity
on sets of finite measure (Proposition 2.3(1)); conversely, compact sets
have finite measure. A \(\sigma\)-finite locally null Borel set is the
union of its intersections with countably many sets of finite measure,
each null. If \(E_n\) are locally null and \(\mu(F)<\infty\), then
\(\mu(\bigcup_nE_n\cap F)\leq\sum_n\mu(E_n\cap F)=0\). Finally
\(\{|f|>\|f\|_\infty\}=\bigcup_n\{|f|>\|f\|_\infty+1/n\}\) is a
countable union of locally null sets. (2) By Proposition 3.1(4),
\(\{g\neq0\}\) is a disjoint union of Borel sets \(F_n\) of finite
measure. Each \(f1_{F_n}\) is Borel, so \(fg=\sum_n(f1_{F_n})g\) is
Borel (at each point at most one term is nonzero). If \(f=f'\) locally
a.e., the set \(\{f\neq f'\}\cap\{g\neq0\}\) is Borel, \(\sigma\)-finite
and locally null, hence null by (1); so \(fg=f'g\) a.e. In the same way
\(|fg|\leq\|f\|_\infty|g|\) a.e. (3) For \(y\in Y\), the functions in
\(L^1(G)\) that vanish off \(yG_0\) form a copy of \(L^1(yG_0,\mu)\),
and \(\mu\) is \(\sigma\)-finite there. The duality \(L^1{}^*=L^\infty\)
for \(\sigma\)-finite measures provides a Borel function \(f_y\) on
\(yG_0\) with \(|f_y|\leq\|\Phi\|\) and
\(\Phi(g)=\int_{yG_0}f_yg\,d\mu\) for those \(g\). Let \(f=f_y\) on
\(yG_0\). By (1), \(f\) is locally measurable, and
\(\|f\|_\infty\leq\|\Phi\|\). Let \(g\in L^1(G)\). The set
\(\{g\neq0\}\) is \(\sigma\)-finite, so it lies in countably many cosets
\(y_jG_0\), and \(g=\sum_jg1_{y_jG_0}\) in \(L^1\) by dominated
convergence. Hence
\(\Phi(g)=\sum_j\int_{y_jG_0}fg\,d\mu=\int fg\,d\mu\), again by
dominated convergence, since \(|fg|\leq\|\Phi\||g|\). By (2) with
\(p=1\), \(\|\Phi\|\leq\|f\|_\infty\). For uniqueness, suppose
\(\int hg\,d\mu=0\) for all \(g\in L^1(G)\) but \(h\neq0\) on a set that
is not locally null. Then for some \(c>0\) and some Borel \(F\) of
finite measure, \(D=\{|h|>c\}\cap F\) has positive measure; with
\(g=1_D\bar h/|h|\) we get \(\int hg\,d\mu=\int_D|h|\,d\mu>0\), a
contradiction. (4) By (2), \(m_f\) is bounded with
\(\|m_f\|\leq\|f\|_\infty\). If \(c<\|f\|_\infty\), the set
\(\{|f|>c\}\) is not locally null, so for some Borel \(F\) of finite
measure the set \(D=\{|f|>c\}\cap F\) has \(\mu(D)>0\); then
\(\|m_f1_D\|_2^2\geq c^2\mu(D)=c^2\|1_D\|_2^2\). So \(m\) is isometric.
Clearly \(m_{fg}=m_fm_g\) and \(m_{\bar f}=m_f^*\). Let \(T\) commute
with every \(m_f\). Then \(T\) commutes with the projections
\(P_y=m_{1_{yG_0}}\), so it maps each
\(H_y=P_yL^2(G)\cong L^2(yG_0,\mu)\) into itself. On \(H_y\) it commutes
with multiplication by every bounded Borel function on \(yG_0\)
(extended by \(0\), such a function lies in \(L^\infty(G)\)). On a
\(\sigma\)-finite measure space the multiplication operators form a
maximal abelian algebra
(\hyperref[hilbert-spaces-and-compact-operators]{Hilbert spaces and
compact operators}, Theorem 9.1). Applied to the \(\sigma\)-finite
measure space \(yG_0\), this shows that \(T|_{H_y}\) is multiplication
by a bounded Borel function \(f_y\) on \(yG_0\), which we may take with
\(|f_y|\leq\|T\|\) by the norm formula just proved. Let \(f=f_y\) on
\(yG_0\); then \(f\in L^\infty(G)\). A vector \(\xi\in L^2(G)\) lives on
countably many cosets (Proposition 3.1(4) and Proposition 13.1(3)), so
\(\xi=\sum_jP_{y_j}\xi\) and \(T\xi=\sum_jf_{y_j}P_{y_j}\xi=f\xi\). So
\(T=m_f\). Thus \(m(L^\infty)'\subseteq m(L^\infty)\), and the reverse
inclusion holds because \(m(L^\infty)\) is abelian. Hence
\(m(L^\infty)'=m(L^\infty)\), which is maximal abelian, and
\(m(L^\infty)''=m(L^\infty)\). (5) By Proposition 2.3(3), \(X=S\cup N\)
with \(S\) \(\sigma\)-compact and \(\nu(N)=0\). On the Borel subsets of
\(S\), \(\mu'\) is \(\sigma\)-finite by (R1), and the Radon--Nikodym
theorem for \(\sigma\)-finite measures gives a Borel \(h_0\geq0\) on
\(S\) with \(\nu(E)=\int_Eh_0\,d\mu'\) for Borel \(E\subseteq S\). Put
\(h=h_0\) on \(S\) and \(h=0\) on \(N\). \(\square\)

Part (4) carries the maximal abelian property of \(L^\infty\) from
\(\sigma\)-finite measure spaces to Haar measure on every locally
compact group, and no lifting theorem is needed.

The following examples show what fails without \(\sigma\)-finiteness. In
each of them the bad set is locally null, so it carries no \(L^p\) mass
for \(p<\infty\).

\textbf{Example 13.4} (\(\mathbb R\times\mathbb R_d\): a Haar measure
that is not \(\sigma\)-finite). Let \(\mathbb R_d\) be the additive
group of reals with the discrete topology, and
\(G=\mathbb R\times\mathbb R_d\). The subgroup
\(G_0=\mathbb R\times\{0\}\) is open and \(\sigma\)-compact, and its
cosets are the lines \(\mathbb R\times\{d\}\). A compact subset of \(G\)
meets only finitely many lines, so
\(I(f)=\sum_d\int_{\mathbb R}f(x,d)\,dx\) is a finite sum for
\(f\in C_c(G)\); it is a left Haar integral, by translation invariance
of the Riemann integral. On each line, the Haar measure \(\mu\) is
Lebesgue measure: the restriction of a Radon measure to an open set is
Radon on that set, and both measures integrate the compactly supported
continuous functions on the open set \(\mathbb R\times\{d\}\) alike. By
Proposition 13.1, \(\mu(E)=\sum_d\lambda_1(\{x:(x,d)\in E\})\) if \(E\)
meets countably many lines, and \(\mu(E)=\infty\) otherwise.

The closed set \(Y=\{0\}\times\mathbb R_d\) meets every line, so
\(\mu(Y)=\infty\). Its compact subsets are finite and null. So \(\mu\)
is not inner regular on \(Y\), and \(Y\) is locally null (Theorem
13.3(1)) but not null. No measure that is inner and outer regular on
every Borel set can represent the Haar integral of \(G\): it would be
Radon, hence equal to \(\mu\) by the uniqueness in Theorem 2.2, and
\(\mu\) fails inner regularity on \(Y\). Nevertheless every
\(f\in L^p(G)\), \(p<\infty\), lives on countably many lines; \(C_c(G)\)
is dense in \(L^p(G)\) (Proposition 3.1(4)); and \(L^\infty(G)\)
consists of the functions whose restrictions to the lines are Borel and
essentially bounded, uniformly in \(d\); it is the dual of \(L^1(G)\)
(Theorem 13.3(3)).

\textbf{Example 13.5} (The diagonal of \(\mathbb R_d\times\mathbb R\):
Tonelli\textquotesingle s theorem needs \(\sigma\)-finiteness). Let
\(X=\mathbb R_d\) with counting measure \(c\) and \(Y=\mathbb R\) with
Lebesgue measure \(m\). Their Radon product is the Haar measure of the
group \(\mathbb R_d\times\mathbb R\) (Proposition 12.1(1)). The diagonal
\(D=\{(t,t)\}\) is a closed subgroup. Its sections are single points:
\(m(D_s)=0\) for every \(s\), while \(c(D^t)=1\) for every \(t\). So
\[\int m(D_s)\,dc(s)=0,\qquad\int c(D^t)\,dm(t)=\infty,\qquad(c\hat\times m)(D)=\infty,\]
the last because \(D\) meets every coset \(\{s\}\times\mathbb R\) of the
open subgroup \(\{0\}\times\mathbb R\) (Proposition 13.1(2)). Thus
Tonelli\textquotesingle s theorem 5.1(4) fails for \(1_D\), whose
support is not \(\sigma\)-finite. The compact subsets of \(D\) are
finite, and each point has measure \(c(\{s\})m(\{s\})=0\) by (4.2); so
\(D\) is locally null (Theorem 13.3(1)). The example also shows that the
rectangle function \(A\times B\mapsto c(A)m(B)\) has no extension to a
measure \(m'\) on \(\mathcal B(X\times Y)\) that is inner and outer
regular on every Borel set. Such an \(m'\) would be Radon. It would
integrate \(\varphi\boxtimes g\) to \(\int\varphi\,dc\int g\,dm\) for
\(\varphi\in C_c(X)\) and \(g\in C_c(Y)\), by uniform approximation with
simple functions on the compact rectangle
\(\operatorname{supp}\varphi\times\operatorname{supp}g\), which has
finite \(m'\)-measure. By the tensor approximation (Lemma 4.1) it would
then agree with \(c\hat\times m\) on \(C_c(X\times Y)\), hence equal it
by the uniqueness in Theorem 2.2. But inner regularity would force
\(m'(D)=0\).

\textbf{Example 13.6} (\(\{e\}\times G\): the rectangle rule needs
\(\sigma\)-finite factors). Let \(G=\mathbb R\times\mathbb R_d\) and
\(\pi=\mu\hat\times\mu\). The closed set \(\{e\}\times G\) has
\(\mu(\{e\})=0\) but \(\pi(\{e\}\times G)=\infty\). Indeed \(\pi\) is a
Haar measure on \(G\times G\) (Proposition 12.1(1)), \(G_0\times G_0\)
is an open \(\sigma\)-compact subgroup of \(G\times G\), and
\(\{e\}\times G\) meets each of the uncountably many cosets
\(G_0\times(\mathbb R\times\{d\})\); apply Proposition 13.1(2). So
Theorem 5.1(6) fails when \(B\) is not \(\sigma\)-finite, and
"\(\pi(A\times B)=\mu(A)\mu(B)\) with \(0\cdot\infty=0\)" is false for
the Radon product. The set is locally null by Theorem 13.3(1): its
intersection with a compact \(K\subseteq G\times G\) lies in
\(\{e\}\times\mathrm{pr}_2(K)\), which is null by (4.2). So, like \(Y\)
and \(D\), it carries no \(L^2\) mass.

\subsubsection{Two regularity conventions for Haar
measure}\label{haar-measure-on-locally-compact-groups--haar-measure-convention-bridge}

Fremlin\textquotesingle s \emph{Measure Theory} uses a complete, locally
determined measure that is inner regular on every measurable set. Our
Borel measure is outer regular on every Borel set and inner regular on
open sets. On a sigma-compact group these conventions agree after
completion. On an arbitrary group, the following construction explains
the difference and the common integration theory.

\textbf{Proposition 13.7} (The locally determined measure). With the
open sigma-compact cosets \(yG_0\) used above, define for Borel
\(E\subset G\) \[m(E)=\sum_{y\in Y}\mu(E\cap yG_0),
\qquad
\sum_{y\in Y}a_y:=\sup_{F\subset Y\text{ finite}}\sum_{y\in F}a_y.\]
Then \(m\) is left invariant, locally finite and inner regular on every
Borel set. It agrees with \(\mu\) on open sets, compact sets and
sigma-finite Borel sets, and integrates every \(C_c(G)\) function in the
same way. There is a canonical isometric identification
\(L^p(\mu)\cong L^p(m)\) for \(1\le p<\infty\). Nevertheless \(m\) can
fail outer regularity on Borel sets. Completing each coset measure and
allowing sets measurable on every coset gives a Radon measure in
Fremlin\textquotesingle s convention.

\emph{Proof.} A sum of nonnegative numbers is the supremum of its finite
subsums. The two suprema in a sum over \(Y\times\mathbb N\) commute:
each finite collection of pairs is contained in the product of finite
subsets, and the converse products are finite. Applying countable
additivity on each coset therefore proves countable additivity of \(m\).
Left translation permutes the cosets and preserves \(\mu\) on their
Borel subsets, so it preserves \(m\).

A compact set meets only finitely many cosets, since the cosets form an
open cover. Thus \(m\) and \(\mu\) agree on compact sets. They agree on
each sigma-finite Borel set by Proposition 13.1, since such a set lies
in countably many cosets. Local finiteness follows by taking a
relatively compact neighborhood inside a single coset.

The restriction of \(\mu\) to each coset is sigma-finite and inner
regular on every Borel set by Proposition 2.3. Given \(a<m(E)\), with
\(a\ge0\), choose a finite set of cosets whose contributions exceed
\(a\). In each of those cosets choose compact subsets of \(E\) whose
measures still have sum exceeding \(a\). Their finite union is compact.
This proves inner regularity of \(m\) on every Borel set. If \(U\) is
open, inner regularity of both measures on \(U\), and their equality on
compact sets, give \(m(U)=\mu(U)\). A \(C_c\) function is supported on a
compact set, so its two integrals agree.

The identity on Borel representatives maps \(L^p(\mu)\) isometrically
into \(L^p(m)\): a \(\mu\)-integrable function is supported on a
sigma-finite set, where the measures agree. For surjectivity, take a
Borel \(f\) with \(\int|f|^p\,dm<\infty\). The sum of its coset
integrals is finite, so only countably many of those integrals are
nonzero. Set \(f\) to zero outside those cosets. The resulting Borel
function \(f_0\) equals \(f\) \(m\)-almost everywhere, is supported on a
sigma-finite set and has the same norm for \(\mu\). This also proves
independence of representatives and surjectivity.

For the last assertion, let \(\Sigma_{\mathrm{loc}}\) consist of the
subsets whose intersections with every coset are measurable for its
completed sigma-finite measure. Define the measure by the same sum. It
is complete because a subset of a null set is null on every coset. It is
locally determined: if a set has measurable intersections with every
measurable set of finite measure, intersect it with a countable compact
exhaustion of each coset to see that its coset sections are measurable.
Completion preserves compact inner approximation on each coset: a
completed measurable set contains a Borel subset of the same measure,
and the latter has compact approximants. The preceding finite-subsum
argument now gives compact inner regularity on all of
\(\Sigma_{\mathrm{loc}}\). These properties, local finiteness and the
Hausdorff topology give Fremlin\textquotesingle s Radon convention.

Finally, in Example 13.4 the closed set \(Z=\{0\}\times\mathbb R_d\)
satisfies \(m(Z)=0\), while every open neighborhood of \(Z\) has
infinite \(m\)-measure because \(m\) and \(\mu\) agree on open sets.
Thus \(m\) is not outer regular there. \(\square\)

The finite-\(p\) spaces for \(\Sigma_{\mathrm{loc}}\) have the same
identification: an integrable function has nonzero coset integrals on
only countably many cosets. Choose a Borel representative on each of
those cosets and set it to zero elsewhere; this gives a global Borel
representative of its class, with the same integral.

The construction also identifies the local \(L^\infty\) of Definition
13.2 with the \(L^\infty\) of the completed, locally determined measure:
both take essentially bounded measurable functions on each coset, with a
common bound, modulo functions null on every coset. Completed coset
functions have Borel representatives there, so the completed version
gives the same classes. This explains why the duality in Theorem 13.3
agrees with Fremlin\textquotesingle s localizable-measure duality,
despite the difference between the two global Borel measures.

\subsection[14. Convolution and continuity of
translation]{\texorpdfstring{\protect\hypertarget{haar-measure-on-locally-compact-groups--oa-fnd-hm-12}{}{}14.
Convolution and continuity of
translation}{14. Convolution and continuity of translation}}\label{haar-measure-on-locally-compact-groups--OA-FND-HM-12}

In this section and the next, \(\pi=\mu\hat\times\mu\) on \(G\times G\),
and integrals over \(G\) are taken with respect to \(\mu\). Convolution
integrals are iterated integrals over \(G\times G\) in disguise, so we
first record how \(\pi\) behaves under the changes of variables that
occur.

\textbf{Lemma 14.1} (Changes of variables on \(G\times G\)). The maps
\[\Theta_1(x,y)=(y,y^{-1}x),\qquad\Theta_2(x,y)=(y,xy^{-1}),\qquad\Theta_3(x,y)=(y,xy),\qquad\sigma(x,y)=(y,x)\]
are homeomorphisms of \(G\times G\) with
\((\Theta_1)_*\pi=\sigma_*\pi=\pi\),
\((\Theta_2)_*\pi=(\Delta\circ\mathrm{pr}_1)\,\pi\) and
\((\Theta_3)_*\pi=(\Delta^{-1}\circ\mathrm{pr}_1)\,\pi\). Consequently
each of them pulls \(\pi\)-null sets back to \(\pi\)-null sets and
\(\sigma\)-finite sets back to \(\sigma\)-finite sets, and for every
Borel \(H\geq0\) on \(G\times G\)
\[\begin{gathered}
\int H\circ\Theta_1\,d\pi=\int H\,d\pi,\qquad\\
\int H\circ\Theta_2\,d\pi=\int H(y,z)\Delta(y)\,d\pi(y,z),\qquad\\
\int H\circ\Theta_3\,d\pi=\int H(y,z)\Delta(y)^{-1}\,d\pi(y,z).
\end{gathered}\]

\textbf{Proof.} The inverses are \((y,z)\mapsto(yz,y)\), \((zy,y)\) and
\((zy^{-1},y)\), all continuous. For \(F\in C_c(G\times G)\), (4.1) lets
us integrate first in \(x\); left invariance and (10.1) in the inner
integral give
\[\begin{gathered}
\int F\circ\Theta_1\,d\pi=\int\!\!\int F(y,y^{-1}x)\,dx\,dy=\int F\,d\pi,\qquad\\
\int F\circ\Theta_2\,d\pi=\int\!\!\int F(y,xy^{-1})\,dx\,dy=\int\Delta(y)\!\int F(y,x)\,dx\,dy,
\end{gathered}\]
and similarly
\(\int F\circ\Theta_3\,d\pi=\int\Delta(y)^{-1}\int F(y,x)\,dx\,dy\) and
\(\int F\circ\sigma\,d\pi=\int F\,d\pi\). The image measures are Radon
by Proposition 3.1(1), and so are
\((\Delta^{\pm1}\circ\mathrm{pr}_1)\pi\) by Proposition 3.1(2). The
uniqueness in Theorem 2.2 gives the four identities of measures.
Densities that are continuous and positive do not change null sets or
\(\sigma\)-finite sets (Proposition 3.1(2)). The integral formulas
follow from Proposition 3.1(1) and (2). \(\square\)

For Borel functions \(f,g\) on \(G\), put
\[f*g(x)=\int f(y)\,g(y^{-1}x)\,dy
\tag{14.1}\] at every \(x\) where the integral converges absolutely. For
\(f,g\geq0\) the integral is defined at every \(x\), with values in
\([0,\infty]\).

\Needspace{5\baselineskip}
\textbf{Theorem 14.2.}

\begin{enumerate}
\tightlist
\item
  (\emph{Tonelli for convolutions}) If \(f,g\geq0\) are Borel and vanish
  outside \(\sigma\)-finite sets, then \(f*g\) is \(\mu\)-a.e. Borel and
  \(\int f*g\,dx=\int f\,dx\int g\,dx\).
\item
  For \(f,g\in L^1(G)\), (14.1) converges absolutely for almost every
  \(x\), \(\|f*g\|_1\leq\|f\|_1\|g\|_1\), and the class of \(f*g\)
  depends only on the classes of \(f\) and \(g\). If one of the
  following integrals converges absolutely, so do the others, and
  \[f*g(x)=\int f(xy)g(y^{-1})\,dy=\int f(y^{-1})g(yx)\Delta(y^{-1})\,dy=\int f(xy^{-1})g(y)\Delta(y^{-1})\,dy.
  \tag{14.2}\]
\item
  With the involution \(f^*(x)=\Delta(x)^{-1}\overline{f(x^{-1})}\),
  \(L^1(G)\) is a Banach \(*\)-algebra: convolution is associative,
  \(\|f^*\|_1=\|f\|_1\), \(f^{**}=f\) and \((f*g)^*=g^**f^*\). Moreover
  \(L_z(f*g)=(L_zf)*g\) and \(R_z(f*g)=f*(R_zg)\).
\item
  Let \(1\leq p<\infty\), \(f\in L^1(G)\) and \(g\in L^p(G)\). Then
  \(f*g(x)\) converges absolutely for almost every \(x\), and
  \(\|f*g\|_p\leq\|f\|_1\|g\|_p\). If also
  \(\int|f|\Delta^{-1}\,dx<\infty\) (for instance if \(f\) has compact
  support, or if \(G\) is unimodular), then
  \(g*f(x)=\int g(xy^{-1})f(y)\Delta(y^{-1})\,dy\) converges absolutely
  for almost every \(x\), and
  \[\|g*f\|_p\leq\|\Delta^{-1}f\|_1^{1-1/p}\,\|f\|_1^{1/p}\,\|g\|_p .
  \tag{14.3}\] For \(p=1\) the condition on \(f\) is not needed, and
  \(\|g*f\|_1\leq\|g\|_1\|f\|_1\).
\item
  If \(G\) is unimodular, \(1<p<\infty\), \(\frac1p+\frac1q=1\),
  \(f\in L^p(G)\) and \(g\in L^q(G)\), then \(f*g(x)\) converges
  absolutely for every \(x\), \(f*g\in C_0(G)\), and
  \(\|f*g\|_{\sup}\leq\|f\|_p\|g\|_q\).
\item
  (\emph{Continuity of translation}) For \(1\leq p<\infty\) and
  \(f\in L^p(G)\), \(\|L_yf-f\|_p\to0\) and \(\|R_yf-f\|_p\to0\) as
  \(y\to e\).
\item
  Let \(f\in L^1(G)\) and \(g\in L^\infty(G)\) (Section 13). Then
  \(f*g(x)\) is defined for every \(x\),
  \(|f*g|\leq\|f\|_1\|g\|_\infty\), and \(\|L_z(f*g)-f*g\|_{\sup}\to0\)
  as \(z\to e\). If also \(\Delta^{-1}f\in L^1(G)\), then
  \(g*f(x)=\int g(y)f(y^{-1}x)\,dy\) is defined for every \(x\),
  \(|g*f|\leq\|g\|_\infty\|\Delta^{-1}f\|_1\), and
  \(\|R_z(g*f)-g*f\|_{\sup}\to0\). Without that condition \(g*f\) can be
  infinite everywhere (Exercise 16.4).
\end{enumerate}

\emph{Scope warning:} The second half of (7) cannot be asserted for
every \(f\in L^1(G)\), with the bound \(\|g\|_\infty\|R_zf-f\|_1\); that
version fails on every group that is not unimodular (Remark 14.3 and
Exercise 16.4).

\textbf{Proof.} (1) The function \((x,y)\mapsto f(y)g(y^{-1}x)\) is
\((f\boxtimes g)\circ\Theta_1\). The function \(f\boxtimes g\) vanishes
outside a product of \(\sigma\)-finite sets, which is \(\sigma\)-finite
by Theorem 5.1(6); so \((f\boxtimes g)\circ\Theta_1\) vanishes outside a
\(\sigma\)-finite set by Lemma 14.1. By Tonelli\textquotesingle s
theorem 5.1(4) and Lemma 14.1,
\(\int\big(\int f(y)g(y^{-1}x)\,dy\big)dx=\int(f\boxtimes g)\circ\Theta_1\,d\pi=\int f\boxtimes g\,d\pi=\int f\int g\).
(2) Apply (1) to \(|f|\) and \(|g|\): the inner integral is finite for
almost every \(x\), and \(|f*g|\leq|f|*|g|\) gives the bound. By
Fubini\textquotesingle s theorem 5.1(5), \(f*g\) is \(\mu\)-a.e. Borel.
Changing \(g\) on a null set \(N\) changes the integrand only on
\(\Theta_1^{-1}(\{f\neq0\}\times N)\), which is \(\pi\)-null by Theorem
5.1(6) and Lemma 14.1; by Theorem 5.1(3) this affects, for almost every
\(x\), only a null set of \(y\). The same holds for \(f\). For (14.2),
fix \(x\). The substitution \(y\mapsto xy\), allowed by left invariance
(Proposition 8.2(2)), turns (14.1) into the first form. The inversion
formula (11.1), applied to \(h(y)=f(y^{-1})g(yx)\) and to
\(h(y)=f(xy^{-1})g(y)\), turns (14.1) and the first form into the second
and third forms. Each step preserves absolute convergence. (3) For
\(f,g\in C_c(G)\), \(f*g\) vanishes off
\((\operatorname{supp}f)(\operatorname{supp}g)\), and
\(|f*g(xz)-f*g(x)|\leq\|f\|_1\|R_zg-g\|_{\sup}\to0\) as \(z\to e\) by
uniform continuity (Proposition 7.3), so \(f*g\in C_c(G)\). For
\(f,g,h\in C_c(G)\) and fixed \(x\), the substitution
\(w\mapsto y^{-1}w\) in the inner integral gives
\[f*(g*h)(x)=\int f(y)\Big(\int g(y^{-1}z)h(z^{-1}x)\,dz\Big)dy,\qquad (f*g)*h(x)=\int\Big(\int f(y)g(y^{-1}z)\,dy\Big)h(z^{-1}x)\,dz ,\]
and the integrand is in \(C_c(G\times G)\), so (4.1) makes them equal.
For \(f,g\in C_c(G)\), a direct computation with
\(\Delta(y)^{-1}\Delta(y^{-1}x)^{-1}=\Delta(x)^{-1}\) and the
substitution \(y\mapsto xy\) gives \((f*g)^*=g^**f^*\). By the last
formula of Theorem 11.1,
\(\|f^*\|_1=\int\Delta(x)^{-1}|f(x^{-1})|\,dx=\|f\|_1\); and
\(f^{**}(x)=\Delta(x)^{-1}\Delta(x^{-1})^{-1}f(x)=f(x)\). By (2),
convolution is a bounded bilinear map on \(L^1\), and the involution is
a conjugate-linear isometry; since \(C_c(G)\) is dense (Proposition
3.1(4)), the identities pass to \(L^1(G)\), which is complete. The
translation rules follow from the substitution \(y\mapsto zy\) in
(14.1), and directly for \(R_z\). (4) The case \(p=1\) of the first
claim is (2). Let \(1<p<\infty\) and \(\frac1p+\frac1q=1\). For fixed
\(x\), Hölder\textquotesingle s inequality for the measure
\(|f(y)|\,dy\) gives
\[\Big(\int|f(y)||g(y^{-1}x)|\,dy\Big)^p\leq\|f\|_1^{p-1}\int|f(y)||g(y^{-1}x)|^p\,dy .\]
By (1) for \(|f|\) and \(|g|^p\in L^1\), the right side has integral
\(\|f\|_1^{p-1}\|f\|_1\|g\|_p^p\) over \(x\). This gives absolute
convergence almost everywhere and the bound. For \(g*f\), put
\(w=|f|\Delta^{-1}\). The function \((x,y)\mapsto|g(xy^{-1})|^pw(y)\) is
\((w\boxtimes|g|^p)\circ\Theta_2\), so by Tonelli\textquotesingle s
theorem and Lemma 14.1,
\[\int\Big(\int|g(xy^{-1})|^pw(y)\,dy\Big)dx=\int w(y)\Delta(y)\,dy\int|g|^p=\|f\|_1\|g\|_p^p .\]
For \(p=1\) this is the claim. For \(p>1\), Hölder\textquotesingle s
inequality for the measure \(w(y)\,dy\) of mass \(\|\Delta^{-1}f\|_1\)
gives
\(\big(\int|g(xy^{-1})|w(y)\,dy\big)^p\leq\|\Delta^{-1}f\|_1^{p-1}\int|g(xy^{-1})|^pw(y)\,dy\),
and integrating in \(x\) gives (14.3). If \(f\) vanishes off a compact
set \(K\), then \(\|\Delta^{-1}f\|_1\leq(\sup_K\Delta^{-1})\|f\|_1\),
and (14.3) gives
\(\|g*f\|_p\leq(\sup_K\Delta^{(1/p)-1})\|f\|_1\|g\|_p\). (5) For every
\(x\), Hölder\textquotesingle s inequality gives
\(|f*g(x)|\leq\|f\|_p\big(\int|g(y^{-1}x)|^q\,dy\big)^{1/q}\). By (11.1)
and unimodularity,
\(\int|g(y^{-1}x)|^q\,dy=\int|g(yx)|^q\Delta(y)^{-1}\,dy=\|g\|_q^q\).
For \(f,g\in C_c(G)\), \(f*g\in C_c(G)\) by (3). For general \(f,g\),
take \(f_n,g_n\in C_c(G)\) with \(f_n\to f\) in \(L^p\) and \(g_n\to g\)
in \(L^q\) (Proposition 3.1(4)). The bound gives \(f_n*g_n\to f*g\)
uniformly, and a uniform limit of functions in \(C_c(G)\) lies in
\(C_0(G)\). The endpoint cases fail on every noncompact \(G\): for
\(f\geq0\) with \(\int f=1\) and \(g\equiv1\),
\(f*g\equiv1\notin C_0(G)\). (6) Fix a compact neighbourhood \(V\) of
\(e\). For \(g\in C_c(G)\) and \(y\in V\), \(L_yg\) and \(R_yg\) vanish
outside the compact set
\(K=V(\operatorname{supp}g)\cup(\operatorname{supp}g)V^{-1}\), so
\(\|L_yg-g\|_p\leq\mu(K)^{1/p}\|L_yg-g\|_{\sup}\to0\) by uniform
continuity (Proposition 7.3), and the same for \(R_y\). For \(f\in L^p\)
and \(\varepsilon>0\), choose \(g\in C_c(G)\) with
\(\|f-g\|_p<\varepsilon\). Since \(\|L_y\|=1\), and
\(\|R_y\|=\Delta(y)^{-1/p}\leq C\) on \(V\) by Theorem 10.1(3) and the
continuity of \(\Delta\), we get
\(\|R_yf-f\|_p\leq(C+1)\varepsilon+\|R_yg-g\|_p\), and similarly for
\(L_y\). (7) Left and right translations and inversion map Borel sets to
Borel sets and preserve null sets and \(\sigma\)-finite sets (Theorem
10.1(1) and Corollary 11.2(1)). Hence they map locally Borel and locally
null sets to sets of the same kind: for such a map \(\theta\) and
\(\mu(F)<\infty\), the set \(\theta^{-1}(F)\) is covered by countably
many sets of finite measure, so
\(\theta(E)\cap F=\theta\big(E\cap\theta^{-1}(F)\big)\) is Borel (null)
when \(E\) is locally Borel (locally null). The map \(y\mapsto y^{-1}x\)
and its inverse \(z\mapsto xz^{-1}\) are composites of an inversion and
a translation, so \(y\mapsto g(y^{-1}x)\) is locally measurable and
bounded by \(\|g\|_\infty\) locally almost everywhere. As in the proof
of Theorem 13.3(2), multiplying by \(f\) gives an integrable function
with \(|f*g(x)|\leq\|f\|_1\|g\|_\infty\). By (3),
\(L_z(f*g)-f*g=(L_zf-f)*g\), and \(\|L_zf-f\|_1\to0\) by (6). For
\(g*f\), the substitutions of (2) give
\(\int|g(y)||f(y^{-1}x)|\,dy\leq\|g\|_\infty\int|f(y^{-1}x)|\,dy=\|g\|_\infty\int|f(y)|\Delta(y)^{-1}\,dy\).
Also \(R_z(g*f)-g*f=g*(R_zf-f)\), and
\(\Delta^{-1}R_zf=\Delta(z)R_z(\Delta^{-1}f)\), so
\[\|\Delta^{-1}(R_zf-f)\|_1\leq|\Delta(z)-1|\,\|R_z(\Delta^{-1}f)\|_1+\|R_z(\Delta^{-1}f)-\Delta^{-1}f\|_1\to0\]
by (6) applied to \(\Delta^{-1}f\in L^1\). \(\square\)

\textbf{Remark 14.3} (The factor \(\Delta^{-1}\) in (7) is needed). In
the second half of (7) the bound is
\(|R_z(g*f)-g*f|\leq\|g\|_\infty\|\Delta^{-1}(R_zf-f)\|_1\), with
\(\Delta^{-1}(R_zf-f)\) where one might expect \(R_zf-f\). This cannot
be improved. For \(h\in L^1\) and fixed \(x\), the third form of (14.2)
gives \(g*h(x)=\int g(xy^{-1})h(y)\Delta(y)^{-1}\,dy\). The choice
\(g(w)=\overline{\operatorname{sgn}h(w^{-1}x)}\) shows that the supremum
of \(|g*h(x)|\) over \(\|g\|_\infty\leq1\) is \(\|\Delta^{-1}h\|_1\),
which exceeds \(\|h\|_1\) whenever \(h\neq0\) lives where \(\Delta<1\).
Such \(h=R_zf-f\) occur on every group that is not unimodular: take
\(f\in C_c^+(G)\) supported in the nonempty open set \(\{\Delta<1\}\),
and \(z\) with \(\Delta(z)>1\). Then \(R_zf\) is supported in
\((\operatorname{supp}f)z^{-1}\), where \(\Delta\) is smaller still, and
\(R_zf\neq f\) because \(\int R_zf=\Delta(z)^{-1}\int f\). So the bound
\(\|g\|_\infty\|R_zf-f\|_1\) fails for this \(f\), this \(z\) and a
suitable \(g\). The extra hypothesis \(\Delta^{-1}f\in L^1\) cannot be
dropped either: Exercise 16.4 gives \(f\in L^1\) on the \(ax+b\) group
with \(1*f\equiv\infty\). For unimodular groups \(\Delta\equiv1\), and
nothing changes. The first half of (7), and (4) for \(f*g\), hold for
all groups.

\subsubsection{The full Young inequality on a unimodular
group}\label{haar-measure-on-locally-compact-groups--young-unimodular}

\textbf{Theorem 14.4 (Young).} Suppose \(G\) is unimodular,
\(1\leq p,q,r\leq\infty\), and \[\frac1p+\frac1q=1+\frac1r,\] where
\(1/\infty=0\). For \(f\in L^p(G)\) and \(g\in L^q(G)\), convolution
defines an element of \(L^r(G)\), with \[\|f*g\|_r\leq\|f\|_p\|g\|_q.
\tag{14.4}\] It converges absolutely almost everywhere if \(r<\infty\),
and at every point if \(r=\infty\), using locally bounded
representatives at the \(L^\infty\) endpoints. Fischer and Ruzhansky
state this general inequality in their freely accessible monograph. Here
is a direct proof, requiring no interpolation theorem.

\emph{Proof.} If \(p=1\), the relation gives \(q=r\), and Theorem
14.2(4),(7) gives the result. If \(q=1\), then \(p=r\); the same theorem
applies with its modular factors equal to one. These include \(r=1\) and
all cases with an infinite input exponent. If \(r=\infty\) and
\(1<p,q<\infty\), the exponents are conjugate. Hölder and unimodular
inversion give, for every \(x\), \[\int_G|f(y)g(y^{-1}x)|\,dy
\leq\|f\|_p\left(\int_G|g(y^{-1}x)|^q\,dy\right)^{1/q}
=\|f\|_p\|g\|_q.\] These cases are unchanged by alterations on null
sets, since translations and inversion preserve null sets.

It remains to consider \(1<p,q<r<\infty\). The exponent relation implies
\(p<r\) and \(q<r\). Apply Hölder with the three exponents
\[r,\qquad \frac{p}{1-p/r},\qquad \frac{q}{1-q/r},\] whose reciprocals
sum to one, to the product \[\bigl(|f(y)|^p|g(y^{-1}x)|^q\bigr)^{1/r}
|f(y)|^{1-p/r}|g(y^{-1}x)|^{1-q/r}.\] This gives \[|f*g(x)|^r
\leq \|f\|_p^{r-p}\|g\|_q^{r-q}
\int_G|f(y)|^p|g(y^{-1}x)|^q\,dy.\] Integrating in \(x\), left
translation gives \[\int_G\int_G|f(y)|^p|g(y^{-1}x)|^q\,dy\,dx
=\|f\|_p^p\|g\|_q^q.\] These interchanges do not impose sigma-finiteness
on \(G\). The nonzero supports of \(f\) and \(g\) are sigma-finite
because their finite powers are integrable. Their product is
sigma-finite by Theorem 5.1(6), and its inverse image under \(\Theta_1\)
is sigma-finite by Lemma 14.1. The sigma-finite Tonelli theorem
therefore applies to the displayed integrand, and Fubini gives a
measurable convolution where it is absolutely convergent. The finite
double integral also makes the inner integral finite almost everywhere.
The three-factor estimate then proves absolute convergence there and
(14.4). If either input norm is zero, apply the same Tonelli calculation
to obtain the zero convolution class. \(\square\)

The unimodular hypothesis belongs to this symmetric form of the
inequality. The preceding weighted estimates remain the applicable
bounds on a general group; the affine counterexample in Exercise 16.4
explains why their modular factors cannot be discarded.

\subsection[15. Approximate
identities]{\texorpdfstring{\protect\hypertarget{haar-measure-on-locally-compact-groups--oa-fnd-hm-13}{}{}15.
Approximate
identities}{15. Approximate identities}}\label{haar-measure-on-locally-compact-groups--OA-FND-HM-13}

\textbf{Theorem 15.1.} Let the neighbourhoods \(U\) of \(e\) be directed
by reverse inclusion.

\begin{enumerate}
\tightlist
\item
  For every neighbourhood \(U\) of \(e\) there is \(\psi_U\in C_c(G)\)
  with \(\psi_U\geq0\), \(\operatorname{supp}\psi_U\subseteq U\),
  \(\int\psi_U=1\) and \(\psi_U(x^{-1})=\psi_U(x)\).
\item
  Let \((\psi_U)\) be any family with \(\psi_U\geq0\), \(\int\psi_U=1\),
  and \(\psi_U\) vanishing outside a compact subset of \(U\) (continuity
  is not required). For \(1\leq p<\infty\) and \(f\in L^p(G)\),
  \[\|\psi_U*f-f\|_p\leq\sup_{y\in U}\|L_yf-f\|_p\longrightarrow0 .
  \tag{15.1}\] If \(f\) is bounded with \(\|L_yf-f\|_{\sup}\to0\) as
  \(y\to e\), then \(\|\psi_U*f-f\|_{\sup}\to0\).
\item
  If moreover \(\psi_U(x^{-1})=\psi_U(x)\), then
  \(\|f*\psi_U-f\|_p\leq\sup_{y\in U}\|R_yf-f\|_p\to0\), and
  \(\|f*\psi_U-f\|_{\sup}\to0\) for bounded \(f\) with
  \(\|R_yf-f\|_{\sup}\to0\).
\item
  (\emph{Weak-integral form}) For \(\psi\in L^1(G)\) and
  \(\xi,\eta\in L^2(G)\),
  \[\langle\psi*\xi,\eta\rangle=\int\psi(s)\,\langle L_s\xi,\eta\rangle\,ds ,
  \tag{15.2}\] where \(s\mapsto\langle L_s\xi,\eta\rangle\) is
  continuous and bounded by \(\|\xi\|\|\eta\|\). So the operator
  \(\xi\mapsto\psi*\xi\) is the weak integral
  \(\int\psi(s)\lambda(s)\,ds\) of the left regular representation, its
  norm is at most \(\|\psi\|_1\), and \(\psi_U*\xi\to\xi\) in \(L^2\)
  for every \(\xi\).
\end{enumerate}

\textbf{Proof.} (1) The set
\(U'=\operatorname{int}U\cap(\operatorname{int}U)^{-1}\) is an open
neighbourhood of \(e\). By (T2) choose \(\varphi\) with
\(\{e\}\prec\varphi\prec U'\), put \(\varphi'=\varphi+\check\varphi\),
which vanishes off
\(\operatorname{supp}\varphi\cup(\operatorname{supp}\varphi)^{-1}\subseteq U'\),
and \(\psi_U=\varphi'/\int\varphi'\); the integral is positive by
Proposition 9.1. (2) Let \(C\) be a compact set outside of which
\(\psi=\psi_U\) vanishes. For almost every \(x\), (14.1) converges
absolutely by Theorem 14.2(4), and since \(\int\psi=1\),
\(\psi*f(x)-f(x)=\int\psi(y)\big(f(y^{-1}x)-f(x)\big)dy\).
Hölder\textquotesingle s inequality for the probability measure
\(\psi(y)\,dy\) gives
\(|\psi*f(x)-f(x)|^p\leq\int\psi(y)|f(y^{-1}x)-f(x)|^p\,dy\). The
function \((x,y)\mapsto\psi(y)|f(y^{-1}x)-f(x)|^p\) is Borel and
vanishes outside
\((\{f\neq0\}\times C)\cup\Theta_1^{-1}(C\times\{f\neq0\})\), which is
\(\sigma\)-finite by Theorem 5.1(6) and Lemma 14.1. By
Tonelli\textquotesingle s theorem,
\[\|\psi*f-f\|_p^p\leq\int\psi(y)\Big(\int|f(y^{-1}x)-f(x)|^p\,dx\Big)dy=\int\psi(y)\,\|L_yf-f\|_p^p\,dy\leq\sup_{y\in U}\|L_yf-f\|_p^p ,\]
which tends to \(0\) by the continuity of translation (Theorem 14.2(6)).
For the uniform statement,
\(|\psi*f(x)-f(x)|\leq\int\psi(y)|L_yf(x)-f(x)|\,dy\leq\sup_{y\in U}\|L_yf-f\|_{\sup}\)
at every \(x\). (3) By the first form of (14.2) and the symmetry of
\(\psi\), \(f*\psi(x)=\int f(xy)\psi(y)\,dy\), so
\(f*\psi(x)-f(x)=\int\psi(y)(R_yf(x)-f(x))\,dy\); the integral converges
absolutely for almost every \(x\) by Theorem 14.2(4), since \(\psi\) has
compact support. Repeat the argument of (2), with \(\Theta_3\) in place
of \(\Theta_1\) and \(\int|f(xy)-f(x)|^p\,dx=\|R_yf-f\|_p^p\). (4) The
function \((x,s)\mapsto\psi(s)\xi(s^{-1}x)\overline{\eta(x)}\) is Borel,
vanishes outside a \(\sigma\)-finite set (as in (2)), and by
Tonelli\textquotesingle s theorem and the Cauchy--Schwarz inequality in
\(x\) it is \(\pi\)-integrable with integral of its absolute value at
most
\(\int|\psi(s)|\,\|L_s\xi\|\,\|\eta\|\,ds=\|\psi\|_1\|\xi\|\|\eta\|\).
Fubini\textquotesingle s theorem 5.1(5) gives (15.2). Continuity:
\(|\langle L_s\xi-L_{s_0}\xi,\eta\rangle|\leq\|L_{s_0^{-1}s}\xi-\xi\|\,\|\eta\|\to0\)
by Theorem 14.2(6). The norm bound is Theorem 14.2(4) with \(p=2\), and
the convergence is (2). \(\square\)

\emph{Units.} If \(G\) is discrete (with counting measure),
\(\delta=1_{\{e\}}\) satisfies \(\delta*f=f*\delta=f\). If \(G\) is not
discrete, \(L^1(G)\) has no unit, which is why approximate identities
are needed. Indeed, suppose \(u*f=f\) for all \(f\in L^1(G)\). Taking
\(f=\psi_U\) symmetric, \(\psi_U=u*\psi_U\to u\) in \(L^1\) by (3). For
a compact \(K\) with \(e\notin K\), \(\int_K|\psi_U|=0\) as soon as
\(U\cap K=\varnothing\), so \(\int_K|u|=0\). By inner regularity on the
\(\sigma\)-finite set \(\{u\neq0\}\) (Proposition 2.3(2)), \(u=0\)
almost everywhere off \(\{e\}\). And \(\mu(\{e\})=0\): otherwise every
point would have the same positive measure (by left invariance), a
compact neighbourhood of \(e\) would be finite by (R1), and \(\{e\}\)
would be open. So \(u=0\) almost everywhere, and \(u*f=0\neq f\) for
\(f\neq0\), a contradiction.

A family as in (1) is a \emph{compactly supported approximate identity}.
Its translates behave well: for \(g\in G\),
\((L_g\psi_U)*\xi=L_g(\psi_U*\xi)\) by Theorem 14.2(3), so convolution
by \(L_g\psi_U\) is \(\lambda(g)\) composed with convolution by
\(\psi_U\).

\subsection{16.
Exercises}\label{haar-measure-on-locally-compact-groups--16-exercises}

\textbf{Exercise 16.1} (A density that vanishes on a closed set). On
\(G=\mathbb R\times\mathbb R_d\) with Haar measure \(\mu\) (Example
13.4), let \(\varphi(x,d)=|x|\) and \(\nu(E)=\int_E\varphi\,d\mu\). Show
that \(\nu\) is a Borel measure, finite on compact sets, with
\(\nu(Y)=0\) for \(Y=\{0\}\times\mathbb R_d\) but \(\nu(U)=\infty\) for
every open \(U\supseteq Y\). Conclude that the positivity of the density
in Proposition 3.1(2) cannot be dropped.

\emph{Solution.} \(\nu\) is a measure because \(\varphi\) is a
nonnegative Borel function. A compact set \(K\) meets finitely many
lines and \(\varphi\) is bounded on it, so
\(\nu(K)\leq\max_K\varphi\cdot\mu(K)<\infty\). Since \(\varphi=0\) on
\(Y\), \(\nu(Y)=0\). If \(U\supseteq Y\) is open, then for each \(d\)
there is \(\delta_d>0\) with
\((-\delta_d,\delta_d)\times\{d\}\subseteq U\), and
\(\nu\big((-\delta_d,\delta_d)\times\{d\}\big)=\int_{-\delta_d}^{\delta_d}|x|\,dx=\delta_d^2>0\),
because \(\mu\) is Lebesgue measure on each line. These sets are
disjoint and uncountably many, so, as in the proof of Proposition
13.1(2), \(\nu(U)=\infty\). So \(\nu\) fails (R2) at \(Y\).

\textbf{Exercise 16.2} (Closed ideals of \(L^1(G)\)). Let
\(\mathcal J\subseteq L^1(G)\) be a closed subspace. Show that
\(\mathcal J\) is a left ideal (\(g*f\in\mathcal J\) for \(g\in L^1\),
\(f\in\mathcal J\)) exactly when \(L_x\mathcal J\subseteq\mathcal J\)
for all \(x\), and a right ideal exactly when
\(R_x\mathcal J\subseteq\mathcal J\) for all \(x\).

\emph{Solution.} \emph{Left, only if.} For \(f\in\mathcal J\) and
\(x\in G\), \(L_x(\psi_U*f)=(L_x\psi_U)*f\in\mathcal J\) by Theorem
14.2(3), and \(L_x(\psi_U*f)\to L_xf\) in \(L^1\) by Theorem 15.1(2),
since \(L_x\) is isometric. As \(\mathcal J\) is closed,
\(L_xf\in\mathcal J\). \emph{Left, if.} Since \(C_c(G)\) is dense in
\(L^1\) and \(\|g*f\|_1\leq\|g\|_1\|f\|_1\), it suffices to take
\(g\in C_c(G)\). Let \(K=\operatorname{supp}g\) and \(\varepsilon>0\).
By Theorem 14.2(6), \(y\mapsto L_yf\) is continuous into \(L^1\), so
each \(y\in K\) has an open neighbourhood \(O_y\) with
\(\|L_{y'}f-L_yf\|_1<\varepsilon\) for \(y'\in O_y\). Cover \(K\) by
\(O_{y_1},\dots,O_{y_n}\), and let
\(E_j=(K\cap O_{y_j})\setminus\bigcup_{i<j}O_{y_i}\). Put
\(h=\sum_j\big(\int_{E_j}g\big)L_{y_j}f\in\mathcal J\). For almost every
\(x\),
\(g*f(x)-h(x)=\sum_j\int_{E_j}g(y)\big(f(y^{-1}x)-f(y_j^{-1}x)\big)dy\).
The integrands vanish outside \(\sigma\)-finite sets (Lemma 14.1), so
Tonelli\textquotesingle s theorem gives
\(\|g*f-h\|_1\leq\sum_j\int_{E_j}|g(y)|\,\|L_yf-L_{y_j}f\|_1\,dy\leq\varepsilon\|g\|_1\).
Hence \(g*f\in\overline{\mathcal J}=\mathcal J\). \emph{Right.} If
\(\mathcal J\) is a right ideal, take symmetric \(\psi_U\) (Theorem
15.1(1)); then \(R_x(f*\psi_U)=f*(R_x\psi_U)\in\mathcal J\), and
\(f*\psi_U\to f\) by Theorem 15.1(3), while \(R_x\) is bounded on
\(L^1\). Conversely, for \(g\in C_c(G)\) the first form of (14.2) reads
\(f*g(x)=\int g(y^{-1})R_yf(x)\,dy\), and the same Riemann-sum argument,
with the continuous map \(y\mapsto R_yf\) and \(\Theta_3\) for the
\(\sigma\)-finiteness, shows \(f*g\in\mathcal J\).

\textbf{Exercise 16.3} (The regular representations). Show that
\(\lambda(g)=L_g\) and \(\rho(g)=\Delta(g)^{1/2}R_g\) are unitary
representations of \(G\) on \(L^2(G)\), continuous for the strong
operator topology, and that \(\lambda(g)\rho(h)=\rho(h)\lambda(g)\).

\emph{Solution.} \(\lambda(g)\) is isometric by left invariance
(Proposition 8.2(2)), and \(\rho(g)\) is isometric because
\(\|R_g\xi\|_2=\Delta(g)^{-1/2}\|\xi\|_2\) (Theorem 10.1(3)). Since
\(L_{gh}=L_gL_h\), \(R_{gh}=R_gR_h\) and \(\Delta\) is multiplicative,
both are homomorphisms, so each \(\lambda(g)\), \(\rho(g)\) has the
inverse \(\lambda(g^{-1})\), \(\rho(g^{-1})\) and is unitary. They
commute: \(L_gR_h\xi(t)=\xi(g^{-1}th)=R_hL_g\xi(t)\). At \(e\),
\(\|\lambda(g)\xi-\xi\|\to0\) by Theorem 14.2(6), and
\(\|\rho(g)\xi-\xi\|\leq|\Delta(g)^{1/2}-1|\,\|R_g\xi\|+\|R_g\xi-\xi\|\to0\)
by the continuity of \(\Delta\) (Theorem 10.1(2)) and Theorem 14.2(6).
At \(g_0\),
\(\|\lambda(g)\xi-\lambda(g_0)\xi\|=\|\lambda(g_0^{-1}g)\xi-\xi\|\), and
likewise for \(\rho\).

\textbf{Exercise 16.4} (The \(ax+b\) group: left against right, and a
failure of \(L^\infty*L^1\)). On the \(ax+b\) group of Example 11.3,
with \(d\mu=a^{-2}da\,db\), \(d\tilde\mu=a^{-1}da\,db\) and
\(\Delta(a,b)=1/a\): (a) find Borel sets \(E,E'\) with
\(\mu(E)<\infty=\tilde\mu(E)\) and \(\tilde\mu(E')<\infty=\mu(E')\), and
conclude that neither of \(L^2(\mu)\), \(L^2(\tilde\mu)\) contains the
other; (b) for \(g\equiv1\in L^\infty(G)\) and \(f=1_E\) with \(E\) from
(a), show \(f\in L^1(\mu)\) and \(g*f(x)=\infty\) for every \(x\).

\emph{Solution.} By Proposition 3.1(2), \(\mu\) and \(\tilde\mu\) are
the measures with densities \(a^{-2}\) and \(a^{-1}\) with respect to
the Radon product of Lebesgue measures on \((0,\infty)\times\mathbb R\):
both sides are Radon and integrate \(C_c\) alike. Measures of regions
are then iterated integrals by Tonelli\textquotesingle s theorem 5.1(4).
(a) Let \(E=\{a>1,\ 0<b<1\}\): \(\mu(E)=\int_1^\infty a^{-2}\,da=1\) and
\(\tilde\mu(E)=\int_1^\infty a^{-1}\,da=\infty\). Let
\(E'=\{0<a<1,\ 0<b<a\}\): \(\tilde\mu(E')=\int_0^1a\cdot a^{-1}\,da=1\)
and \(\mu(E')=\int_0^1a\cdot a^{-2}\,da=\infty\). So
\(1_E\in L^2(\mu)\setminus L^2(\tilde\mu)\) and
\(1_{E'}\in L^2(\tilde\mu)\setminus L^2(\mu)\). Corollary 11.2(3) gives
isometries between the two spaces, but they are different spaces. (b)
\(\|f\|_1=\mu(E)=1\). For every \(x\), left invariance and (11.1) give
\(g*f(x)=\int f(y^{-1}x)\,dy=\int f(y^{-1})\,dy=\int f\Delta^{-1}\,d\mu=\tilde\mu(E)=\infty\).
Here \(\Delta^{-1}f\notin L^1\), which is exactly the extra hypothesis
in Theorem 14.2(7).

\subsection{Results used from other
lessons}\label{haar-measure-on-locally-compact-groups--results-used-from-other-lessons}

\emph{Topology}, from
\hyperref[the-stone-weierstrass-theorem-for-functions-vanishing-at-infinity]{The
Stone--Weierstrass theorem for functions vanishing at infinity} and from
\hyperref[weak-topologies-tychonoff-banach-alaoglu-mazur-bipolars-krein-milman-and-eberlein-smulian]{Weak
topologies: Tychonoff, Banach--Alaoglu, Mazur, bipolars, Krein--Milman
and Eberlein--Šmulian}.

\begin{itemize}
\tightlist
\item
  \textbf{Compact closures}
  (\hyperlink{the-stone-weierstrass-theorem-for-functions-vanishing-at-infinity--oa-fnd-sw-15}{Proposition
  4.2(2) of the Stone--Weierstrass lesson}). If \(K\) is a compact
  subset of an LCH space and \(V\) is a neighbourhood of \(K\), some
  open \(U\supseteq K\) has compact closure contained in \(V\). This is
  (T1).
\item
  \textbf{Urysohn\textquotesingle s lemma}
  (\hyperlink{the-stone-weierstrass-theorem-for-functions-vanishing-at-infinity--oa-fnd-sw-16}{Corollary
  5.2 of the Stone--Weierstrass lesson}). If \(K\) is compact, \(B\) is
  closed and \(K\cap B=\varnothing\) in an LCH space \(X\), some
  \(f\in C_c(X)\) with \(0\leq f\leq1\) equals \(1\) on \(K\) and \(0\)
  on \(B\).
\item
  \textbf{The Stone--Weierstrass theorem}
  (\hyperlink{the-stone-weierstrass-theorem-for-functions-vanishing-at-infinity--oa-fnd-sw-09}{Theorem
  10.1 of that lesson}). Let \(X\) be an LCH space and \(A\) a
  subalgebra of \(C_0(X)\) that is closed under complex conjugation,
  separates the points of \(X\), and contains for each \(x\in X\) a
  function that does not vanish at \(x\). Then \(A\) is dense in
  \(C_0(X)\). It is used only in Example 12.2.
\item
  \textbf{Compact Tietze extension} (Proposition 16.1(2) of the
  Stone--Weierstrass lesson). On a compact Hausdorff space, a continuous
  real function on a closed subset extends continuously with the same
  supremum bound; real and imaginary extensions give the complex
  version. Its proof uses the compact Stone--Weierstrass theorem and a
  uniformly summable correction series. This supplies compact extension
  whenever needed; no normality of an arbitrary LCH space is assumed.
\item
  \textbf{Tychonoff\textquotesingle s theorem} (Theorem 2.3 of the
  lesson on weak topologies). Every product of compact spaces is compact
  in the product topology.
\end{itemize}

\emph{Measure theory}, from
\hyperref[measure-and-hilbert-space-tools]{Measure and Hilbert space
tools for Haar integration}: Theorem 1.1 proves Carathéodory; Section 2
proves measurability of sequential limits; Theorems 2.1--2.2 prove
monotone and dominated convergence; Theorems 3.1--3.2 prove Hölder,
completeness and simple approximation; Theorems 4.1--4.2 prove
Radon--Nikodym and the sigma-finite dual of \(L^1\). Fremlin is a
complementary human source for these classical statements.

\begin{itemize}
\tightlist
\item
  \textbf{Carathéodory\textquotesingle s theorem}
  \hyperref[measure-and-hilbert-space-tools--1-from-an-outer-measure-to-a-measure]{the
  preceding tools lesson, Theorem 1.1}. For an outer measure \(\mu^*\)
  on \(X\), the sets \(A\) with
  \(\mu^*(Y)=\mu^*(Y\cap A)+\mu^*(Y\setminus A)\) for every
  \(Y\subseteq X\) form a \(\sigma\)-algebra on which \(\mu^*\) is
  countably additive. To check the condition it suffices to show
  \(\mu^*(Y)\geq\mu^*(Y\cap A)+\mu^*(Y\setminus A)\) when
  \(\mu^*(Y)<\infty\): the reverse inequality is subadditivity, and the
  case \(\mu^*(Y)=\infty\) is trivial.
\item
  \textbf{Measurable functions} {[}Fremlin{]}. The supremum, the infimum
  and the pointwise limit of a sequence of Borel functions with values
  in \(\mathbb R\) or in \([0,\infty]\) are Borel.
\item
  \textbf{Convergence theorems.} Monotone convergence for functions with
  values in \([0,\infty]\) {[}Fremlin{]}, and for integrable functions
  {[}Fremlin{]}. Dominated convergence {[}Fremlin{]}: if \(f_n\to f\)
  almost everywhere and \(|f_n|\leq g\) with \(g\) integrable, then
  \(\int f_n\,d\mu\to\int f\,d\mu\); applied to \(|f_n-f|\leq2g\) it
  gives \(\int|f_n-f|\,d\mu\to0\).
\item
  \textbf{\(L^p\) spaces.} Hölder\textquotesingle s inequality
  {[}Fremlin{]}. For \(1\leq p<\infty\), \(L^p(\mu)\) is complete
  {[}Fremlin{]}, and the simple functions in \(L^p(\mu)\) are dense in
  it {[}Fremlin{]}.
\item
  \textbf{The dual of \(L^1\)} {[}Fremlin{]}. If \(\mu\) is
  \(\sigma\)-finite, every bounded linear functional \(\Phi\) on
  \(L^1(\mu)\) has the form \(\Phi(g)=\int fg\,d\mu\) for a unique
  \(f\in L^\infty(\mu)\), and \(\|\Phi\|=\|f\|_\infty\).
\item
  \textbf{The Radon--Nikodym theorem} {[}Fremlin{]}. Let \(\mu\) and
  \(\nu\) be measures on the same \(\sigma\)-algebra, \(\nu\) finite and
  \(\mu\) \(\sigma\)-finite, with \(\nu(E)=0\) whenever \(\mu(E)=0\).
  Then \(\nu(E)=\int_Eh\,d\mu\) for an integrable \(h\geq0\). In Theorem
  13.3(5) both measures are \(\sigma\)-finite on the set considered; the
  set is split into countably many Borel pieces of finite measure for
  both, and the theorem is applied on each piece.
\item
  \textbf{Substitution in Riemann integrals.} If
  \(\phi\colon[\alpha,\beta]\to\mathbb R\) is continuously
  differentiable and \(f\) is continuous on \(\phi([\alpha,\beta])\),
  then
  \(\int_{\phi(\alpha)}^{\phi(\beta)}f(u)\,du=\int_\alpha^\beta f(\phi(t))\,\phi'(t)\,dt\).
  The integral identity is proved in
  \hyperref[measure-and-hilbert-space-tools]{Measure and Hilbert space
  tools for Haar integration}, Lemma 6.1, from the elementary
  differential-calculus prerequisites.
\end{itemize}

\emph{Hilbert spaces}, from
\hyperref[hilbert-spaces-and-compact-operators]{Hilbert spaces and
compact operators}.

\begin{itemize}
\tightlist
\item
  \textbf{Hilbert tensor products} (Theorem 8.1). For Hilbert spaces
  \(H\) and \(K\) there is a Hilbert space \(H\otimes K\) with a
  bilinear map \((\xi,\eta)\mapsto\xi\otimes\eta\) such that
  \(\langle\xi\otimes\eta,\xi'\otimes\eta'\rangle=\langle\xi,\xi'\rangle\langle\eta,\eta'\rangle\)
  and the elementary tensors span a dense subspace.
\item
  \textbf{Multiplication operators on \(\sigma\)-finite spaces} (Theorem
  9.1). For a \(\sigma\)-finite measure space \((Z,\nu)\), the operators
  \(m_f\xi=f\xi\) on \(L^2(\nu)\) with \(f\in L^\infty(\nu)\) form an
  algebra that is maximal abelian: every bounded operator on
  \(L^2(\nu)\) that commutes with all \(m_f\) is itself some \(m_f\).
\end{itemize}

\subsection{Where this
leads}\label{haar-measure-on-locally-compact-groups--where-this-leads}

\begin{itemize}
\tightlist
\item
  The lesson \emph{The Plancherel weight and Fourier coefficients of a
  locally compact group}, in the course on modular theory and weights,
  is built on this one and, like it, works for every locally compact
  group. It uses the normalization (10.1) of the modular function and
  the inversion formula (11.1); the regular representations of Exercise
  16.3, whose strong continuity rests on Proposition 7.3 and Theorem
  14.2(6), and the conjugation \(J\) of Corollary 11.2(2); the density
  of \(C_c(G)\) in \(L^2(G)\) (Proposition 3.1(4)) and the core of
  Corollary 11.2(5); the approximate identities and weak integrals of
  Theorem 15.1; the identification
  \(L^2(G)\otimes L^2(G)\cong L^2(G\times G,\mu\hat\times\mu)\) of
  Theorem 6.2, with the operators \(m_a\otimes1\) of part (4), the
  unitary \(W\) of Proposition 12.1, and Tonelli\textquotesingle s
  theorem 5.1(4) for norm computations on elementary tensors; and the
  von Neumann algebra \(L^\infty(G)\) of Theorem 13.3(4). The unitarity
  of \(W\) can also be checked on compactly supported tensors, one \(s\)
  at a time, and then extended by density; Theorem 6.2(3) is what
  justifies working with one \(s\) at a time.
\item
  Tensor products of \(L^\infty(G)\) with other von Neumann algebras are
  spatial tensor products; see the lesson
  \href{https://kokunoyumeto.github.io/open-math-courses/courses/harmonic-analysis-on-locally-compact-groups/reader/programme-reading.html\#ref-c4b6f64f97d8c007}{Spatial
  tensor products of von Neumann algebras}.
\item
  For a Radon measure that is not \(\sigma\)-finite and does not come
  from a group, the maximal abelian property of \(L^\infty\) (Theorem
  13.3(4) for Haar measure) needs a decomposition of the space into
  disjoint compact pieces; for Haar measure, the open \(\sigma\)-compact
  cosets of Section 13 play that role. The lesson on vector-valued
  functions and preduals constructs
  \href{https://kokunoyumeto.github.io/open-math-courses/courses/harmonic-analysis-on-locally-compact-groups/reader/programme-reading.html\#ref-ca782bdbcbfeedb9}{the
  decomposition} and proves
  \href{https://kokunoyumeto.github.io/open-math-courses/courses/harmonic-analysis-on-locally-compact-groups/reader/programme-reading.html\#ref-ca782bdbcbfeedb9}{the
  maximal abelian property for every Radon measure}.
\end{itemize}

\subsection{References}\label{haar-measure-on-locally-compact-groups--references}

\begin{itemize}
\item
  {[}Fremlin{]} D. H. Fremlin, \emph{Measure Theory}, Volumes 1, 2 and
  4,
  \href{https://www1.essex.ac.uk/maths/people/fremlin/mt.htm}{author's
  freely accessible text and editable TeX}. Volume 4 uses complete,
  locally determined Radon measures; Proposition 13.7 relates that
  convention to ours.
\item
  {[}van Doorn{]} F. van Doorn, \emph{Formalized Haar Measure}, 2021,
  \href{https://arxiv.org/abs/2102.07636v1}{arXiv version 1}. The paper
  constructs Haar measure on arbitrary LCH groups; its uniqueness proof
  assumes second countability. The
  \href{https://leanprover-community.github.io/mathlib4_docs/Mathlib/MeasureTheory/Measure/Haar/Unique.html}{current
  mathlib uniqueness development} also distinguishes regularity
  hypotheses.
\item
  {[}Cartan 1940{]} H. Cartan, Sur la mesure de Haar, \emph{Comptes
  Rendus de l\textquotesingle Académie des Sciences de Paris} 211
  (1940), 759--762.
  https://gallica.bnf.fr/ark:/12148/bpt6k3163d/f759.item
\item
  {[}von Neumann 1936{]} J. von Neumann, The uniqueness of
  Haar\textquotesingle s measure, \emph{Matematicheskii Sbornik} (N.S.)
  1(43) (1936), 721--734. https://www.mathnet.ru/eng/sm5481
\item
  {[}Fischer--Ruzhansky{]} V. Fischer and M. Ruzhansky,
  \href{https://biblio.ugent.be/publication/8585474}{\emph{Quantization
  on Nilpotent Lie Groups}}, Birkhäuser, 2016, freely accessible under
  CC BY 4.0. Its Heisenberg and homogeneous-group treatment complements
  Example 12.3; its convolution discussion includes the full Young bound
  of Theorem 14.4. The calculations and proof here are independently
  written.
\end{itemize}

\section{Quotient measures and Weil\textquotesingle s integration
formula}\label{quotient-measures-and-weils-integration-formula}

\emph{Originally written and self-checked by GPT-6 Astra (OpenAI), Ultra
reasoning effort, October 2026. Reconciled, extended and self-checked by
GPT-6.1 Sol (OpenAI), Ultra, October 2026. These are writing-AI checks;
no independent AI or human review of this edition is claimed. Original
exposition is dedicated under CC0. The programme\textquotesingle s CC0
sources are credited at the end.}

Integration over a group can often be separated into integration along a
subgroup and integration over its cosets. For a nonabelian group, an
invariant measure on the coset space need not exist. The precise
obstruction is the mismatch between two modular functions:
\[G/H\text{ has a nonzero invariant Radon measure}
\quad\Longleftrightarrow\quad
\Delta_G|_H=\Delta_H.\] When the obstruction does not vanish, a positive
continuous density still gives a strongly quasi-invariant measure. We
construct it, prove its change-of-variables formula and calculate
examples where the distinction matters.

The lesson works for every locally compact Hausdorff group \(G\) and
every closed subgroup \(H\). Neither is assumed second countable,
\(\sigma\)-compact or unimodular. All interchanges of integrals in the
main construction involve continuous functions with compact support. We
will not infer a Fubini theorem for arbitrary nonnegative Borel
functions on a non-\(\sigma\)-finite product.

\subsection{What to know
first}\label{quotient-measures-and-weils-integration-formula--what-to-know-first}

The preceding programme lesson
\hyperref[haar-measure-on-locally-compact-groups]{Haar measure on
locally compact groups} supplies the following exact results:

\begin{itemize}
\tightlist
\item
  Theorem 2.2: a positive linear functional on \(C_c(X)\) is represented
  by a unique Radon measure.
\item
  Proposition 3.1(2): multiplication by a strictly positive continuous
  function gives a Radon measure in the same measure class.
\item
  Proposition 4.2: iterated integrals of a continuous compactly
  supported function on a product agree.
\item
  Proposition 9.1 and Theorem 9.2: positivity on nonempty open sets and
  uniqueness of Haar measure.
\item
  Theorems 10.1 and 11.1: the modular function and the inversion
  formula.
\end{itemize}

Their elementary measure-theory prerequisites are proved in the
preceding tools lesson. The topology needed specifically for the
quotient construction is proved below, including the cutoff argument
sometimes hidden behind a reference to a partition of unity.

Fix left Haar measures \(dg\) and \(dh\). Our convention is
\[\int_G f(ga)\,dg=\Delta_G(a)^{-1}\int_G f(g)\,dg.
\tag{1}\] Equivalently, the measure of a right translate \(Ea\) is
\(\Delta_G(a)\) times the measure of \(E\). The same convention defines
\(\Delta_H\). Put \(Y=G/H\), the space of cosets \(gH\), and
\(q(g)=gH\).

As in the preceding lesson, a Radon measure here is a Borel measure
finite on compact sets, outer regular on Borel sets and inner regular on
open sets. Equality of Radon measures can therefore be checked on
\(C_c\). Left and right translations of functions mean
\(L_af(g)=f(a^{-1}g)\) and \(R_af(g)=f(ga)\).

\subsection{Local topology and compactly supported
functions}\label{quotient-measures-and-weils-integration-formula--local-topology-and-compactly-supported-functions}

\textbf{Lemma 1.1.} An LCH space has a base of relatively compact open
neighborhoods. If \(K\subset V\), where \(K\) is compact and \(V\) is
open, there is \(u\in C_c(X)\) with \(0\leq u\leq1\), \(u=1\) on \(K\),
and \(\operatorname{supp}u\subset V\).

\emph{Proof.} We first recall the separation argument for a compact
Hausdorff space \(C\). For a point \(x\) and a closed set \(A\) not
containing it, separate \(x\) from each point of \(A\) by disjoint open
neighborhoods. A finite subcover of \(A\) gives an open neighborhood
\(W\) of \(x\) whose closure misses \(A\). For two disjoint closed sets,
do this for each point of the first, then take a finite union of the
resulting neighborhoods. Thus whenever \(A\subset O\), with \(A\) closed
and \(O\) open, there is open \(W\) with
\[A\subset W\subset\overline W\subset O.
\tag{2}\]

This proves the continuous separation needed here. For disjoint closed
\(A,B\subset C\), choose open sets \(U_r\), indexed by dyadic rationals
\(r\in[0,1]\), with \(A\subset U_0\), \(U_1=C\setminus B\), and
\(\overline U_r\subset U_s\) when \(r<s\). Start by shrinking \(A\)
inside \(C\setminus B\); at each finite dyadic stage use (2) to insert
the new intermediate open sets. Set \[v(x)=\inf\{r:x\in U_r\},\] with
value \(1\) when the set is empty. Then \(v=0\) on \(A\) and \(v=1\) on
\(B\). The identity \(\{v<t\}=\bigcup_{r<t}U_r\) and, for \(0\leq t<1\),
\(\{v>t\}=\bigcup_{r>t}(C\setminus\overline U_r)\) prove continuity. In
the second identity, choose a dyadic \(r\) strictly between \(t\) and
\(v(x)\); if \(x\in\overline U_r\), the nesting would put \(x\) in every
\(U_s\) with \(s>r\), contradicting \(v(x)>r\).

Now let \(x\in O\subset X\), with \(O\) open and \(X\) LCH. Choose a
compact neighborhood \(C\) of \(x\). Apply (2) within \(C\) to \(x\) and
\(O\cap\operatorname{int}C\). The resulting neighborhood lies in
\(\operatorname{int}C\), so it is open in \(X\), and its closure is
compact and lies in \(O\). This proves the neighborhood-base assertion.

For compact \(K\subset V\), take finitely many such neighborhoods to
obtain open \(W\) with \(K\subset W\subset\overline W\subset V\) and
\(\overline W\) compact. In \(\overline W\), continuously separate \(K\)
and \(\overline W\setminus W\), taking the value \(1\) on the first and
\(0\) on the second. Extend the function by zero outside \(W\). The two
definitions agree on the boundary, so the pasting lemma on the two
closed sets \(\overline W\) and \(X\setminus W\) proves continuity. This
is the required \(u\). \(\square\)

\textbf{Lemma 1.2.} The quotient map \(q\) is open, \(Y\) is LCH, and
every compact subset \(E\subset Y\) is the image of a compact subset of
\(G\).

\emph{Proof.} For open \(O\subset G\), \(q^{-1}(q(O))=OH\) is open, so
\(q(O)\) is open. If \(xH\neq yH\), then \(y^{-1}x\notin H\). Since
\(H\) is closed, continuity gives neighborhoods \(V\) of \(x\) and \(W\)
of \(y\) with \(W^{-1}V\cap H=\varnothing\). Their images under \(q\)
are disjoint open neighborhoods. Thus \(Y\) is Hausdorff. The image of a
compact neighborhood of \(g\) contains the image of its interior, an
open neighborhood of \(gH\), so \(Y\) is locally compact.

Choose an identity neighborhood \(V\subset G\) with compact closure.
Finitely many \(q(g_iV)\) cover \(E\). Then
\[K=q^{-1}(E)\cap\bigcup_i g_i\overline V\] is compact and \(q(K)=E\).
Here \(q^{-1}(E)\) is closed because \(Y\) is Hausdorff. No continuous
section of \(q\) is being assumed. \(\square\)

\subsection{Averaging along the
subgroup}\label{quotient-measures-and-weils-integration-formula--averaging-along-the-subgroup}

Why not simply push Haar measure forward through \(q\)? Consider
\(G=\mathbb R^2\), with two-dimensional Lebesgue measure, and
\(H=\mathbb R\times\{0\}\), so that \(q(x,y)=y\). The inverse image of a
nondegenerate compact interval \([a,b]\) contains every rectangle
\([-N,N]\times[a,b]\). Its measure is at least \(2N(b-a)\) for every
\(N\), hence is infinite. The pushforward is therefore not finite on
compact intervals and is not the required Radon measure. Nevertheless,
the quotient plainly has the invariant measure \(dy\).

The construction below separates these two operations. First integrate a
compactly supported function along the fiber, rather than measuring the
whole infinite fiber. Then integrate the resulting function over the
base. For the example, this is precisely \(\int f(x,y)\,dx\) followed by
integration in \(y\). A function constant on every fiber is usually not
compactly supported upstairs; the cutoff used below supplies a compactly
supported representative without pretending that the fiber has finite
mass. This is why surjectivity of the averaging map, rather than
existence of a global section or finiteness of subgroup volume, is the
useful starting point.

For \(f\in C_c(G)\), define \[Pf(gH)=\int_H f(gh)\,dh.
\tag{3}\] Left invariance of \(dh\) makes this independent of the
representative \(g\): replacing \(g\) by \(gr\), \(r\in H\), replaces
\(h\) by \(rh\).

\textbf{Lemma 2.1.} The map \(P:C_c(G)\to C_c(Y)\) is positive and
surjective. For each compact \(E\subset Y\), there is \(b\in C_c(G)\),
\(b\geq0\), with \(Pb=1\) on \(E\). A nonnegative function in \(C_c(Y)\)
has a nonnegative lift through \(P\).

\emph{Proof.} Positivity is immediate. For \(g\) in a compact
neighborhood \(A\) of \(g_0\), only
\[h\in H\cap A^{-1}\operatorname{supp}f\] can contribute to (3). This
set is compact. The continuous function \((g,h)\mapsto f(gh)\) varies
uniformly in \(h\) on this compact set as \(g\to g_0\): a finite cover
of that set by continuity neighborhoods proves the assertion. Its
integral therefore varies continuously. The resulting continuous
function on \(G\) is constant on fibers, so descends continuously
through the quotient map. Its support lies in
\(q(\operatorname{supp}f)\), a compact set.

Lift \(E\) to compact \(K\subset G\) by Lemma 1.2, and choose \(a\geq0\)
in \(C_c(G)\) with \(a=1\) on \(K\), by Lemma 1.1. On every fiber over
\(E\), \(a\) is positive on a relatively open \(H\)-neighborhood of a
point of \(K\); hence \(Pa>0\) there by positivity of Haar measure.
Choose \(\chi\in C_c(Y)\), equal to \(1\) on \(E\), with
\(0\leq\chi\leq1\) and \(\operatorname{supp}\chi\subset\{Pa>0\}\). Set
\[b(g)=a(g)\frac{\chi(q(g))}{Pa(q(g))},
\tag{4}\] where the fraction is zero outside \(\{Pa>0\}\). The numerator
has support compactly inside that open set, so this extension is
continuous. Thus \(b\geq0\), \(b\in C_c(G)\) and \(Pb=\chi\).

For \(F\in C_c(Y)\), choose this \(b\) for \(E=\operatorname{supp}F\).
Then \(f=(F\circ q)b\) has compact support and \(Pf=F\). If \(F\geq0\),
so is \(f\). \(\square\)

The map respects left translation and has a simple right-subgroup rule:
\[P(L_a f)(y)=Pf(a^{-1}y),\qquad
P(R_h f)=\Delta_H(h)^{-1}Pf\quad(h\in H).
\tag{5}\] Both follow directly from (3), with the modular convention
(1).

\subsection{Exactly when the quotient measure is
invariant}\label{quotient-measures-and-weils-integration-formula--exactly-when-the-quotient-measure-is-invariant}

\textbf{Theorem 3.1.} The space \(G/H\) has a nonzero \(G\)-invariant
Radon measure if and only if \(\Delta_G(h)=\Delta_H(h)\) for every
\(h\in H\). When this holds, it is unique up to a positive scalar. For
the fixed measures \(dg,dh\), there is exactly one normalization
\(d\mu\) for which \[\int_G f(g)\,dg=\int_Y Pf(y)\,d\mu(y)
\qquad(f\in C_c(G)).
\tag{6}\]

\emph{Proof of necessity.} If \(\mu\) is invariant, the functional
\[I(f)=\int_Y Pf\,d\mu\] is positive, left invariant and nonzero.
Nonzeroness follows from the positive surjectivity of \(P\) and from
Radon representation on \(Y\). By Haar uniqueness,
\(I(f)=c\int_G f\,dg\), for some \(c>0\). For \(h\in H\), (5) gives
\[I(R_h f)=\Delta_H(h)^{-1}I(f),\] while (1) and the description of
\(I\) give the same equation with \(\Delta_G(h)^{-1}\). Choose
nonnegative nonzero \(f\); then \(I(f)>0\), forcing the two modular
functions to agree on \(H\).

\emph{Proof of sufficiency.} Assume the modular functions agree on
\(H\). Fix \(f\in C_c(G)\), and choose \(b\geq0\) in \(C_c(G)\) with
\(Pb=1\) on \(q(\operatorname{supp}f)\). Then \[\begin{aligned}
\int_G f(t)\,dt
&=\int_G\int_H f(t)b(th)\,dh\,dt\\
&=\int_G b(s)\int_H f(sh^{-1})\Delta_G(h)^{-1}\,dh\,ds\\
&=\int_G b(s)Pf(sH)\,ds.
\end{aligned}
\tag{7}\] For the second line use \(s=th\) and (1). For the third use
\(\Delta_G|_H=\Delta_H\) and the inversion formula in \(H\). These
interchanges are justified by compact support: initially
\(t\in\operatorname{supp}f\) and
\(h\in H\cap(\operatorname{supp}f)^{-1}\operatorname{supp}b\). Thus the
integrand is a continuous compactly supported function on \(G\times H\),
where the preceding lesson\textquotesingle s compact-product theorem
applies.

In particular \(Pf=0\) implies \(\int f=0\). Consequently
\[\Lambda(Pf)=\int_G f\,dg\] is a well-defined linear functional on all
of \(C_c(Y)\). It is positive because nonnegative functions have
nonnegative lifts; it is nonzero and left invariant by (5). Radon
representation supplies \(\mu\), and equality of the integrals of
\(C_c(Y)\) functions proves that its translates equal it. This gives
(6).

Finally, any invariant Radon measure produces a scalar multiple of Haar
integration in the necessity argument. Surjectivity of \(P\) then makes
that measure a scalar multiple of \(\mu\). Equation (6) fixes the
scalar. \(\square\)

Two useful consequences require no additional construction. A compact
subgroup always satisfies the criterion: its modular function is one,
and the continuous homomorphism \(\Delta_G:H\to(0,\infty)\) is one on a
compact group. An open subgroup also satisfies it, because restriction
of Haar measure on \(G\) to \(H\) is Haar measure and has the same
right-translation factors. For a discrete subgroup, however, the
criterion is \(\Delta_G|_H=1\), not an automatic assertion that every
discrete-subgroup quotient has an invariant measure.

A compact cutoff in this proof must equal one on
\(q(\operatorname{supp}f)\), not just on the support of \(Pf\). For a
signed function, subgroup integration can cancel to zero on a coset
while \(f\) is nonzero there. The choice above covers those cosets as
well.

\subsection{A cutoff over all
cosets}\label{quotient-measures-and-weils-integration-formula--a-cutoff-over-all-cosets}

For a general subgroup the modular obstruction need not vanish. To
construct a compensating density, we need a continuous nonnegative \(k\)
on \(G\) such that \[Pk=1,\qquad
\operatorname{supp}k\cap q^{-1}(E)\text{ is compact for each compact }E\subset Y.
\tag{8}\] It will not generally have compact support on all of \(G\).

\textbf{Lemma 4.1.} A function \(k\) with properties (8) exists.

\emph{Proof.} Choose a relatively compact symmetric open identity
neighborhood \(V\) in \(G\). The subgroup \(L=\bigcup_{n\geq1}V^n\) is
open and \(\sigma\)-compact: \(V^n\subset(\overline V)^n\subset L\),
with the latter compact, since \(\overline V\subset V^2\). For the last
inclusion, if \(x\in\overline V\), the open set \(xV\) meets \(V\), so
\(x\in VV^{-1}=V^2\).

The \(L\)-orbits in \(Y\) are disjoint open sets, because \(q\) is open;
their complements are unions of other orbits, so they are also closed.
Each orbit \(Y_0\) is the continuous image of \(L\), hence is
\(\sigma\)-compact and LCH.

On such an orbit choose a compact exhaustion
\(K_n\subset\operatorname{int}K_{n+1}\), with union \(Y_0\). To
construct it, start with a countable compact cover; at each step cover
the preceding compact set and the next compact set in that cover by
finitely many relatively compact open neighborhoods, and take the union
of their compact closures. Put \(K_0=K_{-1}=\varnothing\). The compact
shells \[S_n=K_n\setminus\operatorname{int}K_{n-1}\] are contained in
the open sets \(W_n=\operatorname{int}K_{n+1}\setminus K_{n-2}\), and
cover \(Y_0\). By Lemma 1.1 choose \(\chi_n\in C_c(Y_0)\), nonnegative,
equal to one on \(S_n\), with support in \(W_n\). The family is locally
finite: a neighborhood contained in \(\operatorname{int}K_N\) misses
\(W_n\) for \(n\geq N+2\). Therefore
\[\psi_n=\frac{\chi_n}{\sum_j\chi_j}\] are continuous, compactly
supported, locally finite, nonnegative and sum to one on \(Y_0\).

Perform this construction in every orbit and extend the functions by
zero to \(Y\). The resulting family \((\psi_i)\) is still locally finite
because the orbits are open, and each \(\psi_i\) has compact support. By
Lemma 2.1 choose \(b_i\geq0\) in \(C_c(G)\) with \(Pb_i=1\) on
\(\operatorname{supp}\psi_i\). Set \[k(g)=\sum_i \psi_i(q(g))b_i(g).
\tag{9}\] Local finiteness makes this continuous. Averaging gives
\(Pk=\sum_i\psi_i=1\). A compact subset of \(Y\) meets only finitely
many supports of this locally finite family; over that compact set, the
support of \(k\) lies in the union of the corresponding finitely many
compact supports of \(b_i\). It is closed there, so is compact. This
proves (8). \(\square\)

The construction uses neither a countable base for \(G\) nor a global
continuous section of \(G\to G/H\).

\subsection{The modular density and weighted
integration}\label{quotient-measures-and-weils-integration-formula--the-modular-density-and-weighted-integration}

Write \[\chi(h)=\frac{\Delta_G(h)}{\Delta_H(h)}.\] This is a positive
continuous character on \(H\). From the cutoff of Lemma 4.1 define
\[\rho(g)=\int_H k(gh)\chi(h)\,dh.
\tag{10}\]

\textbf{Lemma 5.1.} The function \(\rho\) is continuous, strictly
positive and finite, and
\[\rho(gh)=\frac{\Delta_H(h)}{\Delta_G(h)}\rho(g)
\qquad(g\in G,\ h\in H).
\tag{11}\]

\emph{Proof.} For \(g\) in a compact neighborhood \(A\), property (8)
puts all \(gh\) that contribute into one compact subset \(B\) of \(G\).
Then \(h\) lies in the compact set \(H\cap A^{-1}B\). The integrand is
continuous and varies uniformly on this set; its integral is therefore
finite and continuous in \(g\). Since \(Pk=1\), \(k\) is not identically
zero on any fiber; \(\chi>0\), so \(\rho>0\). For \(r\in H\), substitute
\(t=rh\) using left invariance: \[\rho(gr)=\int_H k(gt)\chi(r^{-1}t)\,dt
=\chi(r)^{-1}\rho(g).\] This is (11). \(\square\)

A continuous positive function satisfying (11) is called a
\(\rho\)-function. It is a density on the group, not usually a function
on the quotient.

\textbf{Theorem 5.2.} For any \(\rho\)-function there is a unique Radon
measure \(\mu_\rho\) on \(Y\) satisfying
\[\int_Y Pf\,d\mu_\rho=\int_G f(g)\rho(g)\,dg
\qquad(f\in C_c(G)).
\tag{12}\] It is positive on nonempty open subsets of \(Y\). The
construction is independent of the cutoff.

\emph{Proof.} Use a cutoff \(k\) as in (8), and define a positive
functional by \[\Lambda(F)=\int_G F(q(g))\,k(g)\rho(g)\,dg
\qquad(F\in C_c(Y)).
\tag{13}\] The integrand has compact support by (8), so this is finite
and Radon representation applies.

For \(F=Pf\), exchange the integrals in (13) and put \(t=gh\). The
translation factor is \(\Delta_G(h)^{-1}\), while (11) gives
\(\rho(th^{-1})=\chi(h)\rho(t)\). Thus \[\begin{aligned}
\Lambda(Pf)
&=\int_G f(t)\rho(t)
\left(\int_H k(th^{-1})\Delta_H(h)^{-1}\,dh\right)dt\\
&=\int_G f(t)\rho(t)Pk(tH)\,dt
=\int_G f(t)\rho(t)\,dt.
\end{aligned}
\tag{14}\] The first integral before substitution is compactly supported
on \(G\times H\): the \(g\)\textquotesingle s lie in
\(\operatorname{supp}k\cap q^{-1}(q(\operatorname{supp}f))\), which is
compact, and the \(h\)\textquotesingle s then lie in a compact set as in
Lemma 2.1. Hence the compact-product theorem justifies the exchange.
Inversion on \(H\) gives the second line.

Uniqueness follows from surjectivity of \(P\) and uniqueness in Radon
representation; consequently no choice of cutoff changes the measure. A
nonempty open subset of \(Y\) has a nonzero nonnegative \(F\in C_c(Y)\)
supported there, by Lemma 1.1. Lift it to nonnegative nonzero
\(f\in C_c(G)\). Then (12), \(\rho>0\) and Haar positivity give
\(\int F\,d\mu_\rho>0\). \(\square\)

If \(\rho_1,\rho_2\) are two \(\rho\)-functions, their ratio is constant
on each coset by (11). It descends to a strictly positive continuous
function \(r\) on \(Y\), because \(q\) is a quotient map. Applying (12)
to \(r(q(g))f(g)\) and using surjectivity of \(P\) proves
\[d\mu_{\rho_2}=r\,d\mu_{\rho_1}.
\tag{15}\] Thus all the measures constructed from continuous positive
\(\rho\)-functions have the same null sets. No assertion about arbitrary
non-regular measures is needed.

\subsection{Quasi-invariance and the direction of
translation}\label{quotient-measures-and-weils-integration-formula--quasi-invariance-and-the-direction-of-translation}

Let \(T_a(y)=ay\). The ratio \[c(a,gH)=\frac{\rho(ag)}{\rho(g)}
\tag{16}\] is well-defined and continuous. Replacing \(g\) by \(gh\)
multiplies both numerator and denominator by the same factor from (11).
Direct cancellation gives \[c(ab,y)=c(a,by)c(b,y).
\tag{17}\]

\textbf{Theorem 6.1.} The measure \(\mu_\rho\) is strongly
quasi-invariant. More precisely,
\[\frac{d[E\mapsto\mu_\rho(aE)]}{d\mu_\rho}(gH)
=\frac{\rho(ag)}{\rho(g)},\qquad
\frac{d(T_a)_*\mu_\rho}{d\mu_\rho}(gH)
=\frac{\rho(a^{-1}g)}{\rho(g)}.
\tag{18}\] Both densities are strictly positive and continuous.

\emph{Proof.} For \(F\in C_c(Y)\), substitute \(g=at\) in (13), using
left Haar invariance: \[\begin{aligned}
\int_Y F(a^{-1}y)\,d\mu_\rho(y)
&=\int_G F(tH)k(at)\rho(at)\,dt\\
&=\int_G F(tH)c(a,tH)k(at)\rho(t)\,dt.
\end{aligned}\] The function \(k_a(t)=k(at)\) is another cutoff:
\(Pk_a(tH)=Pk(atH)=1\), and it has the compact-support-over-compacts
property because \(a\) acts homeomorphically on \(Y\). Cutoff
independence in Theorem 5.2 therefore identifies the last integral with
\(\int_Y F(y)c(a,y)\,d\mu_\rho(y)\).

Both sides represent Radon measures, so equality on \(C_c(Y)\) is
equality of those measures. On the left the measure of \(E\) is
\(\mu_\rho(aE)\). This proves the first formula; replace \(a\) by
\(a^{-1}\) for the ordinary pushforward convention. A strictly positive
density preserves null sets in both directions: if its integral over
\(E\) vanishes, intersect \(E\) with the sets where the density is at
least \(1/n\), then take their countable union. This proves
quasi-invariance. \(\square\)

Notice the inverse in the second formula. The two expressions describe
different measures, not an ambiguity in the modular convention.

\subsubsection{Relative invariance and a
character}\label{quotient-measures-and-weils-integration-formula--relative-invariance-character}

The distinction between invariance and quasi-invariance has a useful
intermediate case. Following the freely accessible treatment by Bekka,
de la Harpe and Valette, call a nonzero Radon measure \(\mu\)
\textbf{relatively invariant with character} \(\chi\) if
\(\chi:G\to\mathbb R_{>0}\) is a continuous homomorphism and
\[\mu(aE)=\chi(a)\mu(E)\] for every Borel \(E\subset Y\).

\textbf{Theorem 6.2.} Such a measure exists if and only if
\[\chi(h)=\frac{\Delta_H(h)}{\Delta_G(h)}\qquad(h\in H).
\tag{6.2}\] For a fixed \(\chi\), it is unique up to a positive scalar.
If it has finite positive total mass, then \(\chi=1\).

\emph{Proof.} If (6.2) holds, \(\rho=\chi\) is a rho-function. Theorem
5.2 constructs \(\mu_\chi\), and Theorem 6.1 gives \(c(a,y)=\chi(a)\),
proving relative invariance.

Conversely, suppose \(\mu\) has this property. Define
\[L(f)=\int_Y P(f/\chi)(y)\,d\mu(y)\qquad(f\in C_c(G)).\] This is a
positive nonzero functional: multiplication by \(\chi^{-1}\) preserves
\(C_c(G)\), and positive surjectivity of \(P\) supplies a nonnegative
test with positive integral against the nonzero Radon measure. Also
\[P((L_af)/\chi)(y)=\chi(a)^{-1}P(f/\chi)(a^{-1}y).\] Relative
invariance cancels these two factors, so \(L(L_af)=L(f)\). Riesz
representation and Haar uniqueness give \(L(f)=C\int_Gf\,dg\), with
\(C>0\). Substituting \(\chi f\) yields
\[\int_Y Pf\,d\mu=C\int_G\chi(g)f(g)\,dg.\] For \(h\in H\), the left
side with \(R_hf\) is multiplied by \(\Delta_H(h)^{-1}\). Right
translation on the right side gives the factor
\(\Delta_G(h)^{-1}\chi(h)^{-1}\). A positive test with nonzero integral
equates the two factors and proves (6.2). The displayed identity and
surjectivity of \(P\) also prove uniqueness up to \(C\). Finally
\(\mu(aY)=\mu(Y)\), so finite positive total mass forces \(\chi(a)=1\)
for all \(a\). \(\square\)

\subsubsection{The quasi-regular
representation}\label{quotient-measures-and-weils-integration-formula--quasi-regular-representation}

\textbf{Theorem 6.3.} For the rho-function measure, define
\[(\pi_\rho(a)\xi)(y)=c(a^{-1},y)^{1/2}\xi(a^{-1}y),
\qquad \xi\in L^2(Y,\mu_\rho).
\tag{6.3}\] Then \(\pi_\rho\) is a strongly continuous unitary
representation of \(G\). If \(\mu_{\rho_2}=r\mu_{\rho_1}\) as in (15),
multiplication by \(r^{-1/2}\) gives a unitary intertwiner from
\(\pi_{\rho_1}\) to \(\pi_{\rho_2}\).

\emph{Proof.} Quasi-invariance makes composition with \(a^{-1}\)
well-defined on equivalence classes. Theorem 6.1, first for nonnegative
simple functions and then by monotone convergence, gives the
change-of-variables rule \[\int_Y\phi(y)\,d\mu_\rho(y)
=\int_Y\phi(az)c(a,z)\,d\mu_\rho(z).\] The cocycle identity says
\(c(a^{-1},az)c(a,z)=1\). Apply the rule to
\(\phi(y)=c(a^{-1},y)|\xi(a^{-1}y)|^2\); it proves
\(\|\pi_\rho(a)\xi\|_2=\|\xi\|_2\). The same identity gives
\[c(a^{-1},y)c(b^{-1},a^{-1}y)=c((ab)^{-1},y),\] so
\(\pi_\rho(a)\pi_\rho(b)=\pi_\rho(ab)\). Its inverse is
\(\pi_\rho(a^{-1})\), hence it is unitary.

For \(\xi\in C_c(Y)\), fix a compact identity neighborhood
\(A\subset G\), with \(e\in A\). All functions \(\pi_\rho(a)\xi\),
\(a\in A\), vanish off the compact set \(A\operatorname{supp}\xi\). The
continuous function \((a,y)\mapsto c(a^{-1},y)^{1/2}\xi(a^{-1}y)\)
converges uniformly on this compact set to \(\xi(y)\) as \(a\to e\):
continuity at every \((e,y)\) and a finite cover of the compact set give
a common identity neighborhood for each prescribed error. Radon
finiteness on that compact set proves convergence in \(L^2\). The
density of \(C_c(Y)\) in \(L^2(Y,\mu_\rho)\), proved by the Haar
lesson\textquotesingle s Radon-measure Proposition 3.1(4), and the
unitary norm bound extend convergence to every \(\xi\). The
representation law gives continuity at every \(a\).

For the last assertion, \(U\xi=r^{-1/2}\xi\) satisfies
\(\int|U\xi|^2\,d\mu_{\rho_2}=\int|\xi|^2\,d\mu_{\rho_1}\) and has
inverse multiplication by \(r^{1/2}\). Since
\[c_2(a^{-1},y)=\frac{r(a^{-1}y)}{r(y)}c_1(a^{-1},y),\] substitution in
(6.3) gives \(U\pi_{\rho_1}(a)=\pi_{\rho_2}(a)U\). \(\square\)

For the affine action on the real line in the following examples,
\(c((a,b),t)=a\). Formula (6.3) consequently becomes
\[(\pi(a,b)\xi)(t)=a^{-1/2}\xi((t-b)/a).\] The square root of the
inverse scaling is exactly the factor needed for unitarity on
\(L^2(\mathbb R,dt)\).

\subsubsection{Finite covolume forces
unimodularity}\label{quotient-measures-and-weils-integration-formula--finite-covolume-unimodularity}

\textbf{Proposition 6.4.} If \(H\) is a closed unimodular subgroup and
\(G/H\) has a nonzero invariant Radon measure of finite total mass, then
\(G\) is unimodular. In particular, every group admitting a lattice is
unimodular.

\emph{Proof.} Theorem 3.1 gives \(\Delta_G|_H=1\), so
\(d(gH)=\Delta_G(g)\) is a well-defined continuous positive function on
\(Y\). Suppose \(\Delta_G(a)=s>1\); if a nontrivial modular value is
smaller than one, replace that element by its inverse. The Borel sets
\[E_j=\{y:s^j\leq d(y)<s^{j+1}\},\qquad j\in\mathbb Z,\] are disjoint
and cover \(Y\). Since \(d(ay)=s\,d(y)\), multiplication by \(a\) sends
\(E_j\) onto \(E_{j+1}\). Invariance makes their measures equal.
Positive total mass forces at least one to have positive measure,
whereas infinitely many disjoint sets of that same positive measure
contradict finite total mass. Thus \(\Delta_G=1\). A discrete subgroup
is unimodular for counting measure. It is also closed: choose an open
identity neighborhood \(V\) with \(V^{-1}V\cap\Gamma=\{e\}\). Every
translate \(xV\) meets \(\Gamma\) in at most one point. If \(x\) lies in
its closure, choose \(\gamma\in xV\cap\Gamma\); if \(x\ne\gamma\), the
open neighborhood \(xV\setminus\{\gamma\}\) would miss \(\Gamma\), a
contradiction. Hence \(x\in\Gamma\). This proves the closedness needed
to apply the proposition and completes the lattice assertion.
\(\square\)

\textbf{Corollary 6.5.} If \(H\) is closed and unimodular and \(G/H\) is
compact, then \(G\) is unimodular and \(G/H\) has a finite nonzero
invariant Radon measure. In particular, a closed discrete cocompact
subgroup is a lattice, often called a \textbf{uniform lattice}.

\emph{Proof.} The character \(\chi=\Delta_G^{-1}\) satisfies (6.2)
because \(\Delta_H=1\). Theorem 6.2 gives a nonzero relatively invariant
Radon measure. Compactness makes its total mass finite; its last
assertion therefore gives \(\chi=1\), or \(\Delta_G=1\). The measure is
invariant and has positive mass. \(\square\)

\subsubsection{Integration of quotient
densities}\label{quotient-measures-and-weils-integration-formula--quotient-density-integral}

Vogan\textquotesingle s freely accessible treatment gives an intrinsic
formulation even when an invariant scalar measure does not exist. It
regards the modular correction as part of the function being integrated.
Here is the formulation and a full derivation from the weighted Weil
formula.

\textbf{Corollary 6.6.} Let \(\mathcal D_c(G/H)\) be the continuous
functions \(F:G\to\mathbb C\), with support compact modulo \(H\),
satisfying \[F(gh)=\frac{\Delta_H(h)}{\Delta_G(h)}F(g).
\tag{6.6}\] Here support compact modulo \(H\) means that the support of
the descended function \(F/\rho\) on \(Y\) is compact; this condition is
independent of the positive rho-function. The map
\[Af(g)=\int_H f(gh)\frac{\Delta_G(h)}{\Delta_H(h)}\,dh\] sends
\(C_c(G)\) onto \(\mathcal D_c(G/H)\). There is a unique positive linear
integral \(J\) on this space with \[J(Af)=\int_G f(g)\,dg.
\tag{6.7}\] It is independent of the choice of rho-function and
invariant under \(F(g)\mapsto F(a^{-1}g)\).

\emph{Proof.} Equation (11) gives \[Af(g)=\rho(g)P(f/\rho)(gH).\] Thus
\(Af\) has the covariance (6.6) and compact quotient support.
Conversely, \(F/\rho\) descends to a member \(\phi\in C_c(Y)\). Choose
\(u\in C_c(G)\) with \(Pu=\phi\) by Lemma 2.1. Then \(f=\rho u\) is
compactly supported and \(Af=F\). Define
\[J(F)=\int_Y\frac{F(g)}{\rho(g)}\,d\mu_\rho(gH).\] The quotient
integrand is well-defined, continuous and compactly supported. Theorem
5.2 gives \(J(Af)=\int(f/\rho)\rho\,dg=\int f\,dg\). This also proves
uniqueness and positivity, using a positive lift when \(F\geq0\).

If \(\rho_2=r\rho_1\) on the quotient, then
\(\mu_{\rho_2}=r\mu_{\rho_1}\), and the factors cancel in this
definition of \(J\). Finally \(A(L_af)(g)=Af(a^{-1}g)\). Surjectivity of
\(A\) and left Haar invariance in (6.7) prove the claimed invariance of
\(J\). The integral on these covariant functions therefore survives even
when Theorem 3.1 rules out an invariant scalar measure. \(\square\)

\subsection{Examples that test the
formula}\label{quotient-measures-and-weils-integration-formula--examples-that-test-the-formula}

\textbf{Translations inside the affine group.} Write
\[G=\{(a,b):a>0,\ b\in\mathbb R\},\qquad
(a,b)(a',b')=(aa',b+ab').\] The preceding Haar lesson, Example 11.3,
calculates \(dg=a^{-2}\,da\,db\) and \(\Delta_G(a,b)=a^{-1}\). For
\(H=\{(1,t):t\in\mathbb R\}\), \(dh=dt\) and \(\Delta_H=1\). The
restriction of \(\Delta_G\) to \(H\) is one, so the quotient has an
invariant measure. It is parametrized by \(a\), and (6) becomes
\[\int_0^\infty\int_{\mathbb R}f(a,b)\,\frac{db\,da}{a^2}
=\int_0^\infty\left(\int_{\mathbb R}f(a,at)\,dt\right)\frac{da}{a}.
\tag{19}\] The substitution \(b=at\) verifies the formula and identifies
the quotient measure \(da/a\).

\textbf{Dilations inside the same group.} Instead take
\(H=\{(a,0):a>0\}\), with \(dh=da/a\). It is abelian, so \(\Delta_H=1\),
but \(\Delta_G|_H=a^{-1}\). There is no invariant quotient measure. The
quotient is parametrized by \(b\), and \(\rho(a,b)=a\) satisfies (11).
Formula (12) gives ordinary \(db\) on \(Y\):
\[\int_{\mathbb R}\int_0^\infty f(a,b)\,\frac{da}{a}\,db
=\int_G f(a,b)\rho(a,b)\,dg.\] The element \((a_0,b_0)\) acts by
\(b\mapsto b_0+a_0b\). Its pullback-on-sets density is \(a_0\), and its
pushforward density is \(a_0^{-1}\), precisely as in (18).

\textbf{Compact subgroups.} If \(H\) is compact, normalize \(dh(H)=1\).
Then \(\rho=1\) is allowed and averaging is ordinary fiberwise
probability averaging. The quotient measure need not itself be finite:
\(\mathbb R/\{0\}\) is already a counterexample. Compactness of the
fiber does not mean compactness of the base.

\textbf{The upper half-plane.} Let \(G=\operatorname{SL}_2(\mathbb R)\)
and \(H=\operatorname{SO}(2)\). The action
\[\begin{pmatrix}a&b\\c&d\end{pmatrix}z=\frac{az+b}{cz+d}\] preserves
\(\mathbb H=\{x+iy:y>0\}\), since direct subtraction of the complex
conjugate gives \[\operatorname{Im}(gz)=\frac{y}{|cz+d|^2}.\] The
denominator never vanishes on \(\mathbb H\). The element
\(s(x+iy)=\begin{pmatrix}\sqrt y&x/\sqrt y\\0&1/\sqrt y\end{pmatrix}\)
sends \(i\) to \(x+iy\). Solving \(gi=i\) gives \(a=d\), \(b=-c\),
\(a^2+b^2=1\), so the stabilizer is exactly \(H\). Thus \(gH\mapsto gi\)
is a continuous bijection \(G/H\to\mathbb H\), and its inverse is the
continuous map \(z\mapsto s(z)H\). This proves the claimed
identification, including its topology.

Here is a direct integral proof of invariance of \[dA(z)=y^{-2}\,dx\,dy.
\tag{20}\] This positive continuous density defines a Radon measure.
Upper triangular matrices act by translations \(z\mapsto z+t\) and
positive dilations \(z\mapsto a^2z\). Iterated one-variable substitution
proves that both preserve \(dA\). The substitution and
integrated-derivative rules are proved in
\hyperref[measure-and-hilbert-space-tools]{Measure and Hilbert space
tools for Haar integration}, Lemma 6.1.

For rotations
\(k_\theta=\begin{pmatrix}\cos\theta&\sin\theta\\ -\sin\theta&\cos\theta\end{pmatrix}\),
their action has infinitesimal vector field
\[v(x,y)=(1+x^2-y^2,2xy),\qquad
\partial_x(y^{-2}v_1)+\partial_y(y^{-2}v_2)=0.\] If
\(F\in C_c^1(\mathbb H)\), put \(J(\theta)=\int F(k_\theta z)\,dA(z)\).
On every compact parameter interval, the supports of these functions and
their derivatives lie in one compact subset of \(\mathbb H\): it is the
image of that interval times \(\operatorname{supp}F\) under
\((\theta,w)\mapsto k_{-\theta}w\). The derivatives are bounded there,
and \(y\) is bounded away from zero. Dominated convergence permits
differentiation. The flow identity gives
\[J'(\theta)=\int v(z)\cdot\nabla_z(F(k_\theta z))\,y^{-2}\,dx\,dy=0.\]
For the last equality, integrate each coordinate derivative of the
compactly supported product \(y^{-2}v_jF(k_\theta z)\). Its integral is
zero by the integrated-derivative rule and Fubini. The remaining term is
minus the displayed divergence, hence zero. Therefore rotations preserve
the integral.

This extends to \(C_c(\mathbb H)\) without assuming a smoothing theorem.
On a rectangle strictly inside \(\mathbb H\), polynomials uniformly
approximate a continuous function by
\hyperref[the-stone-weierstrass-theorem-for-functions-vanishing-at-infinity]{the
compact Stone--Weierstrass theorem}, Theorem 8.2. Multiply them by a
\(C^1\) cutoff equal to one on the function\textquotesingle s support
and supported in that rectangle. Such a cutoff is a product of
one-variable plateaus joined to zero by the cubic \(3t^2-2t^3\), whose
endpoint derivatives vanish. The approximants have one compact support
and converge uniformly; both their original and rotated integrals
converge because the relevant compact sets have finite area.

Finally every \(g\in G\) is \(s(z)k\), as proved by its action on \(i\);
\(s(z)\) is upper triangular. Translations, positive dilations and
rotations therefore generate the action of \(G\). This proves invariance
of (20) under every \(g\). By Theorem 3.1 it is the quotient measure up
to Haar normalization. The next lesson calculates the modular lattice
volume from this measure.

\textbf{An integer lattice.} For \(G=\mathbb R^n\), \(H=\mathbb Z^n\),
Lebesgue and counting measures give \[\int_{\mathbb R^n}f(x)\,dx
=\int_{[0,1)^n}\sum_{m\in\mathbb Z^n}f(x+m)\,dx.
\tag{21}\] For \(f\in C_c(\mathbb R^n)\), only finitely many lattice
translates meet its support while \(x\) ranges over the unit cube.
Translate each summand and use the disjoint half-open tiling of
\(\mathbb R^n\). This proves (21); the boundary faces have Lebesgue
measure zero. The resulting quotient measure has total mass one.

\subsection{Exercises with
solutions}\label{quotient-measures-and-weils-integration-formula--exercises-with-solutions}

\textbf{1. Rescaling the measures.} If \(dg\) is multiplied by \(A>0\)
and \(dh\) by \(B>0\), how does the normalized invariant quotient
measure change?

\emph{Solution.} The new averaging operator is \(BP\). To keep (6), the
quotient measure must be multiplied by \(A/B\). The same rule holds in
(12) when the function \(\rho\) is kept fixed. The modular functions do
not change under Haar rescaling.

\textbf{2. A discrete subgroup with no invariant quotient measure.} In
the affine group let \(H=\{(2^n,0):n\in\mathbb Z\}\). Prove that \(H\)
is closed and discrete, and decide whether \(G/H\) has an invariant
measure.

\emph{Solution.} Any compact interval in the \(a\)-coordinate bounded
away from zero meets \(\{2^n\}\) finitely. Around each group element one
may choose such an interval, so there is no accumulation point in \(G\),
and each point of \(H\) is isolated. Thus \(H\) is closed and discrete.
Its modular function is one, whereas \(\Delta_G(2^n,0)=2^{-n}\). Theorem
3.1 rules out a nonzero invariant quotient measure. This shows why
``discrete'' cannot replace the modular criterion.

\textbf{3. Two choices of density.} Suppose \(\rho_1,\rho_2\) satisfy
(11). Prove both (15) and the transformation rule
\[c_2(a,y)=\frac{r(ay)}{r(y)}c_1(a,y).\]

\emph{Solution.} Their ratio \(r(q(g))=\rho_2(g)/\rho_1(g)\) is
unchanged by right multiplication by \(H\), hence is continuous on the
quotient. For \(F=Pf\), the integral of \(F\) against \(r\mu_{\rho_1}\)
is, by (12),
\(\int_G f(g)r(q(g))\rho_1(g)\,dg=\int_G f(g)\rho_2(g)\,dg\).
Surjectivity of \(P\) and Radon uniqueness give (15). Substitution in
(16) gives the cocycle formula. In particular, changing the density
changes the cocycle by this explicit ratio, without changing the measure
class.

\textbf{4. Quotient averaging is equivariant.} Prove (5) and use it to
show that two invariant quotient measures normalized by (6) coincide.

\emph{Solution.} Directly,
\[P(L_a f)(gH)=\int_Hf(a^{-1}gh)\,dh=Pf(a^{-1}gH).\] For \(r\in H\),
right translation in the integration variable gives
\(\int_H f(ghr)\,dh=\Delta_H(r)^{-1}Pf(gH)\). If two measures satisfy
(6), they integrate \(Pf\) equally for every \(f\). Lemma 2.1 says every
member of \(C_c(Y)\) is such a \(Pf\); Radon uniqueness proves equality.

\textbf{5. No global compact cutoff on a noncompact base.} Show that a
cutoff \(k\) satisfying \(Pk=1\) cannot have compact support in \(G\)
when \(Y\) is noncompact.

\emph{Solution.} If \(k\) had compact support, then
\(\operatorname{supp}Pk\subset q(\operatorname{supp}k)\) would be
compact. But \(Pk=1\) has support all of \(Y\). Therefore \(Y\) would be
compact. The weaker compactness condition in (8) is exactly what permits
the global construction.

\textbf{6. Track one modular factor.} In (14), explain why the product
of the right-translation factor and the density-change factor is
\(\Delta_H(h)^{-1}\), not \(\Delta_H(h)\).

\emph{Solution.} Writing \(t=gh\) in an integral over \(g\) contributes
\(\Delta_G(h)^{-1}\). Equation (11) gives
\(\rho(th^{-1})=\Delta_G(h)\Delta_H(h)^{-1}\rho(t)\). Multiplying
cancels \(\Delta_G(h)\), leaving \(\Delta_H(h)^{-1}\). That is exactly
the factor in the inversion formula
\(\int_H a(h^{-1})\Delta_H(h)^{-1}\,dh=\int_Ha(h)\,dh\).

\subsection{What this
enables}\label{quotient-measures-and-weils-integration-formula--what-this-enables}

The invariant criterion is available to the programme\textquotesingle s
subgroup and quotient, adèle and automorphic courses. Theorem 6.3 turns
the density and cocycle formulas into a strongly continuous
quasi-regular representation, including its unitary equivalence under a
change of density. Proposition 6.4 and Corollary 6.5 explain the
unimodularity and compactness constraints on lattices. Corollary 6.6
gives an invariant integral on quotient densities. The two
inverse-translation conventions in (18) remain explicit in these action
formulas.

The principal formulas above are stated and proved on \(C_c(G)\).
Extending them to arbitrary measurable functions requires the precise
regularity and support conventions of the preceding measure-theory
lessons. The abelian Fourier course already proves its \(L^1\)
extension. This lesson does not quietly replace the general
non-\(\sigma\)-finite measure theory by a second-countable argument.

\subsection{Sources and
credit}\label{quotient-measures-and-weils-integration-formula--sources-and-credit}

\begin{itemize}
\tightlist
\item
  D. H. Fremlin,
  \href{https://www1.essex.ac.uk/maths/people/fremlin/mt4.2013/index.htm}{\emph{Measure
  Theory}, Volume 4}, develops positive subgroup averaging, quotient
  integration and the same modular criterion. His complete, locally
  determined measure convention is related to ours by
  \hyperref[haar-measure-on-locally-compact-groups--haar-measure-convention-bridge]{Proposition
  13.7 of the Haar lesson}. The compactly supported kernel argument
  above explicitly places its cutoff on the projection of the original
  support, including cosets where signed integration cancels.
\item
  David Vogan,
  \href{https://math.mit.edu/~dav/integration.pdf}{\emph{Invariant
  measures on homogeneous spaces}}, gives a complementary treatment of
  averaging and the invariant-measure criterion, and the
  quotient-density viewpoint developed in Corollary 6.6.
\item
  \emph{Haar measure on locally compact groups}, programme lesson
  originally by Claude Opus 5.5, supplies the Haar, Radon-product and
  modular proofs named above. Retained programme expression is CC0.
\item
  \emph{Quotient measure and the subgroup modular correction}, OA-FLOW
  lesson 43, credits OpenAI Codex. Its density construction and sign
  convention inform the rho-function argument here. Its historical exact
  model is not recorded.
\item
  \emph{Subgroups, quotients and annihilators}, HA-LCA lesson by GPT-6.1
  Sol, Ultra, supplies the earlier abelian compact-lift and
  positive-averaging arguments. The present lesson proves their
  general-group versions in full.
\item
  B. Bekka, P. de la Harpe and A. Valette,
  \href{https://perso.univ-rennes1.fr/bachir.bekka/KazhdanTotal.pdf}{\emph{Kazhdan's
  Property (T)}}, freely accessible author preprint dated 23 February
  2007. Its homogeneous-space and lattice treatment motivates Theorems
  6.2--6.3 and the finite-covolume results. The proofs here are
  independently written using this lesson's earlier constructions.
\end{itemize}

\section{Finite covolume and arithmetic
quotients}\label{finite-covolume-and-arithmetic-quotients}

\emph{Written by GPT-6.1 Sol (OpenAI), Ultra reasoning effort, October
2026. Self-checked by the writing AI. Independent AI review remains
separate; no human review is claimed. Original exposition and exercises
are dedicated under CC0.}

A quotient measure can exist and still have infinite total mass. A
\textbf{lattice} in a locally compact group is a discrete subgroup whose
invariant quotient measure has finite positive mass. That mass is its
\textbf{covolume}, once the Haar measures on the group and subgroup have
been fixed. We calculate three arithmetic examples and a nonabelian
Heisenberg lattice, keeping the normalizations visible: the modular
group, the additive rational adèles, and the norm-one rational idèles.

\subsection{What to know
first}\label{finite-covolume-and-arithmetic-quotients--what-to-know-first}

The preceding
\hyperref[quotient-measures-and-weils-integration-formula]{quotient
lesson}, Theorem 3.1, proves the modular criterion and the normalized
Weil formula. Its hyperbolic example proves that \(dx\,dy/y^2\) is
invariant under \(\mathrm{SL}_2(\mathbb R)\).
\hyperref[haar-measure-on-locally-compact-groups]{Haar measure on
locally compact groups}, Theorems 2.2, 5.1, 10.1 and 11.1, supplies
Riesz representation, sigma-finite Tonelli, the modular function and
inversion.

For the last two examples, read these exact programme proofs:

\begin{itemize}
\tightlist
\item
  \hyperref[NT-LOC-02]{Completions, the p-adic numbers and complete
  discretely valued fields}, Theorem 2.1 and Proposition 2.2: the field
  \(\mathbb Q_p\), its compact open ring \(\mathbb Z_p\), and its
  residue quotients. Only the rational case of the construction is
  needed.
\item
  \hyperref[NT-ADL-01]{Restricted products and profinite completions},
  Propositions 1.1--1.2: restricted product topology and its Haar
  measure; the Chinese remainder argument and the ring/residue part of
  Proposition 1.4. Character duality in that lesson is not used here.
\item
  \hyperref[NT-ADL-02]{The adèle ring of a number field}, the rational
  proof of Theorem 2.2: a half-open additive fundamental domain. We
  apply it to integration below.
\item
  \hyperref[NT-ADL-03]{Idèles and the idèle class group}, Proposition
  3.1: the idèle topology. We give the rational norm-one reduction
  explicitly, without using the class number or unit theorem.
\end{itemize}

These are specific internal proofs. The present lesson supplies the
integration and normalization arguments; it does not assume the general
number-field compactness theorem. Elementary integer arithmetic used in
the reductions is recalled with proof next.

\subsection{1. Two integer
reductions}\label{finite-covolume-and-arithmetic-quotients--1-two-integer-reductions}

\textbf{Lemma 1.1.} If \(c,d\in\mathbb Z\) are relatively prime, they
are the bottom row of a matrix in \(\mathrm{SL}_2(\mathbb Z)\). Every
nonzero rational number has a unique expression
\(\varepsilon\prod_p p^{e_p}\), with \(\varepsilon\in\{1,-1\}\), integer
exponents, and only finitely many nonzero exponents. Consequently
\[|r|_\infty\prod_p|r|_p=1.
\tag{1.1}\]

\emph{Proof.} Repeated division with remainder strictly decreases
nonnegative remainders, so the Euclidean algorithm terminates. Back
substitution expresses the greatest common divisor as an integer
combination. For coprime \(c,d\), choose \(a,b\) with \(ad-bc=1\). This
proves the first assertion and Bézout\textquotesingle s identity.

For completeness, factor a positive integer by induction: a composite
integer is a product of smaller positive integers. If a prime \(p\)
divides \(uv\) and does not divide \(u\), Bézout gives \(sp+tu=1\);
multiplying by \(v\) shows that \(p\) divides \(v\). This prime-divisor
rule proves uniqueness by cancelling a prime from two proposed
factorizations and inducting. Applying the result to numerator and
denominator gives the rational factorization. Its real absolute value is
\(\prod p^{e_p}\), and \(|r|_p=p^{-e_p}\), proving (1.1). This also
proves directly that these rational absolute values are multiplicative
and ultrametric: over a common denominator, the valuation of a sum is at
least the lesser valuation. ∎

\subsubsection{Borel fundamental domains and
covolume}\label{finite-covolume-and-arithmetic-quotients--borel-fundamental-domain}

The general quotient facts now have full proofs in
\hyperref[quotient-measures-and-weils-integration-formula--finite-covolume-unimodularity]{the
preceding lesson}: a lattice forces its ambient group to be unimodular,
and a closed discrete cocompact subgroup is a lattice. To calculate its
covolume, the following result turns a Borel set of representatives into
a measure calculation. Bekka, de la Harpe and Valette give a
complementary treatment of fundamental domains; the proof below
explicitly handles Borel representatives.

\textbf{Theorem 1.2.} Let \(G\) be a sigma-compact unimodular LCH group
and \(\Gamma\) a closed discrete subgroup, equipped with counting
measure. Then \(\Gamma\) is countable and there exists a Borel set
\(D\subset G\) meeting each right coset \(g\Gamma\) exactly once. For
the quotient measure \(\nu\) normalized by the Weil formula,
\[\nu(G/\Gamma)=\mu(D),
\tag{1.2}\] allowing infinite values. Any Borel set of unique
right-coset representatives gives the same equality.

\emph{Proof.} A compact set meets \(\Gamma\) in finitely many points:
its intersection with the closed discrete subgroup is a compact discrete
space, whose singleton cover has a finite subcover. A countable compact
cover of \(G\) therefore makes \(\Gamma\) countable.

Choose a relatively compact open identity neighborhood \(V\) with
\(V^{-1}V\cap\Gamma=\{e\}\). It exists by continuity of multiplication
and inversion and discreteness. Sigma-compactness gives a countable
cover \(G=\bigcup_{n\geq1}U_n\), with \(U_n=g_nV\): each set of a
countable compact cover has a finite subcover by translates of \(V\).
The quotient map \(q\) is open, and its restriction to each \(U_n\) is
injective, because \(u^{-1}u'\in\Gamma\) for two points of one coset. It
is consequently a homeomorphism onto the open set \(q(U_n)\).

Define \[D_n=U_n\setminus\bigcup_{m<n}U_m\Gamma,
\qquad D=\bigcup_{n\geq1}D_n.\] The saturated sets \(U_m\Gamma\) are
open, so \(D\) is Borel. Each coset has a first chart which meets it,
and in that chart exactly one point; this point lies in \(D_n\), and
later charts contribute none. Hence \(G\) is the disjoint union of
\(D\gamma\), \(\gamma\in\Gamma\).

The same integration argument works for any Borel transversal \(D\). Put
\(\nu_0(B)=\mu(D\cap q^{-1}(B))\) for Borel \(B\subset G/\Gamma\). This
is a Borel measure. For a Borel \(B\subset q(U_n)\), write
\(U_B=(q|_{U_n})^{-1}(B)\). The sets
\[D\cap U_B\gamma^{-1},\qquad \gamma\in\Gamma,\] partition
\(D\cap q^{-1}(B)\), and their right translates by \(\gamma\) partition
\(U_B\). Right Haar invariance and countable additivity therefore give
\[\nu_0(B)=\mu(U_B).
\tag{1.3}\] In particular, \(\nu_0\) is finite and Radon on every chart,
as the transport of \(\mu|_{U_n}\). For completeness, one can establish
its global Radon property without assuming a regularity theorem for
arbitrary Borel transversals. The functional
\(F\mapsto\int_DF(q(g))\,d\mu(g)\) is positive and finite on
\(C_c(G/\Gamma)\): a finite chart cover of its support bounds its
absolute integral by \(\|F\|_\infty\) times the sum of the finite
\(\mu(U_n)\) in that cover. Riesz representation gives a Radon measure
\(\nu_1\). For tests supported in a chart it agrees with the transported
measure in (1.3), so local Radon uniqueness identifies \(\nu_1\) and
\(\nu_0\) on that chart. Partitioning a Borel set among the first chart
containing each point, using the countable cover \(q(U_n)\), proves
equality on every Borel set. Thus \(\nu_0\) itself is Radon.

Finally, for \(f\in C_c(G)\), countable Tonelli and right invariance
give \[\begin{aligned}
\int_{G/\Gamma}\sum_{\gamma\in\Gamma}f(g\gamma)\,d\nu_0(g\Gamma)
&=\int_D\sum_{\gamma\in\Gamma}f(d\gamma)\,d\mu(d)\\
&=\sum_{\gamma\in\Gamma}\int_{D\gamma}f(u)\,d\mu(u)
=\int_Gf(u)\,d\mu(u).
\end{aligned}\] For complex \(f\), absolute summability follows by
applying the same calculation to \(|f|\), whose integral is finite. The
quotient lesson\textquotesingle s normalized uniqueness identifies
\(\nu_0=\nu\). Taking the measure of the whole quotient proves (1.2).
\(\square\)

\subsection{2. A fundamental region in the upper
half-plane}\label{finite-covolume-and-arithmetic-quotients--2-a-fundamental-region-in-the-upper-half-plane}

Write \(G=\mathrm{SL}_2(\mathbb R)\),
\(\Gamma=\mathrm{SL}_2(\mathbb Z)\), and \(K=\mathrm{SO}(2)\). Integer
matrices form a closed discrete subset of real matrix space, and a
compact set meets them in finitely many points. Hence \(\Gamma\) is a
closed discrete subgroup. Counting measure on \(\Gamma\) is unimodular.

We also need unimodularity of \(G\), rather than an unproved general
assertion about Lie groups. Put
\[U(t)=\begin{pmatrix}1&t\\0&1\end{pmatrix},\quad
L(t)=\begin{pmatrix}1&0\\t&1\end{pmatrix},\quad
D(a)=\begin{pmatrix}a&0\\0&a^{-1}\end{pmatrix}.\] Conjugating \(U(t)\)
by \(D(2)\) gives \(U(4t)=U(t)^4\). A homomorphism into an abelian group
is constant on conjugacy classes. Thus its positive modular value
satisfies \(\Delta(U(t))=\Delta(U(t))^4\), so it is one. Conjugation by
\(D(1/2)\) gives the same conclusion for \(L(t)\). These matrices
generate \(G\). Indeed,
\[W(a)=U(a)L(-1/a)U(a)=\begin{pmatrix}0&a\\-1/a&0\end{pmatrix},
\qquad D(a)=W(a)W(1)^{-1}\] for \(a\ne0\), and a matrix
\(\begin{pmatrix}a&b\\c&d\end{pmatrix}\) with \(a\ne0\) equals
\(L(c/a)D(a)U(b/a)\). If \(a=0\), multiplication on the left by \(U(1)\)
makes the top-left entry nonzero. Therefore \(\Delta_G=1\). The quotient
criterion now applies.

Let \(\mathbb H=\{x+iy:y>0\}\) and
\[\mathcal F=\{z\in\mathbb H:|\operatorname{Re}z|\le1/2,\ |z|\ge1\}.
\tag{2.1}\] The action is \(gz=(az+b)/(cz+d)\), with
\(\operatorname{Im}(gz)=y/|cz+d|^2\).

\textbf{Proposition 2.1.} Every \(\Gamma\)-orbit meets \(\mathcal F\).
If two points of its interior lie in the same orbit, they are equal, and
the transforming matrix is \(I\) or \(-I\).

\emph{Proof.} For fixed \(z=x+iy\), there are only finitely many integer
pairs \((c,d)\) with \(|cz+d|\le R\): the inequality bounds \(|c|y\),
and then \(|cx+d|\). Among relatively prime pairs choose one minimizing
the positive number \(|cz+d|\); the pair \((0,1)\) supplies an initial
bound. Lemma 1.1 extends the minimizing pair to \(\gamma\in\Gamma\). The
point \(w=\gamma z\) has greatest imaginary part in its orbit. An
integer translation moves its real part into \([-1/2,1/2]\) without
changing its imaginary part. If \(|w|<1\), the matrix
\(\begin{pmatrix}0&-1\\1&0\end{pmatrix}\) increases that imaginary part,
a contradiction. Thus \(w\in\mathcal F\).

For an interior point, \(|x|<1/2\) and \(y^2>1-x^2>3/4\). If \(c\ne0\),
then \(|cz+d|>1\). For \(|c|\ge2\), its square is at least \(c^2y^2>3\).
For \(|c|=1,d=0\), use \(|z|>1\). For \(|c|=1,|d|\ge1\), its square is
at least \[(1-|x|)^2+y^2>2-2|x|>1.\] Consequently a matrix with
\(c\ne0\) strictly decreases the imaginary part of an interior point. If
it sends one interior point to another, its inverse has nonzero
bottom-left entry too, and would strictly decrease it back. This is
impossible. If \(c=0\), the determinant and integrality give
\(a=d=\pm1\); the action is an integer translation. Two real parts
strictly between \(-1/2\) and \(1/2\) force that translation to be zero,
leaving \(\gamma=\pm I\). ∎

The boundary consists of two vertical rays and a circular arc. Each has
two-dimensional Lebesgue measure zero: their vertical sections are
finite except at two abscissae, and sigma-finite Tonelli applies. Their
hyperbolic area is also zero because the density \(y^{-2}\) is locally
bounded. Countably many \(\Gamma\)-translates of this boundary still
have area zero. Removing them leaves a full-measure invariant subset on
which the region\textquotesingle s interior gives unique orbit
representatives modulo \(\pm I\).

Its area is finite and explicit: \[\operatorname{area}(\mathcal F)
=\int_{-1/2}^{1/2}\int_{\sqrt{1-x^2}}^\infty\frac{dy\,dx}{y^2}
=\int_{-1/2}^{1/2}\frac{dx}{\sqrt{1-x^2}}
=\frac\pi3.
\tag{2.2}\] The last step uses the one-variable substitution
\(x=\sin t\), with \(-\pi/6\le t\le\pi/6\). The integral identities used
here are proved in the tools lesson, Lemma 6.1.

\subsection{3. From area to group
covolume}\label{finite-covolume-and-arithmetic-quotients--3-from-area-to-group-covolume}

Area in \(\mathbb H\) and Haar volume in \(G\) have different
normalizations. Let
\[s(x+iy)=\begin{pmatrix}\sqrt y&x/\sqrt y\\0&1/\sqrt y\end{pmatrix},
\quad k_\theta=\begin{pmatrix}\cos\theta&\sin\theta\\ -\sin\theta&\cos\theta\end{pmatrix}.\]
Every \(g\in G\) has a unique expression \(s(z)k\): set \(z=gi\), and
then \(s(z)^{-1}g\) fixes \(i\), so belongs to \(K\). These operations
are continuous, giving a homeomorphism \(\mathbb H\times K\to G\). Fix
the Haar normalization \[dg=\frac{dx\,dy}{y^2}\,\frac{d\theta}{2\pi}.
\tag{3.1}\] This really is left Haar measure. For \(a\in G\),
\(a s(z)=s(az)\kappa(a,z)\) for a continuous \(K\)-valued function
\(\kappa\). Left multiplication by \(a\) therefore sends \((z,k)\) to
\((az,\kappa(a,z)k)\). Integrate first over Haar probability on \(K\),
then over invariant hyperbolic area. The integral of every \(C_c(G)\)
function is unchanged. The transported Radon product is consequently
Haar.

\textbf{Theorem 3.1.} With (3.1) and counting measure on \(\Gamma\),
both \(\Gamma\backslash G\) and \(G/\Gamma\) have covolume \(\pi/6\).
The hyperbolic quotient has area \(\pi/3\). In particular, \(\Gamma\) is
a lattice.

\emph{Proof.} On the full-measure set of regular base points from
Proposition 2.1, each left orbit first selects a unique
\(z\in\mathcal F^\circ\), then selects \(k_\theta\) with
\(0\le\theta<\pi\). The remaining ambiguity \(-I\in\Gamma\) changes
\(\theta\) by \(\pi\). Thus the Borel set
\[\mathcal D=\{s(z)k_\theta:z\in\mathcal F^\circ,\ 0\le\theta<\pi\}\]
and its \(\Gamma\)-translates partition a full-Haar-measure subset of
\(G\). Omitted fibers lie over the null boundary translates; Tonelli and
(3.1) show they have Haar measure zero. The group is sigma-compact,
since \(\mathbb H\times K\) is, so ordinary sigma-finite Tonelli is
available. Hence for \(f\in C_c(G)\), \[\int_G f(g)\,dg
=\int_{\mathcal D}\sum_{\gamma\in\Gamma}f(\gamma g)\,dg.
\tag{3.2}\] First prove this for nonnegative \(f\) by the countable
partition; apply it to \(|f|\) to justify complex summation. The
discrete averaging sum is locally finite: on compact subsets, only
finitely many \(\gamma\) can move one compact set to meet another, since
those \(\gamma\) lie in a compact set intersected with \(\Gamma\).

Define a positive functional on \(C_c(\Gamma\backslash G)\) by
integrating its pullback over \(\mathcal D\). It is finite because
\[\mu_G(\mathcal D)=\operatorname{area}(\mathcal F)\cdot\frac\pi{2\pi}=\frac\pi6.
\tag{3.3}\] Riesz representation gives a Radon measure. Formula (3.2),
positive surjectivity of discrete averaging, and right invariance of
\(dg\) show that this measure is right \(G\)-invariant and has the
normalized Weil formula. Its total mass is (3.3): an increasing family
of compactly supported quotient cutoffs tends to one, since the quotient
is sigma-compact; monotone convergence on \(\mathcal D\) proves the
assertion. Such cutoffs are obtained from the compact exhaustion and
bump functions of the preceding lesson.

Finally, inversion identifies \(\Gamma\backslash G\) with \(G/\Gamma\).
The inversion formula has no density factor because \(\Delta_G=1\), so
it preserves the normalization and covolume. Equation (2.2) gives the
hyperbolic area. ∎

\textbf{Normalization example.} The measure \(d\theta\), in place of
\(d\theta/(2\pi)\), multiplies the group covolume by \(2\pi\), giving
\(\pi^2/3\). Passing instead to \(\mathrm{PSL}_2(\mathbb R)\) and
normalizing its compact stabilizer to probability gives covolume
\(\pi/3\). Indeed \(-I\) has disappeared there, so the full compact
fiber is used. A claimed number for a lattice volume is incomplete
without its group and Haar normalization.

\subsection{4. Additive rational
adèles}\label{finite-covolume-and-arithmetic-quotients--4-additive-rational-aduxe8les}

Put \[\mathbb A=\mathbb R\times\prod_p'(\mathbb Q_p,\mathbb Z_p),
\qquad Z=\prod_p\mathbb Z_p.\] Use Lebesgue measure at the real place
and additive Haar measures of mass one on each \(\mathbb Z_p\). The
restricted product provider constructs their compatible Haar measure. On
the open subgroup \(\mathbb R\times Z\) it is Lebesgue measure times
Haar probability on \(Z\).

Theorem 2.2 of the adèle lesson proves, with rational principal parts
and integer reduction, that \[D=[0,1)\times Z
\tag{4.1}\] is a Borel fundamental domain for \(\mathbb A/\mathbb Q\),
with unique representatives, and that this quotient is compact and
Hausdorff. Its proof uses only the local residue construction: subtract
each of the finitely many negative local digit strings, then subtract an
integer. A rational integral at every prime is an integer; two
representatives in (4.1) can therefore differ only by an integer whose
real absolute value is less than one, hence by zero. This also shows why
integer reduction must change the finite components as well as the real
one.

\textbf{Theorem 4.1.} The normalized invariant measure on
\(\mathbb A/\mathbb Q\) is probability, and \[\int_{\mathbb A} f(x)\,dx
=\int_D\sum_{r\in\mathbb Q}f(x+r)\,dx
\qquad(f\in C_c(\mathbb A)).
\tag{4.2}\]

\emph{Proof.} The translates \(D+r\) are a disjoint Borel partition. The
group is sigma-compact: the compact sets \([-n,n]\times(n!)^{-1}Z\),
\(n\ge1\), exhaust it. Countable additivity, translation invariance and
Tonelli give (4.2), with absolute convergence for complex \(f\).
Moreover only finitely many summands occur uniformly on \(D\): a
contributing rational lies in \(\operatorname{supp}f-\overline D\), a
compact set; a compact set meets the closed discrete \(\mathbb Q\)
finitely.

The positive functional \(F\mapsto\int_D F(x+\mathbb Q)\,dx\) on the
compact quotient is represented by a Radon measure. Equation (4.2) and
the quotient theorem identify it with the normalized invariant quotient
measure. Its total mass is \(\mu(D)=1\), since the real interval has
length one and \(Z\) has mass one. ∎

\textbf{Worked reduction.} Take real component \(11/10\), local
components \(1/4\) at \(2\), \(2/9\) at \(3\), and zero at all other
primes. Subtract \(r=1/4+2/9=17/36\). At \(2\) the remainder is
\(-2/9\in\mathbb Z_2\); at \(3\) it is \(-1/4\in\mathbb Z_3\); every
other coordinate is integral. The real remainder is \(113/180\), so it
already lies in (4.1). This is an additive representative, not an idèle
reduction.

\subsection{5. Norm-one rational
idèles}\label{finite-covolume-and-arithmetic-quotients--5-norm-one-rational-iduxe8les}

Let
\[J=\mathbb R^\times\times\prod_p'(\mathbb Q_p^\times,\mathbb Z_p^\times),
\quad |x|=|x_\infty|\prod_p|x_p|_p,
\quad J^1=\ker |\cdot|.\] Use the restricted multiplicative topology, as
proved in Proposition 3.1 of the idèle lesson. The norm is a continuous
homomorphism: on each open chart it is a finite product of continuous
local absolute values. Thus \(J^1\) is a closed LCA group. Lemma 1.1
puts the diagonal \(\mathbb Q^\times\) inside it.

\textbf{Theorem 5.1.} Put \(U=\prod_p\mathbb Z_p^\times\), embedded in
\(J^1\) with real component one. Multiplication gives a topological
group isomorphism
\[\mathbb Q^\times_{\mathrm{discrete}}\times U\longrightarrow J^1.
\tag{5.1}\] Consequently \(J^1/\mathbb Q^\times\simeq U\) is compact.
With counting measure on \(\mathbb Q^\times\) and Haar probability on
\(U\), its covolume is one and \[\int_{J^1}f(x)\,dx
=\int_U\sum_{r\in\mathbb Q^\times}f(ru)\,du.
\tag{5.2}\]

\emph{Proof.} For \(x\in J^1\), set \(e_p=v_p(x_p)\), a finitely
supported integer family. Let \(r_0=\prod_p p^{e_p}>0\). Then each
\(x_p/r_0\) is a unit. The norm-one equation gives \(|x_\infty|=r_0\).
Taking \(r=\operatorname{sign}(x_\infty)r_0\) makes the real component
of \(x/r\) one. Thus \(x=ru\) for \(u\in U\). If a diagonal rational
belongs to \(U\), all its valuations vanish, so it is \(\pm1\); its real
component is one, forcing it to be one. This proves uniqueness and
algebraic bijectivity.

The subgroup \(U\) is compact and open in \(J^1\). Compactness follows
from compactness of each \(\mathbb Z_p^\times\), a closed subset of
\(\mathbb Z_p\), and arbitrary-product compactness. For openness,
intersect \(J^1\) with the open idèle set whose real component is
positive and whose finite components are all units. The norm forces its
real component to be exactly one, so that intersection is \(U\). Its
cosets \(rU\) are therefore open and closed and have their transported
compact product topologies. This proves the topology assertion in (5.1),
including discreteness and closedness of the diagonal factor.

Transport counting measure times Haar probability through (5.1). It is
Haar on \(J^1\); on the open compact subgroup \(U\) it has mass one.
Compact supports meet only finitely many open cosets \(rU\). Finite
summation on such supports proves (5.2). Its quotient measure is Haar
probability on \(U\), establishing the covolume. ∎

\textbf{Worked reduction.} Suppose \(x_\infty=-12/5\), \(x_2=4\),
\(x_3=3\), \(x_5=1/5\), and \(x_p=1\) elsewhere. The norm is
\((12/5)(1/4)(1/3)5=1\). The rational factor in (5.1) is \(-12/5\).
Division by it gives real component one and finite components
\(-5/3,-5/4,-1/12\) at \(2,3,5\), respectively, all local units. At the
remaining primes the component is \(-5/12\), also a unit. The entire
tuple, not just its exceptional coordinates, changes.

\subsubsection{Heisenberg lattice
volume}\label{finite-covolume-and-arithmetic-quotients--heisenberg-lattice-volume}

\textbf{Example 5.2.} In the Heisenberg group \(N\) of
\hyperref[haar-measure-on-locally-compact-groups--heisenberg-haar-volume]{Example
12.3 of the Haar lesson}, the subgroup
\(\Gamma=\{(m,n,k):m,n,k\in\mathbb Z\}\) is a uniform lattice. For
Lebesgue Haar measure and counting measure on \(\Gamma\), its covolume
is one.

\emph{Proof.} The group law and inverse formula preserve integer
coordinates, so \(\Gamma\) is a subgroup. It is closed and discrete in
\(\mathbb R^3\). Let \(D=[0,1)^3\). For \(g=(x,y,z)\), define
\[\begin{aligned}
m&=\lfloor x\rfloor,& n&=\lfloor y\rfloor,\\
u&=x-m,& v&=y-n.
\end{aligned}\] \[k=\lfloor z-un\rfloor,\qquad w=z-un-k.\] Then
\(d=(u,v,w)\in D\), \(\gamma=(m,n,k)\in\Gamma\), and
\[d\gamma=(u+m,v+n,w+k+un)=(x,y,z).\] The half-open conditions first
force \(m,n\), then \(k\), so the factorization is unique. Thus \(D\) is
a Borel right-coset transversal. Its closure is compact and its quotient
image covers \(N/\Gamma\); this proves cocompactness. The group is
sigma-compact and unimodular by the Haar example, so Theorem 1.2 gives
\[\nu(N/\Gamma)=\mu(D)=\int_0^1\int_0^1\int_0^1 dz\,dy\,dx=1.\] There is
also a left transversal with the same cube: in the factorization
\(g=\gamma d\), replace the third-coordinate reduction by
\(k=\lfloor z-mv\rfloor\) and \(w=z-mv-k\). The product has third
coordinate \(k+w+mv\), proving existence and uniqueness in that
direction too. This checks the side of translation explicitly.
\(\square\)

The upper triangular model and its lattice appear in Bekka, de la Harpe
and Valette; Fischer and Ruzhansky explain its relation to homogeneous
groups. The floor reductions above supply the exact representatives and
normalization without a general theorem about nilpotent lattices.

\subsection{6. Exercises with complete
solutions}\label{finite-covolume-and-arithmetic-quotients--6-exercises-with-complete-solutions}

\textbf{1 (easy): a nonunit normalization.} Replace counting measure on
a discrete subgroup by \(b\) times counting measure. Keeping group Haar
fixed, find the covolumes in Theorems 3.1, 4.1 and 5.1.

\emph{Solution.} Averaging is multiplied by \(b\);
Weil\textquotesingle s formula therefore divides the quotient measure by
\(b\). The covolumes become \(\pi/(6b),1/b,1/b\), respectively.
Discreteness does not impose a preferred scale until counting measure is
chosen.

\textbf{2 (medium): another additive reduction.} Take real component
\(-1/5\), components \(3/8\) at \(2\), \(1/5\) at \(5\), and zero
elsewhere. Find its representative in \(D\).

\emph{Solution.} Subtract the principal parts \(q_0=3/8+1/5=23/40\). The
real remainder is \(-31/40\); subtracting the integer \(-1\) moves it to
\(9/40\). Hence the total rational subtraction is \(q=-17/40\). The
finite remainders are \(4/5\) at \(2\), \(5/8\) at \(5\), and \(17/40\)
elsewhere. All are integral at their respective primes. Uniqueness
follows from the rational-integral intersection and interval length in
(4.1).

\textbf{3 (medium): finite-index lattice volumes.} If
\(\Gamma'\le\Gamma\) has finite index \(m\), prove that its covolume in
\(G\), with counting measure, is \(m\pi/6\).

\emph{Solution.} Choose representatives \(t_1,\ldots,t_m\) with
\(\Gamma=\bigsqcup_i\Gamma't_i\). The sets \(t_i\mathcal D\) form a
fundamental domain for the left \(\Gamma'\)-action on the same
full-measure set: an equality of two representatives would contradict
the uniqueness of the \(\Gamma\)-representative and then the distinct
cosets. Left Haar invariance gives total volume \(m\mu(\mathcal D)\).
The proof of Theorem 3.1 applies to their union, and inversion gives the
right quotient too. No assumption that \(-I\in\Gamma'\) is needed; the
index already accounts for any central difference.

\textbf{4 (hard): splitting all rational idèles.} Prove
\(J\simeq\mathbb R_{>0}\times\mathbb Q^\times_{\mathrm{discrete}}\times U\),
where the first coordinate is the idèle norm. Identify
\(J/\mathbb Q^\times\) and decide whether it is compact.

\emph{Solution.} Let \(s(t)\) be the idèle with real component \(t>0\)
and all finite components one. It is a continuous homomorphism with norm
\(t\). Every \(x\in J\) is uniquely \(s(|x|)y\) with \(y\in J^1\); both
directions are continuous, by norm continuity and group operations.
Apply (5.1) to \(y\). The rational diagonal has norm one and occupies
exactly the middle factor, so the quotient is
\(\mathbb R_{>0}\times U\). It is not compact: its continuous norm image
is the noncompact \(\mathbb R_{>0}\). This distinguishes norm-one
compactness from compactness of the full idèle class group.

\subsection{References and relation to the
programme}\label{finite-covolume-and-arithmetic-quotients--references-and-relation-to-the-programme}

The adèle and local-field lessons named above are CC0 programme sources,
written by GPT-6.1 Sol at Ultra and self-checked by their writing AI.
Their local constructions and rational fundamental-domain proof are used
at the exact locators stated. This lesson\textquotesingle s reductions,
volume calculations and exercises were independently authored; their
self-check does not claim a fresh independent review of those providers.

J. S. Milne's
\href{https://www.jmilne.org/math/CourseNotes/MF.pdf}{\emph{Modular
Functions and Modular Forms}, version 1.31}, treats the same modular
region and its boundary identifications. K. Conrad's
\href{https://kconrad.math.uconn.edu/blurbs/gradnumthy/characterQ.pdf}{\emph{The
character group of Q}} gives the same half-open rational adèle
representatives through a pole-part reduction. Here the integer/CRT
reductions, unfolding, central factor and Haar normalizations are proved
explicitly. B. Bekka, P. de la Harpe and A. Valette's
\href{https://perso.univ-rennes1.fr/bachir.bekka/KazhdanTotal.pdf}{\emph{Kazhdan's
Property (T)}}, author preprint dated 23 February 2007, gives a
complementary treatment of lattices, fundamental domains and the
Heisenberg example. V. Fischer and M. Ruzhansky's
\href{https://biblio.ugent.be/publication/8585474}{\emph{Quantization on
Nilpotent Lie Groups}}, 2016, develops the Heisenberg group and its
homogeneous scaling. General number-field adèles and class groups remain
in the adèle course; Fourier duality and Poisson summation remain in the
abelian harmonic-analysis course.

\section{Weak topologies: Tychonoff, Banach--Alaoglu, Mazur, bipolars,
Krein--Milman and
Eberlein--Šmulian}\label{weak-topologies-tychonoff-banach-alaoglu-mazur-bipolars-krein-milman-and-eberlein-smulian}

\emph{Originally written and self-checked by Claude Opus 5.5
(Anthropic), October 2026. GPT-6.1 Sol (OpenAI), at the Ultra setting,
read and self-checked the full lesson and all four solutions, corrected
the zero-space cases and supplied the complex-duality and extreme-point
details, October 2026. Public domain (CC0).}

This lesson proves the facts about weak topologies that the
operator-algebra lessons use:

\begin{itemize}
\tightlist
\item
  the continuous functionals of a weak topology are exactly the
  functionals that define it;
\item
  Tychonoff\textquotesingle s theorem, and from it the Banach--Alaoglu
  theorem;
\item
  Mazur\textquotesingle s theorem: convex sets have the same weak and
  norm closures, and more generally the same closures in all locally
  convex topologies with the same continuous functionals;
\item
  annihilators and bipolars, Goldstine\textquotesingle s theorem, and
  the dual of a quotient;
\item
  second adjoints;
\item
  the Krein--Milman theorem and Milman\textquotesingle s converse;
\item
  the Eberlein--Šmulian theorem, which says that weak compactness of a
  subset of a Banach space can be tested with sequences.
\end{itemize}

It builds on the lesson
\href{https://kokunoyumeto.github.io/open-math-courses/courses/harmonic-analysis-on-locally-compact-groups/reader/programme-reading.html\#ref-5bc8dc4283ec5f38}{Hahn--Banach,
Baire and the basic theorems on Banach spaces}, cited below as \emph{the
previous lesson}.

\subsection{Conventions}\label{weak-topologies-tychonoff-banach-alaoglu-mazur-bipolars-krein-milman-and-eberlein-smulian--conventions}

As in the previous lesson. For a set \(S\) of linear functionals,
\(\bigcap_{f\in S}\ker f\) is written \(\ker S\). A \emph{net} is a
family indexed by a directed set. A subset of a topological space is
\emph{relatively compact} if its closure is compact.

\subsection[1. Weak
topologies]{\texorpdfstring{\protect\hypertarget{weak-topologies-tychonoff-banach-alaoglu-mazur-bipolars-krein-milman-and-eberlein-smulian--oa-fnd-wt-01}{}{}1.
Weak
topologies}{1. Weak topologies}}\label{weak-topologies-tychonoff-banach-alaoglu-mazur-bipolars-krein-milman-and-eberlein-smulian--OA-FND-WT-01}

Let \(X\) be a vector space and \(Y\) a vector space of linear
functionals on \(X\). The \emph{weak topology} \(\sigma(X,Y)\) is the
locally convex topology given by the seminorms \(x\mapsto|y(x)|\),
\(y\in Y\). It is the coarsest topology making every \(y\in Y\)
continuous. It is Hausdorff exactly when \(Y\) separates the points of
\(X\).

For a normed space \(E\):

\begin{itemize}
\tightlist
\item
  the \emph{weak topology} on \(E\) is \(\sigma(E,E^*)\);
\item
  the \emph{weak* topology} on \(E^*\) is \(\sigma(E^*,j(E))\), where
  \(j(x)(\varphi)=\varphi(x)\).
\end{itemize}

A net \(x_i\to x\) weakly when \(\varphi(x_i)\to\varphi(x)\) for every
\(\varphi\in E^*\), and \(\varphi_i\to\varphi\) weak* when
\(\varphi_i(x)\to\varphi(x)\) for every \(x\in E\).

\textbf{Lemma 1.1.} Let \(f,f_1,\dots,f_n\) be linear functionals on a
vector space \(X\) with \(\ker\{f_1,\dots,f_n\}\subseteq\ker f\). Then
\(f\) is a linear combination of \(f_1,\dots,f_n\).

\textbf{Proof.} The map \(\pi:X\to\mathbb K^n\),
\(x\mapsto(f_1(x),\dots,f_n(x))\), has kernel inside \(\ker f\). So
\(g(\pi(x))=f(x)\) defines a linear functional \(g\) on the subspace
\(\pi(X)\subseteq\mathbb K^n\). Extend \(g\) linearly to
\(\mathbb K^n\); then \(g(t)=\sum_ic_it_i\), and \(f=\sum_ic_if_i\).
\(\square\)

\textbf{Theorem 1.2.} The linear functionals on \(X\) that are
continuous for \(\sigma(X,Y)\) are exactly the elements of \(Y\).

\textbf{Proof.} Elements of \(Y\) are continuous by definition. Let
\(f\) be continuous. Then \(\{x:|f(x)|<1\}\) contains a basic
neighbourhood \(\{x:|y_i(x)|<\varepsilon,\ i\le n\}\) with \(y_i\in Y\).
If \(y_i(x)=0\) for all \(i\), then \(tx\) lies in this neighbourhood
for every scalar \(t\), so \(|f(tx)|<1\) for all \(t\), and \(f(x)=0\).
By Lemma 1.1, \(f\) is a linear combination of the \(y_i\), so
\(f\in Y\). \(\square\)

\textbf{Corollary 1.3.} The weak topology of a normed space has
continuous dual \(E^*\). The weak* topology of \(E^*\) has continuous
dual \(j(E)\). Both topologies are Hausdorff: \(E^*\) separates points
of \(E\) by the previous lesson, Corollary 2.3, and \(j(E)\) separates
points of \(E^*\) trivially.

\subsection[2. Tychonoff\textquotesingle s
theorem]{\texorpdfstring{\protect\hypertarget{weak-topologies-tychonoff-banach-alaoglu-mazur-bipolars-krein-milman-and-eberlein-smulian--oa-fnd-wt-02}{}{}2.
Tychonoff\textquotesingle s
theorem}{2. Tychonoff\textquotesingle s theorem}}\label{weak-topologies-tychonoff-banach-alaoglu-mazur-bipolars-krein-milman-and-eberlein-smulian--OA-FND-WT-02}

A \emph{filter} on a set \(S\) is a nonempty family \(\mathcal F\) of
subsets with three properties:

\begin{itemize}
\tightlist
\item
  \(\varnothing\notin\mathcal F\);
\item
  \(A\cap B\in\mathcal F\) whenever \(A,B\in\mathcal F\);
\item
  \(B\in\mathcal F\) whenever \(B\supseteq A\in\mathcal F\).
\end{itemize}

An \emph{ultrafilter} is a filter not properly contained in another. A
filter on a topological space \emph{converges} to \(x\) if it contains
every neighbourhood of \(x\).

\Needspace{5\baselineskip}
\textbf{Lemma 2.1.}

\begin{enumerate}
\tightlist
\item
  Every filter is contained in an ultrafilter.
\item
  A filter \(\mathcal U\) is an ultrafilter iff for every
  \(A\subseteq S\), either \(A\in\mathcal U\) or
  \(S\setminus A\in\mathcal U\).
\item
  If \(f:S\to T\) and \(\mathcal U\) is an ultrafilter on \(S\), then
  \(f_*\mathcal U=\{B\subseteq T:f^{-1}(B)\in\mathcal U\}\) is an
  ultrafilter on \(T\).
\end{enumerate}

\textbf{Proof.} 1. Filters containing the given one, ordered by
inclusion, have unions of chains as upper bounds. Zorn\textquotesingle s
lemma (previous lesson, Theorem 1.1) gives a maximal one.

\begin{enumerate}
\setcounter{enumi}{1}
\item
  Suppose \(\mathcal U\) is an ultrafilter and \(A\notin\mathcal U\). If
  \(A\cap U\neq\varnothing\) for all \(U\in\mathcal U\), then the sets
  containing some \(A\cap U\) form a filter that contains \(\mathcal U\)
  and \(A\), contradicting maximality. So \(A\cap U=\varnothing\) for
  some \(U\in\mathcal U\), and then \(S\setminus A\supseteq U\) lies in
  \(\mathcal U\). Conversely, a filter with this property is maximal: a
  larger filter would contain some \(A\notin\mathcal U\) together with
  \(S\setminus A\in\mathcal U\), hence \(\varnothing\).
\item
  \(f_*\mathcal U\) is a filter, and 2 applies because
  \(f^{-1}(T\setminus B)=S\setminus f^{-1}(B)\). \(\square\)
\end{enumerate}

\textbf{Lemma 2.2.} A topological space \(K\) is compact iff every
ultrafilter on \(K\) converges.

\textbf{Proof.} \emph{Compact implies convergence.} Suppose an
ultrafilter \(\mathcal U\) converges to no point. Every \(x\) then has
an open neighbourhood \(V_x\notin\mathcal U\), so
\(K\setminus V_x\in\mathcal U\) by Lemma 2.1(2). Finitely many
\(V_{x_1},\dots,V_{x_n}\) cover \(K\). Then
\(\bigcap_j(K\setminus V_{x_j})=\varnothing\) lies in \(\mathcal U\),
which is impossible.

\emph{Convergence implies compact.} Let \(\mathcal O\) be an open cover
with no finite subcover. The complements of finite unions of members of
\(\mathcal O\) are nonempty and closed under finite intersections. So
they generate a filter, contained in an ultrafilter \(\mathcal U\)
(Lemma 2.1(1)). Let \(\mathcal U\) converge to \(x\), and pick
\(O\in\mathcal O\) with \(x\in O\). Then \(O\in\mathcal U\) and
\(K\setminus O\in\mathcal U\), so \(\varnothing\in\mathcal U\), a
contradiction. \(\square\)

\textbf{Theorem 2.3} (Tychonoff). A product \(\prod_{i\in I}K_i\) of
compact spaces is compact in the product topology.

\textbf{Proof.} Let \(\mathcal U\) be an ultrafilter on the product and
\(\pi_i\) the projections. By Lemma 2.1(3) each \((\pi_i)_*\mathcal U\)
is an ultrafilter on \(K_i\). By Lemma 2.2 it converges to some \(x_i\);
this uses the axiom of choice to choose the limits.

The point \(x=(x_i)\) is a limit of \(\mathcal U\). A basic
neighbourhood of \(x\) has the form \(\bigcap_{i\in F}\pi_i^{-1}(V_i)\),
with \(F\) finite and \(V_i\) a neighbourhood of \(x_i\). Each
\(\pi_i^{-1}(V_i)\in\mathcal U\), since \(V_i\in(\pi_i)_*\mathcal U\),
and \(\mathcal U\) is closed under finite intersections. By Lemma 2.2
the product is compact. \(\square\)

\subsection[3. The Banach--Alaoglu
theorem]{\texorpdfstring{\protect\hypertarget{weak-topologies-tychonoff-banach-alaoglu-mazur-bipolars-krein-milman-and-eberlein-smulian--oa-fnd-wt-03}{}{}3.
The Banach--Alaoglu
theorem}{3. The Banach--Alaoglu theorem}}\label{weak-topologies-tychonoff-banach-alaoglu-mazur-bipolars-krein-milman-and-eberlein-smulian--OA-FND-WT-03}

\textbf{Theorem 3.1} (Banach--Alaoglu). For every normed space \(E\),
the closed unit ball \(B^*\) of \(E^*\) is weak* compact.

\textbf{Proof.} For \(x\in E\) let
\(D_x=\{\lambda\in\mathbb K:|\lambda|\le\|x\|\}\), and
\(P=\prod_{x\in E}D_x\), compact by Theorem 2.3. The map
\(\Phi:B^*\to P\), \(\varphi\mapsto(\varphi(x))_x\), is injective.

\(\Phi\) is a homeomorphism onto its image, with \(B^*\) carrying the
weak* topology: both topologies are the topology of pointwise
convergence.

The image is closed in \(P\). It is the set of points \((\lambda_x)\)
satisfying \(\lambda_{x+y}=\lambda_x+\lambda_y\) and
\(\lambda_{tx}=t\lambda_x\) for all \(x,y,t\), and each of these
conditions defines a closed set: such a point is a linear functional
bounded by \(\|x\|\) at \(x\), hence an element of \(B^*\).

A closed subset of a compact space is compact. \(\square\)

\textbf{Proposition 3.2.} If \(E\) is separable, the weak* topology on
\(B^*\) is metrizable.

\textbf{Proof.} Let \((x_n)\) be dense in the unit ball of \(E\). Put
\(d(\varphi,\psi)=\sum_n2^{-n}|\varphi(x_n)-\psi(x_n)|\) on \(B^*\).

\begin{itemize}
\tightlist
\item
  \(d\) is a metric: if \(d(\varphi,\psi)=0\), then \(\varphi=\psi\) on
  a dense subset of the ball, so \(\varphi=\psi\).
\item
  The identity map from \((B^*,\text{weak}^*)\) to \((B^*,d)\) is
  continuous: each term is continuous, and the series converges
  uniformly because \(|\varphi(x_n)-\psi(x_n)|\le2\).
\item
  A continuous bijection from a compact space onto a Hausdorff space is
  a homeomorphism. \(\square\)
\end{itemize}

\subsection[4. Mazur\textquotesingle s
theorem]{\texorpdfstring{\protect\hypertarget{weak-topologies-tychonoff-banach-alaoglu-mazur-bipolars-krein-milman-and-eberlein-smulian--oa-fnd-wt-04}{}{}4.
Mazur\textquotesingle s
theorem}{4. Mazur\textquotesingle s theorem}}\label{weak-topologies-tychonoff-banach-alaoglu-mazur-bipolars-krein-milman-and-eberlein-smulian--OA-FND-WT-04}

\textbf{Theorem 4.1} (Mazur). A convex subset \(C\) of a normed space
has the same closure in the weak and in the norm topology. In
particular, norm-closed convex sets are weakly closed. If \(x_n\to x\)
weakly, then some sequence of convex combinations of the \(x_n\)
converges to \(x\) in norm.

\textbf{Proof.} The weak topology is coarser, so the norm closure
\(\overline C\) lies in the weak closure. Conversely, if
\(x\notin\overline C\), the previous lesson, Corollary 6.4(2), gives
\(\varphi\in E^*\) and \(t\) with \(\operatorname{Re}\varphi\le t\) on
\(\overline C\) and \(\operatorname{Re}\varphi(x)>t\). The weakly open
set \(\{\operatorname{Re}\varphi>t\}\) contains \(x\) and misses \(C\),
so \(x\) is not in the weak closure.

For the last statement, apply this to the convex hull of
\(\{x_n:n\ge1\}\). The point \(x\) lies in its weak closure, hence in
its norm closure. \(\square\)

The same argument works in every locally convex space.

\textbf{Theorem 4.2} (Closures of convex sets). Let \((X,\tau)\) be a
locally convex space with continuous dual \(X^*\). A convex set
\(C\subseteq X\) has the same closure for \(\tau\) and for
\(\sigma(X,X^*)\). Consequently, two locally convex topologies on \(X\)
with the same continuous linear functionals have the same closed convex
sets.

\textbf{Proof.} Every \(\varphi\in X^*\) is \(\tau\)-continuous, so
\(\sigma(X,X^*)\subseteq\tau\), and the \(\tau\)-closure \(\overline C\)
lies in the \(\sigma(X,X^*)\)-closure. Let \(x\notin\overline C\). The
set \(\overline C\) is closed and convex. The previous lesson, Corollary
6.4(2), gives \(\varphi\in X^*\) and \(t\) with
\(\operatorname{Re}\varphi\le t\) on \(\overline C\) and
\(\operatorname{Re}\varphi(x)>t\). The \(\sigma(X,X^*)\)-open set
\(\{\operatorname{Re}\varphi>t\}\) contains \(x\) and misses \(C\). For
the last claim, both topologies give the closure of \(C\) that
\(\sigma(X,X^*)\) gives. \(\square\)

\subsection[5. Annihilators, bipolars and Goldstine\textquotesingle s
theorem]{\texorpdfstring{\protect\hypertarget{weak-topologies-tychonoff-banach-alaoglu-mazur-bipolars-krein-milman-and-eberlein-smulian--oa-fnd-wt-05}{}{}5.
Annihilators, bipolars and Goldstine\textquotesingle s
theorem}{5. Annihilators, bipolars and Goldstine\textquotesingle s theorem}}\label{weak-topologies-tychonoff-banach-alaoglu-mazur-bipolars-krein-milman-and-eberlein-smulian--OA-FND-WT-05}

For \(L\subseteq E\) let \(L^\perp=\{\varphi\in E^*:\varphi|_L=0\}\).
For \(N\subseteq E^*\) let
\(N_\perp=\{x\in E:\varphi(x)=0\ \forall\varphi\in N\}\).

\textbf{Theorem 5.1.} Let \(E\) be a normed space.

\begin{enumerate}
\tightlist
\item
  For a subspace \(L\subseteq E\), the norm closure of \(L\) is
  \((L^\perp)_\perp\).
\item
  For a subspace \(N\subseteq E^*\), the weak* closure of \(N\) is
  \((N_\perp)^\perp\).
\end{enumerate}

\textbf{Proof.} 1. The set \((L^\perp)_\perp\) is norm closed and
contains \(L\). If \(x\notin\overline L\), the previous lesson,
Corollary 2.3(3), gives \(\varphi\in L^\perp\) with \(\varphi(x)\ne0\),
so \(x\notin(L^\perp)_\perp\).

\begin{enumerate}
\setcounter{enumi}{1}
\tightlist
\item
  The set \((N_\perp)^\perp\) is weak* closed and contains \(N\). Let
  \(\psi\) lie outside the weak* closure \(\overline N\), a weak* closed
  subspace. The weak* topology is locally convex. By the previous
  lesson, Corollary 6.4(2), there is a weak* continuous functional \(F\)
  with \(\operatorname{Re}F(\psi)>t\ge\operatorname{Re}F(\overline N)\)
  for some \(t\).
\end{enumerate}

Since \(\overline N\) is a subspace, \(F\) vanishes on it: if
\(F(\varphi)\neq0\) for some \(\varphi\in\overline N\), suitable
multiples \(s\varphi\) would make \(\operatorname{Re}F(s\varphi)\) as
large as we like. So \(t\ge0\) and \(\operatorname{Re}F(\psi)>0\).

By Corollary 1.3, \(F=j(x)\) for some \(x\in E\). Then \(x\in N_\perp\)
and \(\psi(x)\ne0\), so \(\psi\notin(N_\perp)^\perp\). \(\square\)

\textbf{Theorem 5.2} (Goldstine). The image \(j(B)\) of the closed unit
ball \(B\) of \(E\) is weak* dense in the closed unit ball of
\(E^{**}\).

\textbf{Proof.} Let \(K\) be the weak* closure of \(j(B)\). It is convex
and lies in the unit ball of \(E^{**}\), which is weak* closed.

Suppose \(\Phi\) is in the unit ball of \(E^{**}\) but not in \(K\). The
space \((E^{**},\sigma(E^{**},E^*))\) is locally convex with continuous
dual \(E^*\) (Theorem 1.2). The previous lesson, Corollary 6.4(2), gives
\(\varphi\in E^*\) and \(t\) with
\(\operatorname{Re}\Phi(\varphi)>t\ge\operatorname{Re}j(x)(\varphi)=\operatorname{Re}\varphi(x)\)
for all \(x\in B\).

Taking the supremum over \(B\), and using that \(B\) is balanced, gives
\(\|\varphi\|\le t\). Then
\(\operatorname{Re}\Phi(\varphi)>t\ge\|\varphi\|\ge\|\Phi\|\,\|\varphi\|\),
which is impossible. \(\square\)

\textbf{Proposition 5.3} (second adjoints). For \(T\in B(E,F)\) let
\(T^*\in B(F^*,E^*)\), \(T^*\psi=\psi\circ T\), and \(T^{**}=(T^*)^*\).

\begin{enumerate}
\tightlist
\item
  \(\|T^*\|=\|T\|\) and \(\|T^{**}\|=\|T\|\).
\item
  \(T^{**}\circ j_E=j_F\circ T\).
\item
  \((ST)^{**}=S^{**}T^{**}\).
\item
  \(T^{**}\) is weak*--weak* continuous.
\end{enumerate}

\textbf{Proof.} 1.
\(\|T^*\psi\|=\sup_{\|x\|\le1}|\psi(Tx)|\le\|\psi\|\|T\|\). Conversely,
for \(\|x\|\le1\) choose \(\psi\) of norm at most \(1\) with
\(\psi(Tx)=\|Tx\|\) (previous lesson, Corollary 2.3(2)); then
\(\|Tx\|\le\|T^*\|\). Apply the same to \(T^*\).

\begin{enumerate}
\setcounter{enumi}{1}
\item
  \[\begin{gathered}
  T^{**}(j_Ex)(\psi)\\
  =j_Ex(T^*\psi)\\
  =\psi(Tx)\\
  =j_F(Tx)(\psi).
  \end{gathered}\]
\item
  \((ST)^*=T^*S^*\), so \((ST)^{**}=S^{**}T^{**}\).
\item
  If \(\Phi_i\to\Phi\) weak*, then \[\begin{gathered}
  T^{**}\Phi_i(\psi)\\
  =\Phi_i(T^*\psi)\to\Phi(T^*\psi)\\
  =T^{**}\Phi(\psi)
  \end{gathered}\] for every \(\psi\in F^*\). \(\square\)
\end{enumerate}

\textbf{Proposition 5.4} (The dual of a quotient). Let \(M\) be a closed
subspace of a normed space \(E\), with quotient map \(q:E\to E/M\) and
quotient norm \(\|x+M\|=\operatorname{dist}(x,M)\). Then
\(g\mapsto g\circ q\) is an isometric isomorphism of \((E/M)^*\) onto
\(M^\perp\).

\textbf{Proof.}

\begin{itemize}
\tightlist
\item
  \(g\circ q\) is bounded and vanishes on \(M\).
\item
  \emph{Isometry.} \(q\) maps the open unit ball of \(E\) onto the open
  unit ball of \(E/M\): if \(\|x+M\|<1\), some \(m\in M\) has
  \(\|x-m\|<1\). Hence \(\|g\circ q\|=\|g\|\).
\item
  \emph{Onto \(M^\perp\).} Let \(\varphi\in M^\perp\). Then
  \(g(x+M)=\varphi(x)\) is well defined and linear. For every
  \(m\in M\), \(|g(x+M)|=|\varphi(x-m)|\le\|\varphi\|\,\|x-m\|\). Taking
  the infimum over \(m\) gives \(|g(x+M)|\le\|\varphi\|\,\|x+M\|\), so
  \(g\in(E/M)^*\) and \(g\circ q=\varphi\). \(\square\)
\end{itemize}

\subsection[6. The Krein--Milman
theorem]{\texorpdfstring{\protect\hypertarget{weak-topologies-tychonoff-banach-alaoglu-mazur-bipolars-krein-milman-and-eberlein-smulian--oa-fnd-wt-06}{}{}6.
The Krein--Milman
theorem}{6. The Krein--Milman theorem}}\label{weak-topologies-tychonoff-banach-alaoglu-mazur-bipolars-krein-milman-and-eberlein-smulian--OA-FND-WT-06}

A point \(e\) of a convex set \(K\) is \emph{extreme} if \(e=tx+(1-t)y\)
with \(x,y\in K\) and \(0<t<1\) forces \(x=y=e\). A nonempty closed
subset \(F\subseteq K\) is a \emph{face} (an \emph{extreme subset}) if
whenever \(tx+(1-t)y\in F\) with \(x,y\in K\) and \(0<t<1\), both
\(x,y\in F\).

\textbf{Theorem 6.1} (Krein--Milman). Let \(X\) be a Hausdorff locally
convex space and \(K\subseteq X\) a nonempty compact convex set. Then
\(K\) has an extreme point, and \(K\) is the closed convex hull of its
extreme points.

\textbf{Proof.} \emph{Faces contain extreme points.} Faces of \(K\),
ordered by reverse inclusion, satisfy the hypothesis of
Zorn\textquotesingle s lemma. The intersection of a chain is nonempty,
by compactness and the finite intersection property, and it is closed
and a face. So there is a minimal face \(F\).

Suppose \(F\) has two points \(x\ne y\). Choose a continuous \(\varphi\)
with \(\operatorname{Re}\varphi(x)\ne\operatorname{Re}\varphi(y)\); this
exists by the previous lesson, Corollary 6.4(1), replacing \(\varphi\)
by \(i\varphi\) if needed. The set
\(F'=\{z\in F:\operatorname{Re}\varphi(z)=\max_F\operatorname{Re}\varphi\}\)
is then:

\begin{itemize}
\tightlist
\item
  nonempty, because \(F\) is compact;
\item
  closed, and a face of \(F\), hence of \(K\): if \(tx'+(1-t)y'\in F'\)
  with \(x',y'\in K\), then \(x',y'\in F\) because \(F\) is a face, and
  both attain the maximum, since a strict average of two values at most
  the maximum equals the maximum only if both do;
\item
  smaller than \(F\), since it cannot contain both \(x\) and \(y\).
\end{itemize}

This contradicts minimality. So \(F=\{e\}\), and \(e\) is an extreme
point of \(K\), because \(\{e\}\) is a face.

\emph{The closed convex hull.} Let \(C\) be the closed convex hull of
the extreme points; it lies in \(K\). If some \(x\in K\setminus C\)
existed, the previous lesson, Theorem 6.3, would give a continuous
\(\varphi\) and \(t\) with \(\operatorname{Re}\varphi<t\) on \(C\) and
\(\operatorname{Re}\varphi(x)>t\).

The set
\(F=\{z\in K:\operatorname{Re}\varphi(z)=\max_K\operatorname{Re}\varphi\}\)
is a face of \(K\) disjoint from \(C\), because its points have
\(\operatorname{Re}\varphi\ge\operatorname{Re}\varphi(x)>t\). By the
first part, applied to the compact convex set \(F\), it contains an
extreme point of \(F\), which is an extreme point of \(K\) because \(F\)
is a face, and so lies in \(C\). This is a contradiction. \(\square\)

\textbf{Theorem 6.2} (Milman). Let \(X\) be a Hausdorff locally convex
space, \(Q\subseteq X\) compact, and suppose that the closed convex hull
\(K\) of \(Q\) is compact. Then every extreme point of \(K\) lies in
\(Q\).

\textbf{Proof.} Let \(e\in K\) be extreme, and suppose \(e\notin Q\).

\emph{A neighbourhood that keeps \(e\) away from \(Q\).} For each
\(q\in Q\), \(e-q\neq0\). Since \(X\) is Hausdorff and locally convex,
there is a convex balanced open neighbourhood \(U_q\) of \(0\) with
\(e-q\notin U_q+U_q\). Finitely many sets \(q_j+U_{q_j}\) cover \(Q\).
Put \(W=\bigcap_jU_{q_j}\). Then \(e\notin Q+W\): if \(e=q+w\) with
\(q\in q_j+U_{q_j}\) and \(w\in W\), then \(e-q_j\in U_{q_j}+U_{q_j}\).

\begin{itemize}
\tightlist
\item
  Choose a convex balanced open neighbourhood \(W'\) of \(0\) with
  \(W'+W'\subseteq W\), and let \(V\) be the closure of \(W'\). Then
  \(V\) is closed and convex, and \(V\subseteq W'+W'\subseteq W\),
  because every point of the closure of \(W'\) lies in \(y+W'\) for some
  \(y\in W'\). So \(e\notin Q+V\).
\end{itemize}

\emph{Splitting \(K\).} Finitely many sets \(p_k+W'\), \(p_k\in Q\),
cover \(Q\). Let \(K_k\) be the closed convex hull of \(Q\cap(p_k+V)\).

\begin{itemize}
\tightlist
\item
  \(K_k\subseteq p_k+V\), because \(p_k+V\) is closed and convex. And
  \(K_k\subseteq K\) is closed, hence compact.
\item
  The convex hull \(K'\) of \(K_1\cup\dots\cup K_n\) is the image of the
  compact set \(K_1\times\dots\times K_n\times\Delta_n\), where
  \(\Delta_n\) is the simplex of weights, under the continuous map
  \((x_1,\dots,x_n,t)\mapsto\sum_kt_kx_k\). So \(K'\) is compact, hence
  closed. It is convex and contains \(Q\), so it contains \(K\).
\end{itemize}

\emph{Conclusion.} So \(e=\sum_kt_kx_k\) with \(x_k\in K_k\) and weights
\(t_k\). Since \(e\) is extreme in \(K\), \(e=x_k\) for every \(k\) with
\(t_k>0\): if \(0<t_1<1\), write \(e=t_1x_1+(1-t_1)y\) with \(y\in K\),
so \(x_1=y=e\), and continue with \(y\). Hence
\(e\in K_k\subseteq p_k+V\subseteq Q+V\) for some \(k\), a
contradiction. \(\square\)

\subsection[7. The Eberlein--Šmulian
theorem]{\texorpdfstring{\protect\hypertarget{weak-topologies-tychonoff-banach-alaoglu-mazur-bipolars-krein-milman-and-eberlein-smulian--oa-fnd-wt-07}{}{}7.
The Eberlein--Šmulian
theorem}{7. The Eberlein--Šmulian theorem}}\label{weak-topologies-tychonoff-banach-alaoglu-mazur-bipolars-krein-milman-and-eberlein-smulian--OA-FND-WT-07}

\textbf{Lemma 7.1.} Let \(E\) be a normed space and
\(M\subseteq E^{**}\) a finite-dimensional subspace. Then there are
finitely many \(\varphi_1,\dots,\varphi_m\in E^*\) of norm \(1\) with
\(\max\{0,|\Phi(\varphi_1)|,\dots,|\Phi(\varphi_m)|\}\ge\frac12\|\Phi\|\)
for every \(\Phi\in M\). The family may be empty when \(M=\{0\}\).

\textbf{Proof.} If \(M=\{0\}\), the empty family works. Otherwise the
unit sphere \(S_M\) of \(M\) is nonempty and compact (previous lesson,
Theorem 7.1). For each \(\Phi\in S_M\), choose \(\varphi_\Phi\) of norm
\(1\) with \(|\Phi(\varphi_\Phi)|>3/4\). The open sets
\(\{\Psi\in S_M:\|\Psi-\Phi\|<1/4\}\) cover \(S_M\), so finitely many
centres \(\Phi_1,\dots,\Phi_m\) suffice. For \(\Psi\in S_M\) with
\(\|\Psi-\Phi_k\|<1/4\),
\(|\Psi(\varphi_{\Phi_k})|\ge|\Phi_k(\varphi_{\Phi_k})|-\|\Psi-\Phi_k\|>1/2\).
Scale for general \(\Phi\in M\). \(\square\)

\textbf{Theorem 7.2} (Eberlein--Šmulian). For a subset \(A\) of a Banach
space \(E\), the following are equivalent:

\begin{enumerate}
\tightlist
\item
  \(A\) is relatively weakly compact;
\item
  every sequence in \(A\) has a subsequence that converges weakly to a
  point of \(E\);
\item
  every sequence in \(A\) has a weak cluster point in \(E\).
\end{enumerate}

\textbf{Proof.} If \(E=\{0\}\) or \(A=\varnothing\), all three
assertions are immediate. Suppose otherwise. \emph{2 implies 3.} The
limit of a weakly convergent subsequence is a weak cluster point.

\emph{1 implies 3.} A sequence in the compact set \(\overline A^{\,w}\),
the weak closure of \(A\), has a cluster point there. The terms
\(\{x_n:n\ge m\}\) form a decreasing family of sets with the finite
intersection property, so their closures meet.

\emph{3 implies 1.}

\textbf{Step 1: \(A\) is bounded.} Otherwise there is \(\varphi\in E^*\)
with \(\sup_{x\in A}|\varphi(x)|=\infty\), by the previous lesson,
Corollary 4.3(2). So there are \(x_n\in A\) with
\(|\varphi(x_n)|\ge n\). A weak cluster point \(x\) of \((x_n)\) would
make \(\varphi(x)\) a cluster point of \((\varphi(x_n))\), which has
none.

\textbf{Step 2: the weak* closure of \(j(A)\) lies in \(j(E)\).} The
weak* closure \(\overline{j(A)}^{w*}\) is weak* compact by Theorem 3.1,
since \(A\) is bounded. Let \(\Phi\) be in it. We build \(x_n\in A\) and
functionals of norm \(1\) inductively.

\begin{itemize}
\tightlist
\item
  Let \(\varphi_1\) be any functional of norm \(1\), and choose
  \(x_1\in A\) with \(|(\Phi-jx_1)(\varphi_1)|<1\).
\item
  Given \(x_1,\dots,x_n\), let
  \(M_n=\operatorname{span}\{\Phi,\Phi-jx_1,\dots,\Phi-jx_n\}\). Lemma
  7.1 gives finitely many functionals of norm \(1\) that norm \(M_n\) up
  to the factor \(\frac12\). Append them to the list
  \(\varphi_1,\varphi_2,\dots\).
\item
  Since \(\Phi\) lies in the weak* closure of \(j(A)\), choose
  \(x_{n+1}\in A\) with \(|(\Phi-jx_{n+1})(\varphi_k)|<1/(n+1)\) for all
  functionals \(\varphi_k\) listed so far.
\end{itemize}

Let \(x\) be a weak cluster point of \((x_n)\).

\textbf{Step 3: \(\Phi=j(x)\).}

\begin{itemize}
\tightlist
\item
  \emph{\(\Psi=\Phi-jx\) lies in \(\overline{\bigcup_nM_n}\).} The point
  \(x\) is in the weak closure of \(\{x_n\}\), hence in the norm-closed
  linear span of \(\{x_n\}\) by Theorem 4.1. So \(jx\) is a norm limit
  of combinations \(\sum c_kjx_k\). Then \(\Psi\) is a norm limit of
  \(\Phi-\sum c_kjx_k\). Each of these lies in \(M_n\) for large \(n\):
  write \[\begin{gathered}
  \Phi-\sum c_kjx_k\\
  =(1-\sum c_k)\Phi+\sum c_k(\Phi-jx_k).
  \end{gathered}\]
\item
  \emph{\(\|\Psi\|\le2\sup_k|\Psi(\varphi_k)|\).} For \(G\in M_n\),
  \(\|G\|\le2\max_k|G(\varphi_k)|\), with \(k\) running over the
  functionals chosen for \(M_n\). Approximating \(\Psi\) by such \(G\)
  gives \(\|\Psi\|\le2\sup_k|\Psi(\varphi_k)|+3\|\Psi-G\|\), and the
  last term tends to \(0\).
\item
  \emph{\(\Psi(\varphi_k)=0\) for every \(k\).} For \(n\) large,
  \(|\Phi(\varphi_k)-\varphi_k(x_n)|<1/n\). Since \(\varphi_k(x)\) is a
  cluster point of \((\varphi_k(x_n))\), we get
  \(\Phi(\varphi_k)=\varphi_k(x)\).
\end{itemize}

Hence \(\Psi=0\), that is, \(\Phi=jx\in j(E)\).

\textbf{Step 4: conclusion.} The map \(j\) is a homeomorphism from
\((E,\text{weak})\) onto \((j(E),\text{weak}^*)\): both topologies are
pointwise convergence on \(E^*\). By Step 2,
\(\overline{j(A)}^{w*}=j(K)\) for some \(K\subseteq E\). This \(K\) is
weakly compact and contains \(A\), so \(A\) is relatively weakly
compact.

\emph{1 implies 2.} Let \((x_n)\) be a sequence in \(A\).

\begin{itemize}
\tightlist
\item
  Let \(E_0\) be the norm-closed span of \(\{x_n\}\), a separable closed
  subspace. By Theorem 4.1 it is weakly closed. The weak closure \(K\)
  of \(\{x_n\}\) is weakly compact by 1 and lies in \(E_0\).
\item
  Let \((y_m)\) be dense in \(E_0\). Choose \(\psi_m\in E^*\) of norm
  \(1\) with \(\psi_m(y_m)=\|y_m\|\). The family \((\psi_m)\) separates
  the points of \(E_0\): if \(z\ne0\) in \(E_0\), choose \(y_m\) with
  \(\|z-y_m\|<\|z\|/3\); then
  \(|\psi_m(z)|\ge\|y_m\|-\|z-y_m\|>\|z\|/3\).
\item
  By 1 implies 3 and Step 1, \(A\) is norm bounded. The weakly closed
  norm ball containing it also contains \(K\), so this series is
  uniformly convergent on \(K\times K\). The topology on \(K\) induced
  by \(d(u,v)=\sum_m2^{-m}|\psi_m(u-v)|\) is Hausdorff and coarser than
  the weak topology. A compact topology admits no strictly coarser
  Hausdorff topology, so the weak topology on \(K\) is this metric
  topology.
\item
  In a compact metric space the closures of the sequence tails have a
  common point \(x\), by the finite intersection property. Every ball
  about \(x\) contains arbitrarily late terms. Choose successively
  increasing indices whose terms lie in the balls of radii \(1/k\); the
  subsequence converges in the metric. Thus \((x_n)\) has a subsequence
  converging weakly in \(K\subseteq E\). \(\square\)
\end{itemize}

\subsection{Exercises}\label{weak-topologies-tychonoff-banach-alaoglu-mazur-bipolars-krein-milman-and-eberlein-smulian--exercises}

\textbf{Exercise 1} (medium). Show that in an infinite-dimensional
normed space the weak closure of the unit sphere \(\{\|x\|=1\}\) is the
closed unit ball.

\emph{Solution.}

\begin{itemize}
\tightlist
\item
  \emph{The closure lies in the ball:} the closed unit ball is convex
  and norm closed, so weakly closed (Theorem 4.1).
\item
  \emph{Every point with \(\|x_0\|<1\) is in the closure.} A basic weak
  neighbourhood \(U=\{x:|\varphi_i(x-x_0)|<\varepsilon,\ i\le n\}\)
  contains \(x_0+\ker\{\varphi_1,\dots,\varphi_n\}\). This kernel is a
  nonzero subspace, because a linear map from an infinite-dimensional
  space to \(\mathbb K^n\) has nonzero kernel. Pick \(y\ne0\) in it. The
  function \(t\mapsto\|x_0+ty\|\) is continuous, equals \(\|x_0\|<1\) at
  \(t=0\), and tends to infinity. Choose \(T>0\) with \(\|x_0+Ty\|>1\).
  Bisect the interval \([0,T]\), retaining endpoint values that bracket
  one. Its nested endpoints converge to a common real number \(t\);
  continuity gives \(\|x_0+ty\|=1\). This produces a point of the sphere
  in \(U\).
\item
  \emph{Every point with \(\|x_0\|=1\) is on the sphere itself.}
  \(\square\)
\end{itemize}

\textbf{Exercise 2} (medium). Let \(E=c_0\), the null sequences with the
supremum norm. Show that \(\delta_n\to0\) weakly but not in norm, and
find convex combinations of the \(\delta_n\) that converge to \(0\) in
norm.

\emph{Solution.}

\begin{itemize}
\tightlist
\item
  \emph{Weak convergence:} put \(a_k=\varphi(\delta_k)\) for
  \(\varphi\in c_0^*\). On a finite set \(F\), choose
  \(\theta_k=\overline{a_k}/|a_k|\) when \(a_k\ne0\), and zero
  otherwise. Then \(\|\sum_{k\in F}\theta_k\delta_k\|_\infty\leq1\), so
  \(\sum_{k\in F}|a_k|\leq\|\varphi\|\). Hence \(a\in\ell^1\). Finite
  cutoffs approximate each \(x\in c_0\) in supremum norm, giving
  \(\varphi(x)=\sum_ka_kx_k\). Conversely this series defines a bounded
  functional for every \(a\in\ell^1\), with norm at most \(\|a\|_1\);
  the same finite phase tests give the reverse inequality. Thus
  \(c_0^*=\ell^1\) isometrically, and \(\varphi_a(\delta_n)=a_n\to0\).
\item
  \emph{No norm convergence:} \(\|\delta_n\|=1\).
\item
  \emph{Convex combinations:} \(\frac1N(\delta_1+\dots+\delta_N)\) has
  norm \(1/N\to0\), as Theorem 4.1 predicts. \(\square\)
\end{itemize}

\textbf{Exercise 3} (medium). Show that the closed unit ball of
\(\ell^1\) has extreme points \(\{\lambda\delta_n:|\lambda|=1\}\), but
that the closed unit ball of \(c_0\) has none. Deduce that \(c_0\) is
not isometrically isomorphic to the dual of any normed space.

\emph{Solution.}

\begin{itemize}
\tightlist
\item
  \emph{\(\ell^1\):} if \(\|x\|_1<1\), choose
  \(0<\varepsilon<1-\|x\|_1\); the distinct vectors
  \(x\pm\varepsilon\delta_1\) belong to the ball and average to \(x\).
  If \(\|x\|_1=1\) and \(x_i,x_j\ne0\) with \(i\ne j\), put
  \(u_i=x_i/|x_i|\), \(u_j=x_j/|x_j|\), and choose
  \(0<\varepsilon<\min(|x_i|,|x_j|)\). The two vectors
  \(x\pm\varepsilon(u_i\delta_i-u_j\delta_j)\) both have norm one, since
  their two altered coordinate magnitudes add to \(|x_i|+|x_j|\). They
  are distinct and average to \(x\). The remaining points are
  \(\lambda\delta_n\), \(|\lambda|=1\), and these are extreme: if
  \(\lambda\delta_n=(x+y)/2\) with \(\|x\|_1,\|y\|_1\leq1\), multiply
  the nth coordinate by \(\bar\lambda\). The real parts of both
  resulting numbers are at most one and average to one, so both numbers
  equal one. Hence \(x_n=y_n=\lambda\), and the norm bounds force every
  other coordinate of \(x,y\) to vanish.
\item
  \emph{\(c_0\):} given \(x\) in the unit ball, \(|x_k|<1/2\) for some
  \(k\). Then \(x\pm\frac12\delta_k\) are in the ball and average to
  \(x\).
\item
  \emph{Deduction:} the unit ball of a dual space is weak* compact and
  convex (Theorem 3.1), so by Theorem 6.1 it has extreme points. A
  linear isometric isomorphism preserves convex combinations and
  therefore extreme points of balls. \(\square\)
\end{itemize}

\textbf{Exercise 4} (easy). Show that a reflexive Banach space (one with
\(j(E)=E^{**}\)) has a weakly compact closed unit ball, and deduce that
every bounded sequence in it has a weakly convergent subsequence.

\emph{Solution.}

\begin{itemize}
\tightlist
\item
  \emph{Weak compactness:} the ball of \(E^{**}=j(E)\) is weak* compact
  by Theorem 3.1, applied to \(E^*\). By Step 4 of the proof of Theorem
  7.2, \(j\) is a homeomorphism from the weak topology to the weak*
  topology. So the unit ball of \(E\) is weakly compact.
\item
  \emph{Sequences:} Theorem 7.2, 1 ⇒ 2, gives the weakly convergent
  subsequence. \(\square\)
\end{itemize}

\subsection{Where this
leads}\label{weak-topologies-tychonoff-banach-alaoglu-mazur-bipolars-krein-milman-and-eberlein-smulian--where-this-leads}

The lesson \hyperref[hilbert-spaces-and-compact-operators]{Hilbert
spaces and compact operators} uses these results for operators on
Hilbert space. The operator topologies on \(B(H)\) are weak topologies
in the sense of Section 1, and the predual of a von Neumann algebra is
studied with Theorems 5.1 and 7.2 in
\href{https://kokunoyumeto.github.io/open-math-courses/courses/harmonic-analysis-on-locally-compact-groups/reader/programme-reading.html\#ref-40074ee66678d763}{Polar
decomposition of functionals and weak compactness in preduals}.

\subsection{References}\label{weak-topologies-tychonoff-banach-alaoglu-mazur-bipolars-krein-milman-and-eberlein-smulian--references}

The results are classical: Tychonoff (1930, 1935), Banach (1932) and
Alaoglu (1940), Mazur (1933), Goldstine (1938), Krein and Milman (1940),
Eberlein (1947) and Šmulian (1940). The proof of the Eberlein--Šmulian
theorem follows R. Whitley, "An elementary proof of the
Eberlein--Šmulian theorem", \emph{Mathematische Annalen} 172 (1967)
116--118, written here in our own words.

\section{Hilbert spaces and compact
operators}\label{hilbert-spaces-and-compact-operators}

\emph{Originally written by Claude Opus 5.5 (Anthropic), October 2026,
and self-checked by that writing AI. GPT-6.1 Sol (OpenAI), at the Ultra
setting, checked the full lesson and its solutions, supplied the metric
compactness tools and repaired the measure prerequisites, October 2026.
Public domain (CC0). GPT-6.1 Sol (OpenAI), at the Ultra setting, also
supplied and self-checked the complete normed-space and inner-product
completion construction and bounded-map extension, October 2026.}

This lesson proves the Hilbert-space facts that the operator-algebra
lessons use. It covers:

\begin{itemize}
\tightlist
\item
  the Cauchy--Schwarz inequality for positive semidefinite forms;
\item
  the projection theorem and the Riesz--Fréchet theorem;
\item
  bounded sesquilinear forms and adjoints;
\item
  orthonormal bases: Bessel\textquotesingle s inequality,
  Parseval\textquotesingle s identity, existence and Hilbert dimension;
\item
  compact operators: they form a closed ideal, closed under adjoints, in
  which the finite-rank operators are dense;
\item
  the spectral theorem for compact self-adjoint operators;
\item
  the Riesz theory of compact operators: for \(\lambda\ne0\),
  \(T-\lambda\) has finite-dimensional kernel and closed range; it is
  invertible when injective; and nonzero points of the spectrum are
  isolated eigenvalues;
\item
  Hilbert tensor products;
\item
  multiplication operators on a \(\sigma\)-finite measure space, which
  form an algebra equal to its own commutant.
\end{itemize}

It builds on
\href{https://kokunoyumeto.github.io/open-math-courses/courses/harmonic-analysis-on-locally-compact-groups/reader/programme-reading.html\#ref-5bc8dc4283ec5f38}{Hahn--Banach,
Baire and the basic theorems on Banach spaces}, cited as \emph{the first
lesson}.

\subsection{Conventions}\label{hilbert-spaces-and-compact-operators--conventions}

Hilbert spaces are complex unless a statement says otherwise. No
separability is assumed. Inner products are linear in the first
variable. \(B(H)\) is the algebra of bounded operators. The
\emph{spectrum} \(\sigma(T)\) of \(T\in B(H)\) is the set of
\(\lambda\in\mathbb C\) for which \(T-\lambda\) has no inverse in
\(B(H)\). For \(\xi,\eta\in H\), \(\theta_{\xi,\eta}\) is the rank-one
operator \(\zeta\mapsto\langle\zeta,\eta\rangle\xi\).
\(L^\perp=\{\xi:\langle\xi,\eta\rangle=0\ \forall\eta\in L\}\).

\subsection[1. Positive semidefinite
forms]{\texorpdfstring{\protect\hypertarget{hilbert-spaces-and-compact-operators--oa-fnd-hs-01}{}{}1.
Positive semidefinite
forms}{1. Positive semidefinite forms}}\label{hilbert-spaces-and-compact-operators--OA-FND-HS-01}

A \emph{sesquilinear form} on a complex vector space \(V\) is a map
\(B:V\times V\to\mathbb C\), linear in the first variable and
conjugate-linear in the second. It is \emph{Hermitian} if
\(B(\eta,\xi)=\overline{B(\xi,\eta)}\), and \emph{positive semidefinite}
if \(B(\xi,\xi)\ge0\) for all \(\xi\).

\Needspace{5\baselineskip}
\textbf{Proposition 1.1.}

\begin{enumerate}
\tightlist
\item
  (Polarization)
  \(4B(\xi,\eta)=\sum_{k=0}^3i^kB(\xi+i^k\eta,\xi+i^k\eta)\).
\item
  A form with \(B(\xi,\xi)\in\mathbb R\) for all \(\xi\) is Hermitian.
\item
  (Cauchy--Schwarz) If \(B\) is positive semidefinite, then
  \(|B(\xi,\eta)|^2\le B(\xi,\xi)B(\eta,\eta)\). Consequently
  \(\{\xi:B(\xi,\xi)=0\}\) is a subspace, and
  \(\xi\mapsto B(\xi,\xi)^{1/2}\) is a seminorm.
\end{enumerate}

\textbf{Proof.}

\begin{enumerate}
\tightlist
\item
  Expand each term by sesquilinearity: the terms \(B(\xi,\xi)\) and
  \(B(\eta,\eta)\) cancel, and the cross terms give \(4B(\xi,\eta)\).
\item
  Apply 1 to \(B(\eta,\xi)\) and to \(\overline{B(\xi,\eta)}\), and
  compare the expressions term by term, using that the diagonal values
  are real.
\item
  By 2, \(B\) is Hermitian. For \(t\in\mathbb R\) and \(\theta\) with
  \(e^{-i\theta}B(\xi,\eta)=|B(\xi,\eta)|\), the quantity
  \[\begin{gathered}
  0\\
  \le B(\xi+te^{i\theta}\eta,\ \xi+te^{i\theta}\eta)\\
  =B(\xi,\xi)+2t|B(\xi,\eta)|+t^2B(\eta,\eta)
  \end{gathered}\] is a nonnegative quadratic polynomial in \(t\). If
  \(B(\eta,\eta)>0\), its discriminant is \(\le0\). If
  \(B(\eta,\eta)=0\), the linear polynomial
  \(B(\xi,\xi)+2t|B(\xi,\eta)|\) is nonnegative for all \(t\), so
  \(|B(\xi,\eta)|=0\).
\end{enumerate}

For the consequence, Cauchy--Schwarz gives
\(B(\xi+\eta,\xi+\eta)\le(B(\xi,\xi)^{1/2}+B(\eta,\eta)^{1/2})^2\).
\(\square\)

An inner product is a positive definite Hermitian form. A \emph{Hilbert
space} is an inner product space that is complete for
\(\|\xi\|=\langle\xi,\xi\rangle^{1/2}\). The \emph{parallelogram law}
\(\|\xi+\eta\|^2+\|\xi-\eta\|^2=2\|\xi\|^2+2\|\eta\|^2\) follows by
expanding.

\subsubsection{Completing normed and inner product
spaces}\label{hilbert-spaces-and-compact-operators--completing-normed-and-inner-product-spaces}

\textbf{Lemma (Completion).} A normed space \(E\) embeds linearly and
isometrically as a dense subspace of a Banach space \(\widehat E\). If
its norm comes from an inner product, that product extends to
\(\widehat E\), which is then a Hilbert space. Every bounded linear map
\(T:E\to F\), with \(F\) a Banach space, has a unique bounded linear
extension to \(\widehat E\), of the same norm. In particular the Hilbert
construction applies to a positive semidefinite form after taking its
quotient by the null subspace.

\emph{Proof.} Let \(\mathscr S\) be the vector space of Cauchy sequences
in \(E\), with termwise operations. Identify \(a=(a_n)\) and \(b=(b_n)\)
when \(\|a_n-b_n\|\to0\). The triangle inequality shows that this is an
equivalence relation compatible with the vector operations. The reverse
triangle inequality gives

\[\begin{gathered}
|\|a_n\|-\|a_m\||\leq\|a_n-a_m\|.
\end{gathered}\]

So these norms have a limit, unchanged on replacing \(a\) by an
equivalent sequence. Set \(\widehat E=\mathscr S/\mathord\sim\) and
\(\|[a]\|=\lim_n\|a_n\|\). Homogeneity and the triangle inequality pass
to this limit. Its value is zero exactly when \(a\) is equivalent to the
zero sequence, so this is a norm. Constant sequences give the asserted
linear isometry \(E\to\widehat E\).

The constants \(a_n\) converge to \([a]\): if
\(\|a_n-a_m\|<\varepsilon\) for all sufficiently large \(n,m\), taking
the limit in \(m\) gives \(\|a_n-[a]\|\leq\varepsilon\). This proves
density. If \((z_k)\) is Cauchy in \(\widehat E\), choose \(y_k\in E\)
with \(\|y_k-z_k\|<1/k\). Then

\[\|y_k-y_l\|\leq1/k+\|z_k-z_l\|+1/l.\]

So \((y_k)\) is Cauchy in \(E\), and represents \(z=[(y_k)]\). The
density argument gives \(y_k\to z\); therefore
\(\|z_k-z\|\leq1/k+\|y_k-z\|\to0\). This proves completeness.

If \(E\) has an inner product, its Cauchy sequences are bounded: a
sufficiently late tail lies within distance \(1\) of one term, and the
remaining initial block is finite. Cauchy--Schwarz therefore shows that
the inner products have limits, using

\[\begin{gathered}
|\langle a_n,b_n\rangle-\langle a_m,b_m\rangle|\\
\leq\|a_n-a_m\|\,\|b_n\|\\
+\|a_m\|\,\|b_n-b_m\|.
\end{gathered}\]

The same estimate makes the limit independent of the representatives.
Sesquilinearity and symmetry pass to the limit; its diagonal value is
\(\lim_n\|a_n\|^2=\|[a]\|^2\), so it is a positive definite inner
product inducing the complete norm.

For the extension put \(\widehat T[a]=\lim_nTa_n\). The limit exists in
\(F\) because \(\|Ta_n-Ta_m\|\leq\|T\|\|a_n-a_m\|\). Equivalent
sequences give the same limit by that bound. Linearity passes to limits,
and \(\|\widehat T[a]\|\leq\|T\|\|[a]\|\). Restriction to the constant
sequences gives the reverse norm inequality. Density gives uniqueness.
Finally Proposition 1.1 shows that a null vector for a positive
semidefinite form is orthogonal to every vector, so the form descends to
a positive definite inner product on its quotient. The construction just
proved then applies. \(\square\)

\subsection[2. The projection theorem and the Riesz--Fréchet
theorem]{\texorpdfstring{\protect\hypertarget{hilbert-spaces-and-compact-operators--oa-fnd-hs-02}{}{}2.
The projection theorem and the Riesz--Fréchet
theorem}{2. The projection theorem and the Riesz--Fréchet theorem}}\label{hilbert-spaces-and-compact-operators--OA-FND-HS-02}

\textbf{Theorem 2.1.} Let \(C\) be a nonempty closed convex subset of a
real or complex Hilbert space \(H\), and \(\xi\in H\). There is exactly
one \(c_0\in C\) with \(\|\xi-c_0\|=\operatorname{dist}(\xi,C)\).

\textbf{Proof.} Let \(d=\operatorname{dist}(\xi,C)\) and \(c_n\in C\)
with \(\|\xi-c_n\|\to d\). The parallelogram law applied to \(\xi-c_n\)
and \(\xi-c_m\), together with \(\frac{c_n+c_m}2\in C\), gives
\[\begin{gathered}
\|c_n-c_m\|^2\\
=2\|\xi-c_n\|^2+2\|\xi-c_m\|^2\\
-4\Big\|\xi-\frac{c_n+c_m}2\Big\|^2\\
\le2\|\xi-c_n\|^2+2\|\xi-c_m\|^2-4d^2\to0.
\end{gathered}\] So \((c_n)\) is Cauchy. Its limit \(c_0\in C\) attains
the distance. Two minimizers form a minimizing sequence when alternated,
so they are equal. \(\square\)

\textbf{Theorem 2.2} (projection theorem). Let \(L\) be a closed
subspace of a real or complex Hilbert space \(H\). Then:

\begin{enumerate}
\tightlist
\item
  \(H=L\oplus L^\perp\);
\item
  the map \(P:\xi\mapsto c_0\) of Theorem 2.1 is the orthogonal
  projection onto \(L\), a bounded self-adjoint idempotent with
  \(\|P\|\le1\);
\item
  for every subspace \(Y\), \(\overline Y=(Y^\perp)^\perp\).
\end{enumerate}

\textbf{Proof.} Let \(c_0\) be the nearest point of \(L\) to \(\xi\).
For \(\eta\in L\) and scalars \(t\),
\(\|\xi-c_0-t\eta\|^2\ge\|\xi-c_0\|^2\). Expanding gives
\(-2\operatorname{Re}(\bar t\langle\xi-c_0,\eta\rangle)+|t|^2\|\eta\|^2\ge0\).
Small \(t\) of suitable phase force \(\langle\xi-c_0,\eta\rangle=0\). So
\(\xi=c_0+(\xi-c_0)\) with \(\xi-c_0\in L^\perp\). The decomposition is
unique because \(L\cap L^\perp=\{0\}\).

The projection is linear and idempotent, with
\(\|P\xi\|^2+\|(1-P)\xi\|^2=\|\xi\|^2\). It is self-adjoint because
\(\langle P\xi,\eta\rangle=\langle P\xi,P\eta\rangle=\langle\xi,P\eta\rangle\).

For 3: \((Y^\perp)^\perp\) is closed and contains \(Y\). By 1, applied
to \(\overline Y\), we have \(H=\overline Y\oplus\overline Y^\perp\) and
\(Y^\perp=\overline Y^\perp\). An element of \((Y^\perp)^\perp\) has
zero component in \(\overline Y^\perp\). \(\square\)

\textbf{Theorem 2.3} (Riesz--Fréchet). Every bounded linear functional
\(f\) on \(H\) is \(f(\xi)=\langle\xi,\eta\rangle\) for exactly one
\(\eta\in H\), and \(\|f\|=\|\eta\|\).

\textbf{Proof.} If \(f=0\), take \(\eta=0\). Otherwise \(\ker f\) is a
closed proper subspace. By Theorem 2.2 there is a unit vector
\(u\perp\ker f\). For every \(\xi\), the vector \(f(\xi)u-f(u)\xi\) lies
in \(\ker f\), so it is orthogonal to \(u\). This gives
\(f(\xi)=f(u)\langle\xi,u\rangle=\langle\xi,\overline{f(u)}u\rangle\).

Uniqueness: \(\langle\xi,\eta-\eta'\rangle=0\) for all \(\xi\) forces
\(\eta=\eta'\). The norm equality is Cauchy--Schwarz together with
\(\xi=\eta\). \(\square\)

\subsection[3. Sesquilinear forms and
adjoints]{\texorpdfstring{\protect\hypertarget{hilbert-spaces-and-compact-operators--oa-fnd-hs-03}{}{}3.
Sesquilinear forms and
adjoints}{3. Sesquilinear forms and adjoints}}\label{hilbert-spaces-and-compact-operators--OA-FND-HS-03}

\textbf{Theorem 3.1.} Let \(B\) be a sesquilinear form on \(H\) with
\(|B(\xi,\eta)|\le C\|\xi\|\|\eta\|\). There is exactly one
\(t\in B(H)\) with \(B(\xi,\eta)=\langle t\xi,\eta\rangle\), and
\(\|t\|=\sup\{|B(\xi,\eta)|:\|\xi\|,\|\eta\|\le1\}\le C\). If
\(B(\xi,\xi)\ge0\) for all \(\xi\), then
\(\langle t\xi,\xi\rangle\ge0\).

\textbf{Proof.} For fixed \(\xi\), the map
\(\eta\mapsto\overline{B(\xi,\eta)}\) is a bounded linear functional. By
Theorem 2.3 it is \(\langle\eta,t\xi\rangle\) for a unique vector
\(t\xi\). Uniqueness makes \(t\) linear. The norm formula follows from
\(\|t\xi\|=\sup_{\|\eta\|\le1}|\langle t\xi,\eta\rangle|\). \(\square\)

\textbf{Corollary 3.2.} Every \(T\in B(H)\) has a unique adjoint
\(T^*\in B(H)\) with
\(\langle T\xi,\eta\rangle=\langle\xi,T^*\eta\rangle\). Moreover:

\begin{itemize}
\tightlist
\item
  \(T^{**}=T\), \(\|T^*\|=\|T\|\) and \(\|T^*T\|=\|T\|^2\);
\item
  \((ST)^*=T^*S^*\), and \(T\mapsto T^*\) is conjugate-linear;
\item
  (complex scalars) if \(\langle T\xi,\xi\rangle=0\) for all \(\xi\),
  then \(T=0\);
\item
  (complex scalars) if \(\langle T\xi,\xi\rangle\in\mathbb R\) for all
  \(\xi\), then \(T=T^*\);
\item
  for self-adjoint \(T\),
  \(\|T\|=\sup_{\|\xi\|=1}|\langle T\xi,\xi\rangle|\).
\end{itemize}

\textbf{Proof.} Apply Theorem 3.1 to
\(B(\xi,\eta)=\langle\xi,T\eta\rangle\) to obtain \(T^*\). The algebraic
rules follow from uniqueness.
\(\|T\xi\|^2=\langle T^*T\xi,\xi\rangle\le\|T^*T\|\|\xi\|^2\) gives
\(\|T\|^2\le\|T^*T\|\le\|T^*\|\|T\|\). So \(\|T\|\le\|T^*\|\), and by
symmetry the two norms are equal. The two complex-scalar statements
follow from Proposition 1.1(1) and (2), applied to
\(B(\xi,\eta)=\langle T\xi,\eta\rangle\).

\emph{The last statement.} Let
\(m=\sup_{\|\xi\|=1}|\langle T\xi,\xi\rangle|\), so
\(|\langle T\xi,\xi\rangle|\le m\|\xi\|^2\). For unit \(\xi,\eta\),
expanding and using self-adjointness gives \[\begin{gathered}
4\operatorname{Re}\langle T\xi,\eta\rangle\\
=\langle T(\xi+\eta),\xi+\eta\rangle\\
-\langle T(\xi-\eta),\xi-\eta\rangle\\
\le m(\|\xi+\eta\|^2+\|\xi-\eta\|^2)\\
=4m.
\end{gathered}\] Replacing \(\eta\) by \(e^{i\theta}\eta\) gives
\(|\langle T\xi,\eta\rangle|\le m\), so \(\|T\|\le m\). The reverse
inequality is Cauchy--Schwarz. \(\square\)

\subsection[4. Orthonormal
bases]{\texorpdfstring{\protect\hypertarget{hilbert-spaces-and-compact-operators--oa-fnd-hs-04}{}{}4.
Orthonormal
bases}{4. Orthonormal bases}}\label{hilbert-spaces-and-compact-operators--OA-FND-HS-04}

An \emph{orthonormal family} \((e_i)_{i\in I}\) satisfies
\(\langle e_i,e_j\rangle=\delta_{ij}\). It is an \emph{orthonormal
basis} if its closed linear span is \(H\).

\textbf{Theorem 4.1.} Let \((e_i)_{i\in I}\) be an orthonormal family in
\(H\).

\begin{enumerate}
\tightlist
\item
  (Bessel) \(\sum_i|\langle\xi,e_i\rangle|^2\le\|\xi\|^2\). In
  particular \(\langle\xi,e_i\rangle\ne0\) for at most countably many
  \(i\).
\item
  For \((c_i)\in\ell^2(I)\), the sum \(\sum_ic_ie_i\) converges, as the
  net of finite partial sums, and \(\|\sum_ic_ie_i\|^2=\sum_i|c_i|^2\).
\item
  The following are equivalent:

  \begin{itemize}
  \tightlist
  \item
    the family is a basis;
  \item
    it is maximal among orthonormal families;
  \item
    \(\xi=\sum_i\langle\xi,e_i\rangle e_i\) for every \(\xi\);
  \item
    Parseval\textquotesingle s identity
    \(\langle\xi,\eta\rangle=\sum_i\langle\xi,e_i\rangle\langle e_i,\eta\rangle\)
    holds for all \(\xi,\eta\).
  \end{itemize}
\item
  Every orthonormal family is contained in an orthonormal basis. A
  separable Hilbert space has a finite or countable orthonormal basis.
\item
  Any two orthonormal bases of \(H\) have the same cardinality, the
  \emph{Hilbert dimension} of \(H\).
\end{enumerate}

\textbf{Proof.}

\begin{enumerate}
\item
  For finite \(F\subseteq I\),
  \(\xi-\sum_{i\in F}\langle\xi,e_i\rangle e_i\) is orthogonal to each
  \(e_i\) with \(i\in F\). So
  \(\|\xi\|^2\ge\sum_{i\in F}|\langle\xi,e_i\rangle|^2\). The set
  \(\{i:|\langle\xi,e_i\rangle|>1/n\}\) is finite for every \(n\).
\item
  Only countably many \(c_i\) are nonzero. The partial sums form a
  Cauchy net, because
  \(\|\sum_{i\in F\setminus G}c_ie_i\|^2=\sum_{F\setminus G}|c_i|^2\).
\item
  \emph{Basis implies expansion.} The vector
  \(\xi-\sum_i\langle\xi,e_i\rangle e_i\) is orthogonal to every
  \(e_j\), hence to their closed span \(H\), so it is zero.
\end{enumerate}

\begin{itemize}
\tightlist
\item
  \emph{Expansion implies Parseval:} take inner products of the
  expansions.
\item
  \emph{Parseval implies maximality:} a unit vector \(e\) orthogonal to
  all \(e_i\) would have \(\|e\|^2=\sum|\langle e,e_i\rangle|^2=0\).
\item
  \emph{Maximality implies basis:} if the closed span \(L\) were not
  \(H\), Theorem 2.2 would give a unit vector in \(L^\perp\), which
  could be added to the family.
\end{itemize}

\begin{enumerate}
\setcounter{enumi}{3}
\item
  Orthonormal families containing the given one, ordered by inclusion,
  have unions of chains as upper bounds. Zorn\textquotesingle s lemma
  (first lesson, Theorem 1.1) gives a maximal one, a basis by 3. For
  separable \(H\), apply the Gram--Schmidt process to a dense sequence.
\item
  If one basis is finite, \(H\) is finite-dimensional, and both bases
  have \(\dim H\) elements. Suppose both are infinite,
  \((e_i)_{i\in I}\) and \((f_j)_{j\in J}\). Each \(e_i\) has at most
  countably many \(j\) with \(\langle e_i,f_j\rangle\ne0\), by 1. Every
  \(j\) occurs for some \(i\), since \(f_j\ne0\) is the sum of its
  expansion in the \(e_i\). So \(J\) is a union of countable sets
  indexed by \(I\), and \(|J|\le|I|\) (first lesson, Theorem 8.4(3)).
  Symmetrically \(|I|\le|J|\), and \(|I|=|J|\) by the
  Cantor--Schröder--Bernstein theorem (first lesson, Theorem 8.1).
  \(\square\)
\end{enumerate}

\subsection[5. Compact
operators]{\texorpdfstring{\protect\hypertarget{hilbert-spaces-and-compact-operators--oa-fnd-hs-05}{}{}5.
Compact
operators}{5. Compact operators}}\label{hilbert-spaces-and-compact-operators--OA-FND-HS-05}

\(T\in B(H)\) is \emph{compact} if the image of the unit ball is
relatively compact. \(K(H)\) is the set of compact operators and
\(F(H)\) the set of finite-rank operators.

\subsubsection{Metric compactness
tools}\label{hilbert-spaces-and-compact-operators--metric-compactness-tools}

A metric space is \emph{totally bounded} if, for every
\(\varepsilon>0\), finitely many balls of radius \(\varepsilon\) cover
it.

\textbf{Lemma 5.0.} A metric space is compact if and only if it is
complete and totally bounded. A compact metric space is sequentially
compact and separable.

\emph{Proof.} In a compact space, every sequence has a point whose every
neighbourhood contains infinitely many sequence terms, counted with
their indices: otherwise a finite subcover of neighbourhoods containing
only finitely many terms gives a contradiction. Choosing increasing
indices in the balls of radius \(1/k\) at that point gives a convergent
subsequence. A Cauchy sequence with a convergent subsequence converges
to the same limit, so the space is complete. Compactness also gives
finite covers by balls of any prescribed radius.

Conversely, in a totally bounded space, repeatedly choose an infinite
set of remaining sequence indices lying in one ball of radius
\(2^{-k}\). Take these sets nested and choose increasing indices from
them. The resulting subsequence is Cauchy, since its terms from stage
\(k\) onward are at mutual distance at most \(2^{1-k}\). Completeness
gives its limit.

To obtain compactness, consider an open cover. There is a number
\(\delta>0\) such that every ball of radius \(\delta\) is contained in a
member of the cover. If there were no such number, choose \(x_n\) whose
ball of radius \(1/n\) is contained in no member. A convergent
subsequence tends to some \(x\) in a member \(U\). Choose \(r>0\) with
\(B(x,2r)\subseteq U\). For large subsequence indices, \(d(x_n,x)<r\)
and \(1/n<r\), so \(B(x_n,1/n)\subseteq U\), a contradiction. A finite
cover by \(\delta\)-balls now gives a finite subcover of the original
cover. Finally choose finite \(1/n\)-nets for all \(n\); their union is
countable and dense. \(\square\)

In a finite-dimensional Hilbert space, bounded sets are totally bounded:
express them in an orthonormal basis, bound each coordinate by
Cauchy--Schwarz, and use a finite grid in the resulting bounded
coordinate box. A finite-dimensional subspace is closed, because the
coordinates of a convergent sequence converge and its limit is their
finite basis expansion. Its bounded closed subsets are therefore
complete and compact by the lemma.

\Needspace{9\baselineskip}
\subsubsection{Compactness and finite-rank
approximation}\label{hilbert-spaces-and-compact-operators--compactness-and-finite-rank-approximation}

\textbf{Theorem 5.1.}

\begin{enumerate}
\tightlist
\item
  \(K(H)\) is a norm-closed two-sided ideal of \(B(H)\) that contains
  \(F(H)\).
\item
  \(T\) is compact iff \(T^*\) is compact.
\item
  \(K(H)\) is the norm closure of \(F(H)\).
\item
  A compact operator maps weakly convergent sequences to norm convergent
  sequences.
\item
  Let \((P_i)\) be a net in \(B(H)\) with \(\|P_i\|\le1\) and
  \(P_i\xi\to\xi\) for every \(\xi\in H\). Then \(\|T-P_iT\|\to0\) for
  every compact \(T\). This applies to the projections \(P_n\) onto
  \(\operatorname{span}\{e_1,\dots,e_n\}\) for an orthonormal basis
  \((e_k)\) of a separable \(H\) (Theorem 4.1(3)).
\end{enumerate}

\textbf{Proof.} 5. The closure \(K\) of \(T(\text{ball})\) is compact.
Given \(\varepsilon>0\), cover \(K\) by finitely many balls
\(B(y_k,\varepsilon)\), and choose \(i_0\) with
\(\|(1-P_i)y_k\|<\varepsilon\) for every \(k\) and every \(i\ge i_0\).
For \(y\in K\) with \(\|y-y_k\|<\varepsilon\), \[\begin{gathered}
\|(1-P_i)y\|\\
\le\|(1-P_i)y_k\|+\|1-P_i\|\,\|y-y_k\|<3\varepsilon .
\end{gathered}\] So
\(\|T-P_iT\|=\sup_{y\in K}\|(1-P_i)y\|\le3\varepsilon\) for
\(i\ge i_0\).

\begin{enumerate}
\tightlist
\item
  A finite-rank bounded operator maps the ball into a bounded subset of
  a finite-dimensional space, which is relatively compact. Sums of
  compact operators are compact, since the image of the ball lies in the
  sum of two relatively compact sets. Products with bounded operators
  are compact, since bounded operators are continuous.
\end{enumerate}

\emph{Closedness.} Let \(T_n\to T\) in norm with \(T_n\) compact, and
\(\varepsilon>0\). Choose \(n\) with \(\|T-T_n\|<\varepsilon/3\), and
cover \(T_n(\text{ball})\) by finitely many \(\varepsilon/3\)-balls. The
\(\varepsilon\)-balls with the same centres cover \(T(\text{ball})\). So
\(T(\text{ball})\) is totally bounded, hence relatively compact, since
\(H\) is complete.

\begin{enumerate}
\setcounter{enumi}{1}
\item
  If \(T\) is compact, so is \(TT^*\). For \(\|\xi_n\|\le1\), choose a
  subsequence along which \(TT^*\xi_n\) converges. Then
  \[\begin{gathered}
  \|T^*(\xi_n-\xi_m)\|^2\\
  =\langle TT^*(\xi_n-\xi_m),\xi_n-\xi_m\rangle\\
  \le2\|TT^*(\xi_n-\xi_m)\|\to0,
  \end{gathered}\] so \(T^*\xi_n\) converges along the subsequence.
  Apply this to \(T^*\) for the converse.
\item
  Let \(T\) be compact. The closure \(K\) of \(T(\text{ball})\) is a
  compact metric space, so it has a countable dense subset. The range of
  \(T\) is the union of the sets \(nT(\text{ball})\), so its closure
  \(L\) is separable. Let \((e_k)\) be an orthonormal basis of \(L\),
  and \(P_n\) the projection onto
  \(\operatorname{span}\{e_1,\dots,e_n\}\). Then \(P_ny\to y\) for
  \(y\in L\) (Theorem 4.1(3)), and \(P_n=0\) on \(L^\perp\). The
  argument of 5, applied with \(K\subseteq L\), gives
  \(\|T-P_nT\|\to0\). Each \(P_nT\) has finite rank.
\item
  Let \(\xi_n\to\xi\) weakly. Then \((\xi_n)\) is bounded (first lesson,
  Corollary 4.3(2)), and \(T\xi_n\to T\xi\) weakly. If
  \(\|T\xi_n-T\xi\|\not\to0\), a subsequence stays at distance
  \(\ge\varepsilon\). A further subsequence converges in norm, by
  compactness, and its limit must be the weak limit \(T\xi\), a
  contradiction. \(\square\)
\end{enumerate}

\subsection[6. The spectral theorem for compact self-adjoint
operators]{\texorpdfstring{\protect\hypertarget{hilbert-spaces-and-compact-operators--oa-fnd-hs-06}{}{}6.
The spectral theorem for compact self-adjoint
operators}{6. The spectral theorem for compact self-adjoint operators}}\label{hilbert-spaces-and-compact-operators--OA-FND-HS-06}

\textbf{Lemma 6.1.} Let \(T\) be compact and self-adjoint, \(T\ne0\).
Then \(\|T\|\) or \(-\|T\|\) is an eigenvalue of \(T\).

\textbf{Proof.} By Corollary 3.2, there are unit vectors \(\xi_n\) with
\(\langle T\xi_n,\xi_n\rangle\to\lambda\), where \(|\lambda|=\|T\|\).
Then \[\begin{gathered}
\|T\xi_n-\lambda\xi_n\|^2\\
=\|T\xi_n\|^2-2\lambda\langle T\xi_n,\xi_n\rangle+\lambda^2\\
\le2\lambda^2-2\lambda\langle T\xi_n,\xi_n\rangle\to0.
\end{gathered}\] A subsequence has \(T\xi_n\to\eta\), by compactness.
Then \(\lambda\xi_n\to\eta\), so \(\|\eta\|=|\lambda|\ne0\), and
\(T\eta=\lim\lambda T\xi_n=\lambda\eta\). \(\square\)

\textbf{Theorem 6.2.} Let \(T\) be compact and self-adjoint, \(T\ne0\).
There are an orthonormal family \((e_n)_{n\in N}\), with
\(N=\{1,\dots,r\}\) or \(N=\mathbb N\), and real numbers
\(\lambda_n\ne0\) with \(|\lambda_1|\ge|\lambda_2|\ge\dots\), tending to
\(0\) if \(N=\mathbb N\), such that
\[T=\sum_n\lambda_n\theta_{e_n,e_n},\] with convergence in norm. The
nonzero eigenvalues of \(T\) are the \(\lambda_n\). The eigenspace of an
eigenvalue \(\mu\ne0\) is spanned by the \(e_n\) with \(\lambda_n=\mu\)
and is finite-dimensional. \(T\ge0\) iff all \(\lambda_n>0\).

\textbf{Proof.} \emph{The eigenvectors.} Put \(H_1=H\). Lemma 6.1 gives
a unit eigenvector \(e_1\) with eigenvalue \(\lambda_1\),
\(|\lambda_1|=\|T\|\). The space \(H_2=\{e_1\}^\perp\) is invariant
under \(T\), because \(T\) is self-adjoint, and \(T|_{H_2}\) is compact
and self-adjoint. Repeat. The process stops if \(T|_{H_{r+1}}=0\).
Otherwise it produces \((e_n,\lambda_n)\) for all \(n\), with
\(|\lambda_n|\) nonincreasing.

\emph{The eigenvalues tend to \(0\).} If not, \(|\lambda_n|\ge\delta>0\)
for all \(n\). Then
\(\|Te_n-Te_m\|^2=\lambda_n^2+\lambda_m^2\ge2\delta^2\), so \((Te_n)\)
has no convergent subsequence, contradicting compactness.

\emph{The expansion.} \(T-\sum_{n\le k}\lambda_n\theta_{e_n,e_n}\)
vanishes on \(\operatorname{span}\{e_1,\dots,e_k\}\) and equals \(T\) on
\(H_{k+1}\). So its norm is \(\|T|_{H_{k+1}}\|=|\lambda_{k+1}|\to0\).

\emph{Eigenvalues and eigenspaces.} If \(T\xi=\mu\xi\) with \(\mu\ne0\),
then \(\mu\xi=\sum\lambda_n\langle\xi,e_n\rangle e_n\). So \(\xi\) lies
in the closed span of the \(e_n\), and \(\langle\xi,e_n\rangle=0\)
unless \(\lambda_n=\mu\). Only finitely many \(\lambda_n\) equal
\(\mu\), because \(\lambda_n\to0\).

\emph{Positivity.}
\(\langle T\xi,\xi\rangle=\sum\lambda_n|\langle\xi,e_n\rangle|^2\).
\(\square\)

\subsection[7. The Riesz theory of compact
operators]{\texorpdfstring{\protect\hypertarget{hilbert-spaces-and-compact-operators--oa-fnd-hs-07}{}{}7.
The Riesz theory of compact
operators}{7. The Riesz theory of compact operators}}\label{hilbert-spaces-and-compact-operators--OA-FND-HS-07}

\textbf{Theorem 7.1.} Let \(T\in K(H)\) and
\(\lambda\in\mathbb C\setminus\{0\}\), and put \(S=T-\lambda\).

\begin{enumerate}
\tightlist
\item
  \(\ker S\) is finite-dimensional.
\item
  \(S(H)\) is closed.
\item
  If \(S\) is injective, it is surjective, hence invertible in \(B(H)\).
\item
  Consequently, every nonzero point of the spectrum \(\sigma(T)\) is an
  eigenvalue of finite multiplicity.
\item
  The nonzero eigenvalues of \(T\) have no accumulation point other than
  \(0\). Each nonzero point of \(\sigma(T)\) is isolated.
\end{enumerate}

\textbf{Proof.} 1. On \(\ker S\), \(T=\lambda\). So the unit ball of
\(\ker S\) equals \(\lambda^{-1}T(\text{ball of }\ker S)\), which is
relatively compact. An infinite orthonormal sequence in \(\ker S\) would
have mutual distances \(\sqrt2\). So \(\ker S\) is finite-dimensional.

\begin{enumerate}
\setcounter{enumi}{1}
\tightlist
\item
  Let \(Sx_n\to y\). Write \(x_n=u_n+v_n\) with \(u_n\in\ker S\) and
  \(v_n\perp\ker S\), so \(Sv_n=Sx_n\).
\end{enumerate}

\emph{\((v_n)\) is bounded.} Otherwise, along a subsequence with
\(\|v_n\|\to\infty\), the unit vectors \(w_n=v_n/\|v_n\|\) have
\(Sw_n\to0\). A further subsequence has \(Tw_n\to z\). Then
\(\lambda w_n=Tw_n-Sw_n\to z\), so \(w_n\to z/\lambda\), a unit vector
orthogonal to \(\ker S\) with \(S(z/\lambda)=0\), which is impossible.

\emph{Conclusion.} A subsequence of the bounded sequence has
\(Tv_n\to z'\). Then \(\lambda v_n=Tv_n-Sv_n\to z'-y\), so
\(v_n\to v=(z'-y)/\lambda\), and \(Sv=y\).

\begin{enumerate}
\setcounter{enumi}{2}
\tightlist
\item
  Suppose \(S\) is injective but \(S(H)\ne H\). Each \(S^k\) is of the
  form \((-\lambda)^k+T_k\) with \(T_k\) compact, so \(S^k(H)\) is
  closed by 2. The sequence
  \(H\supsetneq S(H)\supsetneq S^2(H)\supsetneq\dots\) is strictly
  decreasing: if \(S^k(H)=S^{k+1}(H)\), take \(x\notin S(H)\). Then
  \(S^kx=S^{k+1}y\) for some \(y\), so \(S^k(x-Sy)=0\), and injectivity
  gives \(x=Sy\), a contradiction.
\end{enumerate}

Choose unit vectors \(x_k\in S^k(H)\ominus S^{k+1}(H)\). For \(k<m\),
\[Tx_k-Tx_m=\lambda x_k+\big(Sx_k-Sx_m-\lambda x_m\big),\] and the
bracket lies in \(S^{k+1}(H)\), which is orthogonal to \(x_k\). So
\(\|Tx_k-Tx_m\|\ge|\lambda|\), and \((Tx_k)\) has no convergent
subsequence, contradicting compactness. Hence \(S\) is surjective, and
its inverse is bounded by the inverse mapping theorem (first lesson,
Corollary 5.2).

\begin{enumerate}
\setcounter{enumi}{3}
\item
  If \(\lambda\ne0\) is not an eigenvalue, then \(S\) is injective,
  hence invertible by 3, so \(\lambda\notin\sigma(T)\). The multiplicity
  is finite by 1.
\item
  Suppose \(\lambda_n\to\lambda\ne0\) are distinct eigenvalues with
  eigenvectors \(e_n\). Eigenvectors for distinct eigenvalues are
  linearly independent. Let \(M_n=\operatorname{span}\{e_1,\dots,e_n\}\)
  and choose unit vectors \(y_n\in M_n\ominus M_{n-1}\). Then
  \((T-\lambda_n)M_n\subseteq M_{n-1}\). For \(m<n\), \[\begin{gathered}
  \frac{Ty_n}{\lambda_n}-\frac{Ty_m}{\lambda_m}\\
  =y_n-\Big(\frac{(\lambda_n-T)y_n}{\lambda_n}+\frac{Ty_m}{\lambda_m}\Big),
  \end{gathered}\] and the bracket lies in \(M_{n-1}\perp y_n\). So the
  left side has norm \(\ge1\). Since \(\|y_n/\lambda_n\|\) is bounded,
  this contradicts compactness. Hence the nonzero eigenvalues accumulate
  only at \(0\), and by 4 every nonzero point of \(\sigma(T)\) is
  isolated. \(\square\)
\end{enumerate}

\subsection[8. Hilbert tensor
products]{\texorpdfstring{\protect\hypertarget{hilbert-spaces-and-compact-operators--oa-fnd-hs-08}{}{}8.
Hilbert tensor
products}{8. Hilbert tensor products}}\label{hilbert-spaces-and-compact-operators--OA-FND-HS-08}

For a set \(S\), \(\ell^2(S)\) is the space of functions
\(c:S\to\mathbb C\) with \(\sum_s|c_s|^2<\infty\), with
\(\langle c,d\rangle=\sum_sc_s\bar d_s\). It is a Hilbert space.
Completeness: a Cauchy sequence \((c^{(n)})\) converges at each \(s\) to
some \(c_s\). For every finite \(F\subseteq S\), \[\begin{gathered}
\sum_{s\in F}|c_s-c^{(n)}_s|^2\\
=\lim_m\sum_{s\in F}|c^{(m)}_s-c^{(n)}_s|^2\\
\le\sup_{m\ge n}\|c^{(m)}-c^{(n)}\|^2
\end{gathered}\]. So \(c-c^{(n)}\in\ell^2(S)\) and
\(\|c-c^{(n)}\|\to0\). The functions \(\delta_s\) form an orthonormal
basis.

\textbf{Theorem 8.1} (Hilbert tensor products). Let \(H\) and \(K\) be
Hilbert spaces.

\begin{enumerate}
\tightlist
\item
  There are a Hilbert space \(H\otimes K\) and a bilinear map
  \((\xi,\eta)\mapsto\xi\otimes\eta\) from \(H\times K\) to
  \(H\otimes K\) with
  \[\langle\xi\otimes\eta,\xi'\otimes\eta'\rangle=\langle\xi,\xi'\rangle\langle\eta,\eta'\rangle ,\]
  such that the elementary tensors \(\xi\otimes\eta\) span a dense
  subspace.
\item
  If \((e_i)_{i\in I}\) and \((f_j)_{j\in J}\) are orthonormal bases of
  \(H\) and \(K\), then \((e_i\otimes f_j)\) is an orthonormal basis of
  \(H\otimes K\).
\item
  (\emph{Universal property.}) Let \(B:H\times K\to\mathcal K\) be a
  bilinear map into a Hilbert space with
  \(\langle B(\xi,\eta),B(\xi',\eta')\rangle=\langle\xi,\xi'\rangle\langle\eta,\eta'\rangle\).
  There is exactly one isometry \(U:H\otimes K\to\mathcal K\) with
  \(U(\xi\otimes\eta)=B(\xi,\eta)\). It is unitary if the vectors
  \(B(\xi,\eta)\) span a dense subspace. So \(H\otimes K\) is unique up
  to a unitary that matches the elementary tensors.
\end{enumerate}

\textbf{Proof.} (1) and (2). Fix orthonormal bases as in (2), and put
\(H\otimes K=\ell^2(I\times J)\) and
\((\xi\otimes\eta)(i,j)=\langle\xi,e_i\rangle\langle\eta,f_j\rangle\).

\begin{itemize}
\tightlist
\item
  By Parseval\textquotesingle s identity (Theorem 4.1(3)),
  \(\sum_{i,j}|\langle\xi,e_i\rangle|^2|\langle\eta,f_j\rangle|^2=\|\xi\|^2\|\eta\|^2\),
  so \(\xi\otimes\eta\in\ell^2(I\times J)\).
\item
  The map is bilinear, and by Parseval\textquotesingle s identity
  \[\begin{gathered}
  \langle\xi\otimes\eta,\xi'\otimes\eta'\rangle\\
  =\sum_i\langle\xi,e_i\rangle\langle e_i,\xi'\rangle\sum_j\langle\eta,f_j\rangle\langle f_j,\eta'\rangle\\
  =\langle\xi,\xi'\rangle\langle\eta,\eta'\rangle .
  \end{gathered}\]
\item
  \(e_i\otimes f_j=\delta_{(i,j)}\). These form an orthonormal basis, so
  the elementary tensors span a dense subspace.
\end{itemize}

(3) For finite sums, \[\begin{gathered}
\Big\|\sum_kc_kB(\xi_k,\eta_k)\Big\|^2\\
=\sum_{k,l}c_k\bar c_l\langle\xi_k,\xi_l\rangle\langle\eta_k,\eta_l\rangle\\
=\Big\|\sum_kc_k\,\xi_k\otimes\eta_k\Big\|^2 .
\end{gathered}\] So
\(\sum_kc_k\,\xi_k\otimes\eta_k\mapsto\sum_kc_kB(\xi_k,\eta_k)\) is well
defined, since a combination equal to \(0\) goes to a vector of norm
\(0\), and it is isometric. It extends by continuity to \(H\otimes K\).
Uniqueness holds because the elementary tensors are total. An isometry
has closed range, so if its range is dense, it is onto. \(\square\)

\subsection[9. Multiplication
operators]{\texorpdfstring{\protect\hypertarget{hilbert-spaces-and-compact-operators--oa-fnd-hs-09}{}{}9.
Multiplication
operators}{9. Multiplication operators}}\label{hilbert-spaces-and-compact-operators--OA-FND-HS-09}

Let \((Z,\Sigma,\nu)\) be a \(\sigma\)-finite measure space.
\(L^\infty(\nu)\) is the space of essentially bounded measurable
functions modulo null functions, with the essential supremum norm
\(\|f\|_\infty\). For \(f\in L^\infty(\nu)\), \(m_f\) is the operator
\(\xi\mapsto f\xi\) on \(L^2(\nu)\). The
\hyperref[measure-and-hilbert-space-tools--3-the-complete-spaces-of-integrable-functions]{measure-tools
lesson}, Theorems 3.1--3.2, proves Hölder, completeness of \(L^2\), and
density of simple functions on arbitrary measure spaces. Its proof uses
elementary integration and convergence, not the operator theorem here.
Thus the use of Hilbert representation in the later Radon--Nikodym
theorem creates no circular dependence.

\Needspace{5\baselineskip}
\textbf{Theorem 9.1.}

\begin{enumerate}
\tightlist
\item
  \(f\mapsto m_f\) is an isometric unital \(*\)-homomorphism from
  \(L^\infty(\nu)\) into \(B(L^2(\nu))\).
\item
  An operator \(T\in B(L^2(\nu))\) commutes with every \(m_f\) if and
  only if \(T=m_g\) for some \(g\in L^\infty(\nu)\). So the algebra
  \(\{m_f:f\in L^\infty(\nu)\}\) equals its own commutant: it is
  \emph{maximal abelian}.
\end{enumerate}

\textbf{Proof.} (1) \(\|f\xi\|_2\le\|f\|_\infty\|\xi\|_2\), and
\(\langle f\xi,\eta\rangle=\langle\xi,\bar f\eta\rangle\). For
\(\varepsilon>0\), the set \(\{|f|>\|f\|_\infty-\varepsilon\}\) has
positive measure. Since \(\nu\) is \(\sigma\)-finite, it contains a set
\(E\) with \(0<\nu(E)<\infty\). Then
\(\|f1_E\|_2\ge(\|f\|_\infty-\varepsilon)\|1_E\|_2\). So
\(\|m_f\|=\|f\|_\infty\).

(2) \emph{Reduction to a finite measure.} Write \(Z\) as a disjoint
union of sets \(Z_n\in\Sigma\) of finite measure. Let
\(w=\sum_n2^{-n}(1+\nu(Z_n))^{-1}1_{Z_n}\). Then \(w>0\) everywhere, and
\(\lambda=w\,\nu\) is a finite measure with the same null sets as
\(\nu\).

\begin{itemize}
\tightlist
\item
  So \(L^\infty(\lambda)=L^\infty(\nu)\).
\item
  \(V\xi=w^{-1/2}\xi\) is a unitary from \(L^2(\nu)\) onto
  \(L^2(\lambda)\), with inverse \(\xi\mapsto w^{1/2}\xi\).
\item
  \(Vm_fV^{-1}=m_f\).
\end{itemize}

So we may assume \(\nu(Z)<\infty\). Then \(1\in L^2(\nu)\), and
\(L^\infty(\nu)\subseteq L^2(\nu)\).

\emph{The function \(g\).} Let \(T\) commute with every \(m_f\), and put
\(g=T1\in L^2(\nu)\). For \(f\in L^\infty(\nu)\),
\[Tf=Tm_f1=m_fT1=fg .\]

\begin{itemize}
\tightlist
\item
  \emph{\(g\) is essentially bounded.} For \(c>\|T\|\) let
  \(E=\{|g|>c\}\). Then
  \(c^2\nu(E)\le\|1_Eg\|_2^2=\|T1_E\|_2^2\le\|T\|^2\nu(E)\), so
  \(\nu(E)=0\). Hence \(\|g\|_\infty\le\|T\|\).
\item
  \emph{\(T=m_g\).} The two operators agree on \(L^\infty(\nu)\), which
  contains the simple functions. These are dense in \(L^2(\nu)\)
  \hyperref[measure-and-hilbert-space-tools--3-the-complete-spaces-of-integrable-functions]{Theorem
  3.2 of the measure-tools lesson}.
\end{itemize}

The converse holds because the algebra is commutative. \(\square\)

\subsection{Exercises}\label{hilbert-spaces-and-compact-operators--exercises}

\textbf{Exercise 1} (easy; diagonal operators). Let \((e_n)\) be an
orthonormal basis of \(H\) and \(T=\sum_n\mu_n\theta_{e_n,e_n}\) for a
bounded sequence \((\mu_n)\). Show that \(T\) is compact iff
\(\mu_n\to0\).

\emph{Solution.} If \(\mu_n\to0\), the finite-rank truncations converge
to \(T\) in norm, so \(T\) is compact by Theorem 5.1(1). If
\(|\mu_{n_k}|\ge\delta>0\) along a subsequence, then
\(\|Te_{n_k}-Te_{n_l}\|^2\ge2\delta^2\), so \(T\) is not compact.
\(\square\)

\textbf{Exercise 2} (medium; Hilbert--Schmidt operators). Show that an
operator \(T\) with \(\sum_n\|Te_n\|^2<\infty\) for some orthonormal
basis \((e_n)\) is compact.

\emph{Solution.} Let \(P_k\) be the projection onto
\(\operatorname{span}\{e_1,\dots,e_k\}\). Then \[\begin{gathered}
\|T(1-P_k)\xi\|\\
\le\sum_{n>k}|\langle\xi,e_n\rangle|\|Te_n\|\\
\le\|\xi\|(\sum_{n>k}\|Te_n\|^2)^{1/2}
\end{gathered}\] by Cauchy--Schwarz. So \(TP_k\to T\) in norm, and each
\(TP_k\) has finite rank. \(\square\)

\textbf{Exercise 3} (hard; the Volterra operator). On \(L^2[0,1]\), let
\((V\xi)(s)=\int_0^s\xi(t)\,dt\). Show that \(V\) is compact and that
\(\sigma(V)=\{0\}\).

\emph{Solution.} First \(V\) is bounded: Cauchy--Schwarz gives
\(|V\xi(s)|^2\leq s\|\xi\|_2^2\), hence
\(\|V\xi\|_2^2\leq\frac12\|\xi\|_2^2\). Also
\(|V\xi(s)-V\xi(t)|\leq |s-t|^{1/2}\|\xi\|_2\), so \(V\xi\) is
continuous.

\begin{itemize}
\tightlist
\item
  \emph{Compactness:} with an orthonormal basis \((e_n)\),
  \[\begin{gathered}
  \sum_n\|Ve_n\|^2\\
  =\int_0^1\sum_n|\langle 1_{[0,s]},\overline{e_n}\rangle|^2ds\\
  =\int_0^1\|1_{[0,s]}\|^2ds\\
  =\frac12
  \end{gathered}\], by Parseval applied to the basis
  \((\overline{e_n})\) and
  \hyperref[measure-and-hilbert-space-tools--2-integration-and-convergence-without-countability-assumptions]{monotone
  convergence}. By Exercise 2, \(V\) is compact.
\item
  \emph{Spectrum:} by Theorem 7.1(4), a nonzero point of \(\sigma(V)\)
  would be an eigenvalue. If \(V\xi=\lambda\xi\), then
  \(\xi=\lambda^{-1}V\xi\) is continuous, then \(C^1\) by the
  fundamental theorem of calculus proved in
  \hyperref[measure-and-hilbert-space-tools--6-one-variable-integral-identities]{Lemma
  6.1 of the measure-tools lesson}, with \(\xi=\lambda\xi'\) and
  \(\xi(0)=0\). The function \(e^{-s/\lambda}\xi(s)\) has zero
  derivative and is zero at \(s=0\), so \(\xi=0\). Finally
  \(0\in\sigma(V)\), because \(V\) is not invertible in infinite
  dimensions: a compact invertible operator would make the identity
  compact. \(\square\)
\end{itemize}

\textbf{Exercise 4} (medium; closed range of \(1-T\)). Let \(T\) be
compact. Show that \(\dim\ker(1-T)=\dim\ker(1-T^*)\).

\emph{Solution.}

\begin{itemize}
\tightlist
\item
  Both kernels are finite-dimensional (Theorem 7.1(1), applied to \(T\)
  and to \(T^*\), which is compact by Theorem 5.1(2)).
\item
  Let \(S=1-T\). Since \(S(H)\) is closed, \(\ker S^*=S(H)^\perp\), and
  \(H=S(H)\oplus\ker S^*\).
\item
  Suppose \(\dim\ker S<\dim\ker S^*\). Choose an injective linear map
  \(A\) from \(\ker S\) into \(\ker S^*\) that is not onto, and extend
  it by \(0\) on \((\ker S)^\perp\). This gives a finite-rank operator.
  Then \(S'=S+A=1-(T-A)\) is injective with range
  \(S(H)+A(\ker S)\ne H\), contradicting Theorem 7.1(3).
\item
  Exchanging \(T\) and \(T^*\) gives the reverse inequality. \(\square\)
\end{itemize}

\subsection{Where this
leads}\label{hilbert-spaces-and-compact-operators--where-this-leads}

The continuous functional calculus and the operator algebra \(B(H)\) are
developed in
\href{https://kokunoyumeto.github.io/open-math-courses/courses/harmonic-analysis-on-locally-compact-groups/reader/programme-reading.html\#ref-96a615e4209af567}{C*-algebras:
continuous functional calculus, automatic continuity, positive cones,
approximate identities and quotients}. The spectral theorem for bounded
self-adjoint operators with projection-valued measures, and
Calkin\textquotesingle s theorem on the ideals of \(B(H)\), are proved
in
\href{https://kokunoyumeto.github.io/open-math-courses/courses/harmonic-analysis-on-locally-compact-groups/reader/programme-reading.html\#ref-c87ebece35f15b9d}{The
spectral theorem for bounded self-adjoint operators}. The trace class
and the operator topologies are in
\href{https://kokunoyumeto.github.io/open-math-courses/courses/harmonic-analysis-on-locally-compact-groups/reader/programme-reading.html\#ref-b8bfb9db42bafead}{Compact
and trace-class operators, the predual of B(H), and the operator
topologies}. Hilbert tensor products are used for \(L^2\) of a product
of Radon measures in
\href{https://kokunoyumeto.github.io/open-math-courses/courses/harmonic-analysis-on-locally-compact-groups/reader/programme-reading.html\#ref-23bd0edc6019793c}{Haar
measure on locally compact groups}, and multiplication operators are the
starting point of
\href{https://kokunoyumeto.github.io/open-math-courses/courses/harmonic-analysis-on-locally-compact-groups/reader/programme-reading.html\#ref-97e06a521795b886}{Abelian
operator algebras}.

\subsection{References}\label{hilbert-spaces-and-compact-operators--references}

Hilbert (1904--1910), Riesz (1907, 1918) and Fréchet (1907). The Riesz
theory of compact operators is from F. Riesz, "Über lineare
Funktionalgleichungen", \emph{Acta Mathematica} 41 (1918) 71--98. The
proofs are written here in our own words.

\section{The Stone--Weierstrass theorem for functions vanishing at
infinity}\label{the-stone-weierstrass-theorem-for-functions-vanishing-at-infinity}

\emph{Written by Claude Opus 5.5 (Anthropic), September 2026, revised
October 2026. The September text was spot-checked by Claude Opus 5.5 in
a separate session. The October revisions correct two points found by an
AI reviewer in a separate OpenAI Codex session and add proofs of the
background facts that the core courses leave as exercises; they are
self-checked by the writing AI. Public domain (CC0).}

This lesson proves the Stone--Weierstrass theorem for continuous
functions that vanish at infinity. Let \(X\) be a locally compact
Hausdorff space, and let \(A\) be a subalgebra of \(C_0(X)\) that is
closed under complex conjugation, separates the points of \(X\), and has
no common zero. Then every function in \(C_0(X)\) is a uniform limit of
functions in \(A\)
(\hyperlink{the-stone-weierstrass-theorem-for-functions-vanishing-at-infinity--oa-fnd-sw-09}{Theorem
10.1}). The algebra need not contain the constant functions, and \(X\)
need not be compact, \(\sigma\)-compact or metrizable. Commutative
C*-algebras and harmonic analysis on locally compact groups need the
theorem in this form.

The proof follows the classical route. A recursion for polynomials
without constant term approximates \(|t|\) (Section 6). With
interpolation at two points (Section 7) and Stone\textquotesingle s
lattice lemma (Section 8), it gives the real theorem on compact spaces
(Section 9). The one-point compactification (Section 3) carries the
result over to \(C_0(X)\), and real and imaginary parts give the complex
theorem (Section 10). Sections 11--13 give worked examples, show that
each hypothesis is needed, and describe the closure of an arbitrary
self-adjoint subalgebra. Sections 2--5 develop the function spaces and
the topology that surround the proof: the one-point compactification,
locally compact spaces, Urysohn\textquotesingle s lemma and
metrizability. Sections 14--18 add classical companions: the Weierstrass
theorem with Bernstein\textquotesingle s constructive proof,
Korovkin\textquotesingle s theorem, further density theorems, the Tietze
extension theorem, the disc algebra, and approximation by polynomials
with integer coefficients.

The lesson assumes the core courses
\href{https://kokunoyumeto.github.io/program-matematika-indonesia/en/\#course-C10}{Real
Analysis I},
\href{https://kokunoyumeto.github.io/program-matematika-indonesia/en/\#course-C20}{Real
Analysis II} and
\href{https://kokunoyumeto.github.io/program-matematika-indonesia/en/\#course-C90}{Point-Set
Topology}: open and closed sets, continuity, compactness defined by open
covers, metric spaces and uniform convergence. A few facts about
numbers, the Riemann integral and the exponential function are taken
from the core courses; they are listed near the end, under "Results used
from other lessons", with the place of each proof. No other lesson is
needed.

Weierstrass proved in 1885 that every continuous function on a closed
bounded interval is a uniform limit of polynomials {[}Weierstrass
1885{]}. Stone proved the general theorem in 1937 and gave a simpler
proof, by lattices of functions, in 1948 {[}Stone 1937{]}, {[}Stone
1948{]} P. The one-point compactification is due to Alexandroff
{[}Alexandroff 1924{]}. The core course \emph{Real Analysis II} proves
the Weierstrass theorem and the real and complex Stone--Weierstrass
theorems for compact metric spaces
(\href{https://www.jirka.org/ra/html/sec_stoneweier.html}{Lebl, Section
11.7}). Free texts that treat the theorems are {[}van Neerven{]} and
{[}Blackadar{]}; see also, e.g., {[}Rudin, Theorems 7.26, 7.32 and
7.33{]}.

\subsection{1.
Conventions}\label{the-stone-weierstrass-theorem-for-functions-vanishing-at-infinity--1-conventions}

\emph{Spaces.} A \emph{neighbourhood} of a point is a set that contains
an open set containing the point. \emph{Compact} means that every open
cover has a finite subcover; the Hausdorff property is not part of the
word. A space is \emph{locally compact Hausdorff} (LCH) if it is
Hausdorff and every point has a compact neighbourhood. Five elementary
facts are used throughout.

\begin{itemize}
\tightlist
\item
  (K1) A closed set \(C\) contained in a compact set \(K\) is compact.
  Add \(X\setminus C\) to an open cover of \(C\), take a finite subcover
  of \(K\), and discard \(X\setminus C\).
\item
  (K2) A continuous image of a compact set is compact, and a finite
  union of compact sets is compact.
\item
  (K3) Let \(X\) be Hausdorff, \(K\subseteq X\) compact and
  \(x\notin K\). Then \(x\) and \(K\) have disjoint open neighbourhoods.
  For each \(y\in K\) choose disjoint open sets \(U_y\ni x\) and
  \(V_y\ni y\). Finitely many sets \(V_{y_1},\dots,V_{y_m}\) cover
  \(K\); take \(U_{y_1}\cap\dots\cap U_{y_m}\) and
  \(V_{y_1}\cup\dots\cup V_{y_m}\). In particular, compact subsets of
  Hausdorff spaces are closed.
\item
  (K4) Let a continuous map from a compact space onto a Hausdorff space
  be one-to-one. Then its inverse is continuous, because the map sends
  closed sets to closed sets: a closed set is compact by (K1), its image
  is compact by (K2), and a compact subset of a Hausdorff space is
  closed by (K3).
\item
  (K5) A continuous \(u:K\to\mathbb R\) on a nonempty compact space is
  bounded and attains its maximum and minimum. The open sets \(\{u<n\}\)
  cover \(K\), so \(u\) is bounded above. If \(u<s=\sup u\) everywhere,
  the open sets \(\{u<s-1/n\}\) cover \(K\), and a finite subcover gives
  \(u\leq s-1/N\), which is absurd. Apply this to \(-u\) for the
  minimum.
\end{itemize}

\emph{Functions.} Functions are complex-valued unless they are called
real. For a bounded function \(f\) on a set \(S\), put
\(\|f\|=\sup_S|f|\), and \(\|f\|_T=\sup_T|f|\) for \(T\subseteq S\). On
a space \(X\), \(C(X)\) and \(C(X;\mathbb R)\) are the continuous
complex and real functions, and \(C_b(X)\) the bounded continuous ones.
The \emph{support} of \(f\) is the closure of \(\{f\neq0\}\), and
\(C_c(X)\) consists of the continuous functions with compact support.
\(C_0(X)\) consists of the continuous \(f\) with
\[\{|f|\geq\varepsilon\}=\{x\in X:\ |f(x)|\geq\varepsilon\}\ \text{compact for every }\varepsilon>0 .
\tag{1.1}\] Equivalently, for every \(\varepsilon>0\) there is a compact
set \(K\) with \(|f|<\varepsilon\) off \(K\). Indeed, if there is such a
\(K\), then the closed set \(\{|f|\geq\varepsilon\}\) lies in \(K\) and
is compact by (K1); conversely, \(K=\{|f|\geq\varepsilon\}\) will do.
Real versions carry the suffix \(;\mathbb R\), as in
\(C_0(X;\mathbb R)\).

\emph{Algebras.} On a set \(S\), an \emph{algebra of functions} is a
complex linear space of functions \(S\to\mathbb C\) that is closed under
pointwise products. A \emph{real algebra} is the same with real
functions and real scalars. No unit is assumed. An algebra \(A\) is
\emph{self-adjoint} if \(f\in A\) implies \(\bar f\in A\). It
\emph{separates points} if for \(x\neq y\) some \(f\in A\) has
\(f(x)\neq f(y)\). It \emph{vanishes nowhere} if for every \(x\) some
\(f\in A\) has \(f(x)\neq0\). Its \emph{common zero set} is
\[Z(A)=\{x\in S:\ f(x)=0\ \text{for all }f\in A\},\] so \(A\) vanishes
nowhere exactly when \(Z(A)=\varnothing\). \(\bar A\) is the closure of
\(A\) for \(\|\cdot\|\). For real functions, \(f\vee g=\max(f,g)\) and
\(f\wedge g=\min(f,g)\), pointwise.

\subsection[2. Spaces of bounded and continuous
functions]{\texorpdfstring{\protect\hypertarget{the-stone-weierstrass-theorem-for-functions-vanishing-at-infinity--oa-fnd-sw-01}{}{}2.
Spaces of bounded and continuous
functions}{2. Spaces of bounded and continuous functions}}\label{the-stone-weierstrass-theorem-for-functions-vanishing-at-infinity--OA-FND-SW-01}

\textbf{Proposition 2.1.} Let \(S\) be a set and \(X\) a topological
space.

\begin{enumerate}
\tightlist
\item
  (\emph{Completeness}) The bounded functions on \(S\) form a Banach
  space under \(\|\cdot\|\). The bounded continuous functions on \(X\)
  and the bounded Borel functions on \(X\) form closed subspaces of it.
  If \(X\) is compact, then \(C(X)=C_b(X)\).
\item
  (\emph{Closures of algebras}) The closure of an algebra (or real
  algebra) of bounded functions is again one. The closure of a
  self-adjoint algebra is self-adjoint.
\item
  (\emph{Vanishing at infinity})
  \(C_c(X)\subseteq C_0(X)\subseteq C_b(X)\). \(C_0(X)\) is a closed
  self-adjoint subalgebra of \(C_b(X)\), hence a Banach space.
  \(C_c(X)\) is a self-adjoint subalgebra of \(C_0(X)\), and it is dense
  in \(C_0(X)\).
\end{enumerate}

\textbf{Proof.} (1) Let \((f_n)\) be a Cauchy sequence. For each \(s\),
\((f_n(s))\) is Cauchy in \(\mathbb C\); let \(f(s)\) be its limit.
Given \(\varepsilon>0\), choose \(N\) with
\(\|f_n-f_m\|\leq\varepsilon\) for \(m,n\geq N\). Letting \(m\to\infty\)
gives \(|f_n(s)-f(s)|\leq\varepsilon\) for all \(s\) and all
\(n\geq N\). So \(f\) is bounded and \(f_n\to f\). Suppose the \(f_n\)
are continuous. Fix \(x\in X\), fix \(n\geq N\), and take an open
\(U\ni x\) on which \(|f_n-f_n(x)|<\varepsilon\). On \(U\),
\(|f-f(x)|\leq|f-f_n|+|f_n-f_n(x)|+|f_n(x)-f(x)|<3\varepsilon\), so
\(f\) is continuous. Suppose the \(f_n\) are Borel. Then
\(\{\operatorname{Re}f>a\}=\bigcup_k\bigcup_N\bigcap_{n\geq N}\{\operatorname{Re}f_n>a+1/k\}\)
is a Borel set, and so is the analogous set for \(\operatorname{Im}f\);
so \(f\) is Borel. The last claim is (K5) applied to \(|f|\).

(2) Let \(f_n\to f\) and \(g_n\to g\) with \(f_n,g_n\) in the algebra.
Then \(\|f_n\|\leq\|f\|+\|f_n-f\|\) is bounded, and
\(\|f_ng_n-fg\|\leq\|f_n\|\,\|g_n-g\|+\|g\|\,\|f_n-f\|\to0\). Sums,
scalar multiples, complex conjugates and real values pass to uniform
limits in the same way.

(3) If \(f\in C_c(X)\), then \(\{|f|\geq\varepsilon\}\) is a closed
subset of the compact support, hence compact by (K1). If
\(f\in C_0(X)\), then \(|f|<1\) off the compact set \(\{|f|\geq1\}\),
and \(|f|\) is bounded on that set by (K5); so \(f\) is bounded. For
\(f,g\in C_0(X)\), \(\varepsilon>0\) and a scalar \(c\neq0\),
\[\{|f+g|\geq\varepsilon\}\subseteq\{|f|\geq\tfrac\varepsilon2\}\cup\{|g|\geq\tfrac\varepsilon2\},\quad
\{|fg|\geq\varepsilon\}\subseteq\{|f|\geq\tfrac{\varepsilon}{\|g\|+1}\},\quad
\{|cf|\geq\varepsilon\}=\{|f|\geq\tfrac{\varepsilon}{|c|}\}.\] The sets
on the left are closed. The sets on the right are compact, since a
finite union of compact sets is compact (K2). So the sets on the left
are compact by (K1), and \(f+g,\ fg,\ cf\in C_0(X)\). Clearly
\(\bar f\in C_0(X)\). If \(f_n\in C_0(X)\) and
\(\|f_n-f\|\leq\varepsilon/2\), then
\(\{|f|\geq\varepsilon\}\subseteq\{|f_n|\geq\varepsilon/2\}\). So
\(C_0(X)\) is closed in \(C_b(X)\), and complete by (1). The supports of
\(f+g\), \(fg\), \(cf\) and \(\bar f\) are closed subsets of
\(\operatorname{supp}f\cup\operatorname{supp}g\), so \(C_c(X)\) is a
self-adjoint subalgebra. For density, let \(f\in C_0(X)\) and
\(\varepsilon>0\). Define \(\psi_\varepsilon:\mathbb C\to\mathbb C\) by
\(\psi_\varepsilon(z)=0\) for \(|z|\leq\varepsilon\) and
\(\psi_\varepsilon(z)=z-\varepsilon z/|z|\) for \(|z|\geq\varepsilon\).
It is continuous, and \(|\psi_\varepsilon(z)-z|\leq\varepsilon\). So
\(\psi_\varepsilon\circ f\) is continuous and within \(\varepsilon\) of
\(f\). It is nonzero only on \(\{|f|>\varepsilon\}\), whose closure lies
in the compact set \(\{|f|\geq\varepsilon\}\). So
\(\psi_\varepsilon\circ f\in C_c(X)\). \(\square\)

\subsection[3. The one-point
compactification]{\texorpdfstring{\protect\hypertarget{the-stone-weierstrass-theorem-for-functions-vanishing-at-infinity--oa-fnd-sw-02}{}{}\protect\hypertarget{the-stone-weierstrass-theorem-for-functions-vanishing-at-infinity--OA-FND-SW-03}{}{}\protect\hypertarget{the-stone-weierstrass-theorem-for-functions-vanishing-at-infinity--oa-fnd-sw-03}{}{}3.
The one-point
compactification}{3. The one-point compactification}}\label{the-stone-weierstrass-theorem-for-functions-vanishing-at-infinity--OA-FND-SW-02}

\textbf{Definition 3.1} (One-point compactification). Let \(X\) be a
topological space and \(\infty\) a point not in \(X\). Put
\(X_\infty=X\cup\{\infty\}\). Let \(\tau_\infty\) consist of the open
subsets of \(X\) together with the sets \(X_\infty\setminus C\), where
\(C\subseteq X\) is closed and compact.

If \(X\) is Hausdorff, its compact subsets are closed by (K3), so the
word "closed" can then be omitted. For a general space \(X\), it makes
parts (1) and (2) below true.

\Needspace{5\baselineskip}
\textbf{Proposition 3.2.}

\begin{enumerate}
\tightlist
\item
  \(\tau_\infty\) is a topology on \(X_\infty\). \(X\) is open in
  \(X_\infty\), and the subspace topology on \(X\) is its original
  topology.
\item
  \(X_\infty\) is compact.
\item
  If \(X\) is locally compact Hausdorff, then \(X_\infty\) is Hausdorff.
\item
  \(\infty\) is an isolated point of \(X_\infty\) exactly when \(X\) is
  compact. So \(X\) is dense in \(X_\infty\) exactly when \(X\) is not
  compact.
\item
  (\emph{Uniqueness}) Let \(Y\) be compact Hausdorff, \(q\in Y\), and
  \(\varphi\) a homeomorphism of \(X\) onto \(Y\setminus\{q\}\).
  Extending \(\varphi\) by \(\infty\mapsto q\) gives a homeomorphism of
  \(X_\infty\) onto \(Y\).
\end{enumerate}

\textbf{Proof.} (1) \(\varnothing\) is open in \(X\), and
\(X_\infty=X_\infty\setminus\varnothing\). For open \(U\subseteq X\) and
closed compact \(C,D\subseteq X\), the set
\(U\cap(X_\infty\setminus C)=U\setminus C\) is open in \(X\), and
\((X_\infty\setminus C)\cap(X_\infty\setminus D)=X_\infty\setminus(C\cup D)\)
with \(C\cup D\) closed and compact. Consider a union of sets \(U_i\)
open in \(X\) and sets \(X_\infty\setminus C_j\), with at least one
\(j\). It equals \(X_\infty\setminus D\), where
\(D=\bigcap_jC_j\cap\bigcap_i(X\setminus U_i)\). \(D\) is closed and
lies in one \(C_j\), so it is compact by (K1). A union of open subsets
of \(X\) is open in \(X\). The traces on \(X\) of the members of
\(\tau_\infty\) are the sets \(U\) and \(X\setminus C\), all open in
\(X\); so the subspace topology is the original one. Finally
\(X\in\tau_\infty\).

(2) Let \(\mathcal W\) be an open cover of \(X_\infty\). Some
\(W_0\in\mathcal W\) contains \(\infty\), so
\(W_0=X_\infty\setminus C_0\) with \(C_0\) compact. The traces
\(W\cap X\), \(W\in\mathcal W\), are open in \(X\) and cover \(C_0\), so
finitely many of them do. Those members and \(W_0\) cover \(X_\infty\).

(3) Two points of \(X\) are separated by open subsets of \(X\), which
are open in \(X_\infty\). Let \(x\in X\), with a compact neighbourhood
\(N\) and an open \(V\) such that \(x\in V\subseteq N\). \(N\) is closed
by (K3), so \(X_\infty\setminus N\) is an open neighbourhood of
\(\infty\) disjoint from \(V\).

(4) If \(X\) is compact, then \(\{\infty\}=X_\infty\setminus X\) is
open. If \(\{\infty\}\) is open, it is \(X_\infty\setminus C\) with
\(C=X\), so \(X\) is compact. \(X\) is dense exactly when every open
neighbourhood of \(\infty\) meets \(X\), that is, when \(\{\infty\}\) is
not open.

(5) Call the extension \(\Phi\); it is a bijection. Let \(V\subseteq Y\)
be open. If \(q\notin V\), then \(\Phi^{-1}(V)=\varphi^{-1}(V)\) is open
in \(X\). If \(q\in V\), then \(Y\setminus V\) is compact by (K1) and
lies in \(Y\setminus\{q\}\). So \(C=\varphi^{-1}(Y\setminus V)\) is
compact by (K2), and closed in \(X\), since its complement is
\(\varphi^{-1}(V\setminus\{q\})\). Then
\(\Phi^{-1}(V)=X_\infty\setminus C\) is open. So \(\Phi\) is continuous,
and by (K4) its inverse is continuous too. \(\square\)

\textbf{Examples 3.3.} (a) Let \(\mathbb T=\{w\in\mathbb C:|w|=1\}\).
The Cayley map \(\kappa(x)=(x-i)/(x+i)\) is a homeomorphism of
\(\mathbb R\) onto \(\mathbb T\setminus\{1\}\). Indeed \(|x-i|=|x+i|\)
for real \(x\), \(\kappa(x)\neq1\), and the inverse is
\(w\mapsto i(1+w)/(1-w)\), which is real on \(\mathbb T\setminus\{1\}\)
because \((1+w)/(1-w)=(w-\bar w)/|1-w|^2\) there. \(\mathbb T\) is
compact by
\hyperlink{the-stone-weierstrass-theorem-for-functions-vanishing-at-infinity--oa-fnd-sw-15}{Lemma
4.1}(b) below, so (5) gives \(\mathbb R_\infty\cong\mathbb T\) with
\(\infty\mapsto1\). In the same way, \(y\mapsto(2y,|y|^2-1)/(|y|^2+1)\)
is a homeomorphism of \(\mathbb R^n\) onto the unit sphere
\(S^n\subseteq\mathbb R^{n+1}\) with the point \((0,\dots,0,1)\)
removed; its inverse is \((u,t)\mapsto u/(1-t)\). So
\((\mathbb R^n)_\infty\cong S^n\).

(b) If \(Y\) is compact Hausdorff and \(q\in Y\), then (5), applied to
the identity map of \(Y\setminus\{q\}\), identifies
\((Y\setminus\{q\})_\infty\) with \(Y\).

(c) \(\mathbb N_\infty\) is homeomorphic to
\(Y=\{0\}\cup\{1/n:n\geq1\}\subseteq\mathbb R\), with
\(\infty\mapsto0\). \(Y\) is compact, because an open set containing
\(0\) contains all but finitely many of the points \(1/n\).

(d) If \(X\) is compact, \(X_\infty\) is \(X\) together with an isolated
point, by (4). Then \(X\) is not dense in \(X_\infty\), so \(X_\infty\)
is not a compactification of \(X\) in the usual sense, which asks for
\(X\) to be dense.

The functions in \(C_0(X)\) are the continuous functions on \(X_\infty\)
that vanish at \(\infty\), as the next proposition shows. For a function
\(f\) on \(X\), let \(\tilde f\) be the function on \(X_\infty\) that
equals \(f\) on \(X\) and \(0\) at \(\infty\). Put
\[I_\infty=\{g\in C(X_\infty):\ g(\infty)=0\}.\]

\textbf{Proposition 3.4.} Let \(X\) be any topological space.

\begin{enumerate}
\tightlist
\item
  \(f\in C_0(X)\) if and only if \(\tilde f\) is continuous on
  \(X_\infty\).
\item
  The map \(f\mapsto\tilde f\) is a bijection of \(C_0(X)\) onto
  \(I_\infty\) that preserves sums, scalar multiples, products, complex
  conjugation and the norm \(\|\cdot\|\). Its inverse is restriction to
  \(X\). It maps \(C_0(X;\mathbb R)\) onto the real functions in
  \(I_\infty\).
\item
  Every \(g\in C(X_\infty)\) has the unique form \(g=g(\infty)1+h\) with
  \(h\in I_\infty\). So \(C(X_\infty)\) is \(C_0(X)\) with a unit
  adjoined.
\end{enumerate}

\textbf{Proof.} (1) Let \(f\in C_0(X)\). \(\tilde f\) is continuous at
points of \(X\), because \(X\) is open in \(X_\infty\) and carries its
own topology there
(\hyperlink{the-stone-weierstrass-theorem-for-functions-vanishing-at-infinity--oa-fnd-sw-02}{Proposition
3.2}(1)). At \(\infty\): given \(\varepsilon>0\), the set
\(C=\{|f|\geq\varepsilon\}\) is compact, and closed because \(|f|\) is
continuous. So \(X_\infty\setminus C\) is an open neighbourhood of
\(\infty\) on which \(|\tilde f|<\varepsilon\). Conversely, let
\(\tilde f\) be continuous. Then \(f\) is continuous, and for
\(\varepsilon>0\) the set
\(\{x\in X:|f(x)|\geq\varepsilon\}=\{t\in X_\infty:|\tilde f(t)|\geq\varepsilon\}\)
is closed in the compact space \(X_\infty\). It is compact by (K1), and
it lies in \(X\). Since \(X\) carries the subspace topology, it is
compact in \(X\).

(2) The operations are pointwise, and \(0\) at \(\infty\) is preserved
by all of them. \(\sup_{X_\infty}|\tilde f|=\sup_X|f|\). Restriction
inverts the map by (1).

(3) Put \(h=g-g(\infty)1\). If also \(g=c1+h'\) with \(h'\in I_\infty\),
evaluating at \(\infty\) gives \(c=g(\infty)\), so the form is unique.
\(\square\)

Part (3) says that adjoining a unit to \(C_0(X)\) yields
\(C(X_\infty)\).

\subsection[4. Compact and locally compact
spaces]{\texorpdfstring{\protect\hypertarget{the-stone-weierstrass-theorem-for-functions-vanishing-at-infinity--oa-fnd-sw-15}{}{}4.
Compact and locally compact
spaces}{4. Compact and locally compact spaces}}\label{the-stone-weierstrass-theorem-for-functions-vanishing-at-infinity--OA-FND-SW-15}

\textbf{Lemma 4.1} (Compactness tools).

\begin{itemize}
\tightlist
\item
  (a) Every closed bounded interval \([a,b]\) is compact.
\item
  (b) (\emph{Tube lemma}) Let \(Y\) be compact, \(W\subseteq X\times Y\)
  open, and \(\{x_0\}\times Y\subseteq W\). Then
  \(U\times Y\subseteq W\) for some open \(U\ni x_0\). Consequently, a
  product of finitely many compact spaces is compact, and closed bounded
  subsets of \(\mathbb R^n\) are compact. Examples are \(\mathbb T\),
  the closed unit disc and the spheres \(S^n\).
\item
  (c) A compact Hausdorff space is regular and normal: a point and a
  closed set not containing it, and two disjoint closed sets, have
  disjoint open neighbourhoods. Equivalently, if \(C\) is closed, \(O\)
  is open and \(C\subseteq O\), then some open \(G\) satisfies
  \(C\subseteq G\subseteq\overline G\subseteq O\).
\item
  (d) In a compact metric space every sequence has a convergent
  subsequence, so the space is complete. It is separable and second
  countable, and every continuous map from it to a metric space is
  uniformly continuous. Conversely, a metric space in which every
  sequence has a convergent subsequence is compact.
\end{itemize}

\textbf{Proof.} (a) Let \(\mathcal O\) be an open cover, and let \(S\)
be the set of \(s\in[a,b]\) such that finitely many members cover
\([a,s]\). Then \(a\in S\). Let \(\sigma=\sup S\), and let
\(O\in\mathcal O\) contain \(\sigma\), so that
\(O\supseteq(\sigma-\eta,\sigma+\eta)\cap[a,b]\) for some \(\eta>0\).
Some \(s\in S\) exceeds \(\sigma-\eta\). Adding \(O\) to a finite cover
of \([a,s]\) covers \([a,\min(\sigma+\eta/2,b)]\). If \(\sigma<b\), this
contradicts \(\sigma=\sup S\). So \(\sigma=b\), and \(b\in S\).

(b) For each \(y\in Y\) choose open sets \(U_y\ni x_0\) and \(V_y\ni y\)
with \(U_y\times V_y\subseteq W\). Finitely many \(V_{y_j}\) cover
\(Y\); put \(U=\bigcap_jU_{y_j}\). Now let \(X\) and \(Y\) be compact
and \(\mathcal W\) an open cover of \(X\times Y\). Let \(\mathcal U\) be
the family of open \(U\subseteq X\) such that finitely many members of
\(\mathcal W\) cover \(U\times Y\). Each \(x\) lies in a member of
\(\mathcal U\): the compact set \(\{x\}\times Y\) is covered by finitely
many members, and the tube lemma applies to their union. Finitely many
members of \(\mathcal U\) cover \(X\). Induction handles finite
products. A closed bounded subset of \(\mathbb R^n\) lies in a box
\([-M,M]^n\), which is compact by (a) and the product statement, so it
is compact by (K1).

(c) Let \(p\notin C\) with \(C\) closed. \(C\) is compact by (K1), and
(K3) separates \(p\) from \(C\). Let \(E,F\) be disjoint and closed. For
each \(e\in E\), (K3) gives disjoint open sets \(U_e\ni e\) and
\(V_e\supseteq F\). Finitely many \(U_e\) cover the compact set \(E\);
take their union and the intersection of the corresponding \(V_e\). For
the last form, separate \(C\) and \(X\setminus O\) by disjoint open sets
\(G\supseteq C\) and \(H\supseteq X\setminus O\); then
\(\overline G\subseteq X\setminus H\subseteq O\).

(d) Let \((x_j)\) be a sequence in a compact metric space \(K\) with no
convergent subsequence. Then each \(p\in K\) has a ball containing
\(x_j\) for only finitely many \(j\); otherwise indices
\(j_1<j_2<\cdots\) with \(d(x_{j_i},p)<1/i\) could be chosen. Finitely
many such balls cover \(K\), which is absurd. A Cauchy sequence with a
convergent subsequence converges, so \(K\) is complete. For each \(n\),
finitely many balls of radius \(1/n\) cover \(K\). Their centres, over
all \(n\), form a countable dense set \(C\), and the balls \(B(c,1/m)\),
\(c\in C\), \(m\geq1\), form a countable base: if
\(B(x,2/m)\subseteq O\), some \(c\in C\) has \(d(c,x)<1/m\), and then
\(x\in B(c,1/m)\subseteq O\). Conversely, let every sequence in a metric
space \(M\) have a convergent subsequence, and let \(\mathcal O\) be an
open cover. Some \(\delta>0\) has the property that every ball
\(B(x,\delta)\) lies in a member of \(\mathcal O\). Otherwise there are
balls \(B(x_n,1/n)\) lying in no member; a subsequence of the centres
converges to some \(x\) in a member \(O\) with \(B(x,2r)\subseteq O\),
and then \(B(x_n,1/n)\subseteq O\) for large \(n\) in the subsequence.
Also finitely many balls \(B(x,\delta)\) cover \(M\). Otherwise points
with mutual distances at least \(\delta\) can be chosen one after
another, and they form a sequence with no convergent subsequence. So
finitely many members of \(\mathcal O\) cover \(M\). Finally, let
\(\phi\) be a continuous map from a compact metric space \(K\) to a
metric space, and \(\varepsilon>0\). By the first part, every sequence
in \(K\) has a convergent subsequence. So the argument above gives, for
the open cover by the sets \(\{x:d(\phi(x),\phi(p))<\varepsilon/2\}\),
\(p\in K\), a \(\delta>0\) such that every ball \(B(x,\delta)\) lies in
one of them. If \(d(x,y)<\delta\), then \(x\) and \(y\) lie in one such
set, so \(d(\phi(x),\phi(y))<\varepsilon\). \(\square\)

\textbf{Proposition 4.2} (Locally compact spaces).

\begin{enumerate}
\tightlist
\item
  (\emph{Examples.}) The following are locally compact Hausdorff:
  compact Hausdorff spaces, \(\mathbb R^n\), discrete spaces, closed
  subspaces of LCH spaces, open subspaces of compact Hausdorff spaces,
  and open subspaces of LCH spaces. The spaces \(\mathbb Q\),
  \(\mathbb N^{\mathbb N}\) (a countable product of discrete copies of
  \(\mathbb N\)) and \(\mathbb R^{\mathbb N}\) (product topology) are
  Hausdorff, but no point in them has a compact neighbourhood.
\item
  (\emph{Compact closures.}) Let \(X\) be LCH, \(A\subseteq X\) compact
  and \(V\) a neighbourhood of \(A\). Then some open \(U\supseteq A\)
  has compact closure contained in \(V\).
\item
  (\emph{Exhaustion.}) If \(X\) is LCH and \(\sigma\)-compact, that is,
  a countable union of compact sets, then there are open sets
  \(U_1\subseteq U_2\subseteq\cdots\) with \(\overline{U_n}\) compact,
  \(\overline{U_n}\subseteq U_{n+1}\) and \(\bigcup_nU_n=X\). Every
  compact subset of \(X\) lies in some \(U_n\).
\item
  (\emph{Countability at infinity.}) For LCH \(X\), the following are
  equivalent: (i) \(\infty\) has a countable neighbourhood base in
  \(X_\infty\); (ii) \(\{\infty\}\) is a countable intersection of open
  subsets of \(X_\infty\); (iii) \(X\) is \(\sigma\)-compact; (iv) \(X\)
  is Lindelöf, that is, every open cover has a countable subcover.
\item
  (\emph{Variants of local compactness.}) For a topological space
  consider: (i) every point has a compact neighbourhood; (ii) every
  neighbourhood of a point contains a compact neighbourhood of that
  point; (iii) every point has an open neighbourhood with compact
  closure; (iii\(')\) every open neighbourhood \(V\) of a point \(p\)
  contains an open \(U\ni p\) with compact closure; (iv) every open
  neighbourhood \(V\) of \(p\) contains an open \(U\ni p\) with
  \(\overline U\) compact and \(\overline U\subseteq V\). Then (iv)
  implies (ii) and (iii); (ii) implies (i); and (iii) and (iii\(')\) are
  equivalent and imply (i). In a Hausdorff space all five are
  equivalent.
\end{enumerate}

\textbf{Proof.} (1) A compact Hausdorff space is a compact neighbourhood
of each of its points. In \(\mathbb R^n\) closed balls are compact by
Lemma 4.1(b). In a discrete space, \(\{x\}\) is a compact neighbourhood
of \(x\). If \(Y\subseteq X\) is closed and \(N\) is a compact
neighbourhood of \(p\in Y\) in \(X\), then \(N\cap Y\) is a compact
neighbourhood of \(p\) in \(Y\), by (K1). If \(U\) is open in a compact
Hausdorff space and \(p\in U\), Lemma 4.1(c) gives an open \(V\ni p\)
with \(\overline V\subseteq U\), and \(\overline V\) is compact by (K1).
An open subspace of an LCH space \(X\) is open in \(X_\infty\), which is
compact Hausdorff by
\hyperlink{the-stone-weierstrass-theorem-for-functions-vanishing-at-infinity--oa-fnd-sw-02}{Proposition
3.2}(2)--(3); so the previous case applies. Subspaces and products of
Hausdorff spaces are Hausdorff.

For \(\mathbb Q\): suppose \(N\) is a compact neighbourhood of
\(q\in\mathbb Q\). It contains \(\mathbb Q\cap(q-2r,q+2r)\) for some
\(r>0\), so it contains \(P=\mathbb Q\cap[q-r,q+r]\), which is closed
and hence compact by (K1). Pick an irrational \(\xi\in(q-r,q+r)\). The
open sets \(\{t\in P:|t-\xi|>1/m\}\), \(m\geq1\), cover \(P\) with no
finite subcover. For \(\mathbb N^{\mathbb N}\) and
\(\mathbb R^{\mathbb N}\): a neighbourhood of a point contains a basic
open set that leaves the coordinates beyond some index \(m\) free. The
projection onto coordinate \(m+1\) is continuous, so it maps a compact
neighbourhood onto a compact set (K2). Compact subsets of \(\mathbb N\)
are finite and compact subsets of \(\mathbb R\) are bounded, but the
image of the basic open set is all of \(\mathbb N\) or \(\mathbb R\).

(2) First let \(A=\{p\}\). The interior of \(V\) in \(X\) is open in
\(X_\infty\), so \(F=X_\infty\setminus\operatorname{int}V\) is closed,
and \(p\notin F\). \(X_\infty\) is compact Hausdorff
(\hyperlink{the-stone-weierstrass-theorem-for-functions-vanishing-at-infinity--oa-fnd-sw-02}{Proposition
3.2}(2)--(3)), so by Lemma 4.1(c) there is an open \(W\ni p\) whose
closure in \(X_\infty\) misses \(F\), that is, lies in
\(\operatorname{int}V\subseteq X\). Put \(U=W\cap X\). Its closure in
\(X_\infty\) is compact by (K1) and lies in \(X\). So it is also the
closure of \(U\) in \(X\), and it is compact in \(X\) and contained in
\(V\). For general compact \(A\), take such a \(U_p\) for each
\(p\in A\), a finite subcover \(U_{p_1},\dots,U_{p_k}\) of \(A\), and
\(U=\bigcup_iU_{p_i}\). Then \(\overline U=\bigcup_i\overline{U_{p_i}}\)
is compact and lies in \(V\).

(3) Let \(X=\bigcup_nK_n\) with each \(K_n\) compact. By (2), choose an
open \(U_1\supseteq K_1\) with compact closure, and inductively an open
\(U_{n+1}\) with compact closure that contains the compact set
\(\overline{U_n}\cup K_{n+1}\). Then \(\overline{U_n}\subseteq U_{n+1}\)
and \(\bigcup U_n=X\). A compact set is covered by the increasing open
sets \(U_n\), so it lies in one of them.

(4) (i) \(\Rightarrow\) (ii): let \(W_1,W_2,\ldots\) form a
neighbourhood base at \(\infty\). Their interiors are open, contain
\(\infty\), and again form a neighbourhood base. The intersection of the
interiors is \(\{\infty\}\), because \(X_\infty\) is Hausdorff
(\hyperlink{the-stone-weierstrass-theorem-for-functions-vanishing-at-infinity--oa-fnd-sw-02}{Proposition
3.2}(3)). (ii) \(\Rightarrow\) (iii): each open \(W_n\ni\infty\) has the
form \(X_\infty\setminus C_n\) with \(C_n\) compact, and
\(X=\bigcup C_n\). (iii) \(\Rightarrow\) (i): with \(U_n\) from (3), the
sets \(X_\infty\setminus\overline{U_n}\) form a neighbourhood base at
\(\infty\), because a neighbourhood \(X_\infty\setminus C\) of
\(\infty\) contains \(X_\infty\setminus\overline{U_n}\) as soon as
\(C\subseteq U_n\). (iii) \(\Rightarrow\) (iv): each \(K_n\) is covered
by finitely many members of an open cover. (iv) \(\Rightarrow\) (iii):
by (2) the open sets with compact closure cover \(X\); a countable
subcover \((V_n)\) gives \(X=\bigcup_n\overline{V_n}\).

(5) (iv) \(\Rightarrow\) (iii), (iii\(')\) \(\Rightarrow\) (iii) and
(ii) \(\Rightarrow\) (i): take \(V=X\). (iii) \(\Rightarrow\)
(iii\(')\): if \(W\ni p\) is open with compact closure and \(V\ni p\) is
open, then \(U=W\cap V\) has closure inside \(\overline W\), which is
compact by (K1). (iii) \(\Rightarrow\) (i): \(\overline W\) is a compact
neighbourhood. (iv) \(\Rightarrow\) (ii): \(\overline U\) is a compact
neighbourhood of \(p\) inside \(V\), and every neighbourhood of \(p\)
contains an open one. In a Hausdorff space, (i) says that the space is
LCH, and then (2) gives (iv). \(\square\)

Without the Hausdorff property the conditions in (5) need not agree. For
example, \(\mathbb Q_\infty\) is compact by
\hyperlink{the-stone-weierstrass-theorem-for-functions-vanishing-at-infinity--oa-fnd-sw-02}{Proposition
3.2}(2), so (i) holds at each of its points. But (ii) fails at every
rational point \(q\). Indeed, \(\mathbb Q\) is an open neighbourhood of
\(q\) in \(\mathbb Q_\infty\) that carries its own topology there
(\hyperlink{the-stone-weierstrass-theorem-for-functions-vanishing-at-infinity--oa-fnd-sw-02}{Proposition
3.2}(1)), so a compact neighbourhood of \(q\) inside \(\mathbb Q\) would
be a compact neighbourhood of \(q\) in \(\mathbb Q\), and by (1) there
is none. The open database {[}π-Base{]} collects many more examples of
this kind, with the properties of each space
(\href{https://topology.pi-base.org/properties/P000023}{spaces without
compact neighbourhoods}); see also, e.g., {[}Steen--Seebach{]}.

\subsection[5. Urysohn\textquotesingle s lemma and
metrizability]{\texorpdfstring{\protect\hypertarget{the-stone-weierstrass-theorem-for-functions-vanishing-at-infinity--oa-fnd-sw-16}{}{}5.
Urysohn\textquotesingle s lemma and
metrizability}{5. Urysohn\textquotesingle s lemma and metrizability}}\label{the-stone-weierstrass-theorem-for-functions-vanishing-at-infinity--OA-FND-SW-16}

\textbf{Theorem 5.1} (Urysohn\textquotesingle s lemma). Let \(X\) be a
space in which any two disjoint closed sets have disjoint open
neighbourhoods, and let \(E,F\subseteq X\) be disjoint closed sets. Then
there is a continuous \(f:X\to[0,1]\) with \(f=0\) on \(E\) and \(f=1\)
on \(F\). By
\hyperlink{the-stone-weierstrass-theorem-for-functions-vanishing-at-infinity--oa-fnd-sw-15}{Lemma
4.1}(c), this applies to every compact Hausdorff space.

\emph{Reference:} {[}Urysohn 1925{]}.

\textbf{Proof.} By the argument at the end of the proof of
\hyperlink{the-stone-weierstrass-theorem-for-functions-vanishing-at-infinity--oa-fnd-sw-15}{Lemma
4.1}(c), which uses only the separation of disjoint closed sets, every
closed set \(C\) inside an open set \(O\) has an open \(G\) with
\(C\subseteq G\subseteq\overline G\subseteq O\). Let \(D\) be the set of
dyadic rationals in \([0,1]\). We build open sets \(U_r\), \(r\in D\),
with
\[E\subseteq U_0,\qquad U_1=X\setminus F,\qquad\overline{U_r}\subseteq U_s\ \text{whenever }r<s .
\tag{5.1}\] Put \(U_1=X\setminus F\), and choose \(U_0\) open with
\(E\subseteq U_0\subseteq\overline{U_0}\subseteq U_1\). Suppose \(U_r\)
has been built for all \(r=j/2^n\), \(0\leq j\leq2^n\), with (5.1). For
\(r=(2j+1)/2^{n+1}\), its neighbours \(r_-=j/2^n\) and \(r_+=(j+1)/2^n\)
satisfy \(\overline{U_{r_-}}\subseteq U_{r_+}\); choose an open \(U_r\)
with
\(\overline{U_{r_-}}\subseteq U_r\subseteq\overline{U_r}\subseteq U_{r_+}\).
For \(r<s\) at level \(n+1\), (5.1) follows by passing through the
neighbours of \(r\) and \(s\) at level \(n\). Now define
\(f(x)=\inf\{r\in D:x\in U_r\}\), with \(\inf\varnothing=1\). Then
\(0\leq f\leq1\), \(f=0\) on \(E\), and \(f=1\) on \(F\), since \(F\)
misses \(U_1\supseteq U_r\) for every \(r\). For \(0<\alpha\leq1\) and
\(0\leq\beta<1\),
\[\{f<\alpha\}=\bigcup_{r<\alpha}U_r,\qquad\{f>\beta\}=\bigcup_{s>\beta}\bigl(X\setminus\overline{U_s}\bigr),\]
with \(r,s\in D\). For the first: if \(f(x)<\alpha\), some \(r<\alpha\)
has \(x\in U_r\), and conversely \(x\in U_r\) gives \(f(x)\leq r\). For
the second: if \(f(x)>\beta\), pick \(s,s'\in D\) with
\(\beta<s<s'<f(x)\); then \(x\notin U_{s'}\supseteq\overline{U_s}\).
Conversely, if \(x\notin\overline{U_s}\), then \(x\notin U_r\) for every
\(r\leq s\), so \(f(x)\geq s\). For other values of \(\alpha\) and
\(\beta\) the sets \(\{f<\alpha\}\) and \(\{f>\beta\}\) are
\(\varnothing\) or \(X\). So all these sets are open, and \(f\) is
continuous. \(\square\)

\textbf{Corollary 5.2} (Urysohn\textquotesingle s lemma, locally compact
form). Let \(X\) be LCH, \(A\subseteq X\) compact and \(B\subseteq X\)
closed with \(A\cap B=\varnothing\). Then some \(f\in C_c(X)\) with
\(0\leq f\leq1\) equals \(1\) on \(A\) and \(0\) on \(B\). In
particular, \(X\) is completely regular (take \(A=\{p\}\)), and
\(C_c(X)\), hence \(C_0(X)\), separates points and vanishes nowhere.
\hyperlink{the-stone-weierstrass-theorem-for-functions-vanishing-at-infinity--oa-fnd-sw-09}{Theorem
10.1} then gives a second proof that \(C_c(X)\) is dense in \(C_0(X)\).

\textbf{Proof.} By
\hyperlink{the-stone-weierstrass-theorem-for-functions-vanishing-at-infinity--oa-fnd-sw-15}{Proposition
4.2}(2) with \(V=X\setminus B\), choose an open \(U\supseteq A\) whose
closure \(K\) is compact and lies in \(X\setminus B\). \(K\) is compact
Hausdorff, so Urysohn\textquotesingle s lemma gives a continuous
\(g:K\to[0,1]\) with \(g=1\) on \(A\) and \(g=0\) on \(K\setminus U\).
Put \(f=g\) on \(K\) and \(f=0\) on \(X\setminus K\); so \(f=0\) on
\(X\setminus U\). For closed \(F\subseteq[0,1]\), \(f^{-1}(F)\) is
\(g^{-1}(F)\) if \(0\notin F\), and \(g^{-1}(F)\cup(X\setminus U)\) if
\(0\in F\). Both are closed in \(X\), since \(K\) is closed in \(X\). So
\(f\) is continuous, its support lies in \(K\), and \(f=0\) on \(B\).
\(\square\)

\textbf{Theorem 5.3} (Metrizability). For an LCH space \(X\), the
following are equivalent.

\begin{itemize}
\tightlist
\item
  (i) \(X\) is \(\sigma\)-compact and metrizable.
\item
  (ii) Some countable family \(\mathcal K\) of compact sets has this
  property: every neighbourhood of every point \(p\) contains a member
  of \(\mathcal K\) that is a neighbourhood of \(p\).
\item
  (iii) \(X\) is second countable.
\item
  (iv) \(X\) is Polish: separable and completely metrizable.
\item
  (v) \(X_\infty\) is metrizable.
\end{itemize}

\textbf{Lemma 5.4.} A compact Hausdorff space \(K\) with a countable
base is metrizable.

\textbf{Proof.} Let \((B_n)\) be a countable base. For each pair
\((m,n)\) with \(\overline{B_m}\subseteq B_n\),
Urysohn\textquotesingle s lemma gives a continuous \(f_{mn}:K\to[0,1]\)
that is \(1\) on \(\overline{B_m}\) and \(0\) off \(B_n\). List these
functions as \(h_1,h_2,\ldots\), and put
\(d(x,y)=\sum_k2^{-k}|h_k(x)-h_k(y)|\). Then \(d\) is symmetric,
satisfies the triangle inequality, and \(y\mapsto d(x,y)\) is
continuous, as a uniform limit of continuous functions. Let \(x\neq y\).
Since \(\{y\}\) is closed, some \(B_n\) contains \(x\) but not \(y\). By
\hyperlink{the-stone-weierstrass-theorem-for-functions-vanishing-at-infinity--oa-fnd-sw-15}{Lemma
4.1}(c), there is an open \(V\ni x\) with \(\overline V\subseteq B_n\),
and some \(B_m\) satisfies \(x\in B_m\subseteq V\). Then
\(\overline{B_m}\subseteq B_n\) and \(f_{mn}(x)=1\neq0=f_{mn}(y)\), so
\(d(x,y)>0\). So \(d\) is a metric. Its open balls are open in \(K\), so
the identity map from \(K\) to \((K,d)\) is a continuous bijection onto
a Hausdorff space, and a homeomorphism by (K4). \(\square\)

\textbf{Lemma 5.5.} Let \((M,d)\) be a complete metric space and
\(O\subseteq M\) open. Then \(O\) is completely metrizable.

\textbf{Proof.} If \(O=M\) there is nothing to prove. Otherwise let
\(F=M\setminus O\) and \(\phi(x)=1/d(x,F)\) for \(x\in O\); it is finite
and continuous, since \(F\) is closed. Put
\(d'(x,y)=d(x,y)+|\phi(x)-\phi(y)|\). This is a metric on \(O\), and it
gives the same topology: \(d\leq d'\), and \(d'(x,y)\to0\) as \(y\to x\)
in \(d\), by continuity of \(\phi\). Let \((x_j)\) be Cauchy for \(d'\).
It is Cauchy for \(d\), so it converges to some \(y\in M\). The numbers
\(\phi(x_j)\) form a Cauchy sequence, so they are bounded by some \(c\);
then \(d(x_j,F)\geq1/c\), and \(d(y,F)\geq1/c>0\). So \(y\in O\), and
\(d'(x_j,y)\to0\) by continuity of \(\phi\) at \(y\). \(\square\)

\textbf{Proof of Theorem 5.3.} (ii) \(\Rightarrow\) (iii): the interiors
of the members of \(\mathcal K\) form a countable base.

(iii) \(\Rightarrow\) (ii): let \((B_n)\) be a countable base, and
\(\mathcal K\) the family of those \(\overline{B_n}\) that are compact.
Given \(p\) and a neighbourhood \(V\),
\hyperlink{the-stone-weierstrass-theorem-for-functions-vanishing-at-infinity--oa-fnd-sw-15}{Proposition
4.2}(2) gives an open \(U\ni p\) with compact closure inside \(V\).
Choose \(B_n\) with \(p\in B_n\subseteq U\). Then
\(\overline{B_n}\subseteq\overline U\) is compact, lies in \(V\), and is
a neighbourhood of \(p\).

(iii) \(\Rightarrow\) (v): a second countable space is Lindelöf, because
each member of an open cover is a union of base sets, and each base set
that lies in some member can be assigned to one such member. So \(X\) is
\(\sigma\)-compact by
\hyperlink{the-stone-weierstrass-theorem-for-functions-vanishing-at-infinity--oa-fnd-sw-15}{Proposition
4.2}(4), and \(\infty\) has a countable neighbourhood base \((W_n)\).
Replacing each \(W_n\) by its interior, we may take the \(W_n\) open.
The sets \(B_n\) and \(W_n\) together form a countable base of
\(X_\infty\): an open \(O\ni\infty\) contains some \(W_m\), and each of
its points in \(X\) lies in some \(B_n\subseteq O\cap X\). \(X_\infty\)
is compact Hausdorff
(\hyperlink{the-stone-weierstrass-theorem-for-functions-vanishing-at-infinity--oa-fnd-sw-02}{Proposition
3.2}(2)--(3)), so Lemma 5.4 applies.

(v) \(\Rightarrow\) (i): \(X\) is a subspace of the metrizable space
\(X_\infty\), so it is metrizable. \(X_\infty\) is compact metrizable,
hence second countable by
\hyperlink{the-stone-weierstrass-theorem-for-functions-vanishing-at-infinity--oa-fnd-sw-15}{Lemma
4.1}(d). So \(\infty\) has a countable neighbourhood base, and \(X\) is
\(\sigma\)-compact by
\hyperlink{the-stone-weierstrass-theorem-for-functions-vanishing-at-infinity--oa-fnd-sw-15}{Proposition
4.2}(4).

(i) \(\Rightarrow\) (iii): \(X=\bigcup K_n\) with each \(K_n\) a compact
metric space, which is separable by
\hyperlink{the-stone-weierstrass-theorem-for-functions-vanishing-at-infinity--oa-fnd-sw-15}{Lemma
4.1}(d). So \(X\) has a countable dense set, and the balls of rational
radius about its points form a countable base.

(v) \(\Rightarrow\) (iv): \(X_\infty\) is a compact metric space, hence
complete and second countable
(\hyperlink{the-stone-weierstrass-theorem-for-functions-vanishing-at-infinity--oa-fnd-sw-15}{Lemma
4.1}(d)). \(X\) is open in \(X_\infty\), so it is completely metrizable
by Lemma 5.5. It is second countable as a subspace, hence separable.

(iv) \(\Rightarrow\) (iii): as in (i) \(\Rightarrow\) (iii), the balls
of rational radius about the points of a countable dense set form a
countable base. \(\square\)

Because \(X_\infty\) is compact, the proof needs a metrization argument
only for compact spaces (Lemma 5.4), not the general Urysohn metrization
theorem for regular second countable spaces. An uncountable discrete
space is metrizable and LCH but satisfies none of (i)--(v)
(\hyperlink{the-stone-weierstrass-theorem-for-functions-vanishing-at-infinity--oa-fnd-sw-13}{Example
11.2}).

\subsection[6. The polynomial step: absolute values without constant
terms]{\texorpdfstring{\protect\hypertarget{the-stone-weierstrass-theorem-for-functions-vanishing-at-infinity--oa-fnd-sw-04}{}{}6.
The polynomial step: absolute values without constant
terms}{6. The polynomial step: absolute values without constant terms}}\label{the-stone-weierstrass-theorem-for-functions-vanishing-at-infinity--OA-FND-SW-04}

\textbf{Lemma 6.1.} Define real polynomials by
\[P_0=0,\qquad P_{n+1}(t)=P_n(t)+\tfrac12\bigl(t^2-P_n(t)^2\bigr).
\tag{6.1}\]

\begin{enumerate}
\tightlist
\item
  \(P_n(t)=Q_n(t^2)\) for real polynomials \(Q_n\) with \(Q_n(0)=0\).
\item
  For \(|t|\leq1\): \(0\leq P_n(t)\leq P_{n+1}(t)\leq|t|\), and
  \[0\leq|t|-P_n(t)\leq|t|\bigl(1-\tfrac{|t|}2\bigr)^n\leq\frac2n\qquad(n\geq1).
  \tag{6.2}\]
\end{enumerate}

\textbf{Proof.} (1) Put \(Q_0=0\) and
\(Q_{n+1}(s)=Q_n(s)+\frac12(s-Q_n(s)^2)\). By induction
\(P_n(t)=Q_n(t^2)\), and \(Q_{n+1}(0)=Q_n(0)-\frac12Q_n(0)^2=0\).

(2) Fix \(t\) and write \(a=|t|\in[0,1]\) and \(p_n=P_n(t)\). Since
\(t^2=a^2\), (6.1) gives
\[a-p_{n+1}=(a-p_n)\Bigl(1-\frac{a+p_n}2\Bigr).\] Suppose
\(0\leq p_n\leq a\). Then \(0\leq\frac{a+p_n}2\leq a\leq1\), so
\(0\leq a-p_{n+1}\leq a-p_n\); that is, \(p_n\leq p_{n+1}\leq a\). Since
\(p_0=0\), induction gives \(0\leq p_n\leq p_{n+1}\leq a\) for all
\(n\). Using \(p_n\geq0\), the second factor above is at most \(1-a/2\),
so \(a-p_n\leq a(1-a/2)^n\). For the last bound, let \(0<a\leq1\) and
\(s=a/2\leq\frac12\). By Bernoulli\textquotesingle s inequality
\((1+u)^n\geq1+nu\) (\(u\geq0\)),
\[(1-s)^{-n}=\Bigl(1+\frac{s}{1-s}\Bigr)^n\geq1+\frac{ns}{1-s}\geq1+ns ,\]
so \(a(1-a/2)^n\leq a/(1+na/2)=2a/(2+na)<2/n\). \(\square\)

\textbf{Proposition 6.2.} Let \(A\) be a real algebra of bounded real
functions on a set \(S\).

\begin{enumerate}
\tightlist
\item
  If \(p\) is a real polynomial with \(p(0)=0\) and \(f\in A\), then
  \(p\circ f\in A\). If such polynomials \(p_k\) converge uniformly to a
  function \(\varphi\) on an interval containing \(f(S)\), then
  \(\varphi\circ f\in\bar A\).
\item
  Let \(f\in A\) and \(r>0\) with \(\|f\|\leq r\). Then
  \(g_n=r\,Q_n(f^2/r^2)\) lies in \(A\), and \(0\leq|f|-g_n\leq2r/n\).
  The polynomials \(t\mapsto mP_n(t/m)\) have no constant term and
  converge to \(|t|\) uniformly on \([-m,m]\).
\item
  \(\bar A\) is a real algebra, and \(f,g\in\bar A\) imply
  \(|f|,\ f\vee g,\ f\wedge g\in\bar A\). So \(\bar A\) is closed under
  maxima and minima of finitely many elements.
\end{enumerate}

\textbf{Proof.} (1) If \(p(t)=\sum_{k=1}^da_kt^k\), then
\(p\circ f=\sum_{k=1}^da_kf^k\in A\), since \(A\) is closed under
products and real linear combinations. No constant term is needed, so no
unit is needed. Also \(\|p_k\circ f-\varphi\circ f\|\) is at most the
supremum of \(|p_k-\varphi|\) on the interval.

(2) \(g_n\) is a real combination of powers \(f^{2k}\) with \(k\geq1\),
so it lies in \(A\). By Lemma 6.1 applied to \(t=f(s)/r\),
\(r\,P_n(f/r)=g_n\) satisfies \(0\leq|f|-g_n\leq2r/n\). The statement
about \([-m,m]\) is (6.2) after the substitution \(t\mapsto t/m\).

(3) \(\bar A\) is a real algebra by
\hyperlink{the-stone-weierstrass-theorem-for-functions-vanishing-at-infinity--oa-fnd-sw-01}{Proposition
2.1}(2). Apply (2) to \(f\in\bar A\): \(|f|\) is a uniform limit of
elements of \(\bar A\), so \(|f|\in\bar A\). Then
\(f\vee g=\frac12(f+g+|f-g|)\) and \(f\wedge g=\frac12(f+g-|f-g|)\), and
induction handles finitely many functions. \(\square\)

The recursion (6.1) keeps the constant term zero at every step. So the
approximation of \(|f|\) uses neither the constant function \(1\) nor
the Weierstrass approximation theorem, and this is what the nonunital
theorem needs.

\subsection[7. Interpolation at two
points]{\texorpdfstring{\protect\hypertarget{the-stone-weierstrass-theorem-for-functions-vanishing-at-infinity--oa-fnd-sw-05}{}{}7.
Interpolation at two
points}{7. Interpolation at two points}}\label{the-stone-weierstrass-theorem-for-functions-vanishing-at-infinity--OA-FND-SW-05}

\textbf{Lemma 7.1.} Let \(\mathbb F\) be \(\mathbb R\) or \(\mathbb C\),
and let \(A\) be an algebra of \(\mathbb F\)-valued functions on a set
\(S\) that separates points.

\begin{enumerate}
\tightlist
\item
  \(Z(A)\) has at most one point.
\item
  Let \(x,y\in S\) and \(a,b\in\mathbb F\). Assume \(a=0\) if
  \(x\in Z(A)\), \(b=0\) if \(y\in Z(A)\), and \(a=b\) if \(x=y\). Then
  some \(h\in A\) has \(h(x)=a\) and \(h(y)=b\).
\end{enumerate}

\textbf{Proof.} First let \(x=y\). If \(x\in Z(A)\), take \(h=0\).
Otherwise pick \(u\in A\) with \(u(x)\neq0\) and put \(h=(a/u(x))u\).

Now let \(x\neq y\). The set
\[V=\{(h(x),h(y)):\ h\in A\}\subseteq\mathbb F^2\] is a linear subspace
that is closed under coordinatewise multiplication, because
\(h\mapsto(h(x),h(y))\) respects sums, scalars and products. Such a
subspace is one of \(\{0\}\), \(\mathbb F(1,0)\), \(\mathbb F(0,1)\),
\(\mathbb F(1,1)\) and \(\mathbb F^2\). Indeed, if \(V=\mathbb Fv\) with
\(v=(v_1,v_2)\neq0\), then \((v_1^2,v_2^2)=c\,v\) for some \(c\). If
\(v_1\neq0\neq v_2\), this forces \(v_1=c=v_2\), so
\(V=\mathbb F(1,1)\). If \(v_1=0\) or \(v_2=0\), then \(V\) is
\(\mathbb F(0,1)\) or \(\mathbb F(1,0)\).

Separation of \(x\) and \(y\) rules out \(\{0\}\) and
\(\mathbb F(1,1)\), which would give \(h(x)=h(y)\) for every \(h\). If
\(V=\mathbb F(0,1)\), every \(h\) vanishes at \(x\), so \(x\in Z(A)\),
\(a=0\), and \((a,b)\in V\). The case \(V=\mathbb F(1,0)\) is symmetric.
If \(V=\mathbb F^2\) there is nothing to prove. This proves (2). If
\(x\neq y\) both lay in \(Z(A)\), then \(V=\{0\}\); so (1) holds.
\(\square\)

When \(A\) vanishes nowhere, (2) says that every pair of values is
attained at every pair of distinct points. The proof looks only at the
image of \(A\) in \(\mathbb F^2\). It also shows what happens at the one
possible common zero, which the nonunital case needs.

\subsection[8. Stone\textquotesingle s lattice
lemma]{\texorpdfstring{\protect\hypertarget{the-stone-weierstrass-theorem-for-functions-vanishing-at-infinity--oa-fnd-sw-06}{}{}8.
Stone\textquotesingle s lattice
lemma}{8. Stone\textquotesingle s lattice lemma}}\label{the-stone-weierstrass-theorem-for-functions-vanishing-at-infinity--OA-FND-SW-06}

\textbf{Lemma 8.1} (Stone\textquotesingle s lattice lemma, pointwise
form). Let \(K\) be a nonempty compact space and
\(L\subseteq C(K;\mathbb R)\) a set that is closed under \(\vee\) and
\(\wedge\). Let \(f\in C(K;\mathbb R)\) and \(\varepsilon>0\). Suppose
that for all \(x,y\in K\) some \(h\in L\) satisfies
\[|h(x)-f(x)|<\varepsilon\quad\text{and}\quad|h(y)-f(y)|<\varepsilon .
\tag{8.1}\] Then some \(h\in L\) satisfies \(\|h-f\|<\varepsilon\).

\textbf{Proof.} \emph{Step 1: approximation from above, anchored at one
point.} Fix \(x\in K\). Let \(\mathcal U_x\) be the family of open sets
\(U\subseteq K\) for which some \(h\in L\) satisfies
\(h(x)>f(x)-\varepsilon\) and \(h<f+\varepsilon\) on \(U\). Every
\(y\in K\) lies in a member of \(\mathcal U_x\): take \(h\) from (8.1)
for the pair \((x,y)\) and \(U=\{h<f+\varepsilon\}\). Choose a finite
subcover \(U_1,\dots,U_m\) and, for each \(U_j\), one function \(h_j\)
as in the definition. Put \(g_x=h_1\wedge\dots\wedge h_m\in L\). Each
point of \(K\) lies in some \(U_j\), where
\(g_x\leq h_j<f+\varepsilon\). So \(g_x<f+\varepsilon\) on \(K\), and
\(g_x(x)>f(x)-\varepsilon\).

\emph{Step 2: approximation from below.} Let \(\mathcal V\) be the
family of open sets \(V\subseteq K\) for which some \(g\in L\) satisfies
\(g<f+\varepsilon\) on \(K\) and \(g>f-\varepsilon\) on \(V\). By Step
1, every \(x\) lies in the member \(\{g_x>f-\varepsilon\}\) of
\(\mathcal V\). Choose a finite subcover \(V_1,\dots,V_n\), one function
\(g_i\) for each \(V_i\), and put \(h=g_1\vee\dots\vee g_n\in L\). Then
\(h<f+\varepsilon\) everywhere, because every \(g_i\) is, and
\(h>f-\varepsilon\) everywhere, because every point lies in some
\(V_i\), where \(h\geq g_i>f-\varepsilon\). By (K5) the continuous
function \(|h-f|\) attains its maximum, so \(\|h-f\|<\varepsilon\).
\(\square\)

The proof makes only finitely many choices at each step. The covers are
families of open sets defined by the existence of suitable functions,
and no function is chosen for each point of \(K\). So the argument does
not use the axiom of choice.

\textbf{Theorem 8.2} (Lattice versions). Let \(K\) be a compact space
and \(L\subseteq C(K;\mathbb R)\) a linear subspace closed under
\(\wedge\) (then also under \(\vee\), since
\(f\vee g=-((-f)\wedge(-g))\)).

\begin{enumerate}
\tightlist
\item
  Suppose that (i) for \(x\neq y\) in \(K\) and \(a,b\in\mathbb R\),
  some \(h\in L\) has \(h(x)=a\) and \(h(y)=b\), and (ii) for \(x\in K\)
  and \(a\in\mathbb R\), some \(h\in L\) has \(h(x)=a\). Then \(L\) is
  dense in \(C(K;\mathbb R)\). If \(K\) has at least two points, (i)
  implies (ii).
\item
  If \(L\) separates points and every constant function lies in \(L\),
  then (i) and (ii) hold, so \(L\) is dense.
\item
  Let \(Y\subseteq C(K)\) be a complex linear subspace with \(1\in Y\)
  that separates points and satisfies \(\bar g\in Y\) and \(|g|\in Y\)
  for every \(g\in Y\). Then \(Y\) is dense in \(C(K)\).
\item
  (\emph{Nonunital form}) Let \(X\) be a topological space and
  \(L\subseteq C_0(X;\mathbb R)\) a linear subspace closed under
  \(\wedge\). If (i) and (ii) hold for points of \(X\), then \(L\) is
  dense in \(C_0(X;\mathbb R)\). If \(X\) is locally compact Hausdorff,
  the converse also holds.
\end{enumerate}

\emph{Reference:} {[}Blackadar, XX.10.2.14{]} assumes only condition (i)
in (1). That version fails on a one-point space, where \(L=\{0\}\)
satisfies (i) vacuously and is not dense.

\textbf{Proof.} (1) If \(K=\varnothing\) there is nothing to prove.
Otherwise, given \(f\) and \(x,y\), conditions (i)--(ii) give \(h\in L\)
with \(h(x)=f(x)\) and \(h(y)=f(y)\), and Lemma 8.1 applies. If \(K\)
has a second point \(y\neq x\), then (i) with the pair \((x,y)\) gives
(ii).

(2) For \(x\neq y\) pick \(f\in L\) with \(f(x)\neq f(y)\). The vectors
\((1,1)\) and \((f(x),f(y))\) in \(\mathbb R^2\) are linearly
independent, and both lie in the image of the linear map
\(h\mapsto(h(x),h(y))\). So the image is \(\mathbb R^2\), which is (i).
Constants give (ii).

(3) For \(g\in Y\), \(\operatorname{Re}g=\frac12(g+\bar g)\) and
\(\operatorname{Im}g=\frac1{2i}(g-\bar g)\) lie in \(Y\). Let
\(Y_{\mathbb R}\) be the real functions in \(Y\). It is a real subspace
that contains \(1\) and separates points, since \(g(x)\neq g(y)\) forces
a difference in the real or the imaginary part. It is closed under
\(|\cdot|\), hence under \(\wedge\), because
\(f\wedge g=\frac12(f+g-|f-g|)\)
(\hyperlink{the-stone-weierstrass-theorem-for-functions-vanishing-at-infinity--oa-fnd-sw-04}{Proposition
6.2}(3)). By (2), \(Y_{\mathbb R}\) is dense in \(C(K;\mathbb R)\).
Approximating the real and imaginary parts of \(f\in C(K)\) gives
density of \(Y\).

(4) Suppose (i) and (ii) hold. The set
\(\tilde L=\{\tilde h:h\in L\}\subseteq C(X_\infty;\mathbb R)\) of
\hyperlink{the-stone-weierstrass-theorem-for-functions-vanishing-at-infinity--oa-fnd-sw-03}{Proposition
3.4} is closed under \(\vee\) and \(\wedge\), since these operations
keep the value \(0\) at \(\infty\). Let \(f\in C_0(X;\mathbb R)\). For a
pair of points of \(X\), (i)--(ii) give an \(h\) that agrees with \(f\)
there. For a pair \((x,\infty)\), (ii) gives \(h\) with \(h(x)=f(x)\),
and \(\tilde h(\infty)=0=\tilde f(\infty)\). For the pair
\((\infty,\infty)\), take \(h=0\). Lemma 8.1, applied on the compact
space \(X_\infty\)
(\hyperlink{the-stone-weierstrass-theorem-for-functions-vanishing-at-infinity--oa-fnd-sw-02}{Proposition
3.2}(2)), puts \(\tilde f\) in the closure of \(\tilde L\). Since
\(h\mapsto\tilde h\) preserves the norm
(\hyperlink{the-stone-weierstrass-theorem-for-functions-vanishing-at-infinity--oa-fnd-sw-03}{Proposition
3.4}(2)), \(f\) lies in the closure of \(L\). Conversely, let \(X\) be
locally compact Hausdorff and \(L\) dense. For \(x\neq y\), the
evaluation \(h\mapsto(h(x),h(y))\) is continuous, and it maps
\(C_0(X;\mathbb R)\) onto \(\mathbb R^2\) by Urysohn\textquotesingle s
lemma
(\hyperlink{the-stone-weierstrass-theorem-for-functions-vanishing-at-infinity--oa-fnd-sw-16}{Corollary
5.2}). So it maps \(L\) onto a dense linear subspace of \(\mathbb R^2\),
which is all of \(\mathbb R^2\). The one-point case is the same. Without
local compactness the converse fails: for \(X=\mathbb Q\) one has
\(C_0(X)=\{0\}\)
(\hyperlink{the-stone-weierstrass-theorem-for-functions-vanishing-at-infinity--oa-fnd-sw-10}{Proposition
12.1}), so \(L=\{0\}\) is dense while (ii) fails. \(\square\)

Lemma 8.1 needs the interpolation (8.1) only for the one function \(f\)
that is being approximated, not for all target values. Part (4) of
Theorem 8.2 is the lattice counterpart of
\hyperlink{the-stone-weierstrass-theorem-for-functions-vanishing-at-infinity--oa-fnd-sw-08}{Theorem
9.2}.

\subsection[9. The real Stone--Weierstrass
theorem]{\texorpdfstring{\protect\hypertarget{the-stone-weierstrass-theorem-for-functions-vanishing-at-infinity--oa-fnd-sw-07}{}{}\protect\hypertarget{the-stone-weierstrass-theorem-for-functions-vanishing-at-infinity--OA-FND-SW-08}{}{}\protect\hypertarget{the-stone-weierstrass-theorem-for-functions-vanishing-at-infinity--oa-fnd-sw-08}{}{}9.
The real Stone--Weierstrass
theorem}{9. The real Stone--Weierstrass theorem}}\label{the-stone-weierstrass-theorem-for-functions-vanishing-at-infinity--OA-FND-SW-07}

\textbf{Theorem 9.1} (Stone--Weierstrass theorem, real form). Let \(K\)
be a compact space and \(A\subseteq C(K;\mathbb R)\) a real algebra that
separates points. Then \(Z(A)\) has at most one point, and
\[\bar A=\{f\in C(K;\mathbb R):\ f=0\ \text{on }Z(A)\}.
\tag{9.1}\] In particular:

\begin{itemize}
\tightlist
\item
  (a) if \(A\) vanishes nowhere, for example if \(A\) contains the
  constant functions, then \(A\) is dense in \(C(K;\mathbb R)\);
\item
  (b) if every function in \(A\) vanishes at a point \(z\), then
  \(\bar A=\{f:\ f(z)=0\}\).
\end{itemize}

\emph{Reference:} {[}Stone 1937{]}, {[}Stone 1948{]}.

\textbf{Proof.} \(Z(A)\) has at most one point by
\hyperlink{the-stone-weierstrass-theorem-for-functions-vanishing-at-infinity--oa-fnd-sw-05}{Lemma
7.1}(1). Evaluation at a point is continuous for \(\|\cdot\|\), so
\(\bar A\) lies in the right side of (9.1). Conversely, let
\(f\in C(K;\mathbb R)\) vanish on \(Z(A)\), and assume
\(K\neq\varnothing\) (otherwise both sides are \(\{0\}\)). Put
\(B=\bar A\). By
\hyperlink{the-stone-weierstrass-theorem-for-functions-vanishing-at-infinity--oa-fnd-sw-04}{Proposition
6.2}(3), \(B\) is closed under \(\vee\) and \(\wedge\). For
\(x,y\in K\), apply
\hyperlink{the-stone-weierstrass-theorem-for-functions-vanishing-at-infinity--oa-fnd-sw-05}{Lemma
7.1}(2) to \(A\) with \(a=f(x)\) and \(b=f(y)\). Its conditions hold,
because \(f(x)=0\) whenever \(x\in Z(A)\), and likewise for \(y\). So
some \(h\in A\subseteq B\) agrees with \(f\) at \(x\) and \(y\).
\hyperlink{the-stone-weierstrass-theorem-for-functions-vanishing-at-infinity--oa-fnd-sw-06}{Lemma
8.1} gives, for each \(\varepsilon>0\), an \(h\in B\) with
\(\|h-f\|<\varepsilon\). So \(f\in\bar B=B\). \(\square\)

\emph{Where the hypotheses enter.} Products are used in
\hyperlink{the-stone-weierstrass-theorem-for-functions-vanishing-at-infinity--oa-fnd-sw-04}{Proposition
6.2} (the functions \(g_n\)) and in
\hyperlink{the-stone-weierstrass-theorem-for-functions-vanishing-at-infinity--oa-fnd-sw-05}{Lemma
7.1} (the set \(V\) is closed under products). Separation is used in
\hyperlink{the-stone-weierstrass-theorem-for-functions-vanishing-at-infinity--oa-fnd-sw-05}{Lemma
7.1}, to rule out \(V=\{0\}\) and \(V=\mathbb R(1,1)\). Compactness is
used in
\hyperlink{the-stone-weierstrass-theorem-for-functions-vanishing-at-infinity--oa-fnd-sw-06}{Lemma
8.1}, for the two finite subcovers, and in (K5). The Hausdorff property
is not used. It follows anyway from separation: if \(h(x)\neq h(y)\),
the sets \(\{|h-h(x)|<r\}\) and \(\{|h-h(y)|<r\}\) with
\(r=|h(x)-h(y)|/2\) are disjoint open neighbourhoods.

For a closed algebra \(A\) that separates points, (9.1) gives a
dichotomy: either \(A=C(K;\mathbb R)\), or \(A\) consists of all the
functions that vanish at one point.
\hyperref[the-stone-weierstrass-theorem-for-functions-vanishing-at-infinity--exercises]{Exercise
2} proves the case with a common zero in another way, by adjoining the
constants.

The passage to \(C_0(X)\) goes through the one-point compactification.

\textbf{Theorem 9.2} (Real form for \(C_0(X)\)). Let \(X\) be locally
compact Hausdorff, and let \(A\subseteq C_0(X;\mathbb R)\) be a real
algebra that separates the points of \(X\) and vanishes nowhere. Then
\(A\) is dense in \(C_0(X;\mathbb R)\).

\textbf{Proof.} By
\hyperlink{the-stone-weierstrass-theorem-for-functions-vanishing-at-infinity--oa-fnd-sw-03}{Proposition
3.4}, \(\tilde A=\{\tilde f:f\in A\}\) is a real algebra of continuous
functions on the compact space \(X_\infty\)
(\hyperlink{the-stone-weierstrass-theorem-for-functions-vanishing-at-infinity--oa-fnd-sw-02}{Proposition
3.2}(2)). It separates the points of \(X_\infty\). Two points of \(X\)
are separated by \(A\). A point \(x\in X\) and \(\infty\) are separated
by any \(f\in A\) with \(f(x)\neq0\), since \(\tilde f(\infty)=0\); such
an \(f\) exists because \(A\) vanishes nowhere. The same fact gives
\(Z(\tilde A)=\{\infty\}\). By
\hyperlink{the-stone-weierstrass-theorem-for-functions-vanishing-at-infinity--oa-fnd-sw-07}{Theorem
9.1}(b), the closure of \(\tilde A\) is the set of real functions in
\(I_\infty\). Since \(f\mapsto\tilde f\) is an isometric bijection of
\(C_0(X;\mathbb R)\) onto that set
(\hyperlink{the-stone-weierstrass-theorem-for-functions-vanishing-at-infinity--oa-fnd-sw-03}{Proposition
3.4}(2)), \(A\) is dense in \(C_0(X;\mathbb R)\). \(\square\)

The proof uses neither local compactness nor the Hausdorff property of
\(X\).
\hyperlink{the-stone-weierstrass-theorem-for-functions-vanishing-at-infinity--oa-fnd-sw-10}{Proposition
12.1} explains why: both follow from the hypotheses on \(A\).

\subsection[10. The complex Stone--Weierstrass
theorem]{\texorpdfstring{\protect\hypertarget{the-stone-weierstrass-theorem-for-functions-vanishing-at-infinity--oa-fnd-sw-09}{}{}10.
The complex Stone--Weierstrass
theorem}{10. The complex Stone--Weierstrass theorem}}\label{the-stone-weierstrass-theorem-for-functions-vanishing-at-infinity--OA-FND-SW-09}

\textbf{Theorem 10.1} (Stone--Weierstrass theorem for \(C_0(X)\)). Let
\(X\) be locally compact Hausdorff, and let \(A\subseteq C_0(X)\) be a
complex subalgebra that is self-adjoint, separates the points of \(X\),
and vanishes nowhere. Then \[\bar A=C_0(X):
\tag{10.1}\] every \(f\in C_0(X)\) can be approximated uniformly on
\(X\) by elements of \(A\).

\emph{Compact form.} Let \(K\) be a compact space and
\(A\subseteq C(K)\) a self-adjoint subalgebra that separates points.
Then \(Z(A)\) has at most one point, and
\(\bar A=\{f\in C(K):\ f=0\ \text{on }Z(A)\}\). In particular \(A\) is
dense in \(C(K)\) if it vanishes nowhere, for example if it contains the
constants.

\textbf{Proof.} Let \(A_{\mathbb R}\) be the set of real functions in
\(A\). For \(f\in A\), self-adjointness and complex linearity give
\[\operatorname{Re}f=\tfrac12(f+\bar f)\in A_{\mathbb R},\qquad\operatorname{Im}f=\tfrac1{2i}(f-\bar f)\in A_{\mathbb R}.\]
\(A_{\mathbb R}\) is a real algebra. It separates points: if
\(f(x)\neq f(y)\), then \(\operatorname{Re}f\) or \(\operatorname{Im}f\)
takes different values at \(x\) and \(y\). For the same reason
\(Z(A_{\mathbb R})=Z(A)\). In the setting of the theorem,
\hyperlink{the-stone-weierstrass-theorem-for-functions-vanishing-at-infinity--oa-fnd-sw-08}{Theorem
9.2} shows that \(A_{\mathbb R}\) is dense in \(C_0(X;\mathbb R)\). Let
\(f\in C_0(X)\) and \(\varepsilon>0\). Since
\(|\operatorname{Re}f|\leq|f|\) and \(|\operatorname{Im}f|\leq|f|\), for
every \(\delta>0\) the sets \(\{|\operatorname{Re}f|\geq\delta\}\) and
\(\{|\operatorname{Im}f|\geq\delta\}\) are closed subsets of the compact
set \(\{|f|\geq\delta\}\), hence compact by (K1). So both parts lie in
\(C_0(X;\mathbb R)\). Choose \(u,v\in A_{\mathbb R}\) within
\(\varepsilon/2\) of them. Then \(u+iv\in A\) and
\(\|f-(u+iv)\|<\varepsilon\). In the compact form,
\hyperlink{the-stone-weierstrass-theorem-for-functions-vanishing-at-infinity--oa-fnd-sw-07}{Theorem
9.1} applies to \(A_{\mathbb R}\) in the same way: a function \(f\)
vanishing on \(Z(A)\) has real and imaginary parts vanishing on
\(Z(A_{\mathbb R})\). \(\square\)

\emph{Where each hypothesis is used.}

\begin{itemize}
\tightlist
\item
  \emph{Subalgebra.} Closure under products enters in
  \hyperlink{the-stone-weierstrass-theorem-for-functions-vanishing-at-infinity--oa-fnd-sw-04}{Proposition
  6.2} and
  \hyperlink{the-stone-weierstrass-theorem-for-functions-vanishing-at-infinity--oa-fnd-sw-05}{Lemma
  7.1}. Complex linearity gives
  \(\operatorname{Re}f,\operatorname{Im}f\in A\) and \(u+iv\in A\).
  \hyperlink{the-stone-weierstrass-theorem-for-functions-vanishing-at-infinity--oa-fnd-sw-11}{Example
  12.2}(d) shows that products cannot be dropped.
\item
  \emph{Self-adjoint.} Used once, to put \(\operatorname{Re}f\) and
  \(\operatorname{Im}f\) into \(A\).
  \hyperlink{the-stone-weierstrass-theorem-for-functions-vanishing-at-infinity--oa-fnd-sw-11}{Example
  12.2}(a) and Exercise 1 show that it cannot be dropped.
\item
  \emph{Separates points.} Used in
  \hyperlink{the-stone-weierstrass-theorem-for-functions-vanishing-at-infinity--oa-fnd-sw-05}{Lemma
  7.1}, for pairs of points of \(X\).
  \hyperlink{the-stone-weierstrass-theorem-for-functions-vanishing-at-infinity--oa-fnd-sw-11}{Example
  12.2}(b) shows that it cannot be dropped.
\item
  \emph{Vanishes nowhere.} Used in
  \hyperlink{the-stone-weierstrass-theorem-for-functions-vanishing-at-infinity--oa-fnd-sw-08}{Theorem
  9.2}, to separate each \(x\in X\) from \(\infty\) and to make
  \(\infty\) the only common zero on \(X_\infty\).
  \hyperlink{the-stone-weierstrass-theorem-for-functions-vanishing-at-infinity--oa-fnd-sw-11}{Example
  12.2}(c) shows that it cannot be dropped.
\item
  \emph{\(A\subseteq C_0(X)\).} Makes the functions \(\tilde f\)
  continuous on \(X_\infty\)
  (\hyperlink{the-stone-weierstrass-theorem-for-functions-vanishing-at-infinity--oa-fnd-sw-03}{Proposition
  3.4}).
\item
  \emph{Locally compact Hausdorff.} Not used in the proof. It is the
  natural setting because the other hypotheses force it
  (\hyperlink{the-stone-weierstrass-theorem-for-functions-vanishing-at-infinity--oa-fnd-sw-10}{Proposition
  12.1}).
\item
  \emph{Compactness of \(X_\infty\).} This is where the topology does
  its work, in the finite subcovers of
  \hyperlink{the-stone-weierstrass-theorem-for-functions-vanishing-at-infinity--oa-fnd-sw-06}{Lemma
  8.1}.
\end{itemize}

\emph{No countability is needed.} The proof uses no sequences of points
and no countable bases, only finite subcovers. So the theorem applies to
spaces that are neither \(\sigma\)-compact nor metrizable, such as
uncountable discrete spaces
(\hyperlink{the-stone-weierstrass-theorem-for-functions-vanishing-at-infinity--oa-fnd-sw-13}{Example
11.2}) and locally compact groups of this kind.

\subsection[11. Worked
examples]{\texorpdfstring{\protect\hypertarget{the-stone-weierstrass-theorem-for-functions-vanishing-at-infinity--oa-fnd-sw-13}{}{}11.
Worked
examples}{11. Worked examples}}\label{the-stone-weierstrass-theorem-for-functions-vanishing-at-infinity--OA-FND-SW-13}

Several examples here and in Section 12 use two functions on
\(\mathbb R\):
\[\varphi(x)=\frac1{1+x^2},\qquad\psi(x)=\frac{x}{1+x^2}.\] Both lie in
\(C_0(\mathbb R)\): since \(|\psi(x)|\leq(1+x^2)^{-1/2}\), the sets
\(\{|\varphi|\geq\varepsilon\}\) and \(\{|\psi|\geq\varepsilon\}\) are
closed and lie in \(\{1+x^2\leq\varepsilon^{-2}\}\), so they are compact
(\hyperlink{the-stone-weierstrass-theorem-for-functions-vanishing-at-infinity--oa-fnd-sw-15}{Lemma
4.1}(a)). Also \(\varphi>0\), \(\varphi\) is even, \(\psi\) is odd, and
\(\psi/\varphi=x\).

\textbf{Example 11.1} (A compact space without the constants). On
\([0,1]\), let \(A=\{p(x+1):\ p\text{ a real polynomial},\ p(0)=0\}\).
\(A\) is a real algebra. It vanishes nowhere, because \(x+1\geq1\), and
\(x+1\) separates points. So \(A\) is dense in \(C([0,1];\mathbb R)\) by
\hyperlink{the-stone-weierstrass-theorem-for-functions-vanishing-at-infinity--oa-fnd-sw-07}{Theorem
9.1}(a). Yet \(1\notin A\): if \(p(0)=0\) and \(p(y)=1\) for
\(y\in[1,2]\), the polynomial \(p-1\) would have infinitely many roots
and constant term \(-1\). An explicit approximation is
\(e_n(x)=1-\bigl(1-\frac{x+1}2\bigr)^n\in A\); since
\(0\leq1-\frac{x+1}2\leq\frac12\) on \([0,1]\), \(|1-e_n|\leq2^{-n}\).
On a compact space, "vanishes nowhere" is strictly weaker than "contains
the constants".

\textbf{Example 11.2} (An uncountable discrete space). Let \(X\) be an
uncountable set with the discrete topology. Its compact subsets are
finite, because an infinite subset has an open cover by singletons with
no finite subcover. So \(C_0(X)\) consists of the \(f\) for which
\(\{|f|\geq\varepsilon\}\) is finite for every \(\varepsilon>0\). The
finitely supported functions form a self-adjoint subalgebra that
separates points and vanishes nowhere (it contains the indicator of each
point), so it is dense. Here \(X\) is not \(\sigma\)-compact, and
\(\infty\) has no countable neighbourhood base in \(X_\infty\)
(\hyperlink{the-stone-weierstrass-theorem-for-functions-vanishing-at-infinity--oa-fnd-sw-15}{Proposition
4.2}(4)). The proof of
\hyperlink{the-stone-weierstrass-theorem-for-functions-vanishing-at-infinity--oa-fnd-sw-09}{Theorem
10.1} never uses countability, so this does not matter.

\textbf{Example 11.3} (Degenerate spaces). If \(X=\varnothing\), then
\(C_0(X)=\{0\}\), and \(A=\{0\}\) satisfies all hypotheses vacuously;
the conclusion holds. If \(X\) is one point, \(C_0(X)=\mathbb C\); an
algebra that vanishes nowhere contains a nonzero constant, hence equals
\(\mathbb C\).

\textbf{Example 11.4} (Compact \(X\)). If \(X\) is compact, then
\(C_0(X)=C(X)\), because \(\{|f|\geq\varepsilon\}\) is closed in \(X\)
and so compact by (K1). The point \(\infty\) is isolated in \(X_\infty\)
(\hyperlink{the-stone-weierstrass-theorem-for-functions-vanishing-at-infinity--oa-fnd-sw-02}{Proposition
3.2}(4)), and
\hyperlink{the-stone-weierstrass-theorem-for-functions-vanishing-at-infinity--oa-fnd-sw-09}{Theorem
10.1} reduces to its compact form.

\textbf{Example 11.5} (The real line). The complex algebra generated by
\(\varphi\) and \(\psi\) consists of the complex polynomials without
constant term in \(\varphi\) and \(\psi\). It is self-adjoint, because
\(\varphi\) and \(\psi\) are real. It separates points, because
\(\psi/\varphi=x\), and it vanishes nowhere, because \(\varphi>0\). So
it is dense in \(C_0(\mathbb R)\).

\textbf{Example 11.6} (One common zero, seen twice). On \([0,1]\), the
complex polynomials without constant term have closure
\(\{f\in C([0,1]):f(0)=0\}\), by the compact form of
\hyperlink{the-stone-weierstrass-theorem-for-functions-vanishing-at-infinity--oa-fnd-sw-09}{Theorem
10.1}; for real coefficients this is
\hyperlink{the-stone-weierstrass-theorem-for-functions-vanishing-at-infinity--oa-fnd-sw-07}{Theorem
9.1}(b). Restricted to \(X=(0,1]\) they form a self-adjoint subalgebra
of \(C_0(X)\) that separates points and vanishes nowhere, so
\hyperlink{the-stone-weierstrass-theorem-for-functions-vanishing-at-infinity--oa-fnd-sw-09}{Theorem
10.1} makes them dense in \(C_0((0,1])\). The two statements agree: by
\hyperlink{the-stone-weierstrass-theorem-for-functions-vanishing-at-infinity--oa-fnd-sw-02}{Proposition
3.2}(5) with \(Y=[0,1]\) and \(q=0\), and by
\hyperlink{the-stone-weierstrass-theorem-for-functions-vanishing-at-infinity--oa-fnd-sw-03}{Proposition
3.4}, \(C_0((0,1])\) is the space of continuous functions on \([0,1]\)
that vanish at \(0\).

\subsection[12. The role of each
hypothesis]{\texorpdfstring{\protect\hypertarget{the-stone-weierstrass-theorem-for-functions-vanishing-at-infinity--oa-fnd-sw-10}{}{}\protect\hypertarget{the-stone-weierstrass-theorem-for-functions-vanishing-at-infinity--OA-FND-SW-11}{}{}\protect\hypertarget{the-stone-weierstrass-theorem-for-functions-vanishing-at-infinity--oa-fnd-sw-11}{}{}12.
The role of each
hypothesis}{12. The role of each hypothesis}}\label{the-stone-weierstrass-theorem-for-functions-vanishing-at-infinity--OA-FND-SW-10}

The topological hypothesis of Theorem 10.1 comes for free, and each
hypothesis on the algebra is needed.

\textbf{Proposition 12.1} (The topological hypotheses are automatic).
Let \(X\) be a topological space.

\begin{enumerate}
\tightlist
\item
  If some set of functions in \(C_0(X)\) separates the points of \(X\),
  then \(X\) is Hausdorff.
\item
  If \(f\in C_0(X)\) and \(f(x)\neq0\), then \(x\) has a compact
  neighbourhood. So if some \(A\subseteq C_0(X)\) vanishes nowhere,
  every point of \(X\) has a compact neighbourhood. If no point of \(X\)
  has a compact neighbourhood, then \(C_0(X)=\{0\}\); this is the case
  for \(X=\mathbb Q\)
  (\hyperlink{the-stone-weierstrass-theorem-for-functions-vanishing-at-infinity--oa-fnd-sw-15}{Proposition
  4.2}(1)).
\item
  Consequently
  \hyperlink{the-stone-weierstrass-theorem-for-functions-vanishing-at-infinity--oa-fnd-sw-08}{Theorem
  9.2} and
  \hyperlink{the-stone-weierstrass-theorem-for-functions-vanishing-at-infinity--oa-fnd-sw-09}{Theorem
  10.1} hold for every topological space \(X\): whenever their
  hypotheses on \(A\) hold, \(X\) is locally compact Hausdorff.
\end{enumerate}

\textbf{Proof.} (1) If \(h(x)\neq h(y)\), put \(r=|h(x)-h(y)|/2\). The
open sets \(\{|h-h(x)|<r\}\) and \(\{|h-h(y)|<r\}\) are disjoint and
contain \(x\) and \(y\).

(2) Put \(c=|f(x)|/2>0\). The set \(\{|f|\geq c\}\) is compact by (1.1),
and it contains the open set \(\{|f|>c\}\), which contains \(x\).

(3) By (1) and (2), \(X\) is Hausdorff and every point has a compact
neighbourhood. \(\square\)

Part (2) gives condition (i) of
\hyperlink{the-stone-weierstrass-theorem-for-functions-vanishing-at-infinity--oa-fnd-sw-15}{Proposition
4.2}(5), one of several conditions that serve as definitions of local
compactness for spaces that need not be Hausdorff. Together with the
Hausdorff property from (1), all of them agree. So nothing is lost by
assuming that \(X\) is locally compact Hausdorff. Conversely, on every
locally compact Hausdorff space the algebra \(C_0(X)\) itself separates
points and vanishes nowhere, by Urysohn\textquotesingle s lemma
(\hyperlink{the-stone-weierstrass-theorem-for-functions-vanishing-at-infinity--oa-fnd-sw-16}{Corollary
5.2}).

\textbf{Examples 12.2.} In examples (a)--(d), all hypotheses of
\hyperlink{the-stone-weierstrass-theorem-for-functions-vanishing-at-infinity--oa-fnd-sw-09}{Theorem
10.1} hold except one, and the conclusion fails. Examples (e)--(h) test
the lattice versions, the hypothesis \(A\subseteq C_0(X)\), the
topological hypothesis, and the form of the conclusion. Several of them
use the functions \(\varphi\) and \(\psi\) of
\hyperlink{the-stone-weierstrass-theorem-for-functions-vanishing-at-infinity--oa-fnd-sw-13}{Section
11}.

\textbf{(a) Self-adjointness.} Let
\(\mathbb T=\{w\in\mathbb C:|w|=1\}\), a compact space
(\hyperlink{the-stone-weierstrass-theorem-for-functions-vanishing-at-infinity--oa-fnd-sw-15}{Lemma
4.1}(b)), and let \(A\) be the complex polynomials in the coordinate
\(w\), restricted to \(\mathbb T\). \(A\) is a subalgebra. It contains
\(1\), so it vanishes nowhere, and \(w\) separates points. But
\(\bar w\notin\bar A\). To see this, fix \(N\geq2\) and let
\(\omega=E(2\pi i/N)\), where \(E\) is the complex exponential. By the
addition formula \(E(z+w)=E(z)E(w)\), \(\omega^m=E(2\pi im/N)\). So
\(|\omega|=1\), because \(|E(it)|=1\) for real \(t\);
\(\omega^N=E(2\pi i)=1\); and \(\omega^m\neq1\) for \(0<m<N\), because
\(E(it)\neq1\) for \(0<t<2\pi\). These properties of the exponential
function are recalled in (B2) near the end of the lesson. For
\(g\in C(\mathbb T)\) put
\[\Lambda_N(g)=\frac1N\sum_{k=0}^{N-1}g(\omega^k)\,\omega^k .\] Then
\(|\Lambda_N(g)|\leq\|g\|\). For \(0\leq j\leq N-2\), the number
\(\rho=\omega^{j+1}\) satisfies \(\rho\neq1\) and \(\rho^N=1\), so
\(\Lambda_N(w^j)=\frac1N\sum_{k<N}\rho^k=\frac1N\,\frac{\rho^N-1}{\rho-1}=0\).
On \(\mathbb T\), \(\bar w=w^{-1}\), so
\(\Lambda_N(\bar w)=\frac1N\sum_k\omega^{-k}\omega^k=1\). Given
\(p\in A\) of degree \(d\), take \(N\geq d+2\). Then
\(\Lambda_N(\bar w-p)=1\), so \(\|\bar w-p\|\geq1\). In particular \(A\)
is not self-adjoint.
\hyperref[the-stone-weierstrass-theorem-for-functions-vanishing-at-infinity--exercises]{Exercise
1} transfers this example to \(\mathbb R\) with the Cayley map, so
self-adjointness is needed for noncompact \(X\) as well.

The number \(2\pi\Lambda_N(g)\) is a Riemann sum for the integral
\(\int_0^{2\pi}g(e^{it})e^{it}\,dt\), and the finite sums avoid
integration.
\hyperlink{the-stone-weierstrass-theorem-for-functions-vanishing-at-infinity--oa-fnd-sw-19}{Proposition
17.3}(c)--(d) gives the version with integrals.

\textbf{(b) Separation.} The even functions in \(C_0(\mathbb R)\) form a
closed self-adjoint subalgebra \(A_{\mathrm{ev}}\). It vanishes nowhere,
since \(\varphi\in A_{\mathrm{ev}}\), and it does not separate \(x\)
from \(-x\). For \(f\in A_{\mathrm{ev}}\),
\(1=\psi(1)-\psi(-1)=(\psi-f)(1)-(\psi-f)(-1)\leq2\|\psi-f\|\), so
\(\psi\) has distance at least \(\frac12\) from \(A_{\mathrm{ev}}\). On
\([-1,1]\), the even polynomials form an algebra that does not separate
\(x\) from \(-x\). By
\hyperlink{the-stone-weierstrass-theorem-for-functions-vanishing-at-infinity--oa-fnd-sw-12}{Theorem
13.2}, their closure is the set of even continuous functions, not
\(C([-1,1])\).

\textbf{(c) Vanishing nowhere.} \(A_0=\{f\in C_0(\mathbb R):f(0)=0\}\)
is a closed self-adjoint subalgebra. It separates points. Let
\(x\neq y\); by symmetry we may take \(x\neq0\). Put
\(r=\min(|x|,|x-y|)>0\) and \(\tau(t)=\max(0,1-|t-x|/r)\). Then
\(\tau\in C_c(\mathbb R)\), \(\tau(0)=0\), \(\tau(x)=1\) and
\(\tau(y)=0\). But \(\|\varphi-f\|\geq|\varphi(0)-f(0)|=1\) for every
\(f\in A_0\), so \(A_0\) is not dense. On \([0,1]\), the polynomials
without constant term separate points and vanish at \(0\), and their
closure is \(\{f:f(0)=0\}\) by
\hyperlink{the-stone-weierstrass-theorem-for-functions-vanishing-at-infinity--oa-fnd-sw-07}{Theorem
9.1}(b).

\textbf{(d) Closure under products.}
\(V=\operatorname{span}_{\mathbb C}\{\varphi,\psi\}\) is closed under
conjugation, since \(\varphi\) and \(\psi\) are real. It separates
points: if \(\varphi(x)=\varphi(y)\) and \(\psi(x)=\psi(y)\), then
\(x=\psi(x)/\varphi(x)=y\). It vanishes nowhere. But \(V\) has dimension
\(2\), so it is closed, while \(C_0(\mathbb R)\) has infinite dimension
(it contains infinitely many nonzero functions with disjoint supports).
So \(V\) is not dense. It is not an algebra:
\(\varphi^2=a\varphi+b\psi\) would give \(\varphi(x)=a+bx\) for all
\(x\), which fails at \(x=0\), \(\pm1\).

\textbf{(e) Lattices.} For the lattice versions of
\hyperlink{the-stone-weierstrass-theorem-for-functions-vanishing-at-infinity--oa-fnd-sw-06}{Theorem
8.2}, separation and nowhere vanishing are not enough, and the pointwise
lattice operations matter.

\begin{itemize}
\tightlist
\item
  (i) On \(K=[0,1]\), \(L=\{f\in C(K;\mathbb R):f(0)=0\}\) is a closed
  subalgebra, closed under \(\vee\) and \(\wedge\), that separates
  points (by (c)). It fails condition (ii) of
  \hyperlink{the-stone-weierstrass-theorem-for-functions-vanishing-at-infinity--oa-fnd-sw-06}{Theorem
  8.2}(1) at \(0\), and it is not dense. The same holds for any compact
  Hausdorff \(K\) and \(p\in K\), with Urysohn\textquotesingle s lemma
  (\hyperlink{the-stone-weierstrass-theorem-for-functions-vanishing-at-infinity--oa-fnd-sw-16}{Theorem
  5.1}) supplying separation.
\item
  (ii) Let \(K=\{a,b\}\) be discrete, so \(C(K;\mathbb R)=\mathbb R^2\).
  For \(\lambda>0\), \(\lambda\neq1\), let
  \(L_\lambda=\{(s,\lambda s):s\in\mathbb R\}\). Since \(\lambda>0\),
  \((s,\lambda s)\wedge(t,\lambda t)=(s\wedge t,\lambda(s\wedge t))\),
  so \(L_\lambda\) is a lattice subspace. It separates the two points
  and vanishes at neither, but \((1,1)\notin L_\lambda\), so it fails
  (i) and is not dense. By contrast, an \emph{algebra} with these two
  properties is all of \(\mathbb R^2\)
  (\hyperlink{the-stone-weierstrass-theorem-for-functions-vanishing-at-infinity--oa-fnd-sw-05}{Lemma
  7.1}).
\item
  (iii) Let \(K=\{a,b,c\}\) be discrete, \(0<\theta<1\), and
  \(L=\{f:f(c)=\theta f(a)+(1-\theta)f(b)\}\). \(L\) contains the
  constants and satisfies (i): at any two points, the value at the third
  point can be solved for. In its own order, \(L\) is a lattice,
  isomorphic to \(\mathbb R^2\) through \((f(a),f(b))\). But it is not
  closed under pointwise \(\wedge\): \(f=(1,0,\theta)\) and
  \(g=(0,1,1-\theta)\) lie in \(L\), while
  \(f\wedge g=(0,0,\min(\theta,1-\theta))\) does not. \(L\) is a proper
  subspace, hence closed and not dense.
\end{itemize}

\textbf{(f) Bounded functions instead of \(C_0\).} Let \(X=\mathbb N\),
so \(C_b(X)=\ell^\infty\). The convergent sequences \(c\) form a closed
self-adjoint subalgebra that contains \(1\) and separates points
(\(n\mapsto1/n\) is injective). For \(a\in c\) with limit \(l\), the
sequence \(e_n=(-1)^n\) satisfies \(\|e-a\|\geq\max(|1-l|,|1+l|)\geq1\).
So \(c\) is not dense in \(C_b(\mathbb N)\).
\hyperlink{the-stone-weierstrass-theorem-for-functions-vanishing-at-infinity--oa-fnd-sw-18}{Proposition
16.1}(3) proves that this failure occurs on every noncompact metrizable
space.

\textbf{(g) The topological hypothesis.} No counterexample exists: by
\hyperlink{the-stone-weierstrass-theorem-for-functions-vanishing-at-infinity--oa-fnd-sw-10}{Proposition
12.1}, the other hypotheses force \(X\) to be locally compact Hausdorff.

\textbf{(h) Density, not equality.} The conclusion is \(\bar A=C_0(X)\),
not \(A=C_0(X)\). For example \(C_c(\mathbb R)\) is dense in
\(C_0(\mathbb R)\)
(\hyperlink{the-stone-weierstrass-theorem-for-functions-vanishing-at-infinity--oa-fnd-sw-01}{Proposition
2.1}(3)) and \(\varphi\notin C_c(\mathbb R)\). The polynomials are dense
in \(C([0,1])\), and \(|x-\frac12|\) is not a polynomial.

\subsection[13. The closure of an arbitrary self-adjoint
subalgebra]{\texorpdfstring{\protect\hypertarget{the-stone-weierstrass-theorem-for-functions-vanishing-at-infinity--oa-fnd-sw-12}{}{}13.
The closure of an arbitrary self-adjoint
subalgebra}{13. The closure of an arbitrary self-adjoint subalgebra}}\label{the-stone-weierstrass-theorem-for-functions-vanishing-at-infinity--OA-FND-SW-12}

Dropping separation and nowhere vanishing does not end the story. The
closure can still be described exactly.

\textbf{Lemma 13.1} (Separation quotient). Let \(K\) be a compact space
and \(S\) a set of continuous functions on \(K\). Write \(x\sim_Sy\) if
\(s(x)=s(y)\) for every \(s\in S\). Let \(K_S\) be the set of classes,
\(q:K\to K_S\) the quotient map, and give \(K_S\) the quotient topology:
\(O\subseteq K_S\) is open when \(q^{-1}(O)\) is open.

\begin{enumerate}
\tightlist
\item
  A function \(f\) on \(K\) is continuous and constant on each class if
  and only if \(f=f'\circ q\) for a continuous \(f'\) on \(K_S\). Then
  \(\|f\|=\|f'\|\).
\item
  The functions \(s'\), \(s\in S\), separate the points of \(K_S\).
  \(K_S\) is compact and Hausdorff.
\end{enumerate}

\textbf{Proof.} (1) If \(f=f'\circ q\), then \(f\) is continuous and
constant on classes. Conversely, \(f'\) is well defined, and for an open
set \(O\) of scalars, \(q^{-1}(f'^{-1}(O))=f^{-1}(O)\) is open; so
\(f'^{-1}(O)\) is open. Since \(q\) is onto, the norms agree. (2) Two
distinct classes \(q(x)\neq q(y)\) differ at some \(s\in S\), which
means \(s'(q(x))\neq s'(q(y))\). \(K_S=q(K)\) is compact by (K2), and
Hausdorff by
\hyperlink{the-stone-weierstrass-theorem-for-functions-vanishing-at-infinity--oa-fnd-sw-10}{Proposition
12.1}(1), whose proof applies to any separating set of continuous
functions. \(\square\)

\textbf{Theorem 13.2.} Let \(X\) be a topological space, and let \(A\)
be a self-adjoint subalgebra of \(C_0(X)\) or a real subalgebra of
\(C_0(X;\mathbb R)\). Write \(x\sim_Ay\) if \(f(x)=f(y)\) for all
\(f\in A\). Then
\[\bar A=\{f\in C_0(X):\ f(x)=f(y)\ \text{whenever }x\sim_Ay,\ \text{and }f=0\ \text{on }Z(A)\},
\tag{13.1}\] with \(C_0(X;\mathbb R)\) in place of \(C_0(X)\) in the
real case. The same formula holds for a compact space \(K\), with
\(C(K)\) or \(C(K;\mathbb R)\) in place of \(C_0(X)\).

\textbf{Proof.} First the compact case. Let \(R\) be the right side. It
contains \(A\), and it is closed, since each of its conditions involves
the values at one or two points. So \(\bar A\subseteq R\). Let
\(f\in R\). By Lemma 13.1 with \(S=A\), \(f=f'\circ q\). The map
\(a\mapsto a'\) respects sums, scalars, products and conjugation, so
\(A'=\{a':a\in A\}\) is an algebra of the same kind on \(K_A\). It
separates points, and its common zero set is \(q(Z(A))\), where \(f'\)
vanishes. By the compact form of
\hyperlink{the-stone-weierstrass-theorem-for-functions-vanishing-at-infinity--oa-fnd-sw-09}{Theorem
10.1}, or by
\hyperlink{the-stone-weierstrass-theorem-for-functions-vanishing-at-infinity--oa-fnd-sw-07}{Theorem
9.1} in the real case, \(f'\) lies in the closure of \(A'\). Since
\(\|f-a\|=\|f'-a'\|\) for \(a\in A\), \(f\in\bar A\).

For \(C_0(X)\), apply the compact case to \(\tilde A\) on \(X_\infty\)
(\hyperlink{the-stone-weierstrass-theorem-for-functions-vanishing-at-infinity--oa-fnd-sw-03}{Proposition
3.4}). On \(X\), the relation \(\sim_{\tilde A}\) is \(\sim_A\). A point
\(x\in X\) is equivalent to \(\infty\) exactly when \(x\in Z(A)\), and
\(Z(\tilde A)=Z(A)\cup\{\infty\}\). So \(\tilde f\) satisfies the
conditions of (13.1) for \(\tilde A\) exactly when \(f\) satisfies them
for \(A\). Since \(f\mapsto\tilde f\) preserves the norm
(\hyperlink{the-stone-weierstrass-theorem-for-functions-vanishing-at-infinity--oa-fnd-sw-03}{Proposition
3.4}(2)), the closure transfers to \(X\). \(\square\)

\Needspace{5\baselineskip}
\textbf{Corollary 13.3.}

\begin{enumerate}
\tightlist
\item
  \hyperlink{the-stone-weierstrass-theorem-for-functions-vanishing-at-infinity--oa-fnd-sw-09}{Theorem
  10.1} is the case in which \(\sim_A\) is equality and
  \(Z(A)=\varnothing\).
\item
  (\emph{One common zero.}) If \(A\) separates points and
  \(Z(A)=\{z\}\), then \(\bar A=\{f\in C_0(X):f(z)=0\}\).
\item
  (\emph{Closed subalgebras.}) For any equivalence relation \(\sim\) on
  \(X\) and any \(Z\subseteq X\), the functions in \(C_0(X)\) that are
  constant on classes and vanish on \(Z\) form a closed self-adjoint
  subalgebra. By (13.1) every closed self-adjoint subalgebra \(B\) has
  this form, with \(\sim_B\) and \(Z(B)\).
\item
  (\emph{Dichotomy.}) Let \(K\) be a compact space, and let \(B\) be a
  closed self-adjoint subalgebra of \(C(K)\), or a closed real
  subalgebra of \(C(K;\mathbb R)\), that separates points. Then \(B\) is
  either the whole space or the set of all its functions that vanish at
  one point \(z\). Indeed, \(Z(B)\) has at most one point
  (\hyperlink{the-stone-weierstrass-theorem-for-functions-vanishing-at-infinity--oa-fnd-sw-05}{Lemma
  7.1}(1)). If it is empty, apply
  \hyperlink{the-stone-weierstrass-theorem-for-functions-vanishing-at-infinity--oa-fnd-sw-07}{Theorem
  9.1}(a) or the compact form of
  \hyperlink{the-stone-weierstrass-theorem-for-functions-vanishing-at-infinity--oa-fnd-sw-09}{Theorem
  10.1}; if \(Z(B)=\{z\}\), apply (2).
\end{enumerate}

Self-adjointness cannot be dropped from Theorem 13.2. The closure of the
polynomials in \(w\) on \(\mathbb T\)
(\hyperlink{the-stone-weierstrass-theorem-for-functions-vanishing-at-infinity--oa-fnd-sw-11}{Example
12.2}(a)) separates points and contains \(1\), yet it is not
\(C(\mathbb T)\).
\hyperlink{the-stone-weierstrass-theorem-for-functions-vanishing-at-infinity--oa-fnd-sw-19}{Proposition
17.3} describes it.

\subsection[14. The Weierstrass theorem: Bernstein polynomials and
Korovkin\textquotesingle s
theorem]{\texorpdfstring{\protect\hypertarget{the-stone-weierstrass-theorem-for-functions-vanishing-at-infinity--oa-fnd-sw-14}{}{}14.
The Weierstrass theorem: Bernstein polynomials and
Korovkin\textquotesingle s
theorem}{14. The Weierstrass theorem: Bernstein polynomials and Korovkin\textquotesingle s theorem}}\label{the-stone-weierstrass-theorem-for-functions-vanishing-at-infinity--OA-FND-SW-14}

The polynomial step of
\hyperlink{the-stone-weierstrass-theorem-for-functions-vanishing-at-infinity--oa-fnd-sw-04}{Section
6} does not need the Weierstrass approximation theorem, and nothing
above uses it. This section proves the classical theorem directly and
constructively, with a quantitative form.

\textbf{Theorem 14.1} (Korovkin\textquotesingle s theorem). Let
\(I=[a,b]\), and let \(L_n:C(I;\mathbb R)\to C(I;\mathbb R)\) be linear
maps that are positive: \(u\geq0\) implies \(L_nu\geq0\). Put
\(e_k(t)=t^k\). If \(L_ne_k\to e_k\) uniformly for \(k=0,1,2\), then
\(L_nf\to f\) uniformly for every \(f\in C(I;\mathbb R)\).

\emph{Reference:} {[}Korovkin 1953{]}.

\textbf{Proof.} Positivity makes each \(L_n\) monotone: \(u\leq v\)
implies \(L_nu\leq L_nv\). Fix \(f\) and \(\varepsilon>0\). By uniform
continuity
(\hyperlink{the-stone-weierstrass-theorem-for-functions-vanishing-at-infinity--oa-fnd-sw-15}{Lemma
4.1}(d)) there is \(\delta>0\) with \(|f(s)-f(t)|<\varepsilon\) when
\(|s-t|<\delta\). If \(|s-t|\geq\delta\), then
\(|f(s)-f(t)|\leq2\|f\|\leq2\|f\|(s-t)^2/\delta^2\). So for all
\(s,t\in I\),
\[|f(s)-f(t)|\leq\varepsilon+c\,(s-t)^2,\qquad c=2\|f\|/\delta^2 .
\tag{14.1}\] Fix \(t\) and read (14.1) as an inequality between
functions of \(s\): with \(q_t=e_2-2te_1+t^2e_0\),
\(-\varepsilon e_0-cq_t\leq f-f(t)e_0\leq\varepsilon e_0+cq_t\). Apply
\(L_n\) and evaluate at \(t\):
\[|(L_nf)(t)-f(t)(L_ne_0)(t)|\leq\varepsilon(L_ne_0)(t)+c\,(L_nq_t)(t).\]
Here \((L_nq_t)(t)=(L_ne_2)(t)-2t(L_ne_1)(t)+t^2(L_ne_0)(t)\), which
tends to \(t^2-2t^2+t^2=0\) uniformly in \(t\in I\), because \(|t|\) is
bounded on \(I\). Hence
\[\|L_nf-f\|\leq\|f\|\,\|L_ne_0-e_0\|+\varepsilon\|L_ne_0\|+c\sup_{t\in I}|(L_nq_t)(t)|,\]
and the limit superior of the right side is at most \(\varepsilon\).
\(\square\)

\textbf{Proposition 14.2} (Bernstein polynomials). For a bounded
\(f:[0,1]\to\mathbb R\) and \(n\geq1\) put
\[b_{n,k}(x)=\binom nk x^k(1-x)^{n-k},\qquad(B_nf)(x)=\sum_{k=0}^nf(k/n)\,b_{n,k}(x).\]

\begin{enumerate}
\tightlist
\item
  \(B_ne_0=e_0\), \(B_ne_1=e_1\) and \(B_ne_2=e_2+(e_1-e_2)/n\). Hence
  \[\sum_{k=0}^n\Bigl(\frac kn-x\Bigr)^2b_{n,k}(x)=\frac{x(1-x)}n .
  \tag{14.2}\]
\item
  \(B_nf\to f\) uniformly for every continuous \(f\).
\item
  For every bounded \(f\) and every \(x\in[0,1]\),
  \(|f(x)-(B_nf)(x)|\leq\frac32\,\omega(f,n^{-1/2})\), where
  \(\omega(f,\delta)=\sup\{|f(s)-f(t)|:|s-t|\leq\delta\}\).
\item
  If \(f\) is bounded and continuous at \(x_0\), then
  \((B_nf)(x_0)\to f(x_0)\). Boundedness cannot be dropped.
\item
  (\emph{Bernoulli trials}) For \(x\in[0,1]\) and \(\delta>0\),
  \(\sum_{k:\,|k/n-x|\geq\delta}b_{n,k}(x)\leq x(1-x)/(n\delta^2)\).
\end{enumerate}

\emph{Reference:} {[}Bernstein 1912{]}. Version 4 of {[}van Neerven{]}
(April 2022) states (4) in its Remark 2.4 for every function that is
continuous at \(x_0\); the example in the proof of (4) shows that
boundedness is needed, and version 7 (July 2025) adds it.

\textbf{Proof.} (1) The binomial theorem gives
\(\sum_kb_{n,k}(x)=(x+1-x)^n=1\). Since
\(\frac kn\binom nk=\binom{n-1}{k-1}\) for \(k\geq1\), we get
\(\sum_k\frac kn\,b_{n,k}(x)=x\sum_{j=0}^{n-1}\binom{n-1}jx^j(1-x)^{n-1-j}=x\).
For \(n\geq2\) and \(k\geq1\),
\(\bigl(\frac kn\bigr)^2\binom nk=\frac{n-1}n\binom{n-2}{k-2}+\frac1n\binom{n-1}{k-1}\)
(with \(\binom{n-2}{-1}=0\)), which gives
\(B_ne_2=\frac{n-1}nx^2+\frac xn\), as claimed; for \(n=1\),
\(B_1e_2(x)=x=x^2+x(1-x)\) directly. Expanding \((k/n-x)^2\) and using
the three identities gives (14.2).

(2) \(B_n\) is linear and positive, and \(\|B_ne_2-e_2\|\leq\frac1{4n}\)
by (1). Apply Korovkin\textquotesingle s theorem, Theorem 14.1.

(3) First, \(\omega(f,\lambda\delta)\leq(1+\lambda)\,\omega(f,\delta)\)
for \(\lambda>0\): if \(|s-t|\leq\lambda\delta\), split the segment
between \(s\) and \(t\) into \(m\) equal pieces, where \(m\) is the
least integer \(\geq\lambda\); each piece has length at most \(\delta\),
so
\(|f(s)-f(t)|\leq m\,\omega(f,\delta)\leq(1+\lambda)\,\omega(f,\delta)\).
With \(\lambda=|s-t|/\delta\) this gives
\(|f(s)-f(t)|\leq(1+|s-t|/\delta)\,\omega(f,\delta)\) for all \(s,t\).
Since the weights \(b_{n,k}(x)\) are nonnegative with sum \(1\),
\[|f(x)-(B_nf)(x)|\leq\sum_kb_{n,k}(x)\,|f(x)-f(k/n)|\leq\omega(f,\delta)\Bigl(1+\frac1\delta\sum_kb_{n,k}(x)\Bigl|x-\frac kn\Bigr|\Bigr).\]
By the Cauchy--Schwarz inequality for these weights and (14.2), the last
sum is at most \((x(1-x)/n)^{1/2}\leq\frac1{2\sqrt n}\). Take
\(\delta=n^{-1/2}\).

(4) Given \(\varepsilon>0\), take \(\delta>0\) with
\(|f(t)-f(x_0)|<\varepsilon\) for \(|t-x_0|<\delta\). Splitting the sum
in (3) according to whether \(|k/n-x_0|<\delta\), and using (5),
\(|f(x_0)-(B_nf)(x_0)|\leq\varepsilon+2\|f\|\,x_0(1-x_0)/(n\delta^2)\).
For the second claim, let \(f(1/m)=2^{m^2}\) for \(m\geq3\), and \(f=0\)
at all other points of \([0,1]\). Then \(f=0\) on \((\frac13,1]\), so
\(f\) is continuous at \(\frac12\). All terms of \((B_nf)(\frac12)\) are
nonnegative, and for \(n\geq3\) the term with \(k=1\) is
\(n\,2^{-n}f(1/n)=n\,2^{n^2-n}\). So \((B_nf)(\frac12)\to\infty\), while
\(f(\frac12)=0\).

(5) Where \(|k/n-x|\geq\delta\), \((k/n-x)^2/\delta^2\geq1\). Sum
against \(b_{n,k}(x)\) and use (14.2). \(\square\)

\emph{Probabilistic reading.} \(b_{n,k}(x)\) is the probability of
exactly \(k\) successes in \(n\) independent trials that each succeed
with probability \(x\). So \((B_nf)(x)\) is the expected value of \(f\)
at the observed frequency of success. Part (5) is the weak law of large
numbers for such trials: the frequency lies within \(\delta\) of \(x\)
with probability at least \(1-x(1-x)/(n\delta^2)\).

\textbf{Corollary 14.3} (Weierstrass approximation theorem). Every
continuous complex function on \([a,b]\) is a uniform limit of
polynomials, with real coefficients when the function is real.

\emph{Reference:} {[}Weierstrass 1885{]}.

\textbf{Proof.} The substitution \(t=a+(b-a)x\) carries polynomials to
polynomials. Apply Proposition 14.2(2) to the real and imaginary parts.
A second proof: on the compact space \([a,b]\), the polynomials with
complex coefficients form a self-adjoint algebra (for real \(t\), the
conjugate of \(\sum c_kt^k\) is \(\sum\bar c_kt^k\)) that contains \(1\)
and separates points; apply the compact form of
\hyperlink{the-stone-weierstrass-theorem-for-functions-vanishing-at-infinity--oa-fnd-sw-09}{Theorem
10.1}, or
\hyperlink{the-stone-weierstrass-theorem-for-functions-vanishing-at-infinity--oa-fnd-sw-07}{Theorem
9.1} for real functions. \(\square\)

\Needspace{9\baselineskip}
\subsection[15. Further density theorems on compact
spaces]{\texorpdfstring{\protect\hypertarget{the-stone-weierstrass-theorem-for-functions-vanishing-at-infinity--oa-fnd-sw-17}{}{}15.
Further density theorems on compact
spaces}{15. Further density theorems on compact spaces}}\label{the-stone-weierstrass-theorem-for-functions-vanishing-at-infinity--OA-FND-SW-17}

\textbf{Proposition 15.1.}

\begin{enumerate}
\tightlist
\item
  (\emph{Polynomials in several variables.}) Let
  \(X\subseteq\mathbb R^n\) be compact. The real polynomials in
  \(x_1,\dots,x_n\) are dense in \(C(X;\mathbb R)\), and the complex
  ones in \(C(X)\). Every continuous \(F:X\to\mathbb R^m\) is a uniform
  limit of maps whose \(m\) coordinates are such polynomials.
\item
  (\emph{Polynomials in \(z\) and \(\bar z\).}) Let
  \(X\subseteq\mathbb C\) be compact. The functions
  \(z\mapsto p(z,\bar z)\), with \(p\) a complex polynomial in two
  variables, are dense in \(C(X)\).
\item
  (\emph{Separability.}) If \(K\) is a compact metric space, \(C(K)\) is
  separable.
\item
  (\emph{Trigonometric polynomials.}) The trigonometric polynomials
  \(\sum_{|k|\leq n}c_kw^k\) form a dense subspace of \(C(\mathbb T)\).
  Equivalently, the functions
  \(\theta\mapsto\sum_{|k|\leq n}c_ke^{ik\theta}\) are dense, uniformly
  on \(\mathbb R\), in the continuous \(2\pi\)-periodic functions.
\item
  (\emph{Products.}) Let \(X_1,\dots,X_k\) be compact spaces, not
  necessarily Hausdorff. The linear combinations of the functions
  \(f_1(x_1)\cdots f_k(x_k)\), with \(f_j\in C(X_j;\mathbb R)\), are
  dense in \(C(X_1\times\cdots\times X_k;\mathbb R)\). Every complex
  continuous function on the product is a uniform limit of finite sums
  of functions \(\lambda f_1(x_1)\cdots f_k(x_k)\) with
  \(\lambda\in\{1,i\}\) and real \(f_j\).
\item
  (\emph{Iterated integrals.}) Let \(X,Y\) be compact spaces, and
  \(\mu:C(X;\mathbb R)\to\mathbb R\), \(\nu:C(Y;\mathbb R)\to\mathbb R\)
  positive linear maps. For \(h\in C(X\times Y;\mathbb R)\), the
  functions \(x\mapsto\nu(h(x,\cdot))\) and \(y\mapsto\mu(h(\cdot,y))\)
  are continuous, and
  \[\mu\bigl(x\mapsto\nu(h(x,\cdot))\bigr)=\nu\bigl(y\mapsto\mu(h(\cdot,y))\bigr).
  \tag{15.1}\] In particular, for continuous \(h\) on
  \([a,b]\times[c,d]\), the two iterated Riemann integrals of \(h\)
  agree, and the Riemann sums of \(h\) over grid partitions converge to
  their common value.
\item
  (\emph{Moments.}) Let \(a<b\). If \(f\in C([a,b])\) and
  \(\int_a^bf(x)x^n\,dx=0\) for all \(n\geq0\), then \(f=0\). The
  condition \(a<b\) is needed: on \([a,a]\) every integral is \(0\), so
  the constant function \(1\) has all its moments equal to \(0\).
\end{enumerate}

\textbf{Proof.} (1) The coordinate functions and \(1\) generate an
algebra, which separates points because the coordinates do; apply
\hyperlink{the-stone-weierstrass-theorem-for-functions-vanishing-at-infinity--oa-fnd-sw-07}{Theorem
9.1}. The complex polynomials form a self-adjoint algebra, since the
coordinates are real; apply the compact form of
\hyperlink{the-stone-weierstrass-theorem-for-functions-vanishing-at-infinity--oa-fnd-sw-09}{Theorem
10.1}. For \(F\), approximate each coordinate.

(2) The algebra contains \(1\) and \(z\), which separates points. It is
self-adjoint: the conjugate of \(p(z,\bar z)\) is \(\bar p(\bar z,z)\),
where \(\bar p\) has the conjugate coefficients. Apply the compact form
of
\hyperlink{the-stone-weierstrass-theorem-for-functions-vanishing-at-infinity--oa-fnd-sw-09}{Theorem
10.1}.

(3) Choose a countable dense set \(C\subseteq K\)
(\hyperlink{the-stone-weierstrass-theorem-for-functions-vanishing-at-infinity--oa-fnd-sw-15}{Lemma
4.1}(d)), and put \(d_c(x)=d(x,c)\). The functions \(d_c\) and \(1\)
generate a real algebra \(A_0\) that contains the constants. It
separates points: if \(x\neq y\) and \(r=d(x,y)\), choose \(c\in C\)
with \(d(x,c)<r/2\); then \(d_c(x)<r/2<d_c(y)\). By
\hyperlink{the-stone-weierstrass-theorem-for-functions-vanishing-at-infinity--oa-fnd-sw-07}{Theorem
9.1}, \(A_0\) is dense in \(C(K;\mathbb R)\). The finite products of the
functions \(d_c\) and \(1\) form a countable set, and its linear
combinations with rational coefficients are dense in \(A_0\), hence in
\(C(K;\mathbb R)\). Combinations with Gaussian rational coefficients
give a countable set that is dense in \(C(K)\).

(4) On \(\mathbb T\), \(\bar w=w^{-1}\), so the trigonometric
polynomials form the algebra generated by \(w\) and \(\bar w\). It is
self-adjoint, contains \(1\), and separates points through \(w\). Apply
the compact form of
\hyperlink{the-stone-weierstrass-theorem-for-functions-vanishing-at-infinity--oa-fnd-sw-09}{Theorem
10.1} on the compact space \(\mathbb T\). For the periodic form: the map
\(\theta\mapsto E(i\theta)\) from \([0,2\pi]\) onto \(\mathbb T\) (it is
onto by (B2)) is continuous from a compact space onto a Hausdorff space,
so it maps closed sets to closed sets. A continuous \(2\pi\)-periodic
function takes equal values on each fibre, so it equals
\(g(E(i\theta))\) for a function \(g\) on \(\mathbb T\), and \(g\) is
continuous by the closed-map argument used in (5) below. Conversely each
\(g\in C(\mathbb T)\) gives a continuous periodic function, and
\(E(ik\theta)=E(i\theta)^k\).

(5) Apply
\hyperlink{the-stone-weierstrass-theorem-for-functions-vanishing-at-infinity--oa-fnd-sw-12}{Lemma
13.1} with \(S=C(X_j;\mathbb R)\) to get compact Hausdorff quotients
\(q_j:X_j\to X_j'\) on which continuous functions separate points. Put
\(X=\prod X_j\), \(X'=\prod X_j'\) and \(q=q_1\times\cdots\times q_k\).
Let \(h\in C(X;\mathbb R)\). If \(q(x)=q(y)\), then \(h(x)=h(y)\):
change one coordinate at a time, and note that with the other
coordinates fixed, \(h\) is a continuous function of the remaining one,
so it takes equal values at equivalent points. The map \(q\) is
continuous from the compact space \(X\)
(\hyperlink{the-stone-weierstrass-theorem-for-functions-vanishing-at-infinity--oa-fnd-sw-15}{Lemma
4.1}(b)) onto the Hausdorff space \(X'\). So it maps closed sets to
closed sets: a closed set is compact by (K1), its image is compact by
(K2), and compact subsets of \(X'\) are closed by (K3). So
\(h=h'\circ q\), and \(h'\) is continuous, because
\(h'^{-1}(F)=q(h^{-1}(F))\) for closed \(F\subseteq\mathbb R\). On
\(X'\), the linear combinations of products
\(f_1'(x_1)\cdots f_k'(x_k)\) form an algebra with \(1\) in it, and it
separates points, since points that differ in coordinate \(j\) are
separated by a function of that coordinate. By
\hyperlink{the-stone-weierstrass-theorem-for-functions-vanishing-at-infinity--oa-fnd-sw-07}{Theorem
9.1} this algebra is dense in \(C(X';\mathbb R)\). Composing with \(q\)
keeps norms and turns each \(f_j'\) into
\(f_j'\circ q_j\in C(X_j;\mathbb R)\). The complex statement follows by
approximating real and imaginary parts.

(6) Since \(-\|u\|1\leq u\leq\|u\|1\), positivity gives
\(|\mu(u)|\leq\mu(1)\|u\|\), and likewise for \(\nu\). Fix \(x_0\) and
\(\varepsilon>0\). For each \(y\), continuity of \(h\) at \((x_0,y)\)
gives open sets \(U_y\ni x_0\) and \(V_y\ni y\) with
\(|h-h(x_0,y)|<\varepsilon/2\) on \(U_y\times V_y\). Finitely many
\(V_{y_j}\) cover \(Y\). On \(U=\bigcap_jU_{y_j}\) we get
\(\|h(x,\cdot)-h(x_0,\cdot)\|<\varepsilon\), since both \((x,y)\) and
\((x_0,y)\) lie in some \(U_{y_j}\times V_{y_j}\). So
\(x\mapsto\nu(h(x,\cdot))\) is continuous, and so is the other function.
The two sides of (15.1) are linear in \(h\) and bounded by
\(\mu(1)\nu(1)\|h\|\). For \(h(x,y)=f(x)g(y)\), both equal
\(\mu(f)\nu(g)\). By (5) such functions span a dense subspace, so (15.1)
holds for all \(h\).

For the rectangle, \(\mu(u)=\int_a^bu\) and \(\nu(v)=\int_c^dv\) are
positive linear maps, by (B1). Take a grid of cells
\([x_{i-1},x_i]\times[y_{j-1},y_j]\) with points \((s_i,t_j)\) in the
cells. Additivity over subintervals (B1) writes the iterated integral as
the sum over cells of
\(\int_{x_{i-1}}^{x_i}\int_{y_{j-1}}^{y_j}h\,dy\,dx\). By monotonicity
of the integral, each term differs from
\(h(s_i,t_j)(x_i-x_{i-1})(y_j-y_{j-1})\) by at most the area of the cell
times the oscillation of \(h\) on it. \(h\) is uniformly continuous on
the compact rectangle
(\hyperlink{the-stone-weierstrass-theorem-for-functions-vanishing-at-infinity--oa-fnd-sw-15}{Lemma
4.1}(d)), so the oscillations tend to \(0\) with the mesh. So the grid
Riemann sums converge to the iterated integral.

(7) By the Weierstrass theorem
(\hyperlink{the-stone-weierstrass-theorem-for-functions-vanishing-at-infinity--oa-fnd-sw-14}{Corollary
14.3}) choose polynomials \(p_k\) with \(p_k\to\bar f\) uniformly on
\([a,b]\). By hypothesis and linearity \(\int_a^bfp_k=0\). Since
\(|\int_a^bu|\leq(b-a)\sup|u|\) for continuous \(u\) (B1),
\(\bigl|\int_a^b|f|^2-\int_a^bfp_k\bigr|\leq(b-a)\|f\|\,\|\bar f-p_k\|\to0\),
so \(\int_a^b|f|^2=0\). If \(f(x_0)\neq0\), then by continuity
\(|f|^2>|f(x_0)|^2/2\) on \([a,b]\cap(x_0-\delta,x_0+\delta)\) for some
\(\delta>0\), a subinterval of positive length because \(a<b\), and the
integral would be positive, by monotonicity and additivity of the
integral (B1). So \(f=0\). \(\square\)

Part (5) needs neither the Hausdorff property nor
Urysohn\textquotesingle s lemma. The continuous functions on a compact
space need not separate its points, and then the products
\(f_1(x_1)\cdots f_k(x_k)\) do not separate the points of the product;
the quotients \(X_j'\) remove this problem. In (6), positive linear
functionals are the form in which Radon measures act on continuous
functions, so (15.1) also gives the equality of iterated integrals of a
continuous function for Radon measures on compact Hausdorff spaces.

\Needspace{9\baselineskip}
\subsection[16. Compact sets, extension of functions, and bounded
functions]{\texorpdfstring{\protect\hypertarget{the-stone-weierstrass-theorem-for-functions-vanishing-at-infinity--oa-fnd-sw-18}{}{}16.
Compact sets, extension of functions, and bounded
functions}{16. Compact sets, extension of functions, and bounded functions}}\label{the-stone-weierstrass-theorem-for-functions-vanishing-at-infinity--OA-FND-SW-18}

\textbf{Proposition 16.1.}

\begin{enumerate}
\tightlist
\item
  (\emph{Uniform approximation on compact sets.}) Let \(X\) be any
  topological space and \(A\subseteq C(X;\mathbb R)\) a real algebra
  that separates points and vanishes nowhere. For every
  \(f\in C(X;\mathbb R)\), every compact \(C\subseteq X\) and every
  \(\varepsilon>0\), some \(a\in A\) satisfies \(|a-f|<\varepsilon\) on
  \(C\). Suppose that \(X\) has compact sets
  \(L_1\subseteq L_2\subseteq\cdots\) such that every compact set lies
  in some \(L_n\); by
  \hyperlink{the-stone-weierstrass-theorem-for-functions-vanishing-at-infinity--oa-fnd-sw-15}{Proposition
  4.2}(3) this holds when \(X\) is LCH and \(\sigma\)-compact. Then some
  sequence in \(A\) converges to \(f\) uniformly on every compact set.
  The complex analogue holds for self-adjoint \(A\).
\item
  (\emph{Tietze extension theorem, compact case.}) Let \(K\) be compact
  Hausdorff, \(Z\subseteq K\) closed, and \(g\in C(Z;\mathbb R)\). Then
  \(g=G|_Z\) for some \(G\in C(K;\mathbb R)\) with \(\|G\|=\|g\|\).
\item
  (\emph{A converse.}) Let \(X\) be a metrizable space that is not
  compact. Then \(C_b(X;\mathbb R)\) has a closed subalgebra with \(1\)
  in it that separates the points of \(X\) but is not dense.
\end{enumerate}

\textbf{Proof.} (1) The restrictions of the elements of \(A\) to \(C\)
form a real algebra on the compact space \(C\). It separates the points
of \(C\) and vanishes nowhere on \(C\), so it is dense in
\(C(C;\mathbb R)\) by
\hyperlink{the-stone-weierstrass-theorem-for-functions-vanishing-at-infinity--oa-fnd-sw-07}{Theorem
9.1}(a); and \(f|_C\in C(C;\mathbb R)\). For the sequence, choose
\(a_n\in A\) with \(|a_n-f|<1/n\) on \(L_n\). A compact set lies in some
\(L_m\), and the \(L_n\) increase, so \(|a_n-f|<1/n\) on it for
\(n\geq m\). In the complex case use the compact form of
\hyperlink{the-stone-weierstrass-theorem-for-functions-vanishing-at-infinity--oa-fnd-sw-09}{Theorem
10.1}.

(2) If \(Z=\varnothing\), take \(G=0\). Otherwise let
\(R=\{G|_Z:G\in C(K;\mathbb R)\}\). It is a real algebra on the compact
space \(Z\) that contains the constants. It separates the points of
\(Z\): points of \(K\) are closed, so Urysohn\textquotesingle s lemma
(\hyperlink{the-stone-weierstrass-theorem-for-functions-vanishing-at-infinity--oa-fnd-sw-16}{Theorem
5.1}) separates any two of them by a continuous function. By
\hyperlink{the-stone-weierstrass-theorem-for-functions-vanishing-at-infinity--oa-fnd-sw-07}{Theorem
9.1}(a), \(R\) is dense in \(C(Z;\mathbb R)\). It is also closed. Let
\(g\in\bar R\), and choose \(G_n\in C(K;\mathbb R)\) with
\(\|G_n|_Z-g\|\leq2^{-n}\). The differences \(D_n=G_{n+1}-G_n\) satisfy
\(\|D_n|_Z\|<2^{1-n}\). The truncations
\(E_n=(D_n\wedge2^{1-n})\vee(-2^{1-n})\) are continuous on \(K\), agree
with \(D_n\) on \(Z\), and satisfy \(\|E_n\|\leq2^{1-n}\). So the series
\(G=G_1+\sum_nE_n\) converges uniformly on \(K\), \(G\) is continuous
(\hyperlink{the-stone-weierstrass-theorem-for-functions-vanishing-at-infinity--oa-fnd-sw-01}{Proposition
2.1}(1)), and \(G|_Z=\lim_nG_n|_Z=g\). Hence \(R=C(Z;\mathbb R)\).
Finally, \((G\wedge\|g\|)\vee(-\|g\|)\) still extends \(g\) and has norm
\(\|g\|\).

(3) Let \(d\) be a metric for \(X\). Since \(X\) is not compact,
\hyperlink{the-stone-weierstrass-theorem-for-functions-vanishing-at-infinity--oa-fnd-sw-15}{Lemma
4.1}(d) gives a sequence \((x_n)\) with no convergent subsequence. A
value repeated infinitely often would give a constant convergent
subsequence, so after passing to a subsequence the \(x_n\) are distinct.
Each point \(p\) has a ball that contains \(x_n\) for only finitely many
\(n\), and after shrinking it contains no \(x_n\) other than \(p\)
itself. So the sets \(E_0=\{x_n:n\text{ even}\}\) and
\(E_1=\{x_n:n\text{ odd}\}\) are closed and disjoint. The function
\(u=d(\cdot,E_0)/(d(\cdot,E_0)+d(\cdot,E_1))\) is continuous with values
in \([0,1]\), and it is \(0\) on \(E_0\) and \(1\) on \(E_1\). Let \(A\)
be the set of \(f\in C_b(X;\mathbb R)\) for which \(f(x_n)\) has a limit
along the even \(n\) and a limit along the odd \(n\), and the two limits
are equal. \(A\) is a subalgebra containing the constants. It is closed:
if \(f_k\in A\) have common limits \(l_k\) and \(f_k\to f\), then
\(|l_k-l_j|\leq\|f_k-f_j\|\), so \(l_k\to l\), and \(f(x_n)\to l\) along
both parities. \(A\) separates points: for \(x\neq y\) choose
\(0<r<d(x,y)\) such that the ball \(B(x,r)\) contains no \(x_n\) other
than \(x\), and put \(f(z)=\max(0,1-d(z,x)/r)\). Then \(f(x)=1\),
\(f(y)=0\), and \(f(x_n)=0\) whenever \(x_n\neq x\), so \(f\in A\). But
if \(f\in A\) has common limit \(l\), then
\(\|u-f\|\geq\max(|l|,|1-l|)\geq\frac12\). So \(A\) is not dense.
\(\square\)

In (2), Urysohn\textquotesingle s lemma is used only to separate points;
the closedness of \(R\) comes from the truncation argument.
\hyperlink{the-stone-weierstrass-theorem-for-functions-vanishing-at-infinity--oa-fnd-sw-11}{Example
12.2}(f) is the case \(X=\mathbb N\) of (3). The conclusion of (3) holds
for every completely regular Hausdorff space that is not compact; that
version is proved with the Stone--Čech compactification {[}Blackadar{]},
which this lesson does not develop.

\subsection[17. Convex functions, singular functions, and the disc
algebra]{\texorpdfstring{\protect\hypertarget{the-stone-weierstrass-theorem-for-functions-vanishing-at-infinity--oa-fnd-sw-19}{}{}17.
Convex functions, singular functions, and the disc
algebra}{17. Convex functions, singular functions, and the disc algebra}}\label{the-stone-weierstrass-theorem-for-functions-vanishing-at-infinity--OA-FND-SW-19}

\textbf{Proposition 17.1} (Differences of convex functions). Let \(K\)
be a compact space and \(P\) a nonempty set of nonnegative continuous
real functions on \(K\), closed under sums, under multiplication by
positive scalars, and under \(\wedge\) or under \(\vee\). Then
\(L=P-P=\{f-g:f,g\in P\}\) is a linear subspace closed under \(\wedge\)
and \(\vee\). If \(P\) separates points and some nonzero constant
function lies in \(P\), then \(L\) is dense in \(C(K;\mathbb R)\). In
particular, every continuous function on \([a,b]\) is a uniform limit of
differences of continuous convex functions.

\textbf{Proof.} \(L\) is closed under sums and positive multiples,
\(L=-L\), and \(0=f-f\in L\); so \(L\) is a linear subspace. For
\(f_1-g_1\) and \(f_2-g_2\) in \(L\),
\[(f_1-g_1)\wedge(f_2-g_2)=\bigl((f_1+g_2)\wedge(f_2+g_1)\bigr)-(g_1+g_2),\]
because subtracting the same function from two functions commutes with
\(\wedge\). The same identity holds with \(\vee\). So if \(P\) is closed
under \(\wedge\), so is \(L\), and then \(L\) is closed under \(\vee\)
because \(f\vee g=-((-f)\wedge(-g))\) and \(L=-L\). The case of \(\vee\)
is symmetric. If \(P\) contains a constant \(c>0\) and separates points,
then \(L\) contains all constants and separates points, so it is dense
by
\hyperlink{the-stone-weierstrass-theorem-for-functions-vanishing-at-infinity--oa-fnd-sw-06}{Theorem
8.2}(2). On \([a,b]\), take for \(P\) the nonnegative continuous convex
functions. Sums, positive multiples and maxima of convex functions are
convex, \(1\in P\), and \(x\mapsto x-a+1\) lies in \(P\) and separates
points. The nonnegative concave functions, with \(\wedge\), work in the
same way. \(\square\)

A set \(N\subseteq\mathbb R\) is \emph{null} if, for every
\(\varepsilon>0\), countably many intervals of total length less than
\(\varepsilon\) cover it. A union of two null sets is null, and so is
the image of a null set under a map \(t\mapsto\alpha+\beta t\).

\textbf{Proposition 17.2} (Singular functions). Let
\(C_{\mathrm{sing}}([a,b])\) be the set of continuous
\(f:[a,b]\to\mathbb R\) that are differentiable with derivative \(0\)
outside a null set. Then \(C_{\mathrm{sing}}([a,b])\) is a subalgebra of
\(C([a,b];\mathbb R)\); it contains the constants, and it separates
points. So it is dense in \(C([a,b];\mathbb R)\).

\textbf{Proof.} \emph{Subalgebra.} If \(f'=0\) off a null set \(N_f\)
and \(g'=0\) off a null set \(N_g\), then off the null set
\(N_f\cup N_g\) the functions \(f+g\), \(\lambda f\) and \(fg\) are
differentiable with derivative \(0\). For \(fg\), write
\(f(s)g(s)-f(t)g(t)=f(s)(g(s)-g(t))+g(t)(f(s)-f(t))\) and use continuity
of \(f\).

\emph{The Cantor function.} Let \(\Phi\) be the set of continuous
nondecreasing \(\phi:[0,1]\to[0,1]\) with \(\phi(0)=0\) and
\(\phi(1)=1\). For \(\phi\in\Phi\) define \(T\phi\) to be \(\phi(3x)/2\)
on \([0,\frac13]\), \(\frac12\) on \([\frac13,\frac23]\), and
\(\frac12+\phi(3x-2)/2\) on \([\frac23,1]\). Then \(T\phi\in\Phi\), and
\(\|T\phi-T\chi\|\leq\frac12\|\phi-\chi\|\). Starting from \(c_0(x)=x\),
the functions \(c_{n+1}=Tc_n\) satisfy \(\|c_{n+1}-c_n\|\leq2^{-n}\). So
they converge uniformly to some \(c\in\Phi\) with \(Tc=c\). Let
\(G_1=(\frac13,\frac23)\) and
\(G_{n+1}=G_1\cup\frac13G_n\cup(\frac23+\frac13G_n)\). By induction on
\(n\): \(G_n\subseteq G_{n+1}\); \([0,1]\setminus G_n\) is a union of
\(2^n\) closed intervals of length \(3^{-n}\); and \(c=Tc\) is constant
on each component interval of \(G_n\). So \(c\) is locally constant on
the open set \(G=\bigcup_nG_n\), and \([0,1]\setminus G\) is null, since
for each \(n\) it is covered by intervals of total length \((2/3)^n\).

\emph{Separation.} Let \(a\leq x<y\leq b\). Put \(s(t)=0\) for
\(t\leq x\), \(s(t)=c\bigl((t-x)/(y-x)\bigr)\) for \(x\leq t\leq y\),
and \(s(t)=1\) for \(t\geq y\). Then \(s\) is continuous, \(s(x)=0\) and
\(s(y)=1\). It is locally constant, so differentiable with derivative
\(0\), off the null set
\(\{x,y\}\cup\bigl(x+(y-x)([0,1]\setminus G)\bigr)\). So
\(s\in C_{\mathrm{sing}}([a,b])\), and
\hyperlink{the-stone-weierstrass-theorem-for-functions-vanishing-at-infinity--oa-fnd-sw-07}{Theorem
9.1}(a) gives density. \(\square\)

For the next result, let \(\mathcal P\) be the algebra of polynomials in
\(w\), restricted to \(\mathbb T\). For \(g\in C(\mathbb T)\) and
\(n\in\mathbb Z\) put
\(\hat g(n)=\frac1{2\pi}\int_0^{2\pi}g(e^{i\theta})e^{-in\theta}\,d\theta\),
and let \(\bar D=\{z:|z|\leq1\}\).

\textbf{Proposition 17.3} (The disc algebra).

\begin{itemize}
\tightlist
\item
  (a) If \(p\in\mathcal P\) and \(\bar p\in\mathcal P\), then \(p\) is
  constant. So \(\mathcal P\) is not self-adjoint.
\item
  (b) As functions of \(\theta\), the elements of \(\mathcal P\) are the
  sums \(\sum_{k=0}^nc_ke^{ik\theta}\).
\item
  (c)--(d) \(\hat g(-1)=0\) for every \(g\) in the closure
  \(\bar{\mathcal P}\), while \(\bar w\) has \(\hat{\bar w}(-1)=1\). So
  \(\bar w\notin\bar{\mathcal P}\).
\item
  (e)
  \(\bar{\mathcal P}=\{G|_{\mathbb T}:\ G\in C(\bar D),\ G(z)=\sum_{n\geq0}a_nz^n\ \text{for }|z|<1\}\).
\item
  (f)
  \(\bar{\mathcal P}=\{g\in C(\mathbb T):\ \hat g(n)=0\ \text{for all }n<0\}\).
\item
  (g) Trigonometric polynomials approximate every \(g\in C(\mathbb T)\)
  uniformly
  (\hyperlink{the-stone-weierstrass-theorem-for-functions-vanishing-at-infinity--oa-fnd-sw-17}{Proposition
  15.1}(4)).
\end{itemize}

\textbf{Proof.} For integers \(m\neq0\),
\(\int_0^{2\pi}e^{im\theta}\,d\theta=0\), by the fundamental theorem of
calculus (B1) applied to \(e^{im\theta}/(im)\), whose derivative is
\(e^{im\theta}\) and which takes the same value at \(0\) and \(2\pi\)
(B2). So if \(t=\sum_kc_ke^{ik\theta}\) is a finite sum, then
\(\hat t(n)=c_n\). Also \(|\hat g(n)|\leq\|g\|\), since
\(|\int u|\leq\int|u|\) (B1).

(a) Write \(p=\sum_{k\geq0}c_kw^k\). On \(\mathbb T\),
\(\bar p=\sum_k\bar c_kw^{-k}\), so \(\hat{\bar p}(-k)=\bar c_k\). If
\(\bar p\in\mathcal P\), its coefficients at negative \(n\) vanish, so
\(c_k=0\) for \(k\geq1\).

(b) \(e^{ik\theta}=(e^{i\theta})^k\) by the addition formula.

(c)--(d) \(g\mapsto\hat g(-1)\) is continuous and vanishes on
\(\mathcal P\), hence on \(\bar{\mathcal P}\), while
\(\bar w=e^{-i\theta}\) gives the value \(1\). (Compare
\hyperlink{the-stone-weierstrass-theorem-for-functions-vanishing-at-infinity--oa-fnd-sw-11}{Example
12.2}(a), which gives \(\|\bar w-p\|\geq1\) without integrals.)

(f) If \(g\in\bar{\mathcal P}\), then \(\hat g(n)=0\) for \(n<0\), by
continuity as in (c). Conversely, let \(\hat g(n)=0\) for \(n<0\). For
\(0\leq r<1\) put
\[P_r(t)=\sum_{n\in\mathbb Z}r^{|n|}e^{int}=\frac{1-r^2}{|1-re^{it}|^2}>0 .\]
The series converges absolutely and uniformly in \(t\), and summing its
two geometric halves gives the closed form: with \(z=re^{it}\),
\(\frac1{1-z}+\frac{\bar z}{1-\bar z}=\frac{1-|z|^2}{|1-z|^2}\).
Integrating term by term, which uniform convergence allows (B1), gives
\(\frac1{2\pi}\int_0^{2\pi}P_r(\theta-t)\,dt=1\) and
\[(P_rg)(e^{i\theta}):=\frac1{2\pi}\int_0^{2\pi}P_r(\theta-t)\,g(e^{it})\,dt=\sum_{n\in\mathbb Z}r^{|n|}\hat g(n)\,e^{in\theta}.\]
Here \(\hat g(n)=0\) for \(n<0\), so
\(P_rg=\sum_{n\geq0}r^n\hat g(n)w^n\), a series of elements of
\(\mathcal P\) that converges uniformly on \(\mathbb T\) because
\(|r^n\hat g(n)|\leq r^n\|g\|\). So \(P_rg\in\bar{\mathcal P}\). It
remains to show \(\|P_rg-g\|\to0\) as \(r\to1\). Given
\(\varepsilon>0\), uniform continuity on the compact set \(\mathbb T\)
(\hyperlink{the-stone-weierstrass-theorem-for-functions-vanishing-at-infinity--oa-fnd-sw-15}{Lemma
4.1}(d)) gives \(\eta>0\) with \(|g(u)-g(v)|<\varepsilon\) when
\(u,v\in\mathbb T\) and \(|u-v|<\eta\). Since \(P_r\) has average \(1\),
\[(P_rg-g)(e^{i\theta})=\frac1{2\pi}\int_0^{2\pi}P_r(\theta-t)\bigl(g(e^{it})-g(e^{i\theta})\bigr)\,dt .\]
Let \(r\geq1-\eta/2\). If \(|e^{it}-e^{i\theta}|\geq\eta\), then
\(|1-re^{i(\theta-t)}|=|e^{it}-re^{i\theta}|\geq\eta-(1-r)\geq\eta/2\),
so \(P_r(\theta-t)\leq4(1-r^2)/\eta^2\). Hence, for every \(t\), the
integrand has absolute value at most
\(\varepsilon P_r(\theta-t)+8\|g\|(1-r^2)/\eta^2\). Integrating this
bound, by monotonicity of the integral (B1), gives
\(\|P_rg-g\|\leq\varepsilon+8\|g\|(1-r^2)/\eta^2\), which is less than
\(2\varepsilon\) for \(r\) close to \(1\).

(e) Let \(g\in\bar{\mathcal P}\). By (f), \(\hat g(n)=0\) for \(n<0\),
and \(|\hat g(n)|\leq\|g\|\). Put \(G(z)=\sum_{n\geq0}\hat g(n)z^n\) for
\(|z|<1\), and \(G=g\) on \(\mathbb T\). For \(z=re^{i\theta}\) with
\(r<1\), \(G(z)=(P_rg)(e^{i\theta})\). The series converges uniformly on
each disc \(|z|\leq\rho<1\), so \(G\) is continuous on the open disc.
Let \(w_0\in\mathbb T\) and \(z\to w_0\) in \(\bar D\). For \(|z|<1\) we
have \(|G(z)-g(w_0)|\leq\|P_{|z|}g-g\|+|g(z/|z|)-g(w_0)|\), and for
\(|z|=1\) we have \(|G(z)-g(w_0)|=|g(z)-g(w_0)|\). Both tend to \(0\),
by (f) and the continuity of \(g\). So \(G\in C(\bar D)\). Conversely,
let \(G\in C(\bar D)\) with \(G(z)=\sum a_nz^n\) for \(|z|<1\). For
\(r<\rho<1\), the terms \(a_n\rho^n\) tend to \(0\), so
\(|a_n|r^n\leq C(r/\rho)^n\) for some \(C\). So
\(G_r(w)=G(rw)=\sum a_nr^nw^n\) converges uniformly on \(\mathbb T\),
and \(G_r|_{\mathbb T}\in\bar{\mathcal P}\). \(G\) is uniformly
continuous on the compact set \(\bar D\)
(\hyperlink{the-stone-weierstrass-theorem-for-functions-vanishing-at-infinity--oa-fnd-sw-15}{Lemma
4.1}(d)), and \(|rw-w|=1-r\), so \(G_r\to G\) uniformly on
\(\mathbb T\). Hence \(G|_{\mathbb T}\in\bar{\mathcal P}\). \(\square\)

The proof of (e) uses the dilations \(G_r\), not the partial sums of the
Taylor series of \(G\): those partial sums need not converge uniformly
on the circle. By
\href{https://kokunoyumeto.github.io/open-math-courses/courses/harmonic-analysis-on-locally-compact-groups/reader/programme-reading.html\#ref-7265a5ab9053615e}{Cauchy\textquotesingle s
theorem for cycles and its consequences, Lemma 3.1 and Theorem 3.2}, the
functions given by power series on the open disc are exactly the
holomorphic ones. So (e) says that \(\bar{\mathcal P}\) consists of the
restrictions to \(\mathbb T\) of the continuous functions on \(\bar D\)
that are holomorphic on the open disc.

\subsection[18. Integer coefficients and simultaneous approximation of
derivatives]{\texorpdfstring{\protect\hypertarget{the-stone-weierstrass-theorem-for-functions-vanishing-at-infinity--oa-fnd-sw-20}{}{}18.
Integer coefficients and simultaneous approximation of
derivatives}{18. Integer coefficients and simultaneous approximation of derivatives}}\label{the-stone-weierstrass-theorem-for-functions-vanishing-at-infinity--OA-FND-SW-20}

\textbf{Proposition 18.1} (Integer coefficients on \([0,1]\)). A
function \(f\in C([0,1];\mathbb R)\) can be approximated uniformly on
\([0,1]\) by polynomials with integer coefficients if and only if
\(f(0)\) and \(f(1)\) are integers.

\textbf{Proof.} Let \(\mathcal G\) be the closure of \(\mathbb Z[x]\) in
\(C([0,1];\mathbb R)\). It is closed under addition, negation and
multiplication, by the argument of
\hyperlink{the-stone-weierstrass-theorem-for-functions-vanishing-at-infinity--oa-fnd-sw-01}{Proposition
2.1}(2). \emph{Necessity:} \(p(0),p(1)\in\mathbb Z\) for
\(p\in\mathbb Z[x]\), and a limit of integers is an integer.

\emph{Step 1 (iterates).} Let \(g(x)=2x-3x^2+2x^3=x+x(1-x)(1-2x)\). Then
\(g(0)=0\), \(g(\frac12)=\frac12\), \(g(1)=1\), and
\(g'(x)=6(x-\frac12)^2+\frac12>0\). So \(g\) is an increasing bijection
of \([0,1]\), with \(g(x)>x\) on \((0,\frac12)\) and \(g(x)<x\) on
\((\frac12,1)\). Let \(g_m\) be the \(m\)-th iterate of \(g\); then
\(g_m\in\mathbb Z[x]\) and \(0\leq g_m\leq1\) on \([0,1]\). For
\(0<x\leq\frac12\), the sequence \(g_m(x)\) increases and is bounded by
\(\frac12\), so it converges to a fixed point of \(g\) in
\((0,\frac12]\). The fixed points of \(g\) are \(0\), \(\frac12\) and
\(1\), so the limit is \(\frac12\). The case \([\frac12,1)\) is
symmetric. For \(0<\varepsilon<\frac12\) and
\(x\in[\varepsilon,1-\varepsilon]\), monotonicity gives
\(g_m(\varepsilon)\leq g_m(x)\leq g_m(1-\varepsilon)\), so
\(g_m\to\frac12\) uniformly on \([\varepsilon,1-\varepsilon]\).

\emph{Step 2 (dyadic constants).} For \(n\geq0\) and \(0\leq k\leq2^n\)
there are \(\phi_m\in\mathbb Z[x]\) with \(0\leq\phi_m\leq1\) on
\([0,1]\) and \(\phi_m\to k/2^n\) uniformly on
\([\varepsilon,1-\varepsilon]\) for every \(\varepsilon\in(0,\frac12)\).
For \(k=0\) and \(k=2^n\) take the constants \(0\) and \(1\); this
settles \(n=0\). Suppose the claim holds for \(n\). If \(k\) is even,
\(k/2^{n+1}=(k/2)/2^n\). If \(k=2j+1\), take sequences
\(\phi_m\to j/2^n\) and \(\chi_m\to(j+1)/2^n\), and put
\(\psi_m=g_m\phi_m+(1-g_m)\chi_m\in\mathbb Z[x]\). It takes values in
\([0,1]\), as a convex combination, and
\[\psi_m-\frac{k}{2^{n+1}}=g_m\Bigl(\phi_m-\frac j{2^n}\Bigr)+(1-g_m)\Bigl(\chi_m-\frac{j+1}{2^n}\Bigr)-\Bigl(g_m-\frac12\Bigr)\frac1{2^n}\to0\]
uniformly on \([\varepsilon,1-\varepsilon]\).

\emph{Step 3 (real multiples).} Let \(p\in\mathbb Z[x]\) with
\(p(0)=p(1)=0\). The set
\(\Lambda_p=\{\lambda\in\mathbb R:\lambda p\in\mathcal G\}\) is a closed
additive subgroup of \(\mathbb R\) that contains \(\mathbb Z\). Let
\(\lambda=k/2^n\in[0,1]\) and \(\eta>0\). Choose \(\varepsilon\) with
\(|p|\leq\eta\) on \([0,\varepsilon]\cup[1-\varepsilon,1]\), and then
\(m\) with \(|\phi_m-\lambda|\leq\eta/(1+\|p\|)\) on
\([\varepsilon,1-\varepsilon]\). Since \(|\phi_m-\lambda|\leq1\)
everywhere, \(\|\phi_mp-\lambda p\|\leq\eta\), and
\(\phi_mp\in\mathbb Z[x]\). So \(\lambda\in\Lambda_p\). Adding integers,
\(\Lambda_p\) contains all dyadic rationals, and being closed,
\(\Lambda_p=\mathbb R\).

\emph{Step 4 (conclusion).} A real polynomial \(r\) with \(r(0)=r(1)=0\)
has the form \(x(1-x)s(x)=\sum_js_jx^{j+1}(1-x)\) with real \(s_j\), and
each term lies in \(\mathcal G\) by Step 3; so \(r\in\mathcal G\). These
polynomials form a real algebra. On \(X=(0,1)\) it lies in
\(C_0(X;\mathbb R)\), separates points (the ratio of \(x^2(1-x)\) to
\(x(1-x)\) is \(x\)), and vanishes nowhere. By
\hyperlink{the-stone-weierstrass-theorem-for-functions-vanishing-at-infinity--oa-fnd-sw-08}{Theorem
9.2} it is dense in \(C_0(X;\mathbb R)\). Directly from (1.1),
\(C_0(X;\mathbb R)\) consists of the restrictions of the
\(h\in C([0,1];\mathbb R)\) with \(h(0)=h(1)=0\). So every such \(h\)
lies in \(\mathcal G\). Finally, if \(f(0),f(1)\in\mathbb Z\), then
\(q(x)=f(0)+(f(1)-f(0))x\in\mathbb Z[x]\), and \(f-q\) vanishes at \(0\)
and \(1\); so \(f=q+(f-q)\in\mathcal G\). \(\square\)

\textbf{Proposition 18.2} (Integer coefficients on \([-1,1]\)). A
function \(f\in C([-1,1];\mathbb R)\) can be approximated uniformly on
\([-1,1]\) by polynomials with integer coefficients if and only if
\(f(-1)\), \(f(0)\), \(f(1)\) are integers and
\(f(-1)\equiv f(1)\pmod 2\).

\textbf{Proof.} Let \(\mathcal G'\) be the closure of \(\mathbb Z[x]\)
in \(C([-1,1];\mathbb R)\). \emph{Necessity:} for
\(p=\sum a_kx^k\in\mathbb Z[x]\), \(p(1)-p(-1)=2\sum_{k\text{ odd}}a_k\)
is even. Limits of integers are integers, and a convergent sequence of
integers is eventually constant, so the parity condition passes to
limits. \emph{Sufficiency:} let \(p\in\mathbb Z[x]\) vanish at \(-1\),
\(0\) and \(1\), and let \(\phi_m\) be as in Step 2 of the proof of
Proposition 18.1. The polynomials \(\psi_m(x)=\phi_m(1-x^2)\) have
integer coefficients and values in \([0,1]\) on \([-1,1]\). They
converge to \(k/2^n\) uniformly where
\(\varepsilon\leq1-x^2\leq1-\varepsilon\), that is, where
\(\sqrt\varepsilon\leq|x|\leq\sqrt{1-\varepsilon}\). Near \(0\) and
\(\pm1\), \(p\) is small. Step 3 of that proof, with these changes,
gives \(\lambda p\in\mathcal G'\) for all real \(\lambda\). Every real
polynomial vanishing at \(-1,0,1\) is
\(x(1-x^2)s(x)=\sum_js_jx^{j+1}(1-x^2)\), so it lies in \(\mathcal G'\).
These polynomials form a real algebra which, on \(X=(-1,0)\cup(0,1)\),
separates points (the ratio of \(x^2(1-x^2)\) to \(x(1-x^2)\) is \(x\))
and vanishes nowhere. By
\hyperlink{the-stone-weierstrass-theorem-for-functions-vanishing-at-infinity--oa-fnd-sw-08}{Theorem
9.2} it is dense in \(C_0(X;\mathbb R)\), which by (1.1) consists of the
restrictions of the continuous functions on \([-1,1]\) that vanish at
\(-1\), \(0\) and \(1\). Finally, \(\beta=(f(1)-f(-1))/2\) and
\(\gamma=(f(1)+f(-1))/2-f(0)\) are integers by the parity condition, and
\(q(x)=f(0)+\beta x+\gamma x^2\in\mathbb Z[x]\) agrees with \(f\) at
\(-1\), \(0\) and \(1\). So \(f=q+(f-q)\in\mathcal G'\). \(\square\)

\textbf{Proposition 18.3} (Smooth functions). Let \(I=[a,b]\) and
\(r\geq0\) an integer, and let \(f\) be \(r\) times continuously
differentiable on \(I\), with one-sided derivatives at the ends.

\begin{enumerate}
\tightlist
\item
  For every \(\varepsilon>0\) there is a polynomial \(p\) with
  \(\|f^{(k)}-p^{(k)}\|<\varepsilon\) for \(0\leq k\leq r\).
\item
  If \(f\) is infinitely differentiable, there are polynomials \(p_n\)
  with \(p_n^{(k)}\to f^{(k)}\) uniformly for every \(k\geq0\).
\item
  So the polynomials are dense in \(C^r(I)\) for the norm
  \(\max_{k\leq r}\|f^{(k)}\|\), and in \(C^\infty(I)\) for the topology
  of uniform convergence of each derivative.
\end{enumerate}

\textbf{Proof.} (1) Put \(c=\max(1,b-a)\). By the Weierstrass theorem
(\hyperlink{the-stone-weierstrass-theorem-for-functions-vanishing-at-infinity--oa-fnd-sw-14}{Corollary
14.3}) choose a polynomial \(q_r\) with
\(\|f^{(r)}-q_r\|<\varepsilon/c^r\). Going down, put
\(q_{k-1}(x)=f^{(k-1)}(a)+\int_a^xq_k(t)\,dt\) for \(k=r,\dots,1\), and
\(p=q_0\). Each \(q_{k-1}\) is a polynomial with \(q_{k-1}'=q_k\),
because the indefinite integral of a continuous function is
differentiable with that function as its derivative (B1); so
\(p^{(k)}=q_k\). By the fundamental theorem of calculus (B1),
\(f^{(k-1)}(x)=f^{(k-1)}(a)+\int_a^xf^{(k)}(t)\,dt\). So
\(\|f^{(k-1)}-q_{k-1}\|\leq(b-a)\|f^{(k)}-q_k\|\leq c\,\|f^{(k)}-q_k\|\),
by the bound \(|\int_a^xu|\leq(x-a)\sup|u|\) of (B1). Induction gives
\(\|f^{(k)}-q_k\|<c^{r-k}\varepsilon/c^r\leq\varepsilon\). (2) Apply (1)
with \(r=n\) and \(\varepsilon=1/n\). For fixed \(k\) and \(n\geq k\),
\(\|f^{(k)}-p_n^{(k)}\|<1/n\). (3) restates (1) and (2). \(\square\)

\subsection{Exercises}\label{the-stone-weierstrass-theorem-for-functions-vanishing-at-infinity--exercises}

\textbf{Exercise 1 (self-adjointness is needed on the line).} Let
\(\kappa(x)=(x-i)/(x+i)\) be the Cayley map of
\hyperlink{the-stone-weierstrass-theorem-for-functions-vanishing-at-infinity--oa-fnd-sw-02}{Example
3.3}(a), and \(u=\kappa-1\). Show that
\(A=\{q\circ u:\ q\text{ a complex polynomial},\ q(0)=0\}\) is a
subalgebra of \(C_0(\mathbb R)\) that separates points and vanishes
nowhere, but that \(\|\bar u-a\|\geq1\) for every \(a\in A\).

\emph{Solution.} \(u(x)=-2i/(x+i)\), so \(|u(x)|=2(1+x^2)^{-1/2}\). The
sets \(\{|u|\geq\varepsilon\}\) are closed and bounded, hence compact,
so \(u\in C_0(\mathbb R)\), and \(A\subseteq C_0(\mathbb R)\) by
\hyperlink{the-stone-weierstrass-theorem-for-functions-vanishing-at-infinity--oa-fnd-sw-01}{Proposition
2.1}(3). The condition \(q(0)=0\) survives sums, scalar multiples and
products, so \(A\) is a subalgebra. \(u\) is injective because
\(\kappa\) is, so \(A\) separates points; and \(u(x)\neq0\) because
\(\kappa(x)\neq1\), so \(A\) vanishes nowhere. By
\hyperlink{the-stone-weierstrass-theorem-for-functions-vanishing-at-infinity--oa-fnd-sw-02}{Proposition
3.2}(5) and
\hyperlink{the-stone-weierstrass-theorem-for-functions-vanishing-at-infinity--oa-fnd-sw-03}{Proposition
3.4}, the rule \(G(\kappa(x))=g(x)\), \(G(1)=0\), defines a
norm-preserving bijection \(g\mapsto G\) of \(C_0(\mathbb R)\) onto
\(\{G\in C(\mathbb T):G(1)=0\}\). For \(g=q\circ u\), \(G(w)=q(w-1)\),
also at \(w=1\), since \(q(0)=0\). For \(g=\bar u\), \(G(w)=\bar w-1\).
Take \(\Lambda_N\) from
\hyperlink{the-stone-weierstrass-theorem-for-functions-vanishing-at-infinity--oa-fnd-sw-11}{Example
12.2}(a) with \(N\geq\deg q+2\). Then \(\Lambda_N\) vanishes on the
polynomial \(q(w-1)\) and on \(1\), and \(\Lambda_N(\bar w)=1\). So
\(\|\bar u-q\circ u\|=\|\bar w-1-q(w-1)\|\geq|\Lambda_N(\bar w-1-q(w-1))|=1\).

\textbf{Exercise 2 (the route through the constants).} Let \(K\) be
compact and \(A\subseteq C(K;\mathbb R)\) a real algebra that separates
points and has \(Z(A)=\{z\}\). Let
\(B=\{a+s1:\ a\in\bar A,\ s\in\mathbb R\}\). Show that \(B\) is a closed
algebra, that \(1\in B\), and that \(B\) separates points. Deduce
\(\bar A=\{f:f(z)=0\}\) from the case of algebras with constants.

\emph{Solution.} \(\bar A\) is an algebra
(\hyperlink{the-stone-weierstrass-theorem-for-functions-vanishing-at-infinity--oa-fnd-sw-01}{Proposition
2.1}(2)), and \((a+s)(b+t)=(ab+ta+sb)+st\), so \(B\) is an algebra. Let
\(a_n+s_n\to f\) with \(a_n\in\bar A\). Evaluating at \(z\) gives
\(s_n\to f(z)\), so \(a_n\to f-f(z)1\), which then lies in \(\bar A\);
so \(f\in B\), and \(B\) is closed. \(B\) contains \(1\) and \(A\), so
it separates points. By
\hyperlink{the-stone-weierstrass-theorem-for-functions-vanishing-at-infinity--oa-fnd-sw-04}{Proposition
6.2}(3), \(B\) is closed under \(\wedge\), so
\hyperlink{the-stone-weierstrass-theorem-for-functions-vanishing-at-infinity--oa-fnd-sw-06}{Theorem
8.2}(2) makes it dense, and \(B=C(K;\mathbb R)\). If \(f(z)=0\), write
\(f=a+s1\) with \(a\in\bar A\). Then \(0=a(z)+s=s\), so
\(f=a\in\bar A\). The reverse inclusion is clear.
\hyperlink{the-stone-weierstrass-theorem-for-functions-vanishing-at-infinity--oa-fnd-sw-07}{Theorem
9.1} avoids the detour through \(B\).

\textbf{Exercise 3 (\(C_0\) of a product).} Let \(X\) and \(Y\) be LCH.
Show that the linear span of the functions
\((f\otimes g)(x,y)=f(x)g(y)\), \(f\in C_0(X)\), \(g\in C_0(Y)\), is
dense in \(C_0(X\times Y)\).

\emph{Solution.} \(X\times Y\) is Hausdorff, and a product of compact
neighbourhoods is a compact neighbourhood
(\hyperlink{the-stone-weierstrass-theorem-for-functions-vanishing-at-infinity--oa-fnd-sw-15}{Lemma
4.1}(b)); so it is LCH. If \(|f(x)g(y)|\geq\varepsilon\), then
\(|f(x)|\geq\varepsilon/(\|g\|+1)\) and
\(|g(y)|\geq\varepsilon/(\|f\|+1)\). So
\(\{|f\otimes g|\geq\varepsilon\}\) is closed and lies in a product of
two compact sets, and \(f\otimes g\in C_0(X\times Y)\) by (K1). The span
is a self-adjoint subalgebra, since
\((f\otimes g)(f'\otimes g')=(ff')\otimes(gg')\) and
\(\overline{f\otimes g}=\bar f\otimes\bar g\). Let \((x,y)\neq(x',y')\),
say \(x\neq x'\). By the locally compact form of
Urysohn\textquotesingle s lemma
(\hyperlink{the-stone-weierstrass-theorem-for-functions-vanishing-at-infinity--oa-fnd-sw-16}{Corollary
5.2}), choose \(f\in C_c(X)\) with \(f(x)=1\), \(f(x')=0\), and
\(g\in C_c(Y)\) with \(g=1\) on \(\{y,y'\}\). Then \(f\otimes g\) takes
the values \(1\) and \(0\) at the two points. The case \(y\neq y'\) is
symmetric, and the same functions show that the span vanishes nowhere.
Apply
\hyperlink{the-stone-weierstrass-theorem-for-functions-vanishing-at-infinity--oa-fnd-sw-09}{Theorem
10.1}.

\textbf{Exercise 4 (separability).} Let \(X\) be LCH and second
countable. Show that \(C_0(X)\) is separable.

\emph{Solution.} By the metrizability theorem,
\hyperlink{the-stone-weierstrass-theorem-for-functions-vanishing-at-infinity--oa-fnd-sw-16}{Theorem
5.3}, \(X_\infty\) is a compact metrizable space. By
\hyperlink{the-stone-weierstrass-theorem-for-functions-vanishing-at-infinity--oa-fnd-sw-17}{Proposition
15.1}(3), \(C(X_\infty)\) is separable. A subset of a separable metric
space is separable, since the space has a countable base of balls. So
\(I_\infty\) is separable, and so is \(C_0(X)\), which is isometric to
it
(\hyperlink{the-stone-weierstrass-theorem-for-functions-vanishing-at-infinity--oa-fnd-sw-03}{Proposition
3.4}).

\subsection{Where this
leads}\label{the-stone-weierstrass-theorem-for-functions-vanishing-at-infinity--where-this-leads}

\begin{itemize}
\tightlist
\item
  \emph{Commutative C*-algebras.} In
  \href{https://kokunoyumeto.github.io/open-math-courses/courses/harmonic-analysis-on-locally-compact-groups/reader/programme-reading.html\#ref-a8616a90904d331b}{the
  lesson on C*-algebras and the continuous functional calculus},
  \hyperlink{the-stone-weierstrass-theorem-for-functions-vanishing-at-infinity--oa-fnd-sw-09}{Theorem
  10.1} shows that the Gelfand transform of a commutative C*-algebra
  maps onto \(C_0\) of its character space. Urysohn\textquotesingle s
  lemma and the density theorems of Section 15 are used there as well.
\item
  \emph{States on C*-algebras.} The lesson "Integral representations of
  states" uses Stone\textquotesingle s lattice lemma and the real and
  complex theorems to show that certain spaces of continuous functions
  on compact convex sets and on state spaces are dense.
\item
  \emph{Locally compact groups.} In the course "Modular theory and
  weights", the lesson "The Plancherel weight and Fourier coefficients
  of a locally compact group" shows that the Fourier algebra \(A(G)\) of
  a locally compact group \(G\) lies in \(C_0(G)\), is closed under
  pointwise products and complex conjugation, separates points and
  vanishes nowhere.
  \hyperlink{the-stone-weierstrass-theorem-for-functions-vanishing-at-infinity--oa-fnd-sw-09}{Theorem
  10.1} then makes \(A(G)\) dense in \(C_0(G)\). Since the proof uses no
  countability, this holds for every locally compact group.
\item
  \emph{Algebras that are not self-adjoint.} For these the conclusion
  can fail, as the disc algebra of
  \hyperlink{the-stone-weierstrass-theorem-for-functions-vanishing-at-infinity--oa-fnd-sw-19}{Proposition
  17.3} shows. Bishop\textquotesingle s theorem {[}Bishop 1961{]}
  reduces approximation by a closed subalgebra of \(C(K)\) to its
  maximal antisymmetric sets.
\end{itemize}

\subsection{Results used from other
lessons}\label{the-stone-weierstrass-theorem-for-functions-vanishing-at-infinity--results-used-from-other-lessons}

The following facts are proved in the core courses
\href{https://kokunoyumeto.github.io/program-matematika-indonesia/en/\#course-C10}{Real
Analysis I} and
\href{https://kokunoyumeto.github.io/program-matematika-indonesia/en/\#course-C20}{Real
Analysis II}, whose text is {[}Lebl{]}, and
\href{https://kokunoyumeto.github.io/program-matematika-indonesia/en/\#course-B10}{Proof,
Logic, and Discrete Structures}, whose text is {[}Levin{]}. Each item
gives the place of the proof. A few steps that these texts leave to the
reader are proved here.

\begin{itemize}
\tightlist
\item
  (B1) \emph{The Riemann integral.} Every continuous function
  \(u:[a,b]\to\mathbb R\) is Riemann integrable
  (\href{https://www.jirka.org/ra/html/sec_rintprop.html\#lemma_contint}{Real
  Analysis I, Lemma 5.2.7}). The integral is additive over subintervals:
  \(\int_a^bu=\int_a^cu+\int_c^bu\) for \(a\leq c\leq b\)
  (\href{https://www.jirka.org/ra/html/sec_rintprop.html\#sec_rintprop-3-5}{Proposition
  5.2.2}). It is monotone: if \(u\leq v\), then
  \(\int_a^bu\leq\int_a^bv\)
  (\href{https://www.jirka.org/ra/html/sec_rintprop.html\#sec_rintprop-4-8}{Proposition
  5.2.6}). If \(m\leq u\leq M\), then \(m(b-a)\leq\int_a^bu\leq M(b-a)\)
  (\href{https://www.jirka.org/ra/html/sec_rint.html\#intbound_prop}{Proposition
  5.1.10}). And \(\int_a^b\alpha u=\alpha\int_a^bu\) for \(\alpha\geq0\)
  (\href{https://www.jirka.org/ra/html/sec_rintprop.html\#prop_integrallinear}{Proposition
  5.2.4}). The remaining cases of linearity are proved as follows. For a
  bounded function \(u\) and a partition \(P\), the infimum of \(-u\) on
  each cell of \(P\) is minus the supremum of \(u\) there, so the lower
  and upper sums satisfy \(L(P,-u)=-U(P,u)\) and \(U(P,-u)=-L(P,u)\).
  Taking the supremum and the infimum over \(P\) shows that \(-u\) is
  integrable with \(\int_a^b(-u)=-\int_a^bu\) when \(u\) is. For bounded
  \(u\) and \(v\), on each cell the infimum of \(u+v\) is at least the
  sum of the infima of \(u\) and \(v\), and the supremum of \(u+v\) is
  at most the sum of the suprema, so \(L(P,u)+L(P,v)\leq L(P,u+v)\) and
  \(U(P,u+v)\leq U(P,u)+U(P,v)\). Let \(u\) and \(v\) be integrable, let
  \(P_1\) and \(P_2\) be partitions, and put \(P=P_1\cup P_2\). Refining
  a partition does not lower its lower sums or raise its upper sums
  (\href{https://www.jirka.org/ra/html/sec_rint.html\#prop_refinement}{Real
  Analysis I, Proposition 5.1.7}), and the lower integral is at most the
  upper integral
  (\href{https://www.jirka.org/ra/html/sec_rint.html\#intulbound_prop}{Proposition
  5.1.8}), so
  \[L(P_1,u)+L(P_2,v)\leq L(P,u+v)\leq\underline{\int_a^b}(u+v)\leq\overline{\int_a^b}(u+v)\leq U(P,u+v)\leq U(P_1,u)+U(P_2,v).\]
  Taking the supremum of the left side and the infimum of the right side
  over \(P_1\) and \(P_2\) gives
  \(\int_a^bu+\int_a^bv\leq\underline{\int_a^b}(u+v)\leq\overline{\int_a^b}(u+v)\leq\int_a^bu+\int_a^bv\).
  So \(u+v\) is integrable and \(\int_a^b(u+v)=\int_a^bu+\int_a^bv\).
  For a continuous complex function \(u=u_1+iu_2\) one puts
  \(\int_a^bu=\int_a^bu_1+i\int_a^bu_2\). By the real case, applied to
  real and imaginary parts, this integral is complex linear. Moreover
  \(|\int_a^bu|\leq\int_a^b|u|\leq(b-a)\sup_{[a,b]}|u|\). The second
  inequality is Proposition 5.1.10 applied to the continuous function
  \(|u|\). For the first, let \(\int_a^bu\neq0\) and put
  \(c=|\int_a^bu|/\int_a^bu\). Then \(|c|=1\), and
  \(\int_a^bcu=c\int_a^bu=|\int_a^bu|\) is real, so it equals its real
  part \(\int_a^b\operatorname{Re}(cu)\). Since
  \(\operatorname{Re}(cu)\leq|cu|=|u|\), monotonicity gives
  \(|\int_a^bu|\leq\int_a^b|u|\). A uniformly convergent series of
  continuous functions can be integrated term by term: apply
  \href{https://www.jirka.org/ra/html/sec_liminter.html\#integralinterchange_thm}{Real
  Analysis II, Theorem 6.2.4} to the partial sums, which are continuous
  and so integrable, and to their real and imaginary parts. If \(u\) is
  continuous, then \(x\mapsto\int_a^xu\) is differentiable on \([a,b]\)
  with derivative \(u\)
  (\href{https://www.jirka.org/ra/html/sec_ftc.html\#thm_FTCv2}{Real
  Analysis I, Theorem 5.3.3}). If \(F\) is differentiable on \([a,b]\)
  and \(F'\) is continuous, then \(\int_a^bF'=F(b)-F(a)\)
  (\href{https://www.jirka.org/ra/html/sec_ftc.html\#thm_FTCv1}{Theorem
  5.3.1}, the fundamental theorem of calculus). For complex functions,
  apply both to real and imaginary parts. See also, e.g., {[}Rudin,
  Theorems 6.8, 6.12, 6.13, 6.20, 6.21, 6.25 and 7.16{]}.
\item
  (B2) \emph{The exponential function.} Put \(E(z)=\sum_{n\geq0}z^n/n!\)
  for \(z\in\mathbb C\); the series converges for every \(z\). Write
  \(e^z=E(z)\). The law of exponents \(E(z+w)=E(z)E(w)\) is
  \href{https://www.jirka.org/ra/html/sec_complexexp.html\#sec_complexexp-3-7}{Real
  Analysis II, Proposition 11.4.1}.
  \href{https://www.jirka.org/ra/html/sec_complexexp.html\#sec_complexexp-4-4}{Proposition
  11.4.2} and its proof give the following for real \(x\):
  \(E(ix)=\cos x+i\sin x\) with \(\cos x\) and \(\sin x\) real, and
  \(|E(ix)|=1\); \(E(i\pi/2)=i\) and \(E(2\pi i)=1\), so \(E\) has
  period \(2\pi i\); \(t\mapsto E(it)\) is one-to-one on \([0,2\pi)\),
  so \(E(it)\neq1=E(0)\) for \(0<t<2\pi\); and every \(z\in\mathbb T\)
  with \(\operatorname{Re}z\geq0\) and \(\operatorname{Im}z\geq0\) is
  \(E(ix)\) for some \(x\in[0,\pi/2]\). Two further facts are proved
  here. \emph{Every point of \(\mathbb T\) is \(E(it)\) for some
  \(t\in[0,2\pi)\).} Multiplication by \(-i\) maps \(a+ib\) to \(b-ia\).
  So for \(z\in\mathbb T\), one of \(z\), \(-iz\), \(-z\) and
  \(iz=(-i)^3z\) has nonnegative real and imaginary parts:
  \((-i)^kz=E(ix)\) for some \(k\in\{0,1,2,3\}\) and \(x\in[0,\pi/2]\).
  Then \(z=i^kE(ix)=E\bigl(i(x+k\pi/2)\bigr)\) by the law of exponents,
  and \(x+k\pi/2\in[0,2\pi]\). If \(x+k\pi/2=2\pi\), then
  \(z=E(2\pi i)=1=E(0)\). \emph{For real \(m\), the function
  \(t\mapsto E(imt)\) on \(\mathbb R\) is differentiable with derivative
  \(imE(imt)\).} For complex \(w\) with \(|w|\leq1\),
  \[|E(w)-1-w|\leq\sum_{k\geq2}\frac{|w|^k}{k!}\leq|w|^2\sum_{k\geq2}\frac1{k!}\leq|w|^2,\]
  by the triangle inequality for the partial sums, and since
  \(k!\geq2^{k-1}\) gives
  \(\sum_{k\geq2}1/k!\leq\sum_{k\geq2}2^{1-k}=1\). For real \(h\neq0\)
  with \(|mh|\leq1\), the law of exponents gives
  \(E(im(t+h))-E(imt)-imhE(imt)=E(imt)\bigl(E(imh)-1-imh\bigr)\), and
  \(|E(imt)|=1\). Hence
  \[\Bigl|\frac{E(im(t+h))-E(imt)}h-imE(imt)\Bigr|\leq\frac{|mh|^2}{|h|}=m^2|h|,\]
  which tends to \(0\) as \(h\to0\). See also, e.g., {[}Rudin,
  (8.25)--(8.28), (8.48) and Theorem 8.7{]}.
\item
  (B3) \emph{Numbers.} The real numbers form an ordered field with the
  least upper bound property; \emph{Real Analysis I} starts from this
  (\href{https://www.jirka.org/ra/html/sec_setofreals.html\#sec_setofreals-3-3}{Theorem
  1.2.1}), and a construction is in, e.g., {[}Rudin, Theorem 1.19 and
  the appendix to Chapter 1{]}. The geometric sum
  \(\sum_{k=0}^{N-1}\rho^k=(1-\rho^N)/(1-\rho)\) for \(\rho\neq1\) is
  proved by induction in
  \href{https://www.jirka.org/ra/html/sec_basicset.html\#example_geometricsum}{Real
  Analysis I, Example 0.3.8}. The induction uses only the field
  operations, so the identity holds for complex \(\rho\). If
  \(|\rho|\leq r<1\), then
  \[\Bigl|\sum_{k=0}^{N-1}\rho^k-\frac1{1-\rho}\Bigr|=\frac{|\rho|^N}{|1-\rho|}\leq\frac{r^N}{1-r},\]
  which tends to \(0\) as \(N\to\infty\)
  (\href{https://www.jirka.org/ra/html/sec_factslimsseqs.html\#sec_factslimsseqs-6-6}{Real
  Analysis I, Proposition 2.2.11}). So
  \(\sum_{k\geq0}\rho^k=1/(1-\rho)\), uniformly for \(|\rho|\leq r\).
  The binomial theorem is
  \href{https://kokunoyumeto.github.io/program-matematika-indonesia/en/courses/B10/reader/sec_counting-pascal.html\#thm-binomial}{Proof,
  Logic, and Discrete Structures, Theorem 3.1.9}, and the formula
  \(\binom nk=\frac{n!}{k!\,(n-k)!}\) is
  \href{https://kokunoyumeto.github.io/program-matematika-indonesia/en/courses/B10/reader/sec_counting-combperm.html\#subsec-combinations-9}{Theorem
  3.4.9} there. Bernoulli\textquotesingle s inequality
  \((1+u)^n\geq1+nu\) for \(u\geq0\) follows from the binomial theorem,
  because all its terms are nonnegative. The Cauchy--Schwarz inequality
  \(\bigl(\sum_kw_ka_k\bigr)^2\leq\bigl(\sum_kw_k\bigr)\bigl(\sum_kw_ka_k^2\bigr)\)
  for finitely many real numbers \(a_k\) and weights \(w_k\geq0\) is
  \href{https://www.jirka.org/ra/html/sec_metric.html\#sec_metric-16}{Real
  Analysis II, Lemma 7.1.4}, applied to the vectors \((\sqrt{w_k})_k\)
  and \((\sqrt{w_k}\,a_k)_k\); see also, e.g., {[}Rudin, Theorem
  1.35{]}.
\end{itemize}

\subsection{References}\label{the-stone-weierstrass-theorem-for-functions-vanishing-at-infinity--references}

\begin{itemize}
\tightlist
\item
  {[}Alexandroff 1924{]} P. Alexandroff, Über die Metrisation der im
  Kleinen kompakten topologischen Räume, \emph{Math. Ann.} 92 (1924),
  294--301.
  https://resolver.sub.uni-goettingen.de/purl?PPN235181684\_0092\%7CLOG\_0024
\item
  {[}Bernstein 1912{]} S. Bernstein, Démonstration du théorème de
  Weierstrass fondée sur le calcul des probabilités, \emph{Comm. Soc.
  Math. Kharkow} (2) 13 (1912), 1--2.
\item
  {[}Bishop 1961{]} E. Bishop, A generalization of the
  Stone--Weierstrass theorem, \emph{Pacific J. Math.} 11 (1961),
  777--783. https://doi.org/10.2140/pjm.1961.11.777
\item
  {[}Blackadar{]} B. Blackadar, \emph{Real Analysis}, preliminary
  edition, 2025. https://bruceblackadar.com/mathpubs.html
\item
  {[}Korovkin 1953{]} P. P. Korovkin, On convergence of linear positive
  operators in the space of continuous functions, \emph{Dokl. Akad. Nauk
  SSSR (N.S.)} 90 (1953), 961--964 (in Russian).
\item
  {[}Lebl{]} J. Lebl, \emph{Basic Analysis I} and \emph{Basic Analysis
  II: Introduction to Real Analysis}, version 6.3 (15 May 2026), open
  textbook under CC BY-SA 4.0 and CC BY-NC-SA 4.0,
  \href{https://www.jirka.org/ra/}{www.jirka.org/ra}. It is the text of
  the core courses \emph{Real Analysis I} and \emph{Real Analysis II}.
\item
  {[}Levin{]} O. Levin, \emph{Discrete Mathematics: An Open
  Introduction}, 4th ed., open textbook under CC BY-NC-SA 4.0,
  \href{https://discrete.openmathbooks.org/dmoi4/}{discrete.openmathbooks.org}.
  It is the text of the core course \emph{Proof, Logic, and Discrete
  Structures}, which has an English reader at
  \href{https://kokunoyumeto.github.io/program-matematika-indonesia/en/courses/B10/reader/}{the
  programme site}.
\item
  {[}Rudin{]} W. Rudin, \emph{Principles of Mathematical Analysis}, 3rd
  ed., McGraw-Hill, New York, 1976.
\item
  {[}Steen--Seebach{]} L. A. Steen and J. A. Seebach, Jr.,
  \emph{Counterexamples in Topology}, 2nd ed., Springer, New York, 1978.
\item
  {[}π-Base{]} The π-Base Community, S. Clontz and J. Dabbs,
  \emph{π-Base: a community database of topological counterexamples},
  open database under CC BY 4.0,
  \href{https://topology.pi-base.org/}{topology.pi-base.org}.
\item
  {[}Stone 1937{]} M. H. Stone, Applications of the theory of Boolean
  rings to general topology, \emph{Trans. Amer. Math. Soc.} 41 (1937),
  375--481. https://doi.org/10.1090/S0002-9947-1937-1501905-7
\item
  {[}Stone 1948{]} M. H. Stone, The generalized Weierstrass
  approximation theorem, \emph{Math. Mag.} 21 (1948), 167--184 and
  237--254.
\item
  {[}Urysohn 1925{]} P. Urysohn, Über die Mächtigkeit der
  zusammenhängenden Mengen, \emph{Math. Ann.} 94 (1925), 262--295.
  https://resolver.sub.uni-goettingen.de/purl?PPN235181684\_0094\%7CLOG\_0020
\item
  {[}van Neerven{]} J. van Neerven, \emph{Functional Analysis},
  arXiv:2112.11166, version 4 (28 April 2022) and version 7 (17 July
  2025). https://arxiv.org/abs/2112.11166
\item
  {[}Weierstrass 1885{]} K. Weierstrass, Über die analytische
  Darstellbarkeit sogenannter willkürlicher Functionen einer reellen
  Veränderlichen, \emph{Sitzungsber. Königl. Preuss. Akad. Wiss. Berlin}
  (1885), 633--639 and 789--805.
  https://archive.org/details/sitzungsberichte1885deutsch/page/633/mode/1up
\end{itemize}

\section{Restricted products and profinite completions}\label{NT-ADL-01}

\emph{Written by GPT-6.1 Sol (OpenAI) in Codex, Ultra setting, October
2026. Self-checked by its author. Public domain (CC0).}

A finite adèle records one number in each field \(\mathbb Q_p\), with an
integrality condition at all but finitely many primes. Its topology must
allow those finitely many exceptions while retaining a compact
neighbourhood. Restricted products provide exactly this construction.
They also explain how local integrals and local characters combine into
global ones.

We first build the topology, then integration and characters. We apply
the construction to the profinite integers and finite adèles, and finish
with test functions and the topology of finite idèles. The prerequisites
are point-set topology and measure and integration, together with
\hyperref[haar-measure-on-locally-compact-groups]{Haar measure on
locally compact groups}, \emph{The Pontryagin duality theorem},
\emph{Subgroups, quotients and annihilators}, and \emph{Totally
disconnected groups, the p-adic numbers and the adèles}. The precise
background statements used without proof are listed below.

Basic references are {[}Getz--Hahn{]}, {[}Bost--Connes{]}, {[}Deitmar{]}
and {[}Halter-Koch{]}.

\subsection{A topology with finitely many exceptional
coordinates}\label{NT-ADL-01--a-topology-with-finitely-many-exceptional-coordinates}

Let \(V\) be any index set. Each \(G_v\) is a locally compact Hausdorff
group. Fix a finite set \(E\subseteq V\), and for \(v\notin E\) choose a
compact open subgroup \(K_v\leq G_v\). Coordinates in \(E\) are
unrestricted; this permits real or complex factors. Throughout, ``almost
all'' means all but finitely many.

The \textbf{restricted product} is the group
\[G=\prod_{v\in V}'(G_v,K_v)
=\left\{x\in\prod_{v\in V}G_v:
x_v\in K_v\text{ for almost all }v\notin E\right\}.\] For finite
\(S\supseteq E\), set
\[G_S=\prod_{v\in S}G_v\times\prod_{v\notin S}K_v.\] Give \(G\) the
topology with basis
\[B(S;U)=\prod_{v\in S}U_v\times\prod_{v\notin S}K_v,
\qquad U_v\subseteq G_v\text{ open}.
\tag{1}\] The finite set can always be enlarged: a newly included
coordinate is assigned the open set \(K_v\). This observation proves
that intersections of basic sets are basic sets, allowing empty factors.

\textbf{Proposition 1.1.} The restricted product is a locally compact
Hausdorff topological group. It is the upward-directed union of the open
subgroups \(G_S\), and each \(G_S\) has its product topology. Every
compact subset of \(G\) lies in some \(G_S\).

\textbf{Proof.} Every tuple satisfies the defining restriction outside a
finite set, so \(G=\bigcup_SG_S\). Coordinatewise multiplication and
inversion preserve the condition because the \(K_v\) are subgroups. Thus
\(G\) is an algebraic subgroup of the full product.

The sets \(G_S\) are open by (1). Intersecting (1) with \(G_S\) imposes
open conditions on only finitely many factors of its product
decomposition. Conversely, every product neighbourhood in \(G_S\) is
obtained this way, since its finitely many conditions on tail factors
are open in the corresponding \(G_v\). Hence the induced topology is the
product topology.

For \(x,y\in G\), choose one \(G_S\) containing both. Multiplication is
continuous on \(G_S\times G_S\), and inversion is continuous on \(G_S\).
These are open neighbourhoods of the relevant points in the domain,
proving continuity on \(G\).

Each coordinate projection is continuous: at a point, enlarge \(S\) to
include that coordinate. Distinct tuples therefore have disjoint
neighbourhoods pulled back from a Hausdorff factor. This proves the
Hausdorff property.

For \(x\in G_S\), choose a compact neighbourhood \(C_v\) of \(x_v\) in
each of the finitely many factors indexed by \(S\). The set
\[\prod_{v\in S}C_v\times\prod_{v\notin S}K_v\] is a compact
neighbourhood of \(x\), by the compact-product theorem. This proves
local compactness. Finally, the open cover \((G_S)_S\) of a compact
subset has a finite subcover. The union of the finitely many indexing
sets gives a single chart containing that subset. \(\square\)

The inclusion \(G\hookrightarrow\prod_vG_v\) is continuous. It generally
fails to be a topological embedding: (1) constrains every coordinate
outside \(S\), whereas a product neighbourhood constrains only finitely
many coordinates. Exercise 1 proves the exact criterion. Changing the
chosen subgroups at finitely many indices changes neither the underlying
restricted product nor its topology; include those indices in every
chart.

\subsection{Integrating local
data}\label{NT-ADL-01--integrating-local-data}

Choose left Haar measures \(\mu_v\) on \(G_v\). Enlarge \(E\), if
necessary, so that \[\mu_v(K_v)=1\qquad(v\notin E).
\tag{2}\] The enlargement retains any exceptional local normalization as
an actual finite factor. An infinite product of arbitrary normalization
constants is not part of the construction.

On the compact group \(L_S=\prod_{v\notin S}K_v\), let \(m_S\) be Haar
probability. Its projection to any finite product of tail factors has
the product of their normalized Haar measures: the projected measure is
invariant, has total mass one, and Haar uniqueness applies. This
describes the infinite product as a Radon measure on its compact product
topology.

Use the \textbf{Radon product} for finite products of locally compact
spaces: it is the measure determined by iterated integrals of compactly
supported continuous functions. It avoids imposing a countability
assumption on the Borel sets of a product.

\textbf{Proposition 1.2.} The measures
\[\mu_S=\left(\mathop{\widehat\bigotimes}_{v\in S}\mu_v\right)
\widehat\otimes m_S
\quad\text{on }G_S
\tag{3}\] are compatible under restriction and determine a unique left
Haar measure \(\mu\) on \(G\). If \(f_v\) are Borel integrable functions
and \(f_v=1_{K_v}\) for all \(v\notin S\), where \(S\supseteq E\) is
finite, then \[f(x)=\bigotimes_v f_v(x):=\prod_v f_v(x_v)\] is
integrable and \[\int_Gf\,d\mu=\prod_v\int_{G_v}f_v\,d\mu_v.
\tag{4}\] The product on the right has only finitely many factors
different from one. In particular, (4) applies to compactly supported
continuous local functions.

\textbf{Proof.} Suppose \(S\subseteq T\). Restricting \(\mu_T\) to
\(G_S\) replaces each \(\mu_v\), for \(v\in T\setminus S\), by its
restriction to \(K_v\), of mass one. Their finite product with \(m_T\)
is Haar probability on \(L_S\), hence is \(m_S\). This proves
compatibility in (3).

For \(h\in C_c(G)\), Proposition 1.1 places its support in some \(G_S\).
Define \[I(h)=\int_{G_S}h\,d\mu_S.\] Compatibility makes this
independent of \(S\), and a common larger chart proves linearity and
positivity. Given \(g\in G\), a chart containing both \(g\) and
\(\operatorname{supp}h\) also contains \(g\operatorname{supp}h\). Its
measure is a product of left Haar measures, so
\(I(h(g^{-1}\cdot))=I(h)\).

The Riesz representation theorem gives a unique Radon measure
representing \(I\). It is nonzero and left invariant, so it is Haar.
Functions in \(C_c(G_S)\) extend by zero to \(G\), since an open
subgroup is also closed. Riesz uniqueness then gives
\(\mu|_{G_S}=\mu_S\), establishing the asserted restrictions and
uniqueness.

For (4), \(f\) vanishes off \(G_S\), and on that chart it is a finite
product of local functions, independent of \(L_S\). Each set
\(\{f_v\neq0\}\) is \(\sigma\)-finite for \(\mu_v\): it is the union of
the finite-measure sets \(\{|f_v|>1/j\}\). The rectangle and Tonelli
theorems for Radon products therefore apply to the absolute values of
these factors. They give
\[\int_G|f|\,d\mu=\prod_{v\in S}\int|f_v|\,d\mu_v<\infty.\] Fubini now
gives (4). This argument uses \(\sigma\)-finiteness of the nonzero sets
of the particular functions, not of the ambient groups. \(\square\)

Thus the notation \(\mu=\bigotimes_v\mu_v\) means the compatible Radon
measure just constructed. Its definition does not require \(G\) to be
second countable or \(\sigma\)-compact.

\subsection{Characters see only finitely many integral
tails}\label{NT-ADL-01--characters-see-only-finitely-many-integral-tails}

Now assume all \(G_v\) are abelian and write their law additively. Put
\(\mathbb T=\{z\in\mathbb C:|z|=1\}\). The dual \(\widehat H\) of an LCA
group consists of continuous homomorphisms \(H\to\mathbb T\), with
uniform convergence on compact sets. For a subgroup \(K\leq H\), its
annihilator is
\[K^\perp=\{\chi\in\widehat H:\chi(k)=1\text{ for every }k\in K\}.\] The
subgroup and quotient duality theorems imply that \(K_v^\perp\) is
compact and open when \(K_v\) is compact and open. These are the
distinguished subgroups on the dual side.

We use a small fact about the circle. The arc
\(A=\{z\in\mathbb T:\operatorname{Re}z>0\}\) contains no nontrivial
subgroup. Indeed, write a nonidentity element as \(e^{2\pi it}\) with
\(0<|t|\leq1/2\). If \(|t|\geq1/4\), it is already outside \(A\).
Otherwise choose the least positive integer \(m\) with \(m|t|\geq1/4\).
Then \(1/4\leq m|t|<1/2\), and its \(m\)-th power is outside \(A\).

\textbf{Theorem 1.3.} There is a canonical isomorphism of topological
groups \[\Phi:\prod_v'(\widehat G_v,K_v^\perp)
\xrightarrow{\ \sim\ }\widehat G,
\qquad
\Phi((\chi_v))(x)=\prod_v\chi_v(x_v).
\tag{5}\] The factors indexed by \(E\) remain unrestricted. Every
continuous character of \(G\) has local restrictions trivial on \(K_v\)
for almost all \(v\).

\textbf{Proof.} For a tuple on the left and an element \(x\in G\), both
restrictions hold outside finite sets. Consequently all but finitely
many factors in (5) equal one. On a chart containing the exceptional
coordinates of the character tuple, (5) is a finite product of
continuous local characters. It is therefore a continuous character on
the open cover of \(G\).

Conversely, let \(\chi\in\widehat G\). Continuity provides a basic
neighbourhood \(B(S;U)\) of zero whose image lies in \(A\). The compact
subgroup \[T_S=\{0\}^{S}\times\prod_{v\notin S}K_v\] lies in this
neighbourhood. Since \(\chi(T_S)\) is a subgroup contained in \(A\), it
is trivial. Restricting \(\chi\) to the embedded coordinate group gives
\(\chi_v\in\widehat G_v\), and \(\chi_v\in K_v^\perp\) outside \(S\).
For any \(x\), enlarge \(S\) to contain its nonintegral coordinates.
Decompose \(x\) into a finite-coordinate part and an element of \(T_S\).
The latter contributes one, so \(\chi(x)=\prod_v\chi_v(x_v)\).
Uniqueness follows by evaluating on individual coordinates. This proves
algebraic bijectivity.

We verify both topologies explicitly. Let \(C\subseteq G\) be compact.
Choose \(S\supseteq E\) with \(C\subseteq G_S\), and let \(C_v\) be its
compact coordinate projections for \(v\in S\). Restrict dual tuples to
the open dual chart
\[H_S=\prod_{v\in S}\widehat G_v\times\prod_{v\notin S}K_v^\perp.\]
Their evaluations on \(C\) have only the factors indexed by \(S\).
Requiring each such factor to be uniformly within
\(\varepsilon/\max(1,|S|)\) of one on \(C_v\) gives an open
neighbourhood in \(H_S\). The inequality
\[\left|\prod_{j=1}^rz_j-1\right|\leq\sum_{j=1}^r|z_j-1|\qquad(z_j\in\mathbb T)\]
shows that its image is uniformly within \(\varepsilon\) of one on
\(C\). Thus \(\Phi\) is continuous.

For the inverse, take a basic identity neighbourhood with
finite-coordinate conditions \(\chi_v\in U_v\) for \(v\in S\), and the
condition \(\chi_v\in K_v^\perp\) outside \(S\). Restriction
\(\widehat G\to\widehat G_v\) is continuous, since a compact set in
\(G_v\) embeds as a compact set in \(G\). The remaining conditions say
exactly that \(\chi\) annihilates \(T_S\). Its annihilator is open: the
compact-open condition \(\chi(T_S)\subseteq A\) is equivalent to
triviality there. Intersecting this open annihilator with the finitely
many inverse images of \(U_v\) gives the image of the chosen
neighbourhood. This proves continuity of \(\Phi^{-1}\). \(\square\)

The theorem concerns a restricted product of \textbf{abelian} groups.
Proposition 1.1 and the Haar construction require no commutativity.

\subsection{Integral residues and finite
adèles}\label{NT-ADL-01--integral-residues-and-finite-aduxe8les}

The Chinese remainder theorem used below has a short algebraic proof.
For pairwise coprime positive integers \(m_1,\ldots,m_r\), put
\(M=\prod_i m_i\). Bézout\textquotesingle s identity gives \(b_i\) with
\(b_i(M/m_i)\equiv1\pmod{m_i}\). Given residues \(a_i\), the integer
\(\sum_i a_i b_i(M/m_i)\) realizes all of them. An integer divisible by
every \(m_i\) is divisible by \(M\), proving uniqueness modulo \(M\).

Define the \textbf{profinite integers} by
\[\widehat{\mathbb Z}=\varprojlim_{n\geq1}\mathbb Z/n\mathbb Z,\] where
indices are ordered by divisibility and the transition maps are
reduction maps. A point is a compatible collection of residues. The
topology is inherited from the product of the finite discrete rings.
Compatibility is closed, so this is a compact Hausdorff ring.

For each prime \(p\), we use
\(\mathbb Z_p=\varprojlim_k\mathbb Z/p^k\mathbb Z\) and its fraction
field \(\mathbb Q_p\). Its additive neighbourhoods are
\(p^k\mathbb Z_p\), for \(k\in\mathbb Z\); \(\mathbb Z_p\) is compact
and open. Define \[\mathbb A_f=\prod_p'(\mathbb Q_p,\mathbb Z_p).\]
Addition and multiplication are coordinatewise. This is a topological
ring: on each open chart
\(\prod_{p\in S}\mathbb Q_p\times\prod_{p\notin S}\mathbb Z_p\) both
operations are continuous and preserve the chart. Proposition 1.1 gives
local compactness.

\textbf{Proposition 1.4.} There are canonical isomorphisms
\[\widehat{\mathbb Z}\simeq\prod_p\mathbb Z_p,
\qquad
\widehat{\mathbb Z}^{\times}\simeq\prod_p\mathbb Z_p^{\times},
\tag{6}\] of compact topological rings and of compact topological
groups, respectively. Moreover,
\[\mathbb A_f\simeq\mathbb Q\otimes_{\mathbb Z}\widehat{\mathbb Z}
=\bigcup_{n\geq1}\frac1n\widehat{\mathbb Z}
\tag{7}\] as topological rings, when the tensor product is given the
topology specified below. There are isomorphisms of discrete groups
\[\mathbb A_f/\widehat{\mathbb Z}
\simeq\mathbb Q/\mathbb Z
\simeq\bigoplus_p\mathbb Q_p/\mathbb Z_p,
\qquad
\widehat{\widehat{\mathbb Z}}\simeq\mathbb Q/\mathbb Z.
\tag{8}\] In the last formula, the outer hat denotes the character
group.

\textbf{Proof.} Restriction to the prime-power residues maps
\(\widehat{\mathbb Z}\) to \(\prod_p\mathbb Z_p\). Conversely, given
prime-power residues and \(n=\prod_{p\mid n}p^{a_p}\), the Chinese
remainder theorem supplies a unique residue modulo \(n\). These residues
are compatible as \(n\) varies. This proves (6) for rings. The map is
continuous, and a continuous bijection from a compact space to a
Hausdorff space is a homeomorphism. A tuple in a product ring is a unit
precisely when every component is a unit; componentwise inversion is
continuous on the compact product of the local unit groups. This proves
the second assertion.

We will need the more precise residue calculation
\[\widehat{\mathbb Z}/n\widehat{\mathbb Z}\simeq\mathbb Z/n\mathbb Z.
\tag{9}\] In \(\mathbb Z_p\), the kernel of reduction modulo \(p^a\) is
\(p^a\mathbb Z_p\): a compatible tuple in the kernel can be divided by
\(p^a\) by using its residues modulo \(p^{a+k}\). Multiplication by an
integer prime to \(p\) is invertible, because its inverses modulo every
\(p^k\) are compatible. The kernel of reduction modulo \(n\) in (6) is
therefore exactly \(n\widehat{\mathbb Z}\), and reduction is surjective
by the Chinese remainder theorem. This proves (9). In particular, this
kernel is compact and open and has index \(n\).

Identify \(\widehat{\mathbb Z}\) with the compact open subring
\(\prod_p\mathbb Z_p\) of \(\mathbb A_f\). Given \(x\in\mathbb A_f\),
only finitely many coordinates have negative valuation. Choose
\[n=\prod_p p^{\max(0,-v_p(x_p))},\] where zero coordinates contribute
exponent zero. Then \(nx\in\widehat{\mathbb Z}\). Thus the union in (7)
is all of \(\mathbb A_f\).

Every tensor can be put over a common denominator, hence written as
\((1/n)\otimes z\). The map \(q\otimes z\mapsto(qz_p)_p\) is surjective
by the preceding paragraph. It is injective: if the image of
\((1/n)\otimes z\) is zero, every component of \(z\) is zero. This
identifies the tensor product algebraically with the localization of
\(\widehat{\mathbb Z}\) at the nonzero integers.

Give each stage \(n^{-1}\widehat{\mathbb Z}\) the compact topology
transported by \(z\mapsto z/n\). Give their union the \textbf{stage
topology}: a subset is open precisely when its intersection with each
stage is open there. In \(\mathbb A_f\) the stage is
\[\prod_{p\mid n}p^{-v_p(n)}\mathbb Z_p
\times\prod_{p\nmid n}\mathbb Z_p,\] so it is compact and open with
exactly that transported topology. Since these open stages cover
\(\mathbb A_f\), their stage topology is its topology. This proves the
topological assertion in (7); the algebraic tensor product alone does
not specify a topology.

The diagonal map \(\mathbb Q\to\mathbb A_f/\widehat{\mathbb Z}\) has
kernel \(\mathbb Z\): a rational number integral at every prime has no
prime in its reduced denominator. To prove surjectivity, write
\(x=z/n\). By (9), choose \(a\in\mathbb Z\) with
\(z-a\in n\widehat{\mathbb Z}\). Then \(x-a/n\in\widehat{\mathbb Z}\).
Hence the first isomorphism in (8) follows. The quotient is discrete
because \(\widehat{\mathbb Z}\) is open; give \(\mathbb Q/\mathbb Z\)
its discrete topology.

Coordinate residues give a surjection
\[\mathbb A_f\longrightarrow\bigoplus_p\mathbb Q_p/\mathbb Z_p\] with
kernel \(\widehat{\mathbb Z}\): every finite collection of local cosets
is represented by a tuple with zero in all other coordinates. This
proves the middle identification, with the direct sum also discrete. In
particular \(\mathbb Q_p/\mathbb Z_p\simeq\mathbb Z[1/p]/\mathbb Z\).
For surjectivity of this local map, write \(x=p^{-k}z\) and choose an
integer with the same residue as \(z\) modulo \(p^k\); its kernel
consists of the ordinary integers.

Finally, duality for surjective inverse systems of compact abelian
groups gives \[\widehat{\widehat{\mathbb Z}}
\simeq\varinjlim_n\widehat{\mathbb Z/n\mathbb Z}.\] This is the
compact-limit statement in \emph{Totally disconnected groups, the p-adic
numbers and the adèles}, Proposition 4(b), also {[}Halter-Koch, Theorem
1.7.6.A{]}. A character of \(\mathbb Z/n\mathbb Z\) is determined by an
\(n\)-th root of unity, so its dual identifies with
\(n^{-1}\mathbb Z/\mathbb Z\). For \(n\mid m\), pullback sends \(a/n\)
to \((ma/n)/m\), the same rational class. The limit is
\(\mathbb Q/\mathbb Z\). Explicitly the pairing is
\[\langle z,a/n+\mathbb Z\rangle
=\exp\!\left(2\pi i\frac{a z_n}{n}\right),
\tag{10}\] where \(z_n\) is any integer representative of the residue of
\(z\) modulo \(n\). Compatibility makes it independent of all
representatives. The dual of a compact group is discrete, so this is a
topological isomorphism. \(\square\)

\subsubsection{Worked example: one rational
representative}\label{NT-ADL-01--worked-example-one-rational-representative}

Take \(x_2=1/8\), \(x_3=2/9\), and \(x_p=0\) at all other primes. Then
\(72x\in\widehat{\mathbb Z}\). Its residues modulo \(8\) and \(9\) are
respectively \(1\) and \(7\). The integer \(25\) has both residues, so
\[x+\widehat{\mathbb Z}=25/72+\widehat{\mathbb Z}.\] Indeed, the
difference at \(2\) is \(-2/9\), and the difference at \(3\) is
\(-1/8\), both integral in their respective fields. At every other prime
the denominator \(72\) is a unit. This illustrates how independent local
residues combine into one rational class.

The stages also have a useful countable form:
\[\mathbb A_f=\bigcup_{j\geq1}\frac1{j!}\widehat{\mathbb Z}.
\tag{11}\] They increase because \(j!\mid(j+1)!\), and every denominator
divides some factorial. They are compact open \textbf{additive
subgroups}. For \(n>1\), \(n^{-1}\widehat{\mathbb Z}\) is not a subring:
if \(p\mid n\), the square of the diagonal element \(1/n\) has valuation
\(-2v_p(n)<-v_p(n)\) at \(p\), so it leaves that stage.

\subsection{Finite idèles and the normalization
table}\label{NT-ADL-01--finite-iduxe8les-and-the-normalization-table}

An element \(x\in\mathbb A_f\) is a ring unit exactly when every
\(x_p\neq0\) and \(x_p\in\mathbb Z_p^\times\) for almost all \(p\).
Necessity follows because both \(x\) and its inverse must be integral at
almost all primes; sufficiency follows by taking the componentwise
inverse. Thus, algebraically,
\[\mathbb A_f^\times=\prod_p'(\mathbb Q_p^\times,\mathbb Z_p^\times).\]
We give this \textbf{finite idèle group} the restricted product
topology. Its compact open subgroup is \(\widehat{\mathbb Z}^\times\).
This differs from its additive subspace topology; Exercise 2 identifies
the correct topology through the graph of inversion.

For example, the tuple \((p)_p\) is an adèle with no zero coordinate,
yet is not a unit: its coordinatewise inverse fails integrality at every
prime. In contrast, a nonzero rational number gives a finite idèle,
because its numerator and denominator involve only finitely many primes.
The valuation map identifies
\[\mathbb A_f^\times/\widehat{\mathbb Z}^\times\simeq\bigoplus_p\mathbb Z.\]
It is surjective by taking \(x_p=p^{k_p}\) for a finitely supported
integer tuple \((k_p)\), and its kernel is precisely the unit subgroup.
The quotient is discrete because that subgroup is open.

All measures in this lesson use the following normalizations.

{\def\LTcaptype{none} % do not increment counter
\begin{longtable}[]{@{}
  >{\raggedright\arraybackslash}p{(\linewidth - 6\tabcolsep) * \real{0.2500}}
  >{\raggedright\arraybackslash}p{(\linewidth - 6\tabcolsep) * \real{0.2500}}
  >{\raggedleft\arraybackslash}p{(\linewidth - 6\tabcolsep) * \real{0.2500}}
  >{\raggedright\arraybackslash}p{(\linewidth - 6\tabcolsep) * \real{0.2500}}@{}}
\toprule\noalign{}
\begin{minipage}[b]{\linewidth}\raggedright
Group
\end{minipage} & \begin{minipage}[b]{\linewidth}\raggedright
Distinguished compact open subgroup
\end{minipage} & \begin{minipage}[b]{\linewidth}\raggedleft
Its measure
\end{minipage} & \begin{minipage}[b]{\linewidth}\raggedright
Useful consequence
\end{minipage} \\
\midrule\noalign{}
\endhead
\bottomrule\noalign{}
\endlastfoot
\((\mathbb Q_p,+)\) & \(\mathbb Z_p\) & \(1\) &
\(\mu_p(p^k\mathbb Z_p)=p^{-k}\) \\
\(\mathbb Q_p^\times\) & \(\mathbb Z_p^\times\) & \(1\) &
\(\mu_p^\times(p^k\mathbb Z_p^\times)=1\) \\
\((\mathbb A_f,+)\) & \(\widehat{\mathbb Z}\) & \(1\) &
\(\mu(n^{-1}\widehat{\mathbb Z})=n\) \\
\(\mathbb A_f^\times\) & \(\widehat{\mathbb Z}^\times\) & \(1\) &
multiplicative translates have measure \(1\) \\
\end{longtable}
}

Here \(k\in\mathbb Z\) and \(n\geq1\). The additive local formula
follows by partitioning \(\mathbb Z_p\) into its \(p^k\) residue cosets
for \(k\geq0\); negative \(k\) follow by partitioning the larger ball
into cosets of \(\mathbb Z_p\). The multiplicative formula is Haar
invariance. Equation (9) shows that \(n^{-1}\widehat{\mathbb Z}\) has
\(n\) additive cosets of \(\widehat{\mathbb Z}\), proving its measure
formula. Proposition 1.2 gives both global normalizations.

\subsubsection{Worked example: a restricted
box}\label{NT-ADL-01--worked-example-a-restricted-box}

The compact open set \[B=2^{-2}\mathbb Z_2\times(1+9\mathbb Z_3)
\times\prod_{p\neq2,3}\mathbb Z_p\] has additive measure \(4/9\). The
first factor contributes \(4\), the second contributes \(1/9\), and
every other factor contributes one. Translating the centre of a local
ball leaves its additive measure unchanged. This example also explains
why the tail factors in (2) must have measure one.

\subsection{Test functions made from local
pieces}\label{NT-ADL-01--test-functions-made-from-local-pieces}

On \(\mathbb Q_p\), define a Bruhat--Schwartz function to be a locally
constant, compactly supported complex function. Write this space as
\(\mathcal S(\mathbb Q_p)\). On \(\mathbb A_f\) use the same definition.
Then \[\mathcal S(\mathbb A_f)
=\operatorname{span}\left\{\bigotimes_p f_p:
f_p\in\mathcal S(\mathbb Q_p),\quad
f_p=1_{\mathbb Z_p}\text{ for almost all }p\right\}.
\tag{12}\]

\textbf{Proof of (12).} A tensor on the right is locally constant. Its
support is a finite product of compact local supports with the compact
integral tail, so it is compact. Finite sums retain both properties.

Conversely, let \(f\) be locally constant with compact support. At each
point of its support, choose a compact open restricted rectangle on
which \(f\) is constant. Such rectangles exist because the local balls
\(a+p^k\mathbb Z_p\) form a neighbourhood basis. The rectangles can lie
inside \(\{f\neq0\}\): this set equals the support, since local
constancy makes its complement open. Compactness gives a finite cover by
these rectangles.

Include all their exceptional coordinates in one finite set \(S\). They
now have the same tail \(\prod_{p\notin S}\mathbb Z_p\), and each local
factor in \(S\) is compact and open. At each of these finitely many
primes, partition the union of the local factors into the finitely many
disjoint sets obtained by their intersections and differences. These
sets are compact and open. The resulting finite rectangular partition
makes \(f\) constant on each cell, and zero on cells outside its
support. The indicator of every cell is a tensor of local compact open
indicators with integral tail. Expressing \(f\) as its constant value
times these indicators proves (12). \(\square\)

The same proof works for restricted products of groups whose local
topologies have bases of compact open sets. More generally, given local
test spaces \(\mathcal S(G_v)\) containing \(1_{K_v}\) outside \(E\),
define their restricted tensor space as the span of \(\bigotimes_v f_v\)
with \(f_v\in\mathcal S(G_v)\) and \(f_v=1_{K_v}\) for almost all \(v\).
This definition specifies the local choices; it does not identify the
resulting space with all of \(C_c(G)\).

For arbitrary LCA factors, the intrinsic construction of local
Bruhat--Schwartz spaces uses their structure theory and is beyond this
lesson. Formula (12) proves the full identification needed for finite
adèles. With a Euclidean factor, the adelic test space is defined using
finite sums of products of an ordinary Schwartz function on that factor
and functions in (12).

For example, \[f=1_B-\frac49\,1_{\widehat{\mathbb Z}}\] belongs to
\(\mathcal S(\mathbb A_f)\) and has integral zero. Both terms are
factorizable, and their integrals are the measures already computed.
Such finite local descriptions are what make adelic integration
practical.

\subsection{Exercises}\label{NT-ADL-01--exercises}

\textbf{Exercise 1 (easy).} Prove that the restricted product topology
equals the subspace topology from \(\prod_vG_v\) if and only if
\(K_v=G_v\) for almost all \(v\notin E\).

\textbf{Solution.} If equality holds outside a finite set, include that
set in \(S\). Then \(G\) is the full product, and (1) is its product
topology. Conversely, suppose infinitely many \(K_v\) are proper. The
subgroup \(G_E\) is open in the restricted topology. Any product
neighbourhood of the identity constrains only a finite set \(F\) of
coordinates. Choose \(v\notin E\cup F\) with \(K_v\neq G_v\), and an
element \(g_v\notin K_v\). The tuple with that single nonidentity
coordinate belongs to \(G\) and to the chosen product neighbourhood but
not to \(G_E\). Thus \(G_E\) is not open in the subspace topology, and
the two topologies differ.

\textbf{Exercise 2 (medium).} Show that inversion on
\(\mathbb A_f^\times\), with the subspace topology from \(\mathbb A_f\),
is discontinuous. Show also that \[j:x\longmapsto(x,x^{-1})\] identifies
the restricted product topology with the topology induced from
\(\mathbb A_f\times\mathbb A_f\).

\textbf{Solution.} Enumerate the primes increasingly as \(q_r\). Let
\(u^{(r)}\) have coordinate \(q_r\) at \(q_r\), and coordinate one
elsewhere. Each tuple is an idèle. In an additive neighbourhood of one,
only finitely many coordinates have extra local conditions; outside
them, membership in \(\mathbb Z_p\) suffices. Consequently
\(u^{(r)}\to1\) in \(\mathbb A_f\). But
\((u^{(r)})^{-1}\notin\widehat{\mathbb Z}\) for every \(r\). Since
\(\widehat{\mathbb Z}\) is an additive neighbourhood of one, these
inverses do not converge to one. This proves discontinuity.

Take a basic product neighbourhood of \((x,x^{-1})\) in
\(\mathbb A_f\times\mathbb A_f\), and enlarge its two finite indexing
sets to one \(S\) containing the exceptional coordinates of \(x\) and
\(x^{-1}\). Its inverse image under \(j\) has finite-coordinate
conditions \[y_p\in U_p\cap V_p^{-1}\qquad(p\in S),\] which are open in
\(\mathbb Q_p^\times\). Outside \(S\), both \(y_p\) and \(y_p^{-1}\)
must be integral, equivalently \(y_p\in\mathbb Z_p^\times\). The inverse
image is therefore open in the restricted topology.

Conversely, let a restricted basic neighbourhood of \(x\) prescribe
\(y_p\in W_p\subseteq\mathbb Q_p^\times\) for \(p\in S\), and local
units elsewhere. Each \(W_p\) is open in \(\mathbb Q_p\), because
\(\mathbb Q_p^\times\) is open there. Use the additive basic sets with
factors \(W_p\) and \(\mathbb Q_p\), respectively, at \(p\in S\), and
integral tails for both. Their product pulls back under \(j\) to
precisely that restricted neighbourhood. Hence both induced topologies
agree, and injectivity of \(j\) proves the claimed identification.

\textbf{Exercise 3 (medium).} Compute the multiplicative measure and the
additive measure of \(\widehat{\mathbb Z}^\times\), with the
normalizations in the table.

\textbf{Solution.} Its multiplicative measure is one, by Proposition
1.2. Additively,
\(\mathbb Z_p^\times=\mathbb Z_p\setminus p\mathbb Z_p\) has measure
\(1-1/p\). Let \[C_N=\prod_{p\leq N}\mathbb Z_p^\times
\times\prod_{p>N}\mathbb Z_p.\] These compact sets decrease to
\(\widehat{\mathbb Z}^\times\), within the set \(\widehat{\mathbb Z}\)
of finite additive measure one. Continuity from above and the product
formula give \[\mu(\widehat{\mathbb Z}^\times)
=\lim_{N\to\infty}\prod_{p\leq N}(1-1/p)=0.\] For completeness, the
limit is zero without any asymptotic theorem about primes. Multiplying
the finitely many geometric series gives \[\prod_{p\leq N}(1-1/p)^{-1}
=\sum_{\substack{m\geq1\\\text{every prime divisor of }m\leq N}}\frac1m
\geq\sum_{m=1}^{\lfloor N\rfloor}\frac1m.\] The harmonic sums diverge:
each block \(2^{k-1}<m\leq2^k\) contributes at least \(1/2\). This
proves the zero limit. The different answers arise from two different
Haar measures on two different groups.

\textbf{Exercise 4 (hard).} Reconstruct the proof of Theorem 1.3 using
an open annihilator to establish the topology. Treat the compact
integral tail as a whole, rather than assuming a continuous character is
a finite product in advance.

\textbf{Solution.} For a tuple \((\chi_v)\) in the restricted dual
product, the formula \(\chi(x)=\prod_v\chi_v(x_v)\) has only finitely
many nontrivial factors. It is multiplicative and continuous on every
chart containing the tuple\textquotesingle s exceptional coordinates,
hence on \(G\).

For a continuous \(\chi\) on \(G\), choose \(B(S;U)\) mapped into
\(\{\operatorname{Re}z>0\}\). Its tail subgroup \(T_S\) has trivial
image by the circle argument. The coordinate restrictions consequently
annihilate \(K_v\) outside \(S\). Splitting any element into its
finitely many exceptional coordinates and a tail element recovers the
product formula and proves bijectivity.

For continuity of the product map, put a given compact set \(C\) inside
\(G_S\). Within the dual chart \(H_S\), the tail contributes one on
\(C\). Uniformly small evaluations on the finitely many compact
projections \(C_v\) make the product uniformly small, using the product
inequality in the theorem\textquotesingle s proof.

For continuity in the other direction, the condition of belonging to
\(H_S\) is exactly membership in \(T_S^\perp\). The latter is open in
\(\widehat G\), because its defining triviality condition can be tested
by uniform containment of the image of the compact subgroup \(T_S\) in
the same circle arc. On that open subgroup, any finite-coordinate
neighbourhood is obtained by the continuous restriction maps
\(\widehat G\to\widehat G_v\). Thus every restricted basic neighbourhood
has open image. The two continuity statements prove the topological
isomorphism, including for uncountable index sets and with unrestricted
factors in \(E\).

\subsection{What this lesson does not
prove}\label{NT-ADL-01--what-this-lesson-does-not-prove}

The following are the background results used here.

\begin{itemize}
\tightlist
\item
  The compact-product theorem and the fact that a continuous bijection
  from a compact space to a Hausdorff space is a homeomorphism, from
  point-set topology. The compact-product theorem is used in {[}Deitmar,
  §5.1{]}. Bézout\textquotesingle s identity for coprime integers is
  assumed; the Chinese remainder argument was given above.
\item
  Riesz representation for positive functionals on \(C_c(X)\); existence
  and uniqueness of Haar measure on every locally compact Hausdorff
  group; finite Radon products and their integrable-function Fubini
  theorem. These are {[}Haar measure on locally compact groups, Theorem
  2.2, Definition 4.3, Theorem 5.1, Theorems 8.3 and 9.2, and
  Proposition 12.1{]}. No global \(\sigma\)-finiteness is assumed.
\item
  The dual of an LCA group is LCA; compact groups have discrete duals
  and discrete groups have compact duals. See {[}Halter-Koch, Definition
  and Theorem 1.6.3{]}. Pontryagin duality identifies \(H\) with
  \(\widehat{\widehat H}\) by evaluation, as in \emph{The Pontryagin
  duality theorem} and {[}Halter-Koch, Theorem 1.6.6{]}.
\item
  For closed \(K\leq H\), pullback identifies \(\widehat{H/K}\) with
  \(K^\perp\), and restriction identifies \(\widehat H/K^\perp\) with
  \(\widehat K\), topologically. See \emph{Subgroups, quotients and
  annihilators} and {[}Halter-Koch, Theorems 1.6.5 and 1.6.7{]}.
  Together with compact-discrete duality, these imply that the
  annihilator of a compact open subgroup is compact and open.
\item
  For a surjective inverse system of compact abelian groups, every
  character of the limit factors through a coordinate group, and its
  discrete dual is the direct limit of the coordinate duals under
  pullback. See \emph{Totally disconnected groups, the p-adic numbers
  and the adèles}, Proposition 4(b), and {[}Halter-Koch, Theorem
  1.7.6.A{]}. We use this statement for finite cyclic groups, with
  explicit pullback maps.
\item
  The construction of \(\mathbb Q_p\), its valuation, and the compact
  open valuation ring \(\mathbb Z_p\), whose topology has the balls used
  above. See {[}Halter-Koch, §5.1, items 7 and the valuation-ring
  discussion following it{]}. The elementary residue and unit
  calculations needed for (6)--(9) were proved here.
\end{itemize}

The self-duality of \(\mathbb Q_p\) and of the adèles, the diagonal
lattice in the full adèle ring, and Poisson summation are treated in
\emph{The adèle ring of a number field} and \emph{Additive characters,
self-dual measures and Poisson summation on the adèles}.

\subsection{References}\label{NT-ADL-01--references}

\begin{itemize}
\tightlist
\item
  {[}Getz--Hahn{]} J. R. Getz and H. Hahn, \emph{An Introduction to
  Automorphic Representations: With a View toward Trace Formulae},
  \href{https://sites.duke.edu/jgetz/files/2022/04/Graduate_Text.pdf}{author-hosted
  draft dated 22 April 2022}, §2.3; published as Graduate Texts in
  Mathematics 300, Springer, 2024.
  \href{https://sites.duke.edu/jgetz/graduate-text/}{The
  authors\textquotesingle{} book page} links the published edition and
  errata.
\item
  {[}Bost--Connes{]} J.-B. Bost and A. Connes, \emph{Hecke algebras,
  type III factors and phase transitions with spontaneous symmetry
  breaking in number theory}, Selecta Mathematica (N.S.) 1 (1995),
  411--457, §3, definition of finite adèles.
  \href{https://repo-archives.ihes.fr/FONDS_IHES/I_Prepublications/CONNES/1994-1998/M_95_38/M_95_38_web.pdf}{IHÉS
  preprint}.
\item
  {[}Deitmar{]} A. Deitmar, \emph{Automorphic Forms}, Universitext,
  Springer, 2013, §5.1, especially Lemma 5.1.2.
  \href{https://doi.org/10.1007/978-1-4471-4435-9}{Published edition}.
\item
  {[}Halter-Koch{]} F. Halter-Koch, \emph{Class Field Theory and L
  Functions: Foundations and Main Results}, CRC Press, 2022, §§1.6--1.7
  and §5.2.
  \href{https://www.routledge.com/Class-Field-Theory-and-L-Functions-Foundations-and-Main-Results/Halter-Koch/p/book/9781032202655}{Published
  edition}.
\end{itemize}

\section{The adèle ring of a number field}\label{NT-ADL-02}

\emph{Written by GPT-6.1 Sol (OpenAI) in Codex, Ultra setting, October
2026. Self-checked by its author. Public domain (CC0).}

A number field can be studied in every completion at once. The adèle
ring makes this possible while retaining local compactness. Its topology
expresses an arithmetic distinction: a number field is discrete when all
places are present, but dense after any one place is removed. We will
prove both assertions, construct fundamental domains, calculate their
volume, and obtain simultaneous bounds on a nonzero field element.

We assume the results of \hyperref[NT-ADL-01]{Restricted products and
profinite completions}, \emph{Places of number fields in extensions and
the product formula}, \emph{Local fields: classification and Haar
measure}, and \emph{Lattices, Minkowski\textquotesingle s theorem and
the Minkowski embedding}. In particular, we use integral bases, the
Chinese remainder theorem, the local decomposition of a finite field
extension, and the Minkowski embedding. The exact facts used without
proof are collected near the end. Basic references are {[}Getz{]},
{[}Deitmar 2013{]}, {[}Getz--Hahn draft{]}, {[}Neukirch 1999{]},
{[}Harari 2020{]}, and {[}Halter--Koch 2022{]}.

\subsection{Local information with an integral
tail}\label{NT-ADL-02--local-information-with-an-integral-tail}

Let \(K\) be a number field of degree \(n\), with \(r_1\) real places
and \(r_2\) complex places, so \(n=r_1+2r_2\). At each place \(v\), let
\(K_v\) be its completion. If \(v\) is finite, write \(\mathcal O_v\)
for its valuation ring, \(\mathfrak p_v\) for the corresponding prime
ideal of \(\mathcal O_K\), and \(q_v=\#(\mathcal O_K/\mathfrak p_v)\).
Normalize \(\operatorname{ord}_v\) to have value group \(\mathbf Z\).

Our multiplicative local sizes are

\[|x|_v=
\begin{cases}
q_v^{-\operatorname{ord}_v(x)},&v\text{ finite},\\
|x|,&K_v=\mathbf R,\\
|x|^2,&K_v=\mathbf C.
\end{cases}\]

The last expression is the square of the ordinary complex modulus. It is
useful for the product formula and measure scaling, but does not satisfy
the ordinary triangle inequality. Whenever we use a geometric radius in
\(\mathbf C\), we mean the ordinary modulus. With these conventions,
\(\prod_v|x|_v=1\) for \(x\in K^\times\).

An \textbf{adèle} is a tuple \(a=(a_v)_v\), with \(a_v\in K_v\), such
that \(a_v\in\mathcal O_v\) at all but finitely many finite places. Thus

\[\mathbb A_K=\prod_v'(K_v,\mathcal O_v)
=K_\infty\times\mathbb A_{K,f},
\qquad K_\infty=\prod_{v\mid\infty}K_v
\simeq\mathbf R^{r_1}\times\mathbf C^{r_2}.\]

There is no integral restriction at an infinite place. Addition and
multiplication are coordinatewise. For a finite set \(S\) containing the
infinite places, put

\[\mathbb A_{K,S}=\prod_{v\in S}K_v\times
\prod_{v\notin S}\mathcal O_v.\]

These are open subrings, each with its product topology, and their union
is \(\mathbb A_K\). A basic open set is

\[\prod_{v\in S}U_v\times\prod_{v\notin S}\mathcal O_v,
\tag{1}\]

where every \(U_v\subset K_v\) is open. The requirement on the entire
tail matters: this topology is finer than the topology inherited from
the unrestricted product. An adèle may fail to be integral at finitely
many places, but a neighborhood controls where new failures can occur.

Every element of \(K\) defines a diagonal adèle. Indeed, a nonzero field
element has nonzero valuations at only finitely many finite places. We
identify \(K\) with this diagonal subfield. For later use set

\[\widehat{\mathcal O}_K=\prod_{v<\infty}\mathcal O_v.\]

This is a compact open subring of \(\mathbb A_{K,f}\). It is not open in
the full adèle ring when embedded with zero infinite component.

\subsection{Changing the number
field}\label{NT-ADL-02--changing-the-number-field}

\textbf{Proposition 2.1.} The ring \(\mathbb A_K\) is a Hausdorff
locally compact topological ring. If \(L/K\) is finite, the canonical
map

\[\Phi:L\otimes_K\mathbb A_K\longrightarrow\mathbb A_L,
\qquad
\ell\otimes(a_v)_v\longmapsto(\ell a_v)_{w\mid v}
\tag{2}\]

is an isomorphism of topological rings. On the tensor product, use the
finite product topology from any \(K\)-basis of \(L\). In particular,

\[\mathbb A_K\simeq K\otimes_{\mathbf Q}\mathbb A_{\mathbf Q},
\qquad
\mathbb A_K\simeq\mathbb A_{\mathbf Q}^{\,n}
\tag{3}\]

as topological rings in the first assertion and as topological
\(\mathbb A_{\mathbf Q}\)-modules in the second.

\textbf{Proof.} Each stage \(\mathbb A_{K,S}\) is a product of finitely
many locally compact fields and a compact product of valuation rings. It
is locally compact and Hausdorff. The stages are open and cover the
ring, proving these assertions for the underlying space.

For continuity of addition or multiplication at \((a,b)\), enlarge \(S\)
until both adèles belong to \(\mathbb A_{K,S}\). Operations on that
stage are continuous coordinatewise. Their values stay in the integral
tail because each \(\mathcal O_v\) is a subring. Continuity on these
open neighborhoods proves continuity on the whole ring.

We use the local field-extension identity

\[L\otimes_K K_v\simeq\prod_{w\mid v}L_w.
\tag{4}\]

Its map is the local version of (2), and is a homeomorphism of
finite-dimensional \(K_v\)-vector spaces. To pass to restricted products
we must check the integral lattices.

Here is the completion calculation. For a rational prime \(p\), finite
freeness over \(\mathbf Z\), followed by the Chinese remainder theorem,
gives

\[\begin{aligned}
\mathcal O_K\otimes_{\mathbf Z}\mathbf Z_p
&\simeq\varprojlim_m\mathcal O_K/p^m\mathcal O_K\\
&\simeq\prod_{v\mid p}\varprojlim_m
\mathcal O_K/\mathfrak p_v^{\,m e(v/p)}
\simeq\prod_{v\mid p}\mathcal O_v.
\end{aligned}
\tag{5}\]

The powers \(m e(v/p)\) form a cofinal sequence of powers of
\(\mathfrak p_v\). Localizing before completing does not change these
quotients: every denominator outside \(\mathfrak p_v\) is invertible
modulo \(\mathfrak p_v^k\).

For the relative version, localize \(\mathcal O_L\) at \(\mathfrak p_v\)
as a module over \((\mathcal O_K)_{\mathfrak p_v}\). It is finite and
torsion-free over this discrete valuation ring, hence free. Tensoring
this finite free module with the completed ring equals the inverse limit
of its quotients by \(\mathfrak p_v^m\). Decomposing
\(\mathfrak p_v\mathcal O_L\) into prime powers and applying the Chinese
remainder theorem in this inverse limit gives

\[\mathcal O_L\otimes_{\mathcal O_K}\mathcal O_v
\simeq\prod_{w\mid v}\mathcal O_w
\tag{6}\]

at every finite place. Ramification causes no exception to this identity
of completed integer rings.

Choose a \(K\)-basis \(e_1,\ldots,e_m\) of \(L\) consisting of algebraic
integers. The lattice \(M=\sum_j\mathcal O_K e_j\) lies in
\(\mathcal O_L\), and some nonzero \(d\in\mathcal O_K\) satisfies
\(d\mathcal O_L\subset M\). Indeed, express a finite set of
\(\mathcal O_K\)-module generators of \(\mathcal O_L\) in this basis and
clear denominators. Thus
\(M\otimes\mathcal O_v=\mathcal O_L\otimes\mathcal O_v\) outside the
finitely many primes dividing \(d\).

In these coordinates the local map \(K_v^m\to\prod_{w\mid v}L_w\) takes
\(\mathcal O_v^m\) onto \(\prod_{w\mid v}\mathcal O_w\) at almost every
finite place. The local maps and their inverses therefore take
restricted tuples to restricted tuples. They are homeomorphisms on the
open product stages after adjoining the exceptional places, and assemble
to a global homeomorphism. Formula (2) respects multiplication, so it is
a ring isomorphism.

Changes of \(K\)-basis act by invertible matrices over \(K\) on
\(\mathbb A_K^m\); the matrix and its inverse define continuous maps.
The tensor-product topology is independent of the basis. Taking the base
field to be \(\mathbf Q\) proves (3). The module identification depends
on a basis; the ring map does not. \(\square\)

\subsection{Representatives for the additive
quotient}\label{NT-ADL-02--representatives-for-the-additive-quotient}

Begin with \(\mathbf Q\), and write
\(\widehat{\mathbf Z}=\prod_p\mathbf Z_p\). A rational number belongs to
every \(\mathbf Z_p\) exactly when its reduced denominator has no prime
divisor. Hence

\[\mathbf Q\cap(\mathbf R\times\widehat{\mathbf Z})=\mathbf Z.
\tag{7}\]

Every finite rational adèle can be made integral by subtracting a
rational number. Take its negative \(p\)-adic digits at each of its
finitely many nonintegral places and add the resulting rational
principal parts. A principal part at \(p\) has denominator a power of
\(p\), so is integral at every other prime. Thus

\[\mathbb A_{\mathbf Q}
=\mathbf Q+(\mathbf R\times\widehat{\mathbf Z}).
\tag{8}\]

After this subtraction, subtract an integer to move the real coordinate
into \([0,1)\). The integer subtraction also changes every finite
coordinate; it preserves their integrality.

\textbf{Theorem 2.2.} The diagonal subgroup \(K\) is discrete and closed
in \(\mathbb A_K\), and \(\mathbb A_K/K\) is compact. Every coset in
\(\mathbb A_{\mathbf Q}/\mathbf Q\) has exactly one representative in

\[F_{\mathbf Q}=[0,1)\times\widehat{\mathbf Z}.
\tag{9}\]

\textbf{Proof.} In the rational case, the open neighborhood

\[(-1/2,1/2)\times\widehat{\mathbf Z}\]

meets \(\mathbf Q\) only at zero, by (7). This proves discreteness. A
discrete subgroup of a Hausdorff topological group is closed: choose an
open neighborhood \(V\) of zero whose difference set meets the subgroup
only at zero. Any translate of \(V\) contains at most one subgroup
element. A point in the closure must equal that element, since otherwise
a smaller neighborhood excluding it contradicts the definition of
closure.

Existence of representatives in (9) follows from (8). If two
representatives differ by \(q\in\mathbf Q\), their finite components
imply \(q\in\mathbf Z\). Their real components differ by a number
strictly between \(-1\) and \(1\), so \(q=0\). This proves uniqueness.

The domain (9) is half-open and is not compact. Its compact closure
\([0,1]\times\widehat{\mathbf Z}\) still maps onto the quotient. Thus
the quotient is compact.

Choose a \(\mathbf Q\)-basis of \(K\). Under (3), the diagonal \(K\)
becomes the subgroup \(\mathbf Q^n\subset\mathbb A_{\mathbf Q}^n\). It
is discrete and closed, and its quotient is the finite product of \(n\)
compact rational quotients. This proves the assertions for \(K\).
\(\square\)

For volume calculations we want a domain adapted to the integer rings.
Choose an integral basis \(\omega_1,\ldots,\omega_n\) of
\(\mathcal O_K\), let \(j:K\to K_\infty\) be the Minkowski embedding,
and put

\[P_K=\left\{\sum_{i=1}^n t_i j(\omega_i):0\leq t_i<1\right\}.\]

Equations (3) and (5), applied to this basis and the rational
principal-part construction, give

\[\mathbb A_K=K+(K_\infty\times\widehat{\mathcal O}_K),
\qquad
K\cap(K_\infty\times\widehat{\mathcal O}_K)=\mathcal O_K.
\tag{10}\]

For the intersection, write an element of \(K\) in the integral basis.
By (5), integrality at all finite places says its rational coordinates
belong to every \(\mathbf Z_p\), so they are integers. Reducing the
infinite component modulo \(j(\mathcal O_K)\) now shows that

\[F_K=P_K\times\widehat{\mathcal O}_K
\tag{11}\]

is a Borel fundamental domain with unique representatives. Its closure
is compact. If \(U\subset K_\infty\) is bounded, then

\[K\cap(U\times\widehat{\mathcal O}_K)\]

is finite: it corresponds to the points of the full lattice
\(j(\mathcal O_K)\) in a bounded set.

\textbf{Worked example: rational principal parts.} Take an adèle with
real component \(13/10\), components \(3/8\) at \(2\) and \(7/25\) at
\(5\), and zero at all other primes. Subtract

\[q=3/8+7/25=131/200.\]

At \(2\) the remaining value is \(-7/25\), which is integral; at \(5\)
it is \(-3/8\), also integral. At every other prime the denominator of
\(q\) is a unit. The remaining real component is \(129/200\in[0,1)\), so
this already gives the representative in (9).

\textbf{Worked example: the solenoid quotient.} Inclusion induces an
isomorphism of topological groups

\[(\mathbf R\times\widehat{\mathbf Z})/\mathbf Z
\longrightarrow\mathbb A_{\mathbf Q}/\mathbf Q.
\tag{12}\]

Here \(\mathbf Z\) is embedded diagonally: \(m\mapsto(m,(m)_p)\).
Surjectivity and the kernel follow from (7)--(8). The source is compact,
since \([0,1]\times\widehat{\mathbf Z}\) maps onto it. The target is
Hausdorff by Theorem 2.2. Thus the continuous bijection is a
homeomorphism. This quotient is called the rational solenoid; its
further topological structure is studied in the lesson on solenoids.
Formula (9) supplies measurable representatives, not a continuous
product decomposition: crossing a real endpoint also changes the finite
coordinate by a diagonal integer.

\subsection{How the discriminant measures the
quotient}\label{NT-ADL-02--how-the-discriminant-measures-the-quotient}

Fix additive Haar measures \(dx_v\) by

\[dx_v=dx\quad(K_v=\mathbf R),\qquad
dx_v=2\,dx\,dy=|dz\wedge d\bar z|
\quad(K_v=\mathbf C),\qquad
\mu_v(\mathcal O_v)=1\quad(v<\infty).
\tag{13}\]

The restricted product measure \(dx=\prod_v dx_v\) is characterized on
product sets with integral tails by the product of local measures.
Multiplication by \(a\in K_v^\times\) scales \(dx_v\) by \(|a|_v\). At a
finite place this follows by decomposing \(\mathcal O_v\) into \(q_v\)
translates of its maximal ideal; at a complex place it follows from the
real determinant \(|a|^2\).

Let \(d_K\) be the field discriminant. We give the discrete additive
subgroup \(K\) counting measure. The quotient measure is induced from
\(dx\) with this choice; its total mass is the measure of a Borel
fundamental domain such as (11).

\textbf{Proposition 2.4.} With (13),

\[\operatorname{vol}(\mathbb A_K/K)=\sqrt{|d_K|}.
\tag{14}\]

\textbf{Proof.} In ordinary real coordinates on \(K_\infty\), form the
matrix \(M\) whose columns are \(j(\omega_i)\), recording real and
imaginary parts at each complex place. The Euclidean volume of \(P_K\)
is \(|\det M|\).

Let \(B=(\tau(\omega_i))_{\tau,i}\) be the matrix of all \(n\) complex
embeddings, including both members of each conjugate pair. The trace
matrix is \(B^{\mathsf T}B\), so

\[d_K=\det(B^{\mathsf T}B)=(\det B)^2.\]

Replacing a real-coordinate pair \((u,v)\) by \((u+iv,u-iv)\) has
determinant \(-2i\). Consequently

\[|\det B|=2^{r_2}|\det M|,
\qquad
|\det M|=2^{-r_2}\sqrt{|d_K|}.
\tag{15}\]

Our infinite measure multiplies Euclidean measure by \(2^{r_2}\). Thus
\(\mu_\infty(P_K)=\sqrt{|d_K|}\). The finite factor in (11) has measure
\(1\), proving (14). \(\square\)

Using \(dx\,dy\) at complex places would instead give
\(2^{-r_2}\sqrt{|d_K|}\). A change of integral basis changes \(M\) by a
matrix in \(\mathrm{GL}_n(\mathbf Z)\), whose determinant has absolute
value \(1\), so does not change the answer.

A different finite normalization is used with Fourier transforms. For
the standard trace character
\(\psi_v=\psi_{\mathbf Q_p}\circ\operatorname{Tr}_{K_v/\mathbf Q_p}\),
let \(\mathfrak D_v^{-1}\) be the trace-dual lattice of
\(\mathcal O_v\). Its definition identifies it with the annihilator of
\(\mathcal O_v\). If a self-dual measure gives \(\mathcal O_v\) mass
\(c\), Fourier inversion on its indicator gives

\[1=c^2N(\mathfrak D_v),\qquad
c=N(\mathfrak D_v)^{-1/2}.\]

Thus the self-dual finite measure is \(N(\mathfrak D_v)^{-1/2}dx_v\). At
unramified places the different is the unit ideal, so the measures agree
{[}Neukirch 1999, Chapter III, §2, Theorem (2.6){]}. The infinite
factors in (13) are self-dual for the standard trace characters. Fourier
inversion and its convention are used here from {[}Getz--Hahn draft,
Appendix B, §B.2, equation (B.3){]}; \emph{Additive characters,
self-dual measures and Poisson summation on the adèles} develops that
theory.

\subsection{A box large enough to contain a field
element}\label{NT-ADL-02--a-box-large-enough-to-contain-a-field-element}

An \textbf{idèle} is a tuple \(\alpha=(\alpha_v)\) with
\(\alpha_v\in K_v^\times\) and \(\alpha_v\in\mathcal O_v^\times\) at
almost every finite place. These are precisely the units of
\(\mathbb A_K\): both a tuple and its coordinatewise inverse must have
integral tails. Define

\[\|\alpha\|_{\mathbb A}=\prod_v|\alpha_v|_v.\]

Only finitely many factors differ from \(1\). We need the following
bound before proving strong approximation.

\textbf{Lemma 2.5 (adelic Minkowski lemma).} One may take

\[C_K=\left(\frac2\pi\right)^{r_2}\sqrt{|d_K|}.
\tag{16}\]

For every idèle \(\alpha\) with \(\|\alpha\|_{\mathbb A}\geq C_K\),
there exists \(x\in K^\times\) such that

\[|x|_v\leq|\alpha_v|_v\qquad\text{at every place }v.
\tag{17}\]

\textbf{Proof.} First suppose \(\|\alpha\|_{\mathbb A}>C_K\). Establish
the counting principle behind Blichfeldt\textquotesingle s argument. If
a measurable set \(E\subset\mathbb A_K\) satisfies \(\mu(E)>\mu(F_K)\),
two distinct points of \(E\) differ by a nonzero element of \(K\).
Indeed, the translates \(F_K+k\), for \(k\in K\), partition the ring.
Translation invariance and Tonelli\textquotesingle s theorem give

\[\int_{F_K}\sum_{k\in K}\mathbf1_E(f+k)\,df=\mu(E).
\tag{18}\]

The group \(K\) is countable. If every coset met \(E\) at most once, the
sum would be at most \(1\), contradicting the strict volume inequality.
A coset meeting \(E\) twice gives the asserted difference.

Let \(R_v\) be the ordinary real or complex modulus of \(\alpha_v\) at
an infinite place. Choose

\[E=\prod_{v\text{ real}}[-R_v/2,R_v/2]
\times\prod_{v\text{ complex}}\{z:|z|\leq R_v/2\}
\times\prod_{v<\infty}\alpha_v\mathcal O_v.
\tag{19}\]

This is a compact measurable subset of an open product stage. At real
places its factor has measure \(R_v=|\alpha_v|_v\). At complex places
its measure is

\[2\pi(R_v/2)^2=(\pi/2)|\alpha_v|_v.\]

At finite places its measure is \(|\alpha_v|_v\). Hence

\[\mu(E)=(\pi/2)^{r_2}\|\alpha\|_{\mathbb A}
>\sqrt{|d_K|}=\mu(F_K).\]

Choose distinct \(y,z\in E\) with \(x=y-z\in K^\times\). At a real place
their difference has modulus at most \(R_v\). At a complex place use the
ordinary triangle inequality to obtain \(|y_v-z_v|\leq R_v\), then
square to obtain (17). At a finite place \(\alpha_v\mathcal O_v\) is an
additive subgroup, so the difference stays in it.

Now suppose \(\|\alpha\|_{\mathbb A}=C_K\). Choose an infinite place
\(w\), and for \(j\geq1\) replace \(\alpha_w\) by \((1+1/j)\alpha_w\),
leaving every other component fixed. The resulting idèle has norm
\((1+1/j)^{d_w}C_K>C_K\), where \(d_w=1\) at a real place and \(d_w=2\)
at a complex place. The strict case gives a nonzero \(x_j\in K\) with
the corresponding bounds. All these elements lie in one fixed compact
box: double only the ordinary radius at \(w\), retain the other infinite
radii, and retain every finite subgroup \(\alpha_v\mathcal O_v\). By
Theorem 2.2, this box meets \(K\) in a finite set. Some nonzero \(x\)
therefore occurs for infinitely many \(j\). Its bounds at \(v\ne w\)
already have the required size. Along that unbounded subsequence, its
bound at \(w\) tends to \(|\alpha_w|_w\). Thus (17) also holds at
equality. \(\square\)

The half-radius choice is essential. Applying the counting principle
directly to the full box would only bound the difference by twice its
archimedean radii. Halving affects a complex factor by \(1/4\), because
it has two real dimensions. We need not halve the finite factors, which
are additive subgroups.

Compactness and discreteness extend the strict volume argument to the
closed boundary. The product formula explains why a size hypothesis is
needed: (17) for nonzero \(x\) implies \(1\leq\|\alpha\|_{\mathbb A}\).

\subsection{Removing a place frees
approximation}\label{NT-ADL-02--removing-a-place-frees-approximation}

For a fixed place \(v_0\), write

\[\mathbb A_K^{(v_0)}
=\prod_{v\ne v_0}'(K_v,\mathcal O_v),\]

imposing integral restrictions only at finite places. Projection from
the full ring forgets the \(v_0\)-coordinate.

\textbf{Theorem 2.3 (strong approximation).} For every finite or
infinite place \(v_0\), the diagonal image of \(K\) is dense in
\(\mathbb A_K^{(v_0)}\).

Equivalently, let \(S\) be a finite set disjoint from \(\{v_0\}\),
containing all infinite places other than \(v_0\). Given \(a_v\in K_v\)
and \(\varepsilon_v>0\) for \(v\in S\), there exists \(q\in K\) such
that

\[|q-a_v|_v<\varepsilon_v\quad(v\in S),
\qquad
q\in\mathcal O_v\quad(v<\infty,\ v\notin S\cup\{v_0\}).
\tag{20}\]

\textbf{Proof.} Let \(D=\overline{P_K}\times\widehat{\mathcal O}_K\). By
(11), every adèle is \(k+d\) for \(k\in K\) and \(d\in D\). For each
infinite place set

\[B_v=\max\left(1,\sup_{d\in D}|d_v|_v\right),\]

and for finite places set \(B_v=1\). These are bounds on a fixed compact
set, not norms on \(K_\infty\).

Choose an idèle \(\beta\) with

\[|\beta_v|_v B_v<\varepsilon_v\quad(v\in S),
\qquad
|\beta_v|_v=1\quad(v<\infty,\ v\notin S\cup\{v_0\}).\]

At \(v_0\) choose its component sufficiently large that
\(\|\beta\|_{\mathbb A}>C_K\). Every completion has elements of
arbitrarily large size, so this is possible whether \(v_0\) is finite or
infinite. Lemma 2.5 supplies nonzero \(c\in K\) with
\(|c|_v\leq|\beta_v|_v\) everywhere.

Let \(t\in\mathbb A_K\) have components \(a_v\) for \(v\in S\) and zero
elsewhere, including at \(v_0\). Apply the representative property to
\(c^{-1}t\). There are \(k\in K\) and \(d\in D\) with

\[c^{-1}t=k+d.\]

Set \(q=ck\). Then \(t-q=cd\), so for \(v\in S\),

\[|a_v-q|_v\leq|\beta_v|_v B_v<\varepsilon_v.\]

At a finite place outside \(S\cup\{v_0\}\), both \(c\) and \(d_v\) are
integral and \(t_v=0\), so \(q=-cd_v\) is integral. This proves (20).

A neighborhood of a truncated adèle contains a set as in (20): include
in \(S\) its finitely many nonintegral places, all retained infinite
places, and the places with specified local neighborhoods. Choose
smaller local balls inside those neighborhoods. Thus (20) proves
density. \(\square\)

The omitted coordinate pays for the simultaneous small bounds on the
other coordinates: it allows an arbitrarily large component before
invoking Lemma 2.5. The proof uses neither a unit theorem nor
multiplicative strong approximation.

\textbf{Worked example: rational finite approximation.} Suppose we want

\[q-1/4\in\mathbf Z_2,\qquad q-2/9\in\mathbf Z_3,\]

and integrality at every other finite prime. Set \(q=m/36\). The
congruences become

\[m\equiv1\pmod4,\qquad m\equiv8\pmod9.\]

Choose \(m=17\). Indeed, \(17/36-1/4=2/9\in\mathbf Z_2\), while
\(17/36-2/9=1/4\in\mathbf Z_3\). Its denominator has no other prime
factors. Higher precision gives higher prime-power congruences, solved
by the same Chinese remainder theorem. In particular, \(\mathbf Q\) is
dense in \(\mathbb A_{\mathbf Q,f}\), the additive assertion used in
{[}Bost--Connes 1995, §3, proof of Proposition 12{]}.

This does not give density in the full ring. If a rational \(q\) is
integral at every finite prime, it is an integer. No such \(q\) lies in
\((1/3,2/3)\). Thus \((1/3,2/3)\times\widehat{\mathbf Z}\) contains no
rational number.

\subsection{Exercises}\label{NT-ADL-02--exercises}

\textbf{Exercise 1 (easy).} Prove (7) and (8) directly, without using
Theorem 2.2. Explain why subtracting an arbitrary real number cannot
replace rational subtraction in (8).

\textbf{Exercise 2 (medium).} Exhibit a Borel fundamental domain for
\(\mathbf Q(i)\) in \(\mathbb A_{\mathbf Q(i)}\). Compute its volume
with (13), including the discriminant.

\textbf{Exercise 3 (medium).} Prove strong approximation for
\(\mathbf Q\) directly. Treat both omission of infinity and omission of
a finite prime. In the latter case include approximation at the real
place.

\textbf{Exercise 4 (hard).} Prove Proposition 2.4 from an integral
basis. Explain exactly where the factor \(2^{r_2}\) enters, and why an
arbitrary vector-space identification does not preserve the measure
without a normalization check.

\subsection{Solutions}\label{NT-ADL-02--solutions}

\textbf{Solution 1.} Write \(q=a/b\) with coprime integers \(a,b\) and
\(b>0\). If \(b>1\), some prime \(p\mid b\) has
\(\operatorname{ord}_p(q)<0\), so \(q\notin\mathbf Z_p\). Conversely,
every integer belongs to every \(\mathbf Z_p\). This proves (7).

For \(a=(a_\infty,(a_p))\), let \(T\) be its nonintegral finite places.
Choose \(D=\prod_{p\in T}p^{d_p}\) with \(D a_p\in\mathbf Z_p\) for
\(p\in T\). The Chinese remainder theorem gives an integer \(m\)
satisfying

\[m\equiv D a_p\pmod{p^{d_p}}\qquad(p\in T).\]

Each right side denotes an integral residue class. Then \(q=m/D\)
satisfies \(a_p-q\in\mathbf Z_p\) at primes in \(T\). At other primes
\(D\) is a unit and both \(a_p,q\) are integral. Thus
\(a-q\in\mathbf R\times\widehat{\mathbf Z}\), proving (8). For empty
\(T\), take \(D=1,m=0\). A real number only supplies an infinite
coordinate; it has no specified images in finite completions. Rational
numbers supply the diagonal translation needed here.

\textbf{Solution 2.} First \(\mathcal O_{\mathbf Q(i)}=\mathbf Z[i]\).
If \(a+bi\), with \(a,b\in\mathbf Q\), is integral, traces of it and its
product with \(i\) show \(2a,2b\in\mathbf Z\). Its norm \(a^2+b^2\) is
integral. Writing \(2a=m,2b=n\), the condition \(m^2+n^2\equiv0\pmod4\)
forces both \(m,n\) even. Hence \(a,b\in\mathbf Z\).

With basis \(1,i\), formula (11) gives

\[\{x+iy:0\leq x<1,\ 0\leq y<1\}
\times\prod_{v<\infty}\mathcal O_v.\]

By (10), subtracting a field element makes the finite components
integral. Subtracting a Gaussian integer then moves both real
coordinates into the stated intervals. Two such representatives differ
by a Gaussian integer whose real and imaginary parts lie strictly
between \(-1\) and \(1\), so they agree. The complex square has ordinary
area \(1\), hence measure \(2\); the finite product has measure \(1\).
Its volume is \(2\). Finally,

\[d_{\mathbf Q(i)}=
\det\begin{pmatrix}
\operatorname{Tr}(1)&\operatorname{Tr}(i)\\
\operatorname{Tr}(i)&\operatorname{Tr}(i^2)
\end{pmatrix}
=\det\begin{pmatrix}2&0\\0&-2\end{pmatrix}=-4.\]

The volume is \(\sqrt{|-4|}\), as required.

\textbf{Solution 3.} Finite local conditions may be written
\(q-a_p\in p^{r_p}\mathbf Z_p\), with \(r_p\in\mathbf Z\). Let \(T\) be
the finite set of constrained primes, including all nonintegral target
components. Choose \(d_p\geq0\) so \(D a_p\in\mathbf Z_p\) and
\(d_p+r_p\geq0\), where \(D=\prod_{p\in T}p^{d_p}\).

If the omitted place is infinity, choose

\[m\equiv D a_p\pmod{p^{d_p+r_p}}\qquad(p\in T)\]

and put \(q=m/D\). This solves all conditions and gives integrality at
every other finite prime.

Now omit a finite prime \(\ell\), so \(\ell\notin T\). For every
\(N\geq0\), solve

\[m\equiv D\ell^N a_p\pmod{p^{d_p+r_p}}\qquad(p\in T).\]

All solutions form \(m_N+M\mathbf Z\), where
\(M=\prod_{p\in T}p^{d_p+r_p}\). The rationals

\[q=\frac{m_N+Mk}{D\ell^N}\qquad(k\in\mathbf Z)\]

satisfy the finite conditions. In \(\mathbf R\) they form an arithmetic
progression with mesh \(M/(D\ell^N)\), tending to zero. Given
\(a_\infty\) and \(\varepsilon>0\), choose \(N\) with mesh less than
\(2\varepsilon\), then choose the nearest progression point. Its
distance is at most half the mesh, hence less than \(\varepsilon\). Its
denominator is supported on \(T\cup\{\ell\}\), so it is integral outside
that set. This proves density for every omitted place. For empty \(T\),
use \(D=M=1\).

\textbf{Solution 4.} The Minkowski lattice theorem gives a
parallelepiped for \(j(\mathcal O_K)\) of Euclidean volume
\(2^{-r_2}\sqrt{|d_K|}\). To verify the convention, group the embeddings
into real embeddings and complex conjugate pairs. Replacing real and
imaginary coordinates by a conjugate pair multiplies the
determinant\textquotesingle s absolute value by \(2\) per pair. Thus the
embedding determinant has absolute value \(2^{r_2}\) times the
real-coordinate determinant. Its square is the discriminant, giving the
stated Euclidean volume.

By (10), the parallelepiped times \(\widehat{\mathcal O}_K\) meets every
additive coset exactly once. Counting measure on the diagonal subgroup
identifies its measure with the total quotient mass. Each complex
measure in (13) is twice ordinary area, so the infinite product measure
multiplies Euclidean volume by \(2^{r_2}\). The finite factor has mass
\(1\). The result is

\[2^{r_2}\cdot2^{-r_2}\sqrt{|d_K|}\cdot1=\sqrt{|d_K|}.\]

An arbitrary \(\mathbf Q\)-basis gives a topological module isomorphism,
but its infinite determinant need not have absolute value \(1\), and its
finite maps need not preserve the integer lattices. Haar measures
transform by these local determinants. A topological isomorphism alone
does not preserve their chosen numerical normalization.

\subsection{What this lesson does not
prove}\label{NT-ADL-02--what-this-lesson-does-not-prove}

We use the following prerequisite results with their stated conventions.

\begin{itemize}
\tightlist
\item
  A number field has an integral basis; the Minkowski image of
  \(\mathcal O_K\) is a full lattice. For the canonical infinite
  measure, a nonzero integral ideal \(\mathfrak a\) has lattice covolume
  \(\sqrt{|d_K|}[\mathcal O_K:\mathfrak a]\). {[}Neukirch 1999, Chapter
  I, §5, Propositions (5.1)--(5.2){]}. The determinant calculation
  needed for (14) is given above.
\item
  Each finite completion is locally compact, its valuation ring is
  compact open, and that ring is the inverse limit of its quotients by
  powers of its maximal ideal. {[}Neukirch 1999, Chapter II, §5,
  Proposition (5.1) and its proof{]}.
\item
  For finite separable \(L/K\), the canonical map
  \(L\otimes_K K_v\to\prod_{w\mid v}L_w\) is an isomorphism. Finite
  extensions of number fields are separable. {[}Neukirch 1999, Chapter
  II, §8, Proposition (8.3){]}.
\item
  With real modulus, squared complex modulus, and finite values
  \(q_v^{-\operatorname{ord}_v}\), the product formula holds for
  \(K^\times\). {[}Neukirch 1999, Chapter III, §1, Proposition (1.3);
  Getz--Hahn draft, Proposition 2.1.1{]}.
\item
  The Chinese remainder theorem solves simultaneous conditions for
  distinct prime powers. Its rational approximation application is in
  {[}Deitmar 2013, §5.2, proof of Theorem 5.2.1(c){]}. We also use the
  elementary ideal version of the Chinese remainder theorem in (5).
\item
  Haar measures exist on locally compact additive groups. The restricted
  product measure has the product formula on integral-tail product sets.
  {[}Deitmar 2013, Theorem 5.2.2{]}. The same product construction
  applies to \(K\).
\item
  For a fixed nontrivial local character and perfect pairing, a unique
  self-dual Haar measure satisfies Fourier inversion. The standard
  characteristic-zero characters use local traces. {[}Getz--Hahn draft,
  Appendix B, §§B.1--B.2, equation (B.3){]}. Only the normalization
  comparison after Proposition 2.4 uses this fact.
\item
  An unramified finite local extension has unit different. {[}Neukirch
  1999, Chapter III, §2, Theorem (2.6) and its reduction to
  completions{]}.
\end{itemize}

The solenoid\textquotesingle s classification and duality,
multiplicative idèle class groups, and Poisson summation belong to
subsequent lessons. The quotient proved compact here is the additive
quotient \(\mathbb A_K/K\).

\subsection{References}\label{NT-ADL-02--references}

\begin{itemize}
\tightlist
\item
  \textbf{{[}Getz{]}} J. R. Getz, \emph{An Introduction to Automorphic
  Representations}, open course notes, §1.1.
  \href{https://sites.math.duke.edu/~jgetz/aut_reps.pdf}{Author\textquotesingle s
  notes}.
\item
  \textbf{{[}Bost--Connes 1995{]}} J.-B. Bost and A. Connes, \emph{Hecke
  algebras, type III factors and phase transitions with spontaneous
  symmetry breaking in number theory}, Selecta Mathematica (N.S.) 1
  (1995), 411--457, §3.
  \href{https://doi.org/10.1007/BF01589495}{Publication}.
\item
  \textbf{{[}Deitmar 2013{]}} A. Deitmar, \emph{Automorphic Forms},
  Universitext, Springer, 2013, §5.2.
\item
  \textbf{{[}Getz--Hahn draft{]}} J. R. Getz and H. Hahn, \emph{An
  Introduction to Automorphic Representations: With a View toward Trace
  Formulae},
  \href{https://sites.duke.edu/jgetz/files/2022/04/Graduate_Text.pdf}{author-hosted
  draft dated April 22, 2022}, §§2.1, 2.5--2.6, §3.5, and Appendix B.
  The published book appeared as Graduate Texts in Mathematics 300,
  Springer, 2024.
\item
  \textbf{{[}Neukirch 1999{]}} J. Neukirch, \emph{Algebraic Number
  Theory}, translated by N. Schappacher, Grundlehren der mathematischen
  Wissenschaften 322, Springer, 1999, at the propositions cited above.
\item
  \textbf{{[}Harari 2020{]}} D. Harari, \emph{Galois Cohomology and
  Class Field Theory}, Universitext, Springer, 2020, §12.3.
\item
  \textbf{{[}Halter--Koch 2022{]}} F. Halter-Koch, \emph{Class Field
  Theory and L Functions: Foundations and Main Results}, CRC Press,
  2022, §5.3.
\end{itemize}

\section{Idèles and the idèle class group}\label{NT-ADL-03}

\emph{Written by GPT-6.1 Sol (OpenAI) in Codex, Ultra setting, October
2026. Self-checked by its author. Public domain (CC0).}

An adèle can be integral at every finite place and still fail to have an
adelic inverse. For example, the family whose component at each rational
prime \(p\) is \(p\), and whose real component is \(1\), lies in
\(\mathbb A_{\mathbb Q}\). Its coordinatewise inverse is not integral at
any finite place, so it is not an adèle. The invertible adèles must have
unit components almost everywhere. Their topology must also control the
inverse.

The resulting group connects two kinds of information. Valuations record
fractional ideals. Archimedean absolute values record the logarithms of
units. After quotienting by multiplication by a field element, the part
with total absolute value one is compact. We will prove this compactness
from the adelic box lemma and then obtain both the finiteness of the
ideal class group and Dirichlet\textquotesingle s unit theorem. Neither
of those two conclusions is assumed in the compactness proof.

Throughout, \(K\) is a number field, with \(r_1\) real places and
\(r_2\) complex places. We use the normalized absolute values

\[|x|_v=\begin{cases}
|x|,&v\text{ real},\\
|x|^2,&v\text{ complex},\\
q_v^{-\operatorname{ord}_v(x)},&v\text{ finite},
\end{cases}
\qquad q_v=\#(\mathcal O_v/\mathfrak p_v).
\tag{1}\]

The product formula says \(\prod_v|a|_v=1\) for \(a\in K^\times\). At
complex places, geometric radii and triangle inequalities always use the
ordinary modulus; the square in (1) is a multiplicative normalization.

We use three results from \hyperref[NT-ADL-02]{The adèle ring of a
number field}: \(K\) is discrete and closed in \(\mathbb A_K\) (Theorem
2.2); finite-place integer rings are compact; and if an idèle \(\beta\)
satisfies

\[\prod_v|\beta_v|_v>C_K,
\qquad C_K=(2/\pi)^{r_2}\sqrt{|d_K|},
\tag{2}\]

then there is a nonzero \(a\in K\) with \(|a|_v\leq |\beta_v|_v\) at
every place (Lemma 2.5). That lemma was proved by additive quotient
volume and a difference of two points in an adelic box. We also use the
ordinary ideal theory of the Dedekind domain \(\mathcal O_K\):
fractional ideals have unique finite prime factorization, and the
exponent in \((a)\) is \(\operatorname{ord}_v(a)\). No class-number or
unit-rank assertion enters these prerequisites.

\subsection{A topology that controls multiplication and
inversion}\label{NT-ADL-03--a-topology-that-controls-multiplication-and-inversion}

Define

\[J_K=\mathbb A_K^\times
=\left\{(x_v):x_v\in K_v^\times\text{ for all }v,
\ x_v\in\mathcal O_v^\times\text{ for almost all finite }v\right\}.
\tag{3}\]

The equality follows directly from the definition of the adèle ring. If
both \(x\) and \(x^{-1}\) are adèles, then outside a finite set both
components are integral, which forces \(\operatorname{ord}_v(x_v)=0\).
Conversely, the condition in (3) makes the coordinatewise inverse an
adèle.

For a finite set \(S\) containing all archimedean places, put

\[J_{K,S}=\prod_{v\in S}K_v^\times
\times\prod_{v\notin S}\mathcal O_v^\times.
\tag{4}\]

Equip \(J_K\) with the restricted product topology: the subgroups (4)
are open and have their product topologies. A basic neighbourhood of
\(x\), after \(S\) contains its nonunit components, has the form

\[\prod_{v\in S}W_v\times\prod_{v\notin S}\mathcal O_v^\times,
\qquad x_v\in W_v\subset K_v^\times\text{ open}.
\tag{5}\]

\textbf{Proposition 3.1.} The restricted product topology makes \(J_K\)
a locally compact Hausdorff abelian group. It is also the topology
transported by

\[x\longmapsto (x,x^{-1})
\tag{6}\]

from the subspace \(\{(x,y)\in\mathbb A_K^2:xy=1\}\). The inclusion
\(J_K\to\mathbb A_K\) is continuous, but its image does not have the
idèle topology as its additive subspace topology.

\textbf{Proof.} Each \(\mathcal O_v^\times\) is compact and open in
\(K_v^\times\). Indeed,
\(\mathcal O_v^\times=\mathcal O_v\setminus\mathfrak p_v\), a closed
subset of the compact ring \(\mathcal O_v\), and it is open because
valuations are locally constant away from zero. The restricted product
theorem, Proposition 1.1 of \hyperref[NT-ADL-01]{Restricted products and
profinite completions}, therefore gives all the asserted group and local
compactness properties.

The coordinatewise inclusion into \(\mathbb A_K\) is continuous on each
open chart (4), hence is continuous globally. Inversion is continuous in
the restricted product, so (6) is continuous. It is a bijection onto the
displayed graph because every inverse in a ring is unique.

To check its inverse, take a basic neighbourhood (5). In the first adèle
coordinate impose \(x_v\in W_v\) for \(v\in S\) and
\(x_v\in\mathcal O_v\) outside \(S\). These are additive open
conditions: \(K_v^\times\) is open in \(K_v\), so \(W_v\) is also open
in \(K_v\). In the second adèle coordinate impose \(y_v\in\mathcal O_v\)
outside \(S\), with no restriction at \(S\). On the graph \(y=x^{-1}\),
the two tail conditions are exactly \(x_v\in\mathcal O_v^\times\). Their
intersection pulls back to (5). Thus every restricted product
neighbourhood is open in the graph topology. The graph is closed in
\(\mathbb A_K^2\), since multiplication is continuous and \(\{1\}\) is
closed.

For strictness of the topology comparison, choose pairwise distinct
finite places \(v_n\), and a uniformizer \(\varpi_n\) at \(v_n\). Let
\(x^{(n)}\) have component \(\varpi_n\) at \(v_n\) and component \(1\)
everywhere else. These are idèles. They tend to \(1\) in
\(\mathbb A_K\): any additive neighbourhood restricts only finitely many
places and permits integral components at all remaining finite places.
But none lies in the idèle neighbourhood

\[\prod_{v\mid\infty}K_v^\times\times\prod_{v\nmid\infty}\mathcal O_v^\times
\tag{7}\]

of \(1\). Moreover, \((x^{(n)})^{-1}\) never belongs to the additive
neighbourhood \(K_\infty\times\prod_{v\nmid\infty}\mathcal O_v\), where
\(K_\infty=\prod_{v\mid\infty}K_v\). Inversion is therefore
discontinuous at \(1\) for the additive subspace topology. \(\square\)

Equivalently, the idèle topology is the weakest topology containing the
additive subspace topology for which inversion is continuous. Adding the
inverse images of additive open sets gives precisely the graph topology;
multiplication is then continuous by Proposition 3.1.

\subsection{The norm-one constraint recovers the additive
topology}\label{NT-ADL-03--the-norm-one-constraint-recovers-the-additive-topology}

Only finitely many factors differ from \(1\) in

\[|x|=\prod_v|x_v|_v,
\qquad J_K^1=\{x\in J_K:|x|=1\}.
\tag{8}\]

This number is called the idèle norm, or content. It is not a
vector-space norm.

\textbf{Proposition 3.2.} The map \(|\cdot|:J_K\to\mathbb R_{>0}\) is a
continuous surjective homomorphism. The diagonal subgroup \(K^\times\)
is discrete and closed in \(J_K\) and lies in \(J_K^1\). The subset
\(J_K^1\) is closed in \(\mathbb A_K\), and its subspace topologies from
\(J_K\) and \(\mathbb A_K\) coincide.

\textbf{Proof.} On (4), the norm is the finite product
\(\prod_{v\in S}|x_v|_v\), hence is continuous. Each local absolute
value is multiplicative. Choose an archimedean place \(w\). Define a
homomorphism \(s:\mathbb R_{>0}\to J_K\) by setting all components
except \(w\) equal to \(1\), and setting

\[s(t)_w=\begin{cases}t,&w\text{ real},\\ \sqrt t,&w\text{ complex}.
\end{cases}
\tag{9}\]

It is continuous and \(|s(t)|=t\), proving surjectivity. The product
formula puts \(K^\times\) in the kernel. Since \(K\) is discrete in
\(\mathbb A_K\), some additive open neighbourhood \(U\) of \(1\) meets
\(K\) only at \(1\). Its inverse image in \(J_K\) isolates \(1\) in
\(K^\times\).

A discrete subgroup of a Hausdorff topological group is closed. Here is
the argument we will use. Choose an identity neighbourhood \(V\) with
\(VV^{-1}\) meeting the subgroup only at the identity. Each translate
\(gV\) contains at most one subgroup element. If \(g\) lies in the
closure of the subgroup, such a translate contains one element \(h\).
Were \(g\ne h\), intersecting with a neighbourhood of \(g\) that avoids
\(h\) would contradict closure. Thus every closure point belongs to the
subgroup.

For the topology assertion, fix \(x\in J_K^1\) and a basic idèle
neighbourhood (5), with \(S\) containing all its nonunit components and
all archimedean places. Shrink the finitely many \(W_v\), if necessary,
so that

\[\prod_{v\in S}|y_v|_v<2 \quad\text{whenever }y_v\in W_v.
\tag{10}\]

This is possible because the product at \(x\) equals \(1\). Permit
\(y_v\in\mathcal O_v\) outside \(S\); this gives an additive
neighbourhood \(W\) of \(x\). If \(y\in W\cap J_K^1\) had a nonunit
component at a finite place outside \(S\), that component would have
normalized absolute value at most \(q_v^{-1}\leq 1/2\). All other tail
factors are at most \(1\), so (10) would give \(|y|<1\), a
contradiction. Consequently \(W\cap J_K^1\) lies in (5). The opposite
topology inclusion follows from the continuous inclusion
\(J_K\to\mathbb A_K\).

It remains to establish additive closedness. We give all the cases,
since a closedness assertion in \(J_K\) alone would not justify
intersecting \(J_K^1\) with an additive compact box.

First take \(z\in\mathbb A_K\) with a zero component. Choose a finite
set \(S\) containing that place, every archimedean place and every
finite place where \(z\) is nonintegral. The finite product
\(\prod_{v\in S}|z_v|_v\) is zero. By continuity it is less than \(1\)
throughout some finite-coordinate neighbourhood. Together with integral
tails, this gives an additive neighbourhood containing no norm-one
idèle.

Next suppose every component is nonzero but \(z\) is not an idèle.
Infinitely many integral finite components are nonunits. Starting with
all archimedean and nonintegral places, add enough such nonunit places
to make the finite product less than \(1\); each added factor is at most
\(1/2\). The same neighbourhood argument excludes \(J_K^1\).

Finally suppose \(z\) is an idèle with norm \(c\ne1\). If \(c<1\), take
\(S\) containing all nonunit and archimedean places and again keep its
finite product less than \(1\). If \(c>1\), enlarge this \(S\) to
include all finite places with \(q_v\leq 2c\). There are finitely many:
such a place lies above a rational prime \(p\leq 2c\), and a number
field has finitely many places above each \(p\). Choose
finite-coordinate neighbourhoods so that

\[1<\prod_{v\in S}|y_v|_v<2c,
\qquad y_v\in\mathcal O_v\quad(v\notin S).
\tag{11}\]

If a member of this additive neighbourhood were a norm-one idèle, either
its tail would consist entirely of units, giving norm greater than
\(1\), or some tail component would have absolute value at most
\(q_v^{-1}<(2c)^{-1}\), giving norm less than \(1\). Both are
impossible. Every point outside \(J_K^1\) therefore has an additive
neighbourhood disjoint from it. \(\square\)

The mechanism in (10) is useful: a new tail valuation costs at least a
factor of two, and a sufficiently small change at finitely many places
cannot pay that cost while preserving norm one.

\subsection{Compactness before class numbers and
units}\label{NT-ADL-03--compactness-before-class-numbers-and-units}

The idèle class group and its norm-one part are

\[C_K=J_K/K^\times,
\qquad C_K^1=J_K^1/K^\times.
\tag{12}\]

They have quotient topologies. Quotient maps of topological groups are
open: the inverse image of the image of an open set is the union of its
subgroup translates. Since \(K^\times\) is closed, these quotients are
Hausdorff and locally compact. The restriction of the quotient map to
\(J_K^1\) realizes \(C_K^1\) as the kernel of the induced norm on
\(C_K\), with its subspace topology; this also follows explicitly from
the splitting proved in Proposition 3.5 below.

\textbf{Theorem 3.3 (norm-one compactness).} The group \(C_K^1\) is
compact.

\textbf{Proof.} Choose an idèle \(\eta\) with \(|\eta|>C_K\); the
section (9) supplies one. The additive box

\[B_\eta=\{z\in\mathbb A_K:|z_v|_v\leq |\eta_v|_v\text{ at every }v\}
\tag{13}\]

is compact. Each archimedean factor is a closed bounded interval or
disk; each finite factor is the compact subgroup \(\eta_v\mathcal O_v\);
almost all of the latter equal \(\mathcal O_v\). Thus the box is a
product of compact sets in one additive restricted product chart.

By Proposition 3.2, \(B_\eta\cap J_K^1\) is closed in this box and
compact in the idèle topology. For any \(x\in J_K^1\), apply the adelic
box lemma (2) to \(\eta x^{-1}\). Its norm is \(|\eta|>C_K\), so there
is \(a\in K^\times\) satisfying

\[|a|_v\leq |\eta_vx_v^{-1}|_v,
\qquad |ax_v|_v\leq|\eta_v|_v
\quad\text{for every }v.
\tag{14}\]

The product formula gives \(|ax|=1\). Thus \(ax\in B_\eta\cap J_K^1\),
and it represents the same class as \(x\). The quotient is the
continuous image of this compact set. \(\square\)

This proof uses no choice of ideal-class representatives and no
fundamental domain for a unit lattice. Those two objects will follow
from compactness.

\textbf{Corollary 3.4.} The ideal class group \(\operatorname{Cl}_K\) is
finite. The unit group has the form

\[\mathcal O_K^\times\simeq\mu_K\times\mathbb Z^{r_1+r_2-1},
\tag{15}\]

where \(\mu_K\) is the finite cyclic group of roots of unity in \(K\).
More precisely, the logarithms of units form a full lattice in

\[H=\left\{(t_v)_{v\mid\infty}\in\mathbb R^{r_1+r_2}:
\sum_{v\mid\infty}t_v=0\right\}.
\tag{16}\]

\textbf{Proof.} Let \(I_K\) be the group of nonzero fractional ideals,
with the discrete topology. The map

\[\mathfrak a:J_K\longrightarrow I_K,
\qquad x\longmapsto\prod_{v\nmid\infty}\mathfrak p_v^{\operatorname{ord}_v(x_v)}
\tag{17}\]

is a continuous surjective homomorphism. Its kernel is the open subgroup

\[J_{K,\infty}=K_\infty^\times\times U_f,
\qquad U_f=\prod_{v\nmid\infty}\mathcal O_v^\times.
\tag{18}\]

Surjectivity follows by placing uniformizer powers at the finitely many
primes in an ideal factorization. Continuity follows because the kernel
is open, or directly because valuation vectors are locally constant in
idèle charts. Multiplying a chosen lift by \(s(|x|^{-1})\) changes no
finite valuation and makes its norm one. Hence (17) is still surjective
on \(J_K^1\). It takes a principal idèle \(a\) to \((a)\), and induces a
continuous surjection

\[C_K^1\longrightarrow\operatorname{Cl}_K.
\tag{19}\]

A compact discrete space is finite, proving class-number finiteness.

Set \(J_{K,\infty}^1=J_{K,\infty}\cap J_K^1\). The intersection
\(J_{K,\infty}^1\cap K^\times\) is exactly \(\mathcal O_K^\times\): zero
finite valuations mean that both \(a\) and \(a^{-1}\) are algebraic
integers. Moreover,

\[J_{K,\infty}^1/\mathcal O_K^\times
\simeq \ker\bigl(C_K^1\to\operatorname{Cl}_K\bigr)
\tag{20}\]

as topological groups. To see surjectivity, if \(\mathfrak a(x)\) is
principal, divide \(x\) by a generator; its norm stays one and its
finite valuations become zero. The kernel is the indicated intersection.
For the topology, \(J_{K,\infty}^1\) is open in \(J_K^1\), and the
quotient map is open, so the induced bijection onto its image is open.
The right side of (20) is closed in the compact group \(C_K^1\), hence
the left side is compact.

On \(J_{K,\infty}^1\), define

\[\ell(x)=(\log|x_v|_v)_{v\mid\infty}\in H.
\tag{21}\]

This is a continuous surjection. For any \(t\in H\), take component
\(e^{t_v}\) at a real place and the positive real number \(e^{t_v/2}\)
at a complex place, and take all finite components equal to \(1\). These
components give a lift with norm one. Its kernel is the compact group

\[\{\pm1\}^{r_1}\times(S^1)^{r_2}\times U_f.
\tag{22}\]

Let \(\Gamma=\ell(\mathcal O_K^\times)\). We first prove that bounded
sets meet \(\Gamma\) in finitely many points. For a fixed \(M>0\), all
units whose logarithms satisfy \(|\log|a|_v|\leq M\) lie in the compact
idèle set consisting of finite units and archimedean annuli
\(e^{-M}\leq|x_v|_v\leq e^M\). Its intersection with the discrete closed
subgroup \(K^\times\) is finite. Consequently \(\Gamma\) is discrete and
closed, and the kernel of \(\ell\) on units is finite.

That finite kernel equals \(\mu_K\). Every element of a finite subgroup
has finite order, while every root of unity is an algebraic unit with
all normalized absolute values equal to one. It is cyclic: through any
archimedean embedding it is a finite subgroup of \(\mathbb C^\times\),
contained in the cyclic group of \(N\)-th roots of unity for a common
multiple \(N\) of its element orders.

The map (21) descends to a continuous surjection from the compact group
in (20) onto \(H/\Gamma\), so \(H/\Gamma\) is compact. We finish with
the following elementary lattice argument rather than assuming the unit
theorem.

If a subgroup \(\Gamma\subset\mathbb R^d\) meets bounded sets in
finitely many points, then it has a basis of at most \(d\) linearly
independent vectors over \(\mathbb Z\). For a nonzero group, choose
\(e\ne0\) generating its intersection with the line \(\mathbb Re\); one
can choose a shortest nonzero vector on any line meeting the group, and
subtract integer multiples to prove the cyclic assertion. Orthogonally
project onto \(e^\perp\). The projected group still meets bounded sets
in finitely many points: lift points in a bounded set and subtract
integer multiples of \(e\) so their \(e\)-coordinate lies in one fixed
interval of length \(\|e\|\). All resulting lifts lie in a bounded set,
so there are finitely many. Induction on dimension gives a basis of the
projected group; lifting that basis and adjoining \(e\) gives a basis of
\(\Gamma\). Independence follows by projecting a linear relation and
then using the nonzero vector \(e\).

If \(\mathbb R^d/\Gamma\) is compact, these basis vectors must span
\(\mathbb R^d\). Otherwise a nonzero linear functional annihilating
their span would descend to a continuous surjection
\(\mathbb R^d/\Gamma\to\mathbb R\), impossible for a compact source.
Thus \(\Gamma\) has rank \(d\). Applied to (16), this gives rank
\(r_1+r_2-1\). Choose lifts \(\varepsilon_1,\ldots,\varepsilon_d\) of a
lattice basis in the unit group. Every unit is uniquely a root of unity
times \(\varepsilon_1^{n_1}\cdots\varepsilon_d^{n_d}\), proving (15).
When \(d=0\), the lattice is zero and all units are roots of unity.
\(\square\)

The same framework contains the usual theorem on \(S\)-units. Here are
the details needed to keep track of the extra rank. Let
\(S=S_\infty\cup S_f\), with \(S_f\) a finite set of finite places, and
let

\[E_{K,S}=\{a\in K^\times:\operatorname{ord}_v(a)=0\text{ for finite }v\notin S\}.\]

The valuation map fits into

\[1\longrightarrow\mathcal O_K^\times\longrightarrow E_{K,S}
\longrightarrow\mathbb Z^{S_f}.
\tag{23}\]

Its image has finite index. Indeed, if \(h=\#\operatorname{Cl}_K\), then
\(\mathfrak p_v^h=(a_v)\) for each \(v\in S_f\), and \(a_v\in E_{K,S}\)
maps to \(h\) times the corresponding coordinate vector. Thus the image
contains \(h\mathbb Z^{S_f}\); the lattice argument above shows it is
free of rank \(\#S_f\). Lifting a basis splits (23) over its image and
gives

\[E_{K,S}\simeq\mu_K\times\mathbb Z^{r_1+r_2-1+\#S_f}.
\tag{24}\]

Its logarithms \((\log|a|_v)_{v\in S}\) form a full lattice in the
hyperplane \(\sum_{v\in S}t_v=0\). For a direct check, the finite
coordinates lie in the grid \(\prod_{v\in S_f}\log q_v\,\mathbb Z\), and
their image has finite index in that grid. After reducing these
coordinates by the chosen lifted basis, the remaining logarithms lie in
the archimedean unit lattice of (16). This proves discreteness and gives
a bounded parallelepiped covering the hyperplane modulo the logarithms.
The finite-coordinate grid spans \(\mathbb R^{S_f}\), while the
archimedean unit lattice spans its kernel, so the full span has
dimension \(\#S-1\). The logarithmic kernel is again \(\mu_K\).

\subsection{Separating the norm from the compact
group}\label{NT-ADL-03--separating-the-norm-from-the-compact-group}

\textbf{Proposition 3.5.} A choice of section (9) gives topological
group isomorphisms

\[J_K\simeq J_K^1\times\mathbb R_{>0},
\qquad C_K\simeq C_K^1\times\mathbb R_{>0}.
\tag{25}\]

\textbf{Proof.} The maps

\[(y,t)\longmapsto y\,s(t),
\qquad x\longmapsto\bigl(x\,s(|x|)^{-1},|x|\bigr)
\tag{26}\]

are continuous inverse homomorphisms. The second map sends \(K^\times\)
onto \(K^\times\times\{1\}\). Taking quotients yields (25), since the
product of a quotient map with the identity is open and surjective,
hence a quotient map. Explicitly the class inverse is
\([x]\mapsto([x\,s(|x|)^{-1}],|x|)\), so its continuity also follows
from the defining quotient topology. \(\square\)

The norm homomorphism is canonical; its splitting is a choice. For a
general number field the connected component of \(C_K\) can include
connected groups inside \(C_K^1\). Therefore (25) does not identify the
connected component with its \(\mathbb R_{>0}\) factor in general.

\subsection{Rational
representatives}\label{NT-ADL-03--rational-representatives}

For \(K=\mathbb Q\), write

\[\widehat{\mathbb Z}^{\times}=\prod_p\mathbb Z_p^\times.\]

We embed this group in \(J_{\mathbb Q}\) with real component \(1\),
embed \(\mathbb R_{>0}\) with every finite component \(1\), and embed
\(\mathbb Q^\times\) diagonally.

\textbf{Proposition 3.6.} These embeddings give internal direct products
of topological groups

\[J_{\mathbb Q}=\mathbb Q^\times\times\mathbb R_{>0}
\times\widehat{\mathbb Z}^{\times},
\qquad
\mathbb A_f^\times=\mathbb Q_{>0}^\times\times\widehat{\mathbb Z}^{\times}.
\tag{27}\]

Here the rational factors carry the discrete topology, and the finite
idèles carry their restricted product topology. Consequently

\[C_{\mathbb Q}\simeq\mathbb R_{>0}\times\widehat{\mathbb Z}^{\times},
\qquad C_{\mathbb Q}^1\simeq\widehat{\mathbb Z}^{\times}.
\tag{28}\]

The identity component of \(C_{\mathbb Q}\) is
\(\mathbb R_{>0}\times\{1\}\), and its quotient by that component is
\(\widehat{\mathbb Z}^{\times}\).

\textbf{Proof.} For \(x\in J_{\mathbb Q}\), set

\[q_0=\prod_p p^{\operatorname{ord}_p(x_p)}>0,
\quad q=\operatorname{sgn}(x_\infty)q_0,
\quad t=\frac{x_\infty}{q}>0,
\quad u_p=\frac{x_p}{q}.
\tag{29}\]

The product for \(q_0\) is finite, and
\(\operatorname{ord}_p(q)=\operatorname{ord}_p(x_p)\), so every \(u_p\)
is a unit. Thus \(x=qtu\). Finite valuations determine \(q_0\), and the
real sign determines \(q\); the other two factors then follow. This
proves existence and uniqueness.

The multiplication map from the product in (27) is continuous. Its
inverse is continuous too: near a fixed idèle, the finitely many
relevant finite valuations and the real sign can be kept constant in a
basic neighbourhood (5), so \(q\) is locally constant. Then
\(t=x_\infty/q\) and \(u=x/q\) with the real component removed vary
continuously. Also

\[|x|=|x_\infty|\prod_p p^{-\operatorname{ord}_p(x_p)}
=|x_\infty|/q_0=t.
\tag{30}\]

The same argument without the real sign gives the finite decomposition,
with positive rational factor \(q_0\). Its discreteness can also be seen
from \(\mathbb Q_{>0}^\times\cap\widehat{\mathbb Z}^{\times}=\{1\}\).
Quotienting the first decomposition by its diagonal rational factor
proves (28), and (30) identifies the norm-one factor.

Finally \(\widehat{\mathbb Z}^{\times}\) is totally disconnected.
Reduction modulo \(p^n\) is continuous into a finite discrete unit
group, so a connected subset has constant image in every such quotient.
These reductions separate points in \(\mathbb Z_p^\times\), and the
coordinate projections separate points in the product. Thus every
connected subset is a singleton. In contrast \(\mathbb R_{>0}\) is
connected. Projecting a connected subset of (28) to the second factor
proves the component assertion. \(\square\)

For an explicit representative, take an idèle with real component
\(-7/10\), component \(8\) at \(2\), component \(1/25\) at \(5\), and
component \(1\) at all other primes. Formula (29) gives

\[q_0=8/25,\qquad q=-8/25,\qquad t=35/16,\]

and

\[u_2=-25,\qquad u_5=-1/8,\qquad u_p=-25/8\ (p\ne2,5).\]

Every displayed \(u_p\) is a unit in its own local ring, and the idèle
class is represented by \((35/16,u)\). Its norm is \(35/16\), even
though its real component has absolute value \(7/10\).

For \(K=\mathbb Q(i)\), whose integer ring \(\mathbb Z[i]\) was
identified in Exercise 2 of \emph{The adèle ring of a number field}, the
corresponding ideal-class quotient is trivial:

\[C_K/\operatorname{image}\bigl(\mathbb C^\times\times U_f\bigr)
\simeq\operatorname{Cl}_K=1.
\tag{31}\]

Indeed, \(\mathbb Z[i]\) is Euclidean for \(N(a+bi)=a^2+b^2\). Given a
complex quotient, round each real coordinate to the nearest integer. The
error has squared modulus at most \(1/2<1\), so division by a nonzero
Gaussian integer has a remainder of smaller norm. A nonzero ideal
contains an element of least positive norm; division by that element
shows it generates the ideal. All fractional ideals are principal.
Formula (17) then proves (31). In this case every norm-one class has a
representative with all finite components units; its complex component
must have ordinary modulus one. Hence

\[C_{\mathbb Q(i)}^1\simeq (S^1\times U_f)/\mu_4,
\qquad \mu_4=\{1,-1,i,-i\},
\tag{32}\]

where \(\mu_4\) is embedded diagonally in all components. The unit
calculation follows either by solving \(a^2+b^2=1\), or from the
zero-rank case of (15). Unlike the rational compact factor, (32) visibly
has a connected circle coming from \(S^1\).

There is a different quotient, \(\mathbb A_K/K^\times\), where
multiplication acts on all adèles, including zero. It is not the idèle
class group and is not Hausdorff. For example, the orbit of the adèle
with component \(1\) at a chosen archimedean place and zero at every
other place is not closed: multiplication by positive rational numbers
tending to zero makes it converge to the zero adèle, which is not in
that orbit. If the quotient were Hausdorff, the inverse image of every
singleton would be closed. This observation explains why the topology
assertions above concern invertible adèles.

\subsection{Exercises}\label{NT-ADL-03--exercises}

\begin{enumerate}
\item
  \textbf{Easy.} Starting with the finite valuations and real sign of a
  rational idèle, construct its factors \(q,t,u\) in (27). Show that
  \(t\) is its idèle norm and that all three factors are unique. Carry
  out the construction when \(x_\infty=9\), \(x_3=1/9\), \(x_5=5\), and
  every other finite component is \(1\).
\item
  \textbf{Medium.} For an arbitrary number field, construct idèles
  tending additively to \(1\) whose inverses do not tend additively to
  \(1\). Explain separately why the same idèles fail to converge in the
  idèle topology.
\item
  \textbf{Medium.} Prove directly from the product formula, without
  additive discreteness as a black box, that \(K^\times\) is discrete
  and closed in \(J_K\). Include complex places with the normalization
  (1).
\item
  \textbf{Hard.} Recover Dirichlet\textquotesingle s unit theorem from
  compactness of \(C_K^1\). Identify the compact quotient used in the
  proof, show the logarithmic image is a lattice of full rank, and
  explain why the finite kernel is precisely \(\mu_K\). Then show how
  each finite place allowed in an \(S\)-unit group adds one to the rank.
\end{enumerate}

\subsection{Solutions}\label{NT-ADL-03--solutions}

\textbf{Solution 1.} Finite valuations determine
\(q_0=\prod p^{\operatorname{ord}_p(x_p)}\). Choose \(q\) to have
magnitude \(q_0\) and the sign of \(x_\infty\), put \(t=x_\infty/q\),
and divide each finite component by \(q\). Its valuation becomes zero at
every prime, giving \(u\). The normalized finite product is
\(q_0^{-1}\), so \(t=|x_\infty|/q_0=|x|\). Any proposed decomposition
has these finite valuations and real sign, so yields the same \(q\),
then the same \(t,u\). In the example, \(q=q_0=5/9\), \(t=81/5\),
\(u_3=1/5\), \(u_5=9\), and \(u_p=9/5\) at every other prime. These are
local units. Directly, the norm is \(9\cdot9\cdot(1/5)=81/5\).

\textbf{Solution 2.} Enumerate pairwise distinct finite places \(v_n\).
They exist because there are infinitely many rational primes and each
has at least one place above it. Take a uniformizer at \(v_n\), set that
component of \(x^{(n)}\) equal to it, and all others equal to \(1\). An
additive basic neighbourhood of \(1\) restricts finitely many
coordinates and has integral tails, so eventually contains \(x^{(n)}\).
Each inverse has a nonintegral component at \(v_n\), and therefore none
belongs to the additive open set
\(K_\infty\times\prod_{v\nmid\infty}\mathcal O_v\) containing \(1\).
Also none of the original idèles lies in (7), which is an idèle
neighbourhood. This proves both claims. The example does not contradict
Proposition 3.2: these idèles have norm \(q_{v_n}^{-1}\), not norm one.

\textbf{Solution 3.} At every real place require \(|x_v-1|<1\), at every
complex place require ordinary \(|x_v-1|<1\), and at every finite place
require \(x_v\in\mathcal O_v^\times\). This defines an idèle
neighbourhood \(U\) of \(1\). Suppose a principal idèle \(a\in U\) has
\(a\ne1\). At each finite place, the ultrametric inequality gives
\(|a-1|_v\leq1\). At each archimedean place, \(|a-1|_v<1\); squaring the
ordinary complex inequality preserves this strict inequality. There is
at least one archimedean place, so \(\prod_v|a-1|_v<1\), contradicting
the product formula for \(a-1\ne0\). Thus \(U\cap K^\times=\{1\}\),
proving discreteness. Choose \(V\) with \(VV^{-1}\subset U\). Each
translate \(gV\) contains at most one principal idèle; a point in the
closure must equal that element, because otherwise a smaller
neighbourhood could avoid it. Hence the subgroup is closed.

\textbf{Solution 4.} Use \(J_{K,\infty}^1\), whose diagonal intersection
is \(\mathcal O_K^\times\). The ideal map identifies its quotient by
units with the kernel of \(C_K^1\to\operatorname{Cl}_K\); the
identification is open, and the kernel is closed, so the quotient is
compact. The logarithm map onto the hyperplane (16) has compact kernel
(22). On units, bounded logarithms put the units in a compact idèle set.
Its intersection with discrete closed \(K^\times\) is finite. Thus the
logarithmic subgroup \(\Gamma\) is discrete and closed, with finite
kernel. The kernel is a finite group of units, hence consists of roots
of unity; every root of unity lies in it. A complex embedding proves its
cyclicity.

The compact unit quotient maps onto \(H/\Gamma\). Apply the lattice
induction in Corollary 3.4: project away from a one-dimensional cyclic
subgroup, reduce lifts into a bounded strip to preserve discreteness,
and lift the resulting basis. This gives at most \(\dim H\) independent
generators. A proper real span would allow a nonzero real linear
functional on \(H/\Gamma\) with unbounded image, contradicting
compactness. Thus the rank is exactly \(\dim H=r_1+r_2-1\). Lifting this
basis splits units as \(\mu_K\times\mathbb Z^{r_1+r_2-1}\).

For \(S\)-units, use the finite valuation vector at \(S_f\). Its kernel
is the ordinary unit group. Since the already proved class group has
order \(h\), principal generators of \(\mathfrak p_v^h\) give \(h\)
times each standard vector in the image. The image is therefore a
finite-index subgroup of \(\mathbb Z^{S_f}\), free of rank \(\#S_f\).
Lifting a basis splits off that many further infinite cyclic generators.
The rank is \(r_1+r_2-1+\#S_f\), with the same torsion group \(\mu_K\).

\subsection{What this lesson does not
prove}\label{NT-ADL-03--what-this-lesson-does-not-prove}

The restricted product theorem is Proposition 1.1 of \emph{Restricted
products and profinite completions}. The additive discreteness,
closedness and box lemma are Theorem 2.2 and Lemma 2.5 of \emph{The
adèle ring of a number field}. These are proved there; only their stated
number-field forms are used here.

Unique factorization of fractional ideals in a Dedekind domain and the
correspondence between their exponents and local valuations are
algebraic prerequisites. Precise references are Neukirch,
\emph{Algebraic Number Theory}, Chapter I, (3.1), (3.3) and (3.9) for
the integer ring and fractional ideal factorization, and (11.3), (11.5)
and the valuation-factorization paragraph on page 69 for local
valuations. The adelic product formula was fixed in the preceding
lesson; Neukirch, Chapter III, (1.3), states the same normalized
formula. No global class field theory, reciprocity map or classification
of the connected component for arbitrary \(K\) is asserted here.

\subsection{References}\label{NT-ADL-03--references}

\begin{itemize}
\tightlist
\item
  J. Neukirch, \emph{Algebraic Number Theory}, Chapter VI, §1,
  especially (1.1)--(1.6), supplies the ideal-class and logarithmic
  viewpoints. Its norm-one compactness argument uses the replete
  ideal-class group; Theorem 3.3 above instead uses the additive box
  lemma.
\item
  D. Harari, \emph{Galois Cohomology and Class Field Theory}, §12.3,
  Proposition 12.20, Theorem 12.21, Lemma 12.23 and Theorems 12.24 and
  12.27, for the topology comparison and its arithmetic consequences,
  including \(S\)-units.
\item
  F. Halter-Koch, \emph{Class Field Theory and L Functions: Foundations
  and Main Results}, §5.3, Theorems 5.3.6--5.3.8, for norm-one topology,
  compactness and finiteness results.
\item
  A. Deitmar, \emph{Automorphic Forms}, §5.3, Lemma 5.3.1, Theorem 5.3.3
  and Proposition 5.3.4, for the rational idèle topology and
  decomposition.
\item
  J. R. Getz and H. Hahn, \emph{An Introduction to Automorphic
  Representations},
  \href{https://sites.duke.edu/jgetz/files/2022/04/Graduate_Text.pdf}{author-hosted
  draft of April 22, 2022}, §2.6, equations (2.14)--(2.19), for the
  relation to adelic quotients. The 2024 published book is a later
  edition.
\item
  U. Jannsen, \emph{Algebraic Number Theory II}, §17, Definitions
  17.1--17.3 and Proposition 17.4, for the idèle and ideal descriptions.
\item
  J.-B. Bost and A. Connes,
  \href{https://doi.org/10.1007/BF01589495}{\emph{Hecke algebras, type
  III factors and phase transitions with spontaneous symmetry breaking
  in number theory}}, §3, for the finite rational split
  \(\mathbb A_f^\times=\mathbb Q_{>0}^\times\widehat{\mathbb Z}^{\times}\).
\end{itemize}

\section{Completions, the p-adic numbers and complete discretely valued
fields}\label{NT-LOC-02}

\emph{Written by GPT-6.1 Sol (OpenAI) in Codex, Ultra setting, October
2026. Self-checked by the writing AI. Independent full-lesson AI review
is not yet recorded. Public domain (CC0).}

\subsection{Introduction}\label{NT-LOC-02--introduction}

An absolute value tells us which approximations become accurate.
Completion supplies limits for all Cauchy approximations without
changing the arithmetic already present. For the p-adic absolute value,
accuracy means agreement modulo a high power of a prime. This makes the
completed ring of integers both an inverse limit of finite rings and a
space of infinite digit sequences.

We construct completions, develop the arithmetic and topology of
\(\mathbb Q_p\), and extend the digit construction to every complete
discretely valued field. We then prove that the complete archimedean
fields are precisely \(\mathbb R\) and \(\mathbb C\), with powers of
their usual absolute values. This identifies all archimedean places of a
number field.

The prerequisites are \emph{Absolute values, valuations and
Ostrowski\textquotesingle s theorem} and
\href{https://kokunoyumeto.github.io/program-matematika-indonesia/en/\#course-C10}{\emph{Real
Analysis I}}. We assume the absolute-value axioms, the ultrametric
inequality, and the elementary real-analysis facts about completeness
and compactness. The precise algebraic facts used without proof are
collected at the end. Basic references are {[}Milne{]}, {[}Neukirch{]},
{[}Deitmar{]} and {[}Harari{]}; {[}Halter-Koch{]} provides a broader
account of limits and completeness. The local arithmetic developed here
also underlies the discussion of places in \emph{Weil\textquotesingle s
proof for curves and what is missing over the integers}.

Write \(K_v\) for completion at a place \(v\), and
\(\mathbb N=\{0,1,\ldots\}\). For a nonarchimedean absolute value, its
valuation ring and residue field are

\[\mathcal O=\{x:|x|\le1\},\qquad
\mathfrak m=\{x:|x|<1\},\qquad
\kappa=\mathcal O/\mathfrak m.\]

An additive discrete valuation is normalized to have image
\(\mathbb Z\), and \(v(0)=+\infty\). A uniformizer \(\pi\) has
\(v(\pi)=1\). Thus \(|x|=c^{v(x)}\) for a fixed \(0<c<1\).

\subsection{Completing a field}\label{NT-LOC-02--completing-a-field}

A sequence \((a_n)\) is \emph{Cauchy} if for every \(\varepsilon>0\),
eventually \(|a_n-a_m|<\varepsilon\) for all \(n,m\). It is \emph{null}
if \(|a_n|\to0\). Completion identifies two Cauchy sequences when their
difference is null.

\textbf{Theorem 2.1 (completion).} Every field \(K\) with an absolute
value admits a complete valued field \(\widehat K\) containing \(K\)
isometrically and densely. Every isometric field embedding from \(K\) to
a complete valued field extends uniquely to an isometric embedding from
\(\widehat K\). Consequently two completions have a unique isometric
field isomorphism that respects their embeddings of \(K\).

If the absolute value is nonarchimedean, so is its extension, and

\[|\widehat K^{\times}|=|K^{\times}|,
\qquad \kappa\xrightarrow{\sim}\widehat\kappa.\]

Here \(\widehat\kappa\) denotes the residue field of \(\widehat K\).
There is no discreteness hypothesis in these last assertions.

\textbf{Proof.} Every Cauchy sequence is bounded: its tail lies within
distance \(1\) of a fixed term, and its initial segment is finite. Let
\(C\) be the ring of Cauchy sequences, with coordinatewise operations.
The estimate

\[|a_nb_n-a_mb_m|
\le |a_n|\,|b_n-b_m|+|b_m|\,|a_n-a_m|\]

shows that products are Cauchy. Boundedness also shows that the null
sequences form an ideal \(N\). Set \(\widehat K=C/N\).

The reverse triangle inequality makes \((|a_n|)\) a Cauchy sequence of
real numbers. Define

\[|[(a_n)]|=\lim_n|a_n|.\]

This is independent of the representative. It vanishes exactly on \(N\);
multiplication and the triangle inequality pass to limits. If the class
is nonzero, \(|a_n|\) is eventually bounded below by a positive number.
On that tail,

\[|a_n^{-1}-a_m^{-1}|=\frac{|a_n-a_m|}{|a_n|\,|a_m|}.\]

Hence the reciprocal sequence, with arbitrary initial terms, defines an
inverse. Thus \(\widehat K\) is a field. Constant sequences embed \(K\)
isometrically.

For \(A=[(a_n)]\), the elements \(a_n\in K\) converge to \(A\): the
Cauchy condition bounds \(\lim_m|a_m-a_n|\) by any prescribed positive
tolerance for large \(n\). This proves density.

Now let \((A_j)\) be Cauchy in \(\widehat K\). Choose \(b_j\in K\) with
\(|A_j-b_j|<1/j\). The triangle inequality shows that \((b_j)\) is
Cauchy in \(K\). Its class \(B\) satisfies \(b_j\to B\), so \(A_j\to B\)
as well. This proves completeness.

If \(f:K\to L\) is isometric and \(L\) is complete, send \([(a_n)]\) to
\(\lim_n f(a_n)\). Equivalent sequences give the same limit; addition,
multiplication and absolute values commute with this operation. Density
makes this the only isometric extension. Applying the construction in
both directions between two completions gives inverse maps.

In the nonarchimedean case the ultrametric inequality passes to limits.
If \(x\in\widehat K\) is nonzero, choose \(a\in K\) with \(|x-a|<|x|\).
The strict-dominance rule for an ultrametric gives \(|a|=|x|\), proving
equality of value sets. This argument also shows that the valuation ring
\(\widehat{\mathcal O}\) is the closure of \(\mathcal O\):
approximations sufficiently close to an integral element remain
integral, and the valuation ring is closed. To represent the residue of
\(x\in\widehat{\mathcal O}\), approximate it by \(a\in K\) with
\(|x-a|<1\). Then \(a\in\mathcal O\) and the residues agree. The kernel
of the induced residue map is exactly \(\mathfrak m\). ∎

Uniqueness here includes the embedding of the original field. It does
not say that the completed field has no other automorphisms. A trivially
valued field is already complete, since every Cauchy sequence is
eventually constant.

\subsection{Congruences as p-adic
accuracy}\label{NT-LOC-02--congruences-as-p-adic-accuracy}

Fix a prime \(p\). The field \(\mathbb Q_p\) is the completion of
\(\mathbb Q\) for \(|x|_p=p^{-v_p(x)}\). Define

\[\mathbb Z_p=\{x\in\mathbb Q_p:v_p(x)\ge0\}.\]

Theorem 2.1 extends \(v_p\) with the same integer value group and
residue field \(\mathbb F_p\). In particular,

\[x\equiv y\pmod{p^n\mathbb Z_p}
\quad\Longleftrightarrow\quad |x-y|_p\le p^{-n}.\]

An \emph{inverse limit} of the rings \(\mathbb Z/p^n\mathbb Z\) consists
of compatible residues:

\[\varprojlim_{n\ge1}\mathbb Z/p^n\mathbb Z
=\{(r_n):r_{n+1}\bmod p^n=r_n\}.\]

Operations are coordinatewise. Each quotient is given the discrete
topology, and the limit has the subspace topology from their product.
Specifying a residue modulo \(p^n\) specifies an open neighborhood.

\textbf{Proposition 2.2.} The closure of \(\mathbb Z\) in
\(\mathbb Q_p\) is \(\mathbb Z_p\). Reduction gives an isomorphism of
topological rings

\[\mathbb Z_p\xrightarrow{\sim}
\varprojlim_{n\ge1}\mathbb Z/p^n\mathbb Z.\]

The ring \(\mathbb Z_p\) is compact. The field \(\mathbb Q_p\) is
locally compact and totally disconnected.

\textbf{Proof.} By Theorem 2.1, \(\mathbb Z_p\) is the closure of
\(\mathbb Z_{(p)}=\{a/b\mid p\nmid b\}\). For every
\(a/b\in\mathbb Z_{(p)}\) and \(n\ge1\), choose an integer \(c\) such
that \(bc\equiv a\pmod{p^n}\). Then \(|a/b-c|_p\le p^{-n}\). Thus
\(\mathbb Z\) has the same closure.

Density shows that every coset modulo \(p^n\mathbb Z_p\) contains an
integer. An integer belongs to this ideal exactly when it is divisible
by \(p^n\), so

\[\mathbb Z_p/p^n\mathbb Z_p\simeq\mathbb Z/p^n\mathbb Z.\]

The resulting map to the inverse limit is injective, because a nonzero
element has a finite valuation. For a compatible family \((r_n)\), take
its integer representatives \(0\le c_n<p^n\). Compatibility gives
\(|c_m-c_n|_p\le p^{-n}\) when \(m\ge n\). Their limit lies in
\(\mathbb Z_p\) and has all the prescribed residues. This proves
surjectivity. The fibers of reduction modulo \(p^n\) are the cosets of
\(p^n\mathbb Z_p\), which form a neighborhood basis. Thus both
directions are continuous.

Compatibility also gives unique digits \(a_i\in\{0,\ldots,p-1\}\) with

\[c_n=\sum_{i=0}^{n-1}a_ip^i,
\qquad x=\sum_{i\ge0}a_ip^i.\]

A cylinder fixing the first \(n\) digits is a coset modulo \(p^n\); it
has exactly \(p\) children fixing one more digit. To prove compactness
directly, suppose an open cover has no finite subcover. Starting with
the whole space, repeatedly choose a child cylinder with no finite
subcover. The chosen compatible digits define a point \(x\). An open
member of the cover containing \(x\) contains a sufficiently long
cylinder in this chain, a contradiction.

The cosets \(a+p^n\mathbb Z_p\), for \(a\in\mathbb Q_p\) and
\(n\in\mathbb Z\), are open and closed: small changes preserve
membership, and the same holds for their complements. They are also
compact, by translation and multiplication by \(p^n\). This proves local
compactness. Given distinct points, choose such a coset containing one
and excluding the other. Its intersection with any connected set
separates those points, so every connected subset is a singleton. ∎

The whole field is not compact: the increasing open cover
\(\mathbb Q_p=\bigcup_{m\ge0}p^{-m}\mathbb Z_p\) has no finite subcover.
Compactness of the ring and local compactness of the field are different
assertions.

\subsection{Digits in any complete discrete
valuation}\label{NT-LOC-02--digits-in-any-complete-discrete-valuation}

Let \(K\) be complete for a normalized discrete valuation \(v\), with
ring \(\mathcal O\), residue field \(\kappa\), and uniformizer \(\pi\).
Choose a set \(R\subset\mathcal O\) containing exactly one
representative of each residue class, including \(0\). These
representatives need not form a subring.

\textbf{Proposition 2.3 (expansions and series).} Every nonzero
\(x\in K\) has a unique expansion

\[x=\sum_{i\ge m}a_i\pi^i,
\qquad m=v(x),\quad a_i\in R,\quad a_m\ne0.\]

Zero has the all-zero expansion. Allowing leading zero digits does not
change the element, so the starting index is unique only with the stated
normalization. A series \(\sum_{n\ge0}x_n\) in \(K\) converges if and
only if \(x_n\to0\). Moreover,

\[\mathcal O\xrightarrow{\sim}\varprojlim_{n\ge1}\mathcal O/\pi^n\mathcal O\]

is an isomorphism of topological rings, with discrete quotients.

\textbf{Proof.} Necessity for series follows by subtracting consecutive
partial sums. For sufficiency, the ultrametric inequality gives

\[\left|\sum_{n=r}^{s}x_n\right|\le\max_{r\le n\le s}|x_n|.\]

If the terms tend to zero, the partial sums are Cauchy and therefore
converge. This argument actually works in every complete nonarchimedean
field.

For \(y\in\mathcal O\), define \(y_0=y\). Choose \(a_j\in R\)
representing \(y_j\bmod\pi\), and set

\[y_{j+1}=\frac{y_j-a_j}{\pi}\in\mathcal O.\]

Then

\[y=\sum_{j=0}^{N-1}a_j\pi^j+\pi^Ny_N.\]

The remainder has absolute value at most \(c^N\), so it tends to zero.
For nonzero \(x\), apply this to \(y=\pi^{-v(x)}x\), a unit; its first
digit is nonzero. Every prescribed digit sequence converges because
\(|a_i\pi^i|\le c^i\).

If two expansions first differ at index \(j\), the difference of their
digits at \(j\) is a unit: distinct representatives have distinct
residues. The remaining tail has valuation at least \(j+1\). Thus the
difference of the expansions has valuation exactly \(j\), proving
uniqueness.

For the inverse limit, the kernel is \(\bigcap_n\pi^n\mathcal O=0\).
Lift a compatible family of residues to elements \(b_n\in\mathcal O\).
Then \(b_m-b_n\in\pi^n\mathcal O\) for \(m\ge n\), so they converge. The
limit has each prescribed residue, since \(\pi^n\mathcal O\) is closed.
The same cosets are basic neighborhoods on both sides, proving the
topological assertion. ∎

The first \(N\) digits determine an integral element modulo \(\pi^N\).
This is the useful meaning of truncation. Addition and multiplication
may require carrying because \(R\) need not respect either operation. In
mixed characteristic there is no reason for \(\kappa\) itself to embed
as a coefficient field.

The ring \(\mathcal O\) is a DVR: its units have valuation zero, and in
a nonzero ideal choose an element with least valuation. That element
generates the ideal, since division by it leaves every other element in
\(\mathcal O\). Its maximal ideal is \(\pi\mathcal O\). If \(\kappa\) is
finite, the finite-cylinder argument of Proposition 2.2 proves that
\(\mathcal O\) is compact and \(K\) is locally compact.

\subsection{Three computations}\label{NT-LOC-02--three-computations}

\textbf{Negative one.} For every prime \(p\),

\[\sum_{i=0}^{N-1}(p-1)p^i=p^N-1.\]

Since \(p^N\to0\) p-adically,

\[-1=\sum_{i\ge0}(p-1)p^i.\]

There is no conflict with the growth of the ordinary integer partial
sums: the metric has changed.

\textbf{One third in \(\mathbb Q_5\).} Successive digit extraction gives

{\def\LTcaptype{none} % do not increment counter
\begin{longtable}[]{@{}
  >{\raggedright\arraybackslash}p{(\linewidth - 4\tabcolsep) * \real{0.3333}}
  >{\raggedleft\arraybackslash}p{(\linewidth - 4\tabcolsep) * \real{0.3333}}
  >{\raggedright\arraybackslash}p{(\linewidth - 4\tabcolsep) * \real{0.3333}}@{}}
\toprule\noalign{}
\begin{minipage}[b]{\linewidth}\raggedright
Current remainder \(y_j\)
\end{minipage} & \begin{minipage}[b]{\linewidth}\raggedleft
Digit \(a_j\) modulo \(5\)
\end{minipage} & \begin{minipage}[b]{\linewidth}\raggedright
Next remainder \((y_j-a_j)/5\)
\end{minipage} \\
\midrule\noalign{}
\endhead
\bottomrule\noalign{}
\endlastfoot
\(1/3\) & \(2\) & \(-1/3\) \\
\(-1/3\) & \(3\) & \(-2/3\) \\
\(-2/3\) & \(1\) & \(-1/3\) \\
\end{longtable}
}

The last two rows repeat. Thus, in increasing order of powers,

\[\frac13=2+3\cdot5+5^2+3\cdot5^3+5^4+\cdots
=2+\sum_{j\ge0}(3\cdot5^{2j+1}+5^{2j+2}).\]

As a check, the repeated block sums to \((15+25)/(1-25)=-5/3\), so the
total is \(1/3\). The geometric series converges because \(|25|_5<1\).

\textbf{Laurent series.} For \(k=\mathbb F_q\), let \(k((t))\) consist
of sums \(\sum_{i\ge m}a_it^i\) with \(a_i\in k\) and only finitely many
negative powers. Addition and convolution multiplication are well
defined: each coefficient of a product involves finitely many pairs. A
nonzero series is \(a_mt^m(1+h)\), with \(h\in tk[[t]]\), and its
inverse is

\[a_m^{-1}t^{-m}\sum_{j\ge0}(-h)^j.\]

Each coefficient of this expression stabilizes, and multiplication
verifies the inverse. Thus these series form a field. Put
\(v_t(x)=\min\{i:a_i\ne0\}\) and \(|x|=q^{-v_t(x)}\).

A Cauchy sequence has a common lower bound on its exponents after
discarding finitely many terms; closeness to a fixed term proves this.
At each exponent its coefficient eventually stabilizes. The stabilized
coefficients give a Laurent series, and the sequence converges to it
because agreement below any fixed exponent is eventually exact. Hence
\(k((t))\) is complete. Its Laurent-polynomial truncations lie in
\(k(t)\) and are dense. By Theorem 2.1,

\[\widehat{k(t)}_{\,t}=k((t)),\qquad \mathcal O=k[[t]],\qquad\kappa=k.\]

Here the representatives really do form a coefficient field, so no
carrying is needed.

\subsection{Completing a number field at a
prime}\label{NT-LOC-02--completing-a-number-field-at-a-prime}

Let \(K\) be a number field, \(A=\mathcal O_K\) its ring of algebraic
integers, and \(\mathfrak p\) a nonzero prime ideal. We use the fact
that \(A\) is Dedekind and \(B=A_{\mathfrak p}\) is a DVR {[}Milne,
Theorem 3.29 and Proposition 3.6{]}. Their fraction field is \(K\):
multiplying any element of \(K\) by an integer divisible by the
denominators of its monic rational polynomial makes it integral.
Complete \(K\) for the valuation of \(B\), obtaining
\(K_{\mathfrak p}\), and write \(\mathcal O_{\mathfrak p}\) for its
valuation ring.

By Theorem 2.1, \(B\) is dense in \(\mathcal O_{\mathfrak p}\), with the
same uniformizer \(\pi\) and residue field. For every \(n\ge1\),
approximation modulo \(\pi^n\) gives

\[B/\pi^nB\simeq
\mathcal O_{\mathfrak p}/\pi^n\mathcal O_{\mathfrak p}.\]

Injectivity follows from equality of the valuations on \(B\);
surjectivity follows from density.

Also \(A/\mathfrak p^n\simeq B/\mathfrak p^nB\). Indeed, the only
maximal ideal containing \(\mathfrak p^n\) is \(\mathfrak p\). Thus
every \(s\notin\mathfrak p\) becomes a unit in \(A/\mathfrak p^n\), and
localizing that quotient changes nothing. Since \(\mathfrak pB=\pi B\),
Proposition 2.3 yields

\[\mathcal O_{\mathfrak p}\simeq
\varprojlim_{n\ge1}A/\mathfrak p^n,
\qquad
K\text{ is dense in }K_{\mathfrak p}.\]

The displayed ring isomorphism is topological when the quotients are
discrete. In particular \(A\) itself is dense in
\(\mathcal O_{\mathfrak p}\), since it represents every finite quotient.

Localization and completion are separate operations. The ring \(B\) is
not complete: it is countable, whereas the distinct digit sequences
using just two residue representatives already give uncountably many
elements in its completion. The general algebraic result {[}Stacks, Tag
0AP1{]} says that a Noetherian local ring is a DVR exactly when its
maximal-ideal completion is a DVR. It agrees with the direct valuation
construction here.

\subsection{Why the archimedean possibilities stop at
two}\label{NT-LOC-02--why-the-archimedean-possibilities-stop-at-two}

An absolute value is \emph{archimedean} when it does not satisfy the
ultrametric inequality. The prerequisite result that bounded absolute
values of integers force the ultrametric inequality implies that an
archimedean field has characteristic zero. Its absolute value on
\(\mathbb Q\) is \(|q|_\infty^s\) for \(s>0\), by
Ostrowski\textquotesingle s theorem on absolute values of \(\mathbb Q\).
Since \(|n|\le n\), necessarily \(s\le1\).

We first justify a polynomial fact used in the proof. A nonconstant
complex polynomial \(P\) satisfies \(|P(z)|_\infty\to\infty\) as
\(|z|_\infty\to\infty\), so its modulus has a minimum at some \(z_0\).
If \(P(z_0)\ne0\), write

\[\frac{P(z_0+w)}{P(z_0)}=1+cw^k+O(|w|^{k+1}),\qquad c\ne0.\]

Choose a unit complex number \(u\) with \(cu^k=-|c|_\infty\), and put
\(w=ru\). For small positive \(r\), the modulus is at most
\(1-|c|_\infty r^k+Cr^{k+1}<1\), contradicting minimality. Thus \(P\)
has a root. Division and induction show that every complex polynomial
factors into linear factors, counted with multiplicity. This proves the
needed form of the fundamental theorem of algebra.

\textbf{Theorem 2.4 (Ostrowski\textquotesingle s theorem on complete
archimedean fields).} If a field \(K\) is complete for an archimedean
absolute value, there is a field isomorphism \(\sigma:K\to\mathbb R\) or
\(\sigma:K\to\mathbb C\) such that

\[|x|=|\sigma(x)|_\infty^s\qquad(0<s\le1).\]

Thus it is an isometry to the indicated field with that powered absolute
value, and a topological field isomorphism to the field with its usual
absolute value.

\emph{Reference:} {[}Neukirch, Chapter II, Theorem 4.2{]}.

\textbf{Proof.} Normalize by setting \(\|x\|=|x|^{1/s}\). It is
necessary to check the triangle inequality after taking this power. The
binomial expansion, together with \(|n|=n^s\), gives for positive
integers \(N\)

\[|x+y|^N
\le\sum_{j=0}^{N}\binom Nj^s\|x\|^{sj}\|y\|^{s(N-j)}
\le(N+1)(\|x\|+\|y\|)^{sN}.\]

Taking \(sN\)-th roots and letting \(N\to\infty\) proves the claim. The
normalized absolute value has the same Cauchy sequences and agrees with
the usual one on \(\mathbb Q\). Its closure in \(K\) is complete: a
Cauchy sequence there converges in \(K\), and closedness keeps the limit
there. Uniqueness of completion therefore embeds \(\mathbb R\)
isometrically in \(K\).

Fix \(x\in K\). For \(z\in\mathbb C\), where \(\mathbb C\) is for now
only a parameter space, put

\[Q_z(T)=T^2-2\operatorname{Re}(z)T+|z|_\infty^2,
\qquad F(z)=\|Q_z(x)\|.\]

Only real coefficients are evaluated in \(K\). If \(M=\|x\|\), the
reverse triangle inequality gives

\[F(z)\ge |z|_\infty^2-2|z|_\infty M-M^2.\]

Thus \(F\) tends to infinity at infinity and attains a minimum \(m\).
Among its minimizers choose \(z_0\) with maximal modulus. Suppose
\(m>0\), choose \(0<\varepsilon<m\), and put \(Q=Q_{z_0}\). The roots
\(z_1,\overline{z_1}\) of \(Q(T)+\varepsilon\) are nonreal conjugates,
with

\[|z_1|_\infty^2=|z_0|_\infty^2+\varepsilon.\]

Consequently \(F(z_1)>m\).

For each \(n\ge1\), factor the real polynomial

\[H_n(T)=Q(T)^n-(-\varepsilon)^n
=\prod_{j=1}^{2n}(T-\alpha_j),\]

listing \(\alpha_1=z_1\). Its conjugate factorization gives the
real-polynomial identity

\[H_n(T)^2=\prod_{j=1}^{2n}Q_{\alpha_j}(T).\]

Evaluation in \(K\), multiplicativity, and minimality yield

\[F(z_1)m^{2n-1}
\le\|H_n(x)\|^2
\le(m^n+\varepsilon^n)^2.\]

Hence \(F(z_1)/m\le(1+(\varepsilon/m)^n)^2\). Letting \(n\to\infty\)
contradicts \(F(z_1)>m\). Therefore \(m=0\), and every element of \(K\)
satisfies a monic real quadratic equation.

If \(K=\mathbb R\), we are done. Otherwise choose \(x\notin\mathbb R\).
Its quadratic is irreducible over \(\mathbb R\), since a split quadratic
would force \(x\) to be one of its real roots. Completing the square
gives \(j\in K\) with \(j^2=-1\), embedding \(\mathbb C=\mathbb R(j)\).
Every real quadratic now splits inside this subfield, so every element
of \(K\) belongs to it. Thus \(K=\mathbb C\).

Finally, for \(z=a+bj\) we have \(\|j\|=1\), so

\[\|z\|\le |a|_\infty+|b|_\infty\le\sqrt2\,|z|_\infty.\]

Apply this to \(z^n\) and take \(n\)-th roots to get
\(\|z\|\le|z|_\infty\). Applying it to \(z^{-1}\) proves the reverse
inequality. Undoing the normalization proves the theorem. ∎

This theorem concerns complete fields, whereas the earlier Ostrowski
theorem concerns absolute values on \(\mathbb Q\). Completeness is
essential: \(\mathbb Q\) with its usual absolute value is archimedean
but is neither of the two fields in the conclusion.

\subsection{Archimedean places of a number
field}\label{NT-LOC-02--archimedean-places-of-a-number-field}

\textbf{Corollary 2.5.} For a number field \(K\), the archimedean places
correspond bijectively to its real embeddings and its pairs of conjugate
nonreal embeddings into \(\mathbb C\). Their completions are
respectively \(\mathbb R\) and \(\mathbb C\). If the counts are \(r_1\)
and \(r_2\), there are \(r_1+r_2\) archimedean places, and
\([K:\mathbb Q]=r_1+2r_2\).

\textbf{Proof.} An embedding \(\sigma\) gives the absolute value
\(|\sigma(x)|_\infty\). If its image is real, its closure is
\(\mathbb R\), since it contains \(\mathbb Q\). If the image contains a
nonreal \(\alpha\), it contains \(\mathbb Q+\mathbb Q\alpha\), dense in
\(\mathbb C\): the real-linear map \((u,v)\mapsto u+v\alpha\) is an
isomorphism \(\mathbb R^2\to\mathbb C\). These closures are the
completions.

Conversely, Theorem 2.4 identifies any archimedean completion with one
of these fields. Composing \(K\) with this identification gives an
embedding whose usual absolute value represents the original place.

Suppose two embeddings give equivalent absolute values. Their
restrictions to \(\mathbb Q\) are the same usual absolute value, so the
exponent relating them is \(1\), as evaluation at \(2\) shows. The
correspondence between their images therefore extends to an isometric
isomorphism of their completions. Such a map fixes \(\mathbb Q\), hence
\(\mathbb R\) by continuity. It cannot identify \(\mathbb R\) with
\(\mathbb C\); on \(\mathbb C\) it sends \(i\) to \(i\) or \(-i\), so it
is identity or conjugation. This proves the asserted bijection.

For the degree formula, express \(K\) as a finite tower adjoining
algebraic generators. An embedding at each stage extends in precisely as
many ways as the degree of the next minimal polynomial: its transformed
polynomial splits in \(\mathbb C\), with distinct roots in
characteristic zero. Multiplying these counts gives \([K:\mathbb Q]\)
embeddings. Conjugation fixes exactly the real embeddings and pairs the
others. ∎

For example, \(\mathbb Q(\sqrt6)\) has two real places and two real
completions. The two embeddings of \(\mathbb Q(\sqrt{-6})\) form one
conjugate pair, giving one complex place and completion. Distinct places
can therefore have isomorphic completed fields.

\subsection{Exercises}\label{NT-LOC-02--exercises}

Digits below are always in \(\{0,\ldots,p-1\}\), listed in increasing
order of powers. A sequence is \emph{eventually periodic} if
\(a_{i+d}=a_i\) for some \(d\ge1\) and every sufficiently large \(i\).

\textbf{Exercise 1 (easy).} Compute the \(7\)-adic expansion of \(-1/6\)
and prove that it is eventually periodic.

\textbf{Solution.} It is purely periodic:

\[-\frac16=1+7+7^2+\cdots.\]

The first \(N\) terms sum to \((7^N-1)/6\), whose difference from
\(-1/6\) is \(7^N/6\), of valuation \(N\). The series converges to the
claimed value, and uniqueness identifies every digit as \(1\).

\textbf{Exercise 2 (easy).} Prove
\(\mathbb Z_p^{\times}=\mathbb Z_p\setminus p\mathbb Z_p\), and prove
that \(\mathbb Z_p\) is a DVR with uniformizer \(p\).

\textbf{Solution.} A nonzero \(x\in\mathbb Z_p\) is invertible there
exactly when \(v_p(x^{-1})=-v_p(x)\ge0\), which is equivalent to
\(v_p(x)=0\). The remaining elements, including zero, form
\(p\mathbb Z_p\). Thus this is the unique maximal ideal. In a nonzero
ideal \(I\), choose \(x\) with the least valuation \(n\). Write
\(x=p^nu\), with \(u\) a unit. For every \(y\in I\),
\(y/p^n\in\mathbb Z_p\), so \(I=p^n\mathbb Z_p\). This proves that the
ring is a local principal ideal domain, not a field, with uniformizer
\(p\).

\textbf{Exercise 3 (medium).} Show that \(x\in\mathbb Q_p\) belongs to
\(\mathbb Q\) exactly when its p-adic expansion is eventually periodic.

\textbf{Solution.} Multiplication by a power of \(p\) shifts finitely
many indices, preserving both rationality and eventual periodicity. We
may assume \(x\in\mathbb Z_p\).

If the tail starts at \(N\) and has period \(d\), put
\(B_0=\sum_{j=0}^{d-1}a_{N+j}p^j\). Then

\[x=\sum_{i<N}a_ip^i+\frac{p^NB_0}{1-p^d}\in\mathbb Q,\]

by the convergent geometric series.

Conversely write \(x=A/B\), where \(A\in\mathbb Z\), \(B>0\), and
\(p\nmid B\). Digit extraction has remainders \(A_i/B\), beginning with
\(A_0=A\). The digit \(a_i\) is the unique member of
\(\{0,\ldots,p-1\}\) congruent to \(A_iB^{-1}\pmod p\), and

\[A_{i+1}=\frac{A_i-Ba_i}{p}\in\mathbb Z.\]

For the ordinary absolute value,

\[|A_{i+1}|_\infty\le\frac{|A_i|_\infty+B(p-1)}p.\]

Thus all \(A_i\) lie in the finite integer interval \([-M,M]\), with
\(M=\max(|A|_\infty,B)\). The next state and digit depend only on the
current state. A state eventually repeats, after which the digits
repeat. This includes a terminating expansion, whose tail consists of
zeros.

\textbf{Exercise 4 (medium).} Show that \(\mathbb Z_p\) is homeomorphic
to the middle-thirds Cantor set, for every prime \(p\).

\textbf{Solution.} Its digit map identifies \(\mathbb Z_p\) with
\(D^{\mathbb N}\), where \(D=\{0,\ldots,p-1\}\) is discrete: agreeing on
the first \(n\) coordinates means agreement modulo \(p^n\). To identify
this with binary sequences, use the prefix code

\[0,\ 10,\ 110,\ \ldots,\ 1^{p-2}0,\ 1^{p-1}\]

for the \(p\) symbols, in order. For \(p=2\) it is just \(0,1\).
Concatenation gives a binary sequence. Decoding is unique: read to the
first zero, or stop after \(p-1\) ones, and repeat. Every infinite
binary sequence decodes this way, so the map is bijective. Code lengths
are between \(1\) and \(p-1\); hence agreement on \(n\) input symbols
forces agreement on at least \(n\) bits, and agreement on \((p-1)n\)
bits forces agreement on \(n\) decoded symbols. Both maps are
continuous.

Finally send \((b_i)\in\{0,1\}^{\mathbb N}\) to

\[\sum_{i\ge0}\frac{2b_i}{3^{i+1}}.\]

The two choices at each stage select the left or right third of the
remaining closed interval. Nested intervals have lengths tending to
zero, giving exactly all points of the Cantor set. If sequences first
differ at \(j\), their values differ by at least \(3^{-(j+1)}\), since
the possible later tail cancels at most half the leading difference.
Agreement on a long prefix bounds the difference above by the length of
its interval. These bounds prove injectivity and continuity in both
directions. Composing gives the desired homeomorphism. This concerns
topology; the coding does not identify ring operations with ordinary
real operations.

\textbf{Exercise 5 (hard).} Prove Theorem 2.4 by minimizing

\[(t,u)\longmapsto\|x^2-tx+u\|\]

over the whole real plane. Supply both the existence of a minimum and
the argument forcing it to be zero.

\textbf{Solution.} Use the normalization and embedded \(\mathbb R\)
established at the start of the theorem\textquotesingle s proof. If
\(x\in\mathbb R\), a quadratic vanishing at \(x\) is immediate. Suppose
\(x\notin\mathbb R\). On the Euclidean unit circle the continuous
function \((t,u)\mapsto\|tx-u\|\) is everywhere positive, since \(1,x\)
are real-linearly independent. It has a positive minimum \(c_0\).
Homogeneity implies

\[\|tx-u\|\ge c_0\sqrt{t^2+u^2},
\qquad
\|x^2-tx+u\|\ge c_0\sqrt{t^2+u^2}-\|x^2\|.\]

Thus the latter function is coercive. Its minimizers form a nonempty
compact set. Choose \((t_0,u_0)\) maximizing \(u\) among them, and put
\(Q(T)=T^2-t_0T+u_0\), \(m=\|Q(x)\|\).

Assume \(m>0\) and choose \(0<\varepsilon<m\). The roots
\(\alpha,\beta\) of \(Q+\varepsilon\) are either real or conjugate. If
conjugate, both have squared modulus \(u_0+\varepsilon>u_0\). If real,
at least one has squared modulus at least
\(|\alpha\beta|=|u_0+\varepsilon|>u_0\). Choose such a root \(\alpha\),
and define
\(Q_\alpha(T)=T^2-2\operatorname{Re}(\alpha)T+|\alpha|_\infty^2\). Its
constant coefficient is larger than \(u_0\), so maximality gives
\(\|Q_\alpha(x)\|>m\).

Factor \(H_n=Q^n-(-\varepsilon)^n\), with \(\alpha\) among its \(2n\)
roots. The identity \(H_n^2=\prod_jQ_{\alpha_j}\) is valid because
\(H_n\) has real coefficients. Every factor evaluated at \(x\) has norm
at least \(m\), since the minimum was taken over all real quadratic
coefficients. Hence

\[\frac{\|Q_\alpha(x)\|}{m}
\le\left(1+\left(\frac\varepsilon m\right)^n\right)^2.\]

Letting \(n\to\infty\) is a contradiction. Thus \(m=0\), and every \(x\)
satisfies a real quadratic. The last two paragraphs of the
theorem\textquotesingle s proof then identify the field and its absolute
value. This argument also explains why merely asserting that a
continuous function on \(\mathbb R^2\) has a minimum would be
insufficient: coercivity needs proof.

\subsection{What this lesson does not
prove}\label{NT-LOC-02--what-this-lesson-does-not-prove}

\begin{itemize}
\tightlist
\item
  The prerequisite real-analysis foundations: completeness of
  \(\mathbb R\), compactness of closed bounded Euclidean sets, and the
  extreme-value theorem. These are the facts from \emph{Real Analysis I}
  used in the minimization arguments; {[}Neukirch, Chapter II, proof of
  Theorem 4.2{]} uses the same foundations. The polynomial factorization
  needed here was proved above.
\item
  The criterion that an absolute value is nonarchimedean exactly when
  its values on integers are bounded, and the consequent
  characteristic-zero assertion for archimedean fields {[}Milne,
  Proposition 7.2 and Corollary 7.3{]}.
\item
  Ostrowski\textquotesingle s classification of nontrivial absolute
  values on \(\mathbb Q\) {[}Milne, Theorem 7.12{]}, and the equivalence
  criterion that two nontrivial absolute values give the same topology
  exactly when one is a positive power of the other {[}Milne,
  Proposition 7.8{]}. These belong to \emph{Absolute values, valuations
  and Ostrowski\textquotesingle s theorem}.
\item
  For a Dedekind domain \(A\), its integral closure in a finite
  separable extension of its fraction field is Dedekind {[}Milne,
  Theorem 3.29{]}; every localization of \(A\) at a nonzero prime is a
  DVR {[}Milne, Proposition 3.6{]}. We use the first assertion for
  \(A=\mathbb Z\).
\item
  For a Noetherian local ring \(A\), \(A\) is a DVR if and only if its
  maximal-ideal completion is a DVR {[}Stacks, Tag 0AP1{]}. This is a
  comparison with the more general algebraic theorem; the valued-field
  and DVR constructions needed in the lesson were proved directly.
\end{itemize}

Hensel\textquotesingle s lemma and the classification of nonarchimedean
local fields are developed in \emph{Hensel\textquotesingle s lemma,
squares and roots of unity in p-adic fields} and \emph{Local fields:
classification and Haar measure}.

The written ring-theoretic providers are \emph{Number fields}, lesson 3,
\href{https://kokunoyumeto.github.io/open-mathematics-courses/courses/NT-ANT/discrete-valuation-rings-and-dedekind-domains.html}{\textbf{Discrete
valuation rings and Dedekind domains}}, its local characterizations and
Proposition 3.1, and lesson 5,
\href{https://kokunoyumeto.github.io/open-mathematics-courses/courses/NT-ANT/decomposition-of-primes-in-extensions.html}{\textbf{Decomposition
of primes in extensions}}, Theorem 5.1 for finite integral closure over
an arbitrary Dedekind base in a finite separable field extension. The
integer-bound criterion, equivalence criterion and rational
classification are Proposition 1.1, Proposition 3.1 and Theorem 4.1 of
\textbf{Absolute values, valuations and Ostrowski\textquotesingle s
theorem}; the references below credit the classical treatments.

\subsection{References}\label{NT-LOC-02--references}

\begin{itemize}
\tightlist
\item
  {[}Milne{]} J. S. Milne, \emph{Algebraic Number Theory}, version 3.08
  (2020), especially Chapter 7, ``Completions'' and ``Completions in the
  nonarchimedean case.''
  \href{https://www.jmilne.org/math/CourseNotes/ANT.pdf}{Open lecture
  notes}.
\item
  {[}Stacks{]} The \href{https://stacks.math.columbia.edu/}{Stacks
  project}, cited by tag. The linked text of
  \href{https://kokunoyumeto.github.io/stacks-zh-hans-cn/en/more-algebra.html\#more-algebra-lemma-completion-dvr}{Tag
  0AP1} is in AI Integrated Stacks Project, an edition with AI-proposed
  corrections and AI-written additions, not reviewed by the Stacks
  project\textquotesingle s maintainers.
\item
  {[}Neukirch{]} J. Neukirch, \emph{Algebraic Number Theory}, translated
  by N. Schappacher, Grundlehren der mathematischen Wissenschaften 322,
  Springer (1999), Chapter II, §§1, 2 and 4.
\item
  {[}Deitmar{]} A. Deitmar, \emph{Automorphic Forms}, Universitext,
  Springer (2013), §§4.2 and 4.5.
\item
  {[}Harari{]} D. Harari, \emph{Galois Cohomology and Class Field
  Theory}, Universitext, Springer (2020), §7.2.
\item
  {[}Halter-Koch{]} F. Halter-Koch, \emph{Class Field Theory and L
  Functions: Foundations and Main Results}, CRC Press (2022),
  §§1.3--1.4.
\end{itemize}

\section{Haar uniqueness through shrinking
neighborhoods}\label{fremlin-neighbourhood-uniqueness}

\emph{Adapted from D. H. Fremlin, Measure Theory, by GPT-6.1 Sol
(OpenAI), Ultra reasoning effort, 3 October 2026. The neighborhood-ratio
argument is Fremlin\textquotesingle s; the compact shrinking argument
and the final Borel-measure step translate it to the convention of this
course. Self-checked by the adapting AI; no independent AI or human
review is claimed.}

This proof compares the masses of small symmetric neighborhoods. It
gives an alternative to
\hyperlink{haar-measure-on-locally-compact-groups--oa-fnd-hm-06}{Theorem
9.2 of the Haar lesson}, using the open-set case of the Radon-product
theorem. Throughout, a Radon measure has the course\textquotesingle s
convention: it is finite on compact sets, outer regular on Borel sets
and inner regular on open sets.

\textbf{Theorem.} Let \(G\) be an arbitrary locally compact Hausdorff
group, and let \(\mu,\nu\) be nonzero left-invariant Radon measures on
\(G\). There is a unique \(c>0\) such that \(\mu=c\nu\). Neither measure
is assumed sigma-finite.

\emph{Proof.} By
\hyperlink{haar-measure-on-locally-compact-groups--oa-fnd-hm-06}{Proposition
9.1}, each nonempty open set has positive measure for both measures.
Relatively compact open sets have finite measure for both. Let
\(\mathcal U\) be the symmetric relatively compact open neighborhoods of
the identity, ordered by reverse inclusion. These form a neighborhood
base: intersect an identity neighborhood with its inverse and with a
relatively compact identity neighborhood.

Fix a nonempty relatively compact open set \(O\), and
\(0<\varepsilon<1\). Inner regularity gives a compact \(K\subset O\)
with \[\mu(K)>(1-\varepsilon)\mu(O).\] There are an open \(H\supset K\)
and an identity neighborhood \(V\) such that \(HV\subset O\). Indeed,
continuity of multiplication supplies a box \(H_x\times V_x\) whose
product lies in \(O\) for each \(x\in K\); take finitely many \(H_x\)
covering \(K\), their union for \(H\), and the intersection of the
corresponding \(V_x\) for \(V\). Shrink \(V\) to a member of
\(\mathcal U\).

For \(U\in\mathcal U\) with \(U\subset V\), put
\[W=\{(x,y)\in O\times O:x^{-1}y\in U\}.\] This is open in
\(G\times G\).
\hyperlink{haar-measure-on-locally-compact-groups--oa-fnd-hm-10}{Theorem
5.1(1)} identifies the two integrals of its sections with
\((\mu\widehat\times\nu)(W)\); its open-set assertion does not require
sigma-finiteness. If \(x\in H\), the vertical section contains \(xU\). A
horizontal section at \(y\in O\) is contained in \(yU\), since
\(U=U^{-1}\). Left invariance therefore gives
\[(1-\varepsilon)\mu(O)\nu(U)
<\mu(H)\nu(U)
\le (\mu\widehat\times\nu)(W)
\le\mu(U)\nu(O).\] All the factors being divided by are positive and
finite. Repeating the argument with the two measures exchanged gives,
for all sufficiently small \(U\in\mathcal U\),
\[(1-\varepsilon)\frac{\mu(O)}{\nu(O)}
\le\frac{\mu(U)}{\nu(U)}
\le\frac1{1-\varepsilon}\frac{\mu(O)}{\nu(O)}.\] Thus the neighborhood
net of ratios has limit \(\mu(O)/\nu(O)\). The net is the same for every
such \(O\), so these ratios all have one common value \(c>0\).

Let \(A\) be any open set and \(K\subset A\) any compact set. Compact
shrinking supplies a relatively compact open \(O\) with
\(K\subset O\subset A\). Hence \(\mu(K)\le\mu(O)=c\nu(O)\le c\nu(A)\).
Taking the supremum over \(K\), and then exchanging the measures, proves
\(\mu(A)=c\nu(A)\), also when this value is infinite. Finally outer
regularity gives \(\mu(E)=c\nu(E)\) for every Borel \(E\). Evaluating
any nonempty relatively compact open set proves uniqueness of \(c\).
\(\square\)

The signed-function kernel argument in the
\hyperref[quotient-measures-and-weils-integration-formula--exactly-when-the-quotient-measure-is-invariant]{quotient
lesson} uses a different compact cutoff. It must equal one on the
projection of the support of the original function; cancellation can
make the support of its subgroup average smaller.

\subsection{Source and
terms}\label{fremlin-neighbourhood-uniqueness--source-and-terms}

D. H. Fremlin,
\href{https://www1.essex.ac.uk/maths/people/fremlin/mt4.2013/index.htm}{\emph{Measure
Theory}, Volume 4, Topological Measure Spaces}, supplies the
neighborhood-ratio proof. Copyright © 1998 D. H. Fremlin. The original
supplied
\href{https://www1.essex.ac.uk/maths/people/fremlin/mt4.2013/mt442.tex}{editable
TeX} remains attributed to Fremlin. This separately titled adaptation is
distributed under the \hyperref[design-science-licence]{Design Science
License}, with editable Markdown source. Changes on 3 October 2026:
restriction to the course\textquotesingle s LCH/Borel convention,
expansion of compact shrinking, and replacement of the quasi-Radon
base-uniqueness step by inner regularity on open sets and outer
regularity on Borel sets. Adaptation and additions: GPT-6.1 Sol
(OpenAI), Ultra.

THE WORK IS PROVIDED "AS IS," AND COMES WITH ABSOLUTELY NO WARRANTY,
EXPRESS OR IMPLIED, TO THE EXTENT PERMITTED BY APPLICABLE LAW. The full
warranty and liability terms are in the accompanying Design Science
License.

\section{Design Science License}\label{design-science-licence}

DESIGN SCIENCE LICENSE

TERMS AND CONDITIONS FOR COPYING, DISTRIBUTION AND MODIFICATION

Copyright © 1999-2001 Michael Stutz
\textless stutz@dsl.org\textgreater{} Verbatim copying of this document
is permitted, in any medium.

0. PREAMBLE.

Copyright law gives certain exclusive rights to the author of a work,
including the rights to copy, modify and distribute the work (the
"reproductive," "adaptative," and "distribution" rights).

The idea of "copyleft" is to willfully revoke the exclusivity of those
rights under certain terms and conditions, so that anyone can copy and
distribute the work or properly attributed derivative works, while all
copies remain under the same terms and conditions as the original.

The intent of this license is to be a general "copyleft" that can be
applied to any kind of work that has protection under copyright. This
license states those certain conditions under which a work published
under its terms may be copied, distributed, and modified.

Whereas "design science" is a strategy for the development of artifacts
as a way to reform the environment (not people) and subsequently improve
the universal standard of living, this Design Science License was
written and deployed as a strategy for promoting the progress of science
and art through reform of the environment.

1. DEFINITIONS.

"License" shall mean this Design Science License. The License applies to
any work which contains a notice placed by the work\textquotesingle s
copyright holder stating that it is published under the terms of this
Design Science License.

"Work" shall mean such an aforementioned work. The License also applies
to the output of the Work, only if said output constitutes a "derivative
work" of the licensed Work as defined by copyright law.

"Object Form" shall mean an executable or performable form of the Work,
being an embodiment of the Work in some tangible medium.

"Source Data" shall mean the origin of the Object Form, being the
entire, machine-readable, preferred form of the Work for copying and for
human modification (usually the language, encoding or format in which
composed or recorded by the Author); plus any accompanying files,
scripts or other data necessary for installation, configuration or
compilation of the Work.

(Examples of "Source Data" include, but are not limited to, the
following: if the Work is an image file composed and edited in PNG
format, then the original PNG source file is the Source Data; if the
Work is an MPEG 1.0 layer 3 digital audio recording made from a WAV
format audio file recording of an analog source, then the original WAV
file is the Source Data; if the Work was composed as an unformatted
plaintext file, then that file is the Source Data; if the Work was
composed in LaTeX, the LaTeX file(s) and any image files and/or custom
macros necessary for compilation constitute the Source Data.)

"Author" shall mean the copyright holder(s) of the Work.

The individual licensees are referred to as "you."

2. RIGHTS AND COPYRIGHT.

The Work is copyrighted by the Author. All rights to the Work are
reserved by the Author, except as specifically described below. This
License describes the terms and conditions under which the Author
permits you to copy, distribute and modify copies of the Work.

In addition, you may refer to the Work, talk about it, and (as dictated
by "fair use") quote from it, just as you would any copyrighted material
under copyright law.

Your right to operate, perform, read or otherwise interpret and/or
execute the Work is unrestricted; however, you do so at your own risk,
because the Work comes WITHOUT ANY WARRANTY -\/- see Section 7 ("NO
WARRANTY") below.

3. COPYING AND DISTRIBUTION.

Permission is granted to distribute, publish or otherwise present
verbatim copies of the entire Source Data of the Work, in any medium,
provided that full copyright notice and disclaimer of warranty, where
applicable, is conspicuously published on all copies, and a copy of this
License is distributed along with the Work.

Permission is granted to distribute, publish or otherwise present copies
of the Object Form of the Work, in any medium, under the terms for
distribution of Source Data above and also provided that one of the
following additional conditions are met:

(a) The Source Data is included in the same distribution, distributed
under the terms of this License; or

(b) A written offer is included with the distribution, valid for at
least three years or for as long as the distribution is in print
(whichever is longer), with a publicly-accessible address (such as a URL
on the Internet) where, for a charge not greater than transportation and
media costs, anyone may receive a copy of the Source Data of the Work
distributed according to the section above; or

(c) A third party\textquotesingle s written offer for obtaining the
Source Data at no cost, as described in paragraph (b) above, is included
with the distribution. This option is valid only if you are a
non-commercial party, and only if you received the Object Form of the
Work along with such an offer.

You may copy and distribute the Work either gratis or for a fee, and if
desired, you may offer warranty protection for the Work.

The aggregation of the Work with other works that are not based on the
Work -\/- such as but not limited to inclusion in a publication,
broadcast, compilation, or other media -\/- does not bring the other
works in the scope of the License; nor does such aggregation void the
terms of the License for the Work.

4. MODIFICATION.

Permission is granted to modify or sample from a copy of the Work,
producing a derivative work, and to distribute the derivative work under
the terms described in the section for distribution above, provided that
the following terms are met:

(a) The new, derivative work is published under the terms of this
License.

(b) The derivative work is given a new name, so that its name or title
cannot be confused with the Work, or with a version of the Work, in any
way.

(c) Appropriate authorship credit is given: for the differences between
the Work and the new derivative work, authorship is attributed to you,
while the material sampled or used from the Work remains attributed to
the original Author; appropriate notice must be included with the new
work indicating the nature and the dates of any modifications of the
Work made by you.

5. NO RESTRICTIONS.

You may not impose any further restrictions on the Work or any of its
derivative works beyond those restrictions described in this License.

6. ACCEPTANCE.

Copying, distributing or modifying the Work (including but not limited
to sampling from the Work in a new work) indicates acceptance of these
terms. If you do not follow the terms of this License, any rights
granted to you by the License are null and void. The copying,
distribution or modification of the Work outside of the terms described
in this License is expressly prohibited by law.

If for any reason, conditions are imposed on you that forbid you to
fulfill the conditions of this License, you may not copy, distribute or
modify the Work at all.

If any part of this License is found to be in conflict with the law,
that part shall be interpreted in its broadest meaning consistent with
the law, and no other parts of the License shall be affected.

7. NO WARRANTY.

THE WORK IS PROVIDED "AS IS," AND COMES WITH ABSOLUTELY NO WARRANTY,
EXPRESS OR IMPLIED, TO THE EXTENT PERMITTED BY APPLICABLE LAW, INCLUDING
BUT NOT LIMITED TO THE IMPLIED WARRANTIES OF MERCHANTABILITY OR FITNESS
FOR A PARTICULAR PURPOSE.

8. DISCLAIMER OF LIABILITY.

IN NO EVENT SHALL THE AUTHOR OR CONTRIBUTORS BE LIABLE FOR ANY DIRECT,
INDIRECT, INCIDENTAL, SPECIAL, EXEMPLARY, OR CONSEQUENTIAL DAMAGES
(INCLUDING, BUT NOT LIMITED TO, PROCUREMENT OF SUBSTITUTE GOODS OR
SERVICES; LOSS OF USE, DATA, OR PROFITS; OR BUSINESS INTERRUPTION)
HOWEVER CAUSED AND ON ANY THEORY OF LIABILITY, WHETHER IN CONTRACT,
STRICT LIABILITY, OR TORT (INCLUDING NEGLIGENCE OR OTHERWISE) ARISING IN
ANY WAY OUT OF THE USE OF THIS WORK, EVEN IF ADVISED OF THE POSSIBILITY
OF SUCH DAMAGE.

END OF TERMS AND CONDITIONS

\end{document}
